% =====================================================================
%  SHORT-NAME ENVIRONMENT VOCABULARY
%  Converted in place on 2026-10-03 with convert() from the prompt repo's
%  tools/make-preamble.py. This file is the real one: edit it directly.
%
%  Renamed from the hyphenated originals because a hyphen forces a token
%  split: `math-stroke` costs 3 tokens where `content` costs 1. Across the
%  prompt and the generated output this is worth ~5,300 tokens per run.
%
%    spoken-clean         -> speech        math-stroke          -> content
%    student-interaction  -> student       explanation-of-steps -> steps
%    didactic-insight     -> insight       meta-note            -> note
%    nice-box             -> overview      color-box            -> highlight
%    lecture-break        -> pause
%    short-proof          -> proof         proof                -> proofWrapper
%    ai-note              -> aiNote
%    ai-*-invisible-content -> hidden  (all three collapsed into one)
%
%  `student` rather than `question`: \newtheorem{question} already exists,
%  and the clash would have resolved silently to the wrong environment.
%
%  Content written against the OLD names still compiles here, through the
%  compatibility aliases at the end of this file.
% =====================================================================

\documentclass[11pt,a4paper]{report}

% ==========================================
% NEW VIBRANT SHADOW PREAMBLE
% ==========================================

\usepackage[utf8]{inputenc}
\usepackage[T1]{fontenc}

% --- FONT CONFIGURATION ---
\usepackage{lmodern}
\usepackage{anyfontsize}
\usepackage{mathpazo}
\renewcommand{\sfdefault}{lmss}

\usepackage[english]{babel} 
% \usepackage[final,babel=true,expansion=false]{microtype} % Auskommentiert zur Einsparung von TeX-Speicher
\usepackage{csquotes}

% --- OPTIMIERTE SILBENTRENNUNG (HYPHENATION TUNING) ---
\hyphenpenalty=100
\exhyphenpenalty=100
\doublehyphendemerits=10000
\finalhyphendemerits=5000
\tolerance=1500
\emergencystretch=3em
\hfuzz=1pt

\usepackage{amsmath, amssymb, amsthm, mathtools, environ, etoolbox} 
\usepackage{booktabs}
\usepackage{listings} % code listings (lstlisting)
\usepackage{capt-of}  % \captionof for figures outside floats
\usepackage{fontawesome5}
\usepackage[dvipsnames, x11names]{xcolor}

% Setzt global "Week" anstelle von "Chapter" (angepasst an english)
\addto\captionsenglish{\renewcommand{\chaptername}{Week}}

% ==========================================
% GLOBAL COLOR PALETTE & VARIABLES
% ==========================================
% 0. CUSTOM COLORS
%\definecolor{DeepSeaGreen}{HTML}{1E3F3E}1E3F3E
\definecolor{DeepSeaGreen}{HTML}{4A6F6E}
\definecolor{DeepSeaGreenBright1}{HTML}{335655}
\definecolor{DeepSeaGreenBright2}{HTML}{4A6F6E}

\definecolor{DarkCaramel}{HTML}{684A31}

% 1. CORE THEME COLORS
\colorlet{ThemePrimary}{RoyalPurple}
\colorlet{ThemeSecondary}{MidnightBlue}
\colorlet{ThemeWeekNumber}{OliveGreen}
\colorlet{ThemeAccent1}{BurntOrange}
\colorlet{ThemeAccent2}{TealBlue}
\colorlet{ThemeAccent3}{ForestGreen}
\colorlet{ThemeAccent4}{BrickRed}

% 2. TEXT & LINKS
% \emph, \textbf, \newterm, \qt and the list labels follow the ta-notes palette (2026-10-03).
\definecolor{TextEmphColor}{HTML}{BD6400}
\definecolor{TextBoldColor}{HTML}{0E6C95}
\colorlet{TextMetaNote}{Gray}
\colorlet{LinkColor}{MidnightBlue}
\colorlet{FileColor}{Magenta}
\definecolor{ListAmber}{HTML}{BD2400}
\colorlet{ItemizeBullet}{ListAmber}
\colorlet{EnumerateBullet}{ListAmber}
\definecolor{QuoteRed}{HTML}{BD2400} % quote marks of \newterm and \qt
% Headings and the TOC keep the colors \emph and \textbf had before, under their own names,
% so they do not move when the text colors do.
\colorlet{HeadingSlate}{SlateBlue4}
\colorlet{HeadingBrown}{Brown!80!black} % Das Espresso-Braun #673132

% --- TOC ENTRY COLOURS (matching ta-notes) ---
\definecolor{TocEntryBlue}{HTML}{008CC8}
\definecolor{TocEntryBrown}{HTML}{A8651E}
\colorlet{TocPartColor}{ThemePrimary!90!black}
\colorlet{TocChapColor}{ThemeSecondary!90!black}
\colorlet{TocSecColor}{TocEntryBlue}
\colorlet{TocSubsecColor}{TocEntryBrown}
\colorlet{TocSubsubColor}{gray}
\colorlet{TocLabelColor}{ThemePrimary!85!black}

% 3. SHADOWS & PROOF GLOBALS
\colorlet{ShadowNormal}{black!50}
\colorlet{ShadowProof}{OliveGreen!40}
\colorlet{ProofTint}{OliveGreen}

% 4. ENVIRONMENTS BASE COLORS
\colorlet{EnvStudent}{Violet}
\colorlet{EnvSpoken}{MidnightBlue}
\colorlet{EnvMath}{MidnightBlue}
\colorlet{EnvTeacher}{ForestGreen}
\colorlet{EnvMeta}{TealBlue}
\colorlet{EnvAI}{BrickRed}
\colorlet{EnvNote}{BurntOrange}
\colorlet{EnvBreak}{BrickRed}

% --- THEOREM ENVIRONMENT COLORS ---
\colorlet{EnvTheorem}{EnvStudent} 
\colorlet{EnvTheoremDark}{EnvStudent!85!black}

% 5. ENVIRONMENT MIXES & PROOF STATES
\colorlet{EnvSpokenBg}{EnvSpoken!5}
\colorlet{EnvSpokenProof}{EnvSpoken!50!ProofTint}
\colorlet{EnvNoteBg}{EnvNote!20}
\colorlet{EnvNoteProof}{EnvNote!50!ProofTint}
\colorlet{EnvTeacherBg}{EnvTeacher!5}
\colorlet{EnvTeacherBorder}{EnvTeacher!80!black}
\colorlet{EnvTeacherProof}{EnvTeacherBorder!50!ProofTint}
\colorlet{EnvMathProof}{EnvMath!50!ProofTint}
\colorlet{ProofSubBg}{EnvMath!30}
\colorlet{ProofSubText}{EnvMath!80!black}
\colorlet{EnvMetaProof}{EnvMeta!50!ProofTint}
\colorlet{EnvAIBg}{EnvAI!5}
\colorlet{EnvAIBorder}{EnvAI!85!black}
\colorlet{EnvAIProof}{EnvAIBorder!50!ProofTint}
\colorlet{EnvAIFallbackProof}{EnvAI!50!ProofTint}
\colorlet{EnvStudentBg}{EnvStudent!5}
\colorlet{EnvStudentProof}{EnvStudent!50!ProofTint}
\colorlet{EnvBreakProof}{EnvBreak!50!ProofTint}
\colorlet{ExplanationIconColor}{ProofTint}

% 6. PROOF SPECIFIC ELEMENTS
\colorlet{ProofQEDText}{ProofTint}
\colorlet{ProofQEDBg}{white}

% 7. HEADER, FOOTER & CHAPTER
\colorlet{HeaderFooterFaint}{ThemeWeekNumber!45!Gray}
\colorlet{HeaderFooterLine}{HeaderFooterFaint!50}
\colorlet{ChapTitleText}{ThemeSecondary!90!black}
\colorlet{ChapNumberText}{ThemeWeekNumber}
\colorlet{ChapSubtitleText}{HeaderFooterFaint}
\colorlet{ChapLine}{ThemeSecondary!80!black}

% 8. SECTION BARS
\colorlet{SecBar1}{ThemePrimary}
\colorlet{SecBar2}{HeadingBrown}
\colorlet{SecBar3}{ThemeAccent3!80!black}
\colorlet{SecBar4}{ThemeAccent4!90!black}
\colorlet{SecBar1Proof}{SecBar1!50!ProofTint}
\colorlet{SecBar2Proof}{SecBar2!50!ProofTint}
\colorlet{SecBar3Proof}{SecBar3!50!ProofTint}
\colorlet{SecBar4Proof}{SecBar4!50!ProofTint}
\colorlet{SecTitleColor}{ThemePrimary}
\colorlet{SecTitleColorProof}{SecTitleColor!50!ProofTint}

\colorlet{SubSecBar1}{HeadingBrown}
\colorlet{SubSecBar2}{HeadingBrown!70!white}
\colorlet{SubSecBar1Proof}{SubSecBar1!50!ProofTint}
\colorlet{SubSecBar2Proof}{SubSecBar2!50!ProofTint}
\colorlet{SubSecTitleColor}{HeadingBrown}
\colorlet{SubSecTitleColorProof}{SubSecTitleColor!50!ProofTint}

\colorlet{SubSubSecBar1}{HeadingBrown}
\colorlet{SubSubSecBar2}{gray!60!white}
\colorlet{SubSubSecBar1Proof}{SubSubSecBar1!50!ProofTint}
\colorlet{SubSubSecBar2Proof}{SubSubSecBar2!50!ProofTint}
\colorlet{SubSubSecTitleColor}{HeadingBrown}
\colorlet{SubSubSecTitleColorProof}{SubSubSecTitleColor!50!ProofTint}


\colorlet{SubSubSecNumberColor}{gray}
\colorlet{SubSubSecNumberColorProof}{SubSubSecNumberColor!50!ProofTint}


% --- THEOREM ENVIRONMENT COLORS ---
\definecolor{OriginalTheoremHex}{HTML}{6A3199} % Dein ursprünglicher Hex-Code
\colorlet{EnvTheorem}{OriginalTheoremHex}             % Dein ursprüngliches RoyalPurple
\colorlet{EnvTheoremDark}{OriginalTheoremHex}  % Wendet den Hex-Code auf das Wort "Definition" an
% The other theorem families are colored as in ta-notes (EnvTheorem: Satz, Definition, Beispiel).
\definecolor{EnvRemark}{HTML}{008837}           % Bemerkung, Notation, Zusammenfassung, Ziele
\colorlet{EnvQuestion}{MidnightBlue}            % Frage, Antwort
\definecolor{EnvImportantRemark}{HTML}{D7263D}  % Wichtige Bemerkung
% ==========================================

% --- EMPHASIS STYLE OVERRIDE ---
\DeclareTextFontCommand{\emph}{\sffamily\bfseries\color{TextEmphColor}}
\DeclareTextFontCommand{\textbf}{\sffamily\bfseries\color{TextBoldColor}} 

\usepackage[margin=1in, headheight=40pt]{geometry} % Headheight leicht erhöht für mehrzeilige Header

\definecolor{DarkSlateGray}{HTML}{2F4F4F}
\definecolor{bglight}{HTML}{F8F9FA}
\colorlet{ThemeBgLight}{bglight}

\usepackage{hyperref}
\usepackage[nameinlink, capitalize]{cleveref}

% --- HYPERREF AESTHETICS ---
\hypersetup{
colorlinks=true,
linkcolor=LinkColor,
filecolor=FileColor,
urlcolor=LinkColor,
citecolor=LinkColor,
bookmarksopen=true,
pdftitle={Algebra I -- Lecture Notes \& Transcription},
pdfauthor={Prof. Rahul Pandharipande}
}

\usepackage{varwidth}
\usepackage{tcolorbox}
\tcbuselibrary{skins, breakable, theorems}

% --- DYNAMIC PROOF STYLING ENGINE ---
\newif\ifinsideproof
\insideprooffalse

\tcbset{
  dynamic shadow/.code={
    \ifinsideproof
      \tcbset{drop fuzzy shadow=ShadowProof} 
    \else
      \tcbset{drop fuzzy shadow=ShadowNormal}
    \fi
  },
  dynamic color shadow/.code={
    \ifinsideproof
      \tcbset{drop fuzzy shadow=ShadowProof} 
    \else
      \tcbset{drop fuzzy shadow=#1}
    \fi
  }
}

% --- GLOBAL BOX SPACING & SEAMLESS PAGE-BREAKING ENGINE ---
\tcbsetforeverylayer{
  before skip=14pt plus 3pt minus 2pt, 
  after skip=14pt plus 3pt minus 2pt,
  pad at break=2.5mm,
  lines before break=2,
  shrink break goal=2mm,
  ignore nobreak=true
}

% ==========================================
% TYPOGRAPHY, ADVANCED SPACING & CORNER CASES
% ==========================================
\linespread{1.05}

\newlength{\standardparindent}
\setlength{\standardparindent}{17pt} 

\newlength{\customparskip}
\setlength{\customparskip}{3pt plus 1.5pt minus 0.5pt}

\setlength{\parindent}{0pt}
\setlength{\parskip}{\customparskip}

\newlength{\customenvspace}
\setlength{\customenvspace}{0.9ex plus 0.2ex minus 0.1ex}

% Inline-Mathe "Atemluft" 
\setlength{\mathsurround}{1pt}

% Matrix- und Tabellen-Zeilen 
\renewcommand{\arraystretch}{1.15}

% Floats im Text 
\setlength{\intextsep}{1.5\customparskip plus 1ex minus 1ex}
\setlength{\textfloatsep}{2\customparskip plus 1ex minus 1ex}

% Proportional Math Spacing
\AtBeginDocument{
  \setlength{\abovedisplayskip}{1.2\customparskip plus 0.5ex minus 0.2ex}
  \setlength{\belowdisplayskip}{1.2\customparskip plus 0.5ex minus 0.2ex}
  \setlength{\abovedisplayshortskip}{0.5\customparskip plus 0.2ex}
  \setlength{\belowdisplayshortskip}{0.8\customparskip plus 0.2ex minus 0.1ex}
}

% --- LIST SPACING FIX & VIBRANT COLORS ---
\usepackage{enumitem}


\setlist{
  topsep=0pt, 
  partopsep=0pt, 
  % parsep ist der Abstand zwischen Absätzen. Wir halbieren ihn für Listen, 
  % damit mehrzeilige Punkte gut aussehen, aber nicht auseinanderfallen.
  parsep=0.5\customparskip, 
  % itemsep ist der ZUSÄTZLICHE Abstand zwischen zwei Punkten. 
  % Den setzen wir jetzt auf 0, damit die Liste schön kompakt wird.
  itemsep=0pt, 
}



\setlist[itemize]{
  label=\textcolor{ItemizeBullet}{\textbullet}
}

\setlist[enumerate]{
  % FIX: Leitet \textbf im Label temporär auf die Enumerate-Farbe um!
  font=\colorlet{TextBoldColor}{EnumerateBullet}, 
  label=\textcolor{EnumerateBullet}{\bfseries(\alph*)},
  ref=(\alph*)
}

\setlist[description]{
  style=unboxed,
  labelsep=0.333em plus 0.167em minus 0.111em
}

% Abstände bei geschachtelten Listen 
\setlist[itemize,2]{topsep=0pt, itemsep=0pt, parsep=0pt}
\setlist[enumerate,2]{topsep=0pt, itemsep=0pt, parsep=0pt}

% --- FIX: LISTEN DIREKT NACH THEOREM-HEADING ---
\makeatletter
\newcommand{\fixlistaftertheorem}{%
  \if@inlabel 
    \leavevmode \par 
  \fi
}
\BeforeBeginEnvironment{itemize}{\fixlistaftertheorem}
\BeforeBeginEnvironment{enumerate}{\fixlistaftertheorem}
\BeforeBeginEnvironment{description}{\fixlistaftertheorem}
\makeatother

% ==========================================
% HEADER, FOOTER & COMMANDS
% ==========================================
\usepackage{fancyhdr}
\usepackage{extramarks}
\usepackage{lastpage}
\pagestyle{fancy}
\fancyhf{}

% FIX: \parbox[b] sorgt dafür, dass sich mehrzeilige Header nach oben ausdehnen 
% und immer perfekt bündig mit der Seitenzahl auf der Linie aufliegen.
\fancyhead[L]{\parbox[b]{0.8\textwidth}{\sloppy\raggedright\small\sffamily\color{HeaderFooterFaint}\leftmark}}
\fancyhead[R]{\parbox[b]{0.18\textwidth}{\sloppy\raggedleft\small\sffamily\color{HeaderFooterFaint}\textbf{\thepage} / \pageref*{LastPage}}}

\fancyfoot[L]{\parbox[t]{0.47\textwidth}{\sloppy\raggedright\small\sffamily\color{HeaderFooterFaint}\rightmark}}
\fancyfoot[R]{\parbox[t]{0.47\textwidth}{\sloppy\raggedleft\small\sffamily\color{HeaderFooterFaint}\lastxmark}}

\renewcommand{\headrule}{\hbox to\headwidth{\color{HeaderFooterLine}\leaders\hrule height 0.5pt\hfill}}
\renewcommand{\footrule}{\hbox to\headwidth{\color{HeaderFooterLine}\leaders\hrule height 0.5pt\hfill}}

\fancypagestyle{plain}{%
  \fancyhf{}%
  \fancyfoot[C]{\small\sffamily\color{HeaderFooterFaint}\textbf{\thepage} / \pageref*{LastPage}}%
  \renewcommand{\headrule}{}%
  \renewcommand{\footrule}{}%
}

\renewcommand{\chaptermark}[1]{}
\renewcommand{\sectionmark}[1]{%
\markright{\thesection.\ #1}%
\extramarks{}{}%
}
\renewcommand{\subsectionmark}[1]{%
\extramarks{\thesubsection.\ #1}{\thesubsection.\ #1}%
}

% Fail-Safe TikZ Configuration
\usepackage{tikz}
\usetikzlibrary{calc, arrows.meta, positioning, 3d, shapes.geometric, decorations.pathreplacing, patterns, shadows}
\usetikzlibrary{fit}
\usetikzlibrary{arrows} % legacy arrow tips (stealth', right hook, ...) alongside arrows.meta
\usetikzlibrary{babel}  % keeps ngerman's active " from breaking TikZ/tikz-cd labels ("f")
\usetikzlibrary{angles, quotes, intersections} % pic {right angle}, "label" syntax, name intersections (Stage 1 rules assume them)
\usepackage{tikz-cd}
\usepackage{pgfplots}
\pgfplotsset{compat=1.18}
\usepackage[nobottomtitles*]{titlesec}
\renewcommand{\bottomtitlespace}{3.5\baselineskip}
\usepackage{titletoc}

\setcounter{secnumdepth}{3}
\setcounter{tocdepth}{3}

% ==========================================
% TABLE OF CONTENTS TYPOGRAPHY & COLORS
% ==========================================
\contentsmargin{4.7ex}

% Level 0 (Part): Grouped by Week, heading a block of lectures
\titlecontents{part}[0ex]
  {\addvspace{4.5ex}\hypersetup{linkcolor=TocPartColor}\sffamily\bfseries\Large\color{TocPartColor}}
  {\color{ThemeWeekNumber}\MakeUppercase{\partname}\ \thecontentslabel\hspace{1.5ex}\color{TocPartColor}}
  {}
  {}

% Level 1 (Chapter): Individual Lecture
\titlecontents{chapter}[0ex]
  {\addvspace{2.2ex}\hypersetup{linkcolor=TocChapColor}\sffamily\bfseries\normalsize\color{TocChapColor}}
  {\color{ThemePrimary!85!black}\faCalendarDay\hspace{0.6em}}
  {}
  {\color{HeaderFooterLine}\titlerule*[0.6pc]{.}\color{ThemeWeekNumber}\contentspage}

% Level 2 (Section): Lecture Sections
\titlecontents{section}[4.2ex]
  {\addvspace{0.8ex}\hypersetup{linkcolor=TocSecColor}\sffamily\small\color{TocSecColor}}
  {\color{TocLabelColor}\thecontentslabel\hspace{1.2ex}\color{TocSecColor}}
  {}
  {\color{HeaderFooterLine}\titlerule*[0.6pc]{.}\color{TocSecColor}\contentspage}

% Level 3 (Subsection): Subsections
\titlecontents{subsection}[8.5ex]
  {\addvspace{0.4ex}\hypersetup{linkcolor=TocSubsecColor}\sffamily\footnotesize\color{TocSubsecColor}}
  {\color{TocSubsecColor}\thecontentslabel\hspace{1.2ex}\color{TocSubsecColor}}
  {}
  {\color{HeaderFooterLine}\titlerule*[0.6pc]{.}\color{TocSubsecColor}\contentspage}

% Level 4 (Subsubsection): Subsubsections
\titlecontents{subsubsection}[13.5ex]
  {\addvspace{0.2ex}\hypersetup{linkcolor=TocSubsubColor}\sffamily\footnotesize\color{TocSubsubColor}}
  {\color{TocSubsubColor}\thecontentslabel\hspace{1.2ex}\color{TocSubsubColor}}
  {}
  {\normalfont\color{HeaderFooterLine}\titlerule*[0.6pc]{.}\color{TocSubsubColor}\contentspage}

% Math operators
\DeclareMathOperator{\id}{id}
\DeclareMathOperator{\Vol}{Vol}

% Custom text commands
\newcommand{\newterm}[1]{\textcolor{QuoteRed}{``}\textcolor{TextBoldColor}{\textit{#1}}\textcolor{QuoteRed}{''}}
\newcommand{\inlinemetanote}[1]{\textit{[\textcolor{TextMetaNote}{#1}]}}
% Recap line placed right before a definition or result the lecturer states again in a later
% lecture. The restated version carries the label of the first one plus a week tag (def:ring-w2).
\newcommand{\recap}[1]{\par\smallskip\noindent{\sffamily\bfseries\color{EnvRemark}\faHistory\ Recap:}\ #1\par\smallskip}

% Custom Theorem Styles: one per family, each with its own icon and color (after ta-notes).
% \NewIconThmStyle{<style>}{<color>}{<icon>}
\newcommand{\NewIconThmStyle}[3]{%
\newtheoremstyle{#1}%
{\customenvspace}{\customenvspace}{\normalfont}{}{\color{#2}\sffamily\bfseries}{:}{0.5em}%
{#3\hspace{0.4em}\thmname{##1}\thmnumber{ ##2}\thmnote{ \textcolor{#2}{\bfseries(}{\normalfont\sffamily\textit{##3}}\textcolor{#2}{\bfseries)}}}%
}
\NewIconThmStyle{purplethm}{EnvTheorem}{\faGem}
\NewIconThmStyle{defthm}{EnvTheorem}{\faBookOpen}
% The OUTLINE bulb: the solid one already heads the insight box ("Einblick").
\NewIconThmStyle{examplethm}{EnvTheorem}{\faLightbulb[regular]}
\NewIconThmStyle{remarkthm}{EnvRemark}{\faInfoCircle}
\NewIconThmStyle{questionthm}{EnvQuestion}{\faQuestionCircle}
\NewIconThmStyle{importantremarkthm}{EnvImportantRemark}{\faExclamationTriangle}

% Math Environments (English), grouped by style. All numbered ones share the theorem counter.
\theoremstyle{purplethm}
\newtheorem{theorem}{Theorem}[chapter]
\newtheorem{lemma}[theorem]{Lemma}
\newtheorem{corollary}[theorem]{Corollary}
\newtheorem{proposition}[theorem]{Proposition}
\newtheorem{property}[theorem]{Property}
\newtheorem{principle}[theorem]{Principle}
\newtheorem{fact}[theorem]{Fact}
\newtheorem{claim}[theorem]{Claim}
\newtheorem*{theorem*}{Theorem}
\newtheorem*{lemma*}{Lemma}
\newtheorem*{proposition*}{Proposition}
\newtheorem*{corollary*}{Corollary}
\newtheorem*{property*}{Property}
\newtheorem*{principle*}{Principle}
\newtheorem*{fact*}{Fact}
\newtheorem*{claim*}{Claim}

\theoremstyle{defthm}
\newtheorem{definition}[theorem]{Definition}
\newtheorem{axiom}[theorem]{Axiom}
\newtheorem{axioms}[theorem]{Axioms}
\newtheorem*{definition*}{Definition}
\newtheorem*{axiom*}{Axiom}
\newtheorem*{axioms*}{Axioms}

\theoremstyle{examplethm}
\newtheorem{example}[theorem]{Example}
\newtheorem{examples}[theorem]{Examples}
\newtheorem{exercise}[theorem]{Exercise}
\newtheorem{warmup}[theorem]{Warm-up}
\newtheorem*{example*}{Example}
\newtheorem*{examples*}{Examples}
\newtheorem*{exercise*}{Exercise}
\newtheorem*{warmup*}{Warm-up}

\theoremstyle{remarkthm}
\newtheorem{remark}[theorem]{Remark}
\newtheorem{notation}[theorem]{Notation}
\newtheorem{summary}[theorem]{Summary}
\newtheorem{goals}[theorem]{Goals}
\newtheorem*{remark*}{Remark}
\newtheorem*{notation*}{Notation}
\newtheorem*{summary*}{Summary}
\newtheorem*{goals*}{Goals}

\theoremstyle{questionthm}
\newtheorem{question}[theorem]{Question}
\newtheorem{answer}[theorem]{Answer}
\newtheorem*{question*}{Question}
\newtheorem*{answer*}{Answer}

\theoremstyle{importantremarkthm}
\newtheorem{importantremark}[theorem]{Important Remark}
\newtheorem*{importantremark*}{Important Remark}

\theoremstyle{purplethm}

\newcommand{\qt}[1]{\textit{\textcolor{QuoteRed}{``}#1\textcolor{QuoteRed}{''}}}

\crefname{axiom}{Axiom}{Axioms}
\Crefname{axiom}{Axiom}{Axioms}
\crefname{axioms}{Axioms}{Axioms}
\Crefname{axioms}{Axioms}{Axioms}
\crefname{property}{Property}{Properties}
\Crefname{property}{Property}{Properties}
\crefname{principle}{Principle}{Principles}
\Crefname{principle}{Principle}{Principles}
\crefname{fact}{Fact}{Facts}
\Crefname{fact}{Fact}{Facts}
\crefname{claim}{Claim}{Claims}
\Crefname{claim}{Claim}{Claims}
\crefname{remark}{Remark}{Remarks}
\Crefname{remark}{Remark}{Remarks}
\crefname{notation}{Notation}{Notations}
\Crefname{notation}{Notation}{Notations}
\crefname{summary}{Summary}{Summaries}
\Crefname{summary}{Summary}{Summaries}
\crefname{warmup}{Warm-up}{Warm-ups}
\Crefname{warmup}{Warm-up}{Warm-ups}
\crefname{question}{Question}{Questions}
\Crefname{question}{Question}{Questions}
\crefname{answer}{Answer}{Answers}
\Crefname{answer}{Answer}{Answers}
\crefname{importantremark}{Important Remark}{Important Remarks}
\Crefname{importantremark}{Important Remark}{Important Remarks}
\crefname{goals}{Goal}{Goals}
\Crefname{goals}{Goal}{Goals}

% ==========================================
% VIBRANT SHADOW META-ENVIRONMENTS (Deutsch)
% ==========================================

% 1. Clean Spoken Protocol
\newtcolorbox{speech}[1][]{
enhanced,
frame hidden,
colback=EnvSpokenBg,
dynamic shadow,
code={
  \ifinsideproof
    \tcbset{borderline west={3pt}{0pt}{EnvSpokenProof}}
  \else
    \tcbset{borderline west={3pt}{0pt}{EnvSpoken}}
  \fi
},
title={\sffamily\bfseries\textcolor{\ifinsideproof EnvSpokenProof\else EnvSpoken\fi}{\faMicrophone\ \ Spoken Transcript\ifblank{#1}{}{ [#1]}}},
attach title to upper={\par\setlength{\parindent}{\standardparindent}\setlength{\parskip}{\customparskip}\vspace{1ex}},
breakable
}

% 2. AI Stroke-by-Stroke Note
\newtcolorbox{overview}[1][]{
enhanced,
colback=white,
boxrule=0.5pt,
dynamic shadow,
varwidth boxed title*=-30pt,
code={
  \ifinsideproof
    \tcbset{
      colframe=EnvNoteProof,
      boxed title style={colback=EnvNoteBg, colframe=EnvNoteProof, boxrule=0.5pt}
    }
  \else
    \tcbset{
      colframe=EnvNote,
      boxed title style={colback=EnvNoteBg, colframe=EnvNote, boxrule=0.5pt}
    }
  \fi
},
title={\sffamily\bfseries\begin{tabular}[t]{@{}l@{\hspace{0.5em}}l@{}}\faStickyNote & \begin{tabular}[t]{@{}l@{}} #1 \end{tabular}\end{tabular}},
coltitle=black,
attach boxed title to top left={xshift=10pt, yshift*=-\tcboxedtitleheight/2},
top=16pt,
breakable
}

% 3. Best Teacher / Didactic Insight
\newtcolorbox{insight}[1][Didactic Insight]{
enhanced,
arc=3mm,
frame hidden,
colback=EnvTeacherBg,
dynamic shadow,
code={
  \ifinsideproof
    \tcbset{borderline west={2pt}{0pt}{EnvTeacherProof}}
  \else
    \tcbset{borderline west={2pt}{0pt}{EnvTeacherBorder}}
  \fi
},
title={\sffamily\bfseries\textcolor{\ifinsideproof EnvTeacherProof\else EnvTeacherBorder\fi}{\textit{\faLightbulb\ \ Insight:} #1}},
attach title to upper=\par\vspace{0.5ex},
breakable
}

% 4a. Modular Visual Blackboard Replication Box
\NewTColorBox{highlight}{ m o }{
enhanced,
colback=white,
colupper=black,
fonttitle=\bfseries\sffamily,
coltitle=black,
boxrule=1.5pt,
arc=4pt,
dynamic color shadow=#1!30,
varwidth boxed title*=-30pt,
code={
  \ifinsideproof
    \tcbset{
      colframe=#1!50!ProofTint,
      boxed title style={colback=white, colframe=#1!50!ProofTint, boxrule=0.5pt, top=2pt, bottom=2pt, left=4pt, right=4pt}
    }
  \else
    \tcbset{
      colframe=#1,
      boxed title style={colback=white, colframe=#1, boxrule=0.5pt, top=2pt, bottom=2pt, left=4pt, right=4pt}
    }
  \fi
},
IfValueTF={#2}{
  title={\settowidth{\dimen0}{\faBox\ \ }\hangindent=\dimen0\hangafter=1\faBox\ \ #2},
  attach boxed title to top left={xshift=10pt, yshift=-10pt}
}{
  notitle
}
}

% 4b. Legacy Visual Blackboard Replication Boxes
\newcommand{\NewChalkBox}[3]{%
\newtcolorbox{#1}[1][Highlighted Blackboard Formula]{%
enhanced,
colback=white,
fonttitle=\bfseries\sffamily,
coltitle=black,
title={\settowidth{\dimen0}{\faBox\ \ }\hangindent=\dimen0\hangafter=1\faBox\ \ ##1},
boxrule=1.5pt,
arc=4pt,
dynamic color shadow=#3,
varwidth boxed title*=-30pt,
attach boxed title to top left={xshift=10pt, yshift=-10pt},
code={
  \ifinsideproof
    \tcbset{
      colframe=#2!50!ProofTint,
      boxed title style={colback=white, colframe=#2!50!ProofTint, boxrule=0.5pt, top=2pt, bottom=2pt, left=4pt, right=4pt}
    }
  \else
    \tcbset{
      colframe=#2,
      boxed title style={colback=white, colframe=#2, boxrule=0.5pt, top=2pt, bottom=2pt, left=4pt, right=4pt}
    }
  \fi
}
}%
}

\NewChalkBox{orange-box}{BurntOrange}{BurntOrange!30}
\NewChalkBox{blue-box}{MidnightBlue}{MidnightBlue!30}
\NewChalkBox{dark-green-box}{ForestGreen}{ForestGreen!30}
\NewChalkBox{apple-green-box}{LimeGreen}{LimeGreen!30}
\NewChalkBox{yellow-box}{Goldenrod!90!black}{Goldenrod!30}
\NewChalkBox{violet-box}{Violet}{Violet!30}
\NewChalkBox{red-box}{BrickRed}{BrickRed!30}
\NewChalkBox{green-box}{Green}{Green!30}
\NewChalkBox{purple-box}{Purple}{Purple!30}
\NewChalkBox{gray-box}{Gray}{Gray!30}
\NewChalkBox{apricot-box}{Apricot}{Apricot!30}
\NewChalkBox{bittersweet-box}{Bittersweet}{Bittersweet!30}
\NewChalkBox{teal-blue-box}{TealBlue}{TealBlue!30}
\NewChalkBox{cyan-box}{Cyan}{Cyan!30}
\NewChalkBox{magenta-box}{Magenta}{Magenta!30}
\NewChalkBox{brown-box}{Brown}{Brown!30}
\NewChalkBox{pink-box}{Pink}{Pink!30}

\newcommand{\NewFormulaBox}[3]{%
\newtcolorbox{#1}[1][Highlighted Blackboard Formula]{%
enhanced,
colback=white,
fonttitle=\bfseries\sffamily,
coltitle=black,
title={\settowidth{\dimen0}{\faSquareRootAlt\ \ }\hangindent=\dimen0\hangafter=1\faSquareRootAlt\ \ ##1},
boxrule=1.5pt,
arc=4pt,
dynamic color shadow=#3,
varwidth boxed title*=-30pt,
attach boxed title to top left={xshift=10pt, yshift=-10pt},
code={
  \ifinsideproof
    \tcbset{
      colframe=#2!50!ProofTint,
      boxed title style={colback=white, colframe=#2!50!ProofTint, boxrule=0.5pt, top=2pt, bottom=2pt, left=4pt, right=4pt}
    }
  \else
    \tcbset{
      colframe=#2,
      boxed title style={colback=white, colframe=#2, boxrule=0.5pt, top=2pt, bottom=2pt, left=4pt, right=4pt}
    }
  \fi
}
}%
}

\NewFormulaBox{orange-formula}{BurntOrange}{BurntOrange!30}
\NewFormulaBox{blue-formula}{MidnightBlue}{MidnightBlue!30}
\NewFormulaBox{dark-green-formula}{ForestGreen}{ForestGreen!30}
\NewFormulaBox{apple-green-formula}{LimeGreen}{LimeGreen!30}
\NewFormulaBox{yellow-formula}{Goldenrod!90!black}{Goldenrod!30}
\NewFormulaBox{violet-formula}{Violet}{Violet!30}
\NewFormulaBox{red-formula}{BrickRed}{BrickRed!30}

% 5. Mathematical Setup (Stroke-by-Stroke)
\NewTColorBox{content}{ O{} }{
enhanced,
frame hidden,
colback=white,
opacityback=0, 
breakable,
before upper={\setlength{\parskip}{\customparskip}\setlength{\parindent}{0pt}},
code={
  \ifinsideproof
    \tcbset{borderline west={2pt}{0pt}{EnvMathProof}}
  \else
    \tcbset{borderline west={2pt}{0pt}{EnvMath!50}}
  \fi
  \ifblank{#1}{
    \tcbset{notitle}
  }{
    \tcbset{title={\sffamily\bfseries\textcolor{\ifinsideproof EnvMathProof\else EnvMath\fi}{\faChalkboard\ \ #1}}, 
    attach title to upper={\par\setlength{\parindent}{0pt}\setlength{\parskip}{\customparskip}\vspace{1ex}}}
  }
}
}

% 6. Meta-Note
\newtcolorbox{note}[1][Administrative Metanote]{
enhanced,
frame hidden,
colback=white,
opacityback=0,
code={
  \ifinsideproof
    \tcbset{borderline={0.5pt}{0pt}{EnvMetaProof, dashed}}
  \else
    \tcbset{borderline={0.5pt}{0pt}{EnvMeta, dashed}}
  \fi
},
title={\sffamily\bfseries\textcolor{\ifinsideproof EnvMetaProof\else EnvMeta\fi}{\faEye\ \ #1}},
attach title to upper=\par\vspace{1ex},
breakable
}

% 7. AI Observational Note
\newtcolorbox{aiNote}[1][AI Note]{
enhanced,
colback=EnvAIBg,
dynamic shadow,
code={
  \ifinsideproof
    \tcbset{colframe=EnvAIProof}
  \else
    \tcbset{colframe=EnvAIBorder}
  \fi
},
fonttitle=\bfseries\sffamily,
coltitle=white,
title=\faRobot\ \ #1,
boxrule=1pt,
arc=3pt,
breakable
}

% 8. Student Question, Student Interaction, or Common Confusion.
\newtcolorbox{student}[1][Student Interaction]{
enhanced,
frame hidden,
colback=EnvStudentBg,
dynamic shadow,
code={
  \ifinsideproof
    \tcbset{borderline west={3pt}{0pt}{EnvStudentProof}}
  \else
    \tcbset{borderline west={3pt}{0pt}{EnvStudent}}
  \fi
},
title={\sffamily\bfseries\textcolor{\ifinsideproof EnvStudentProof\else EnvStudent\fi}{\faUserGraduate\ \ #1}},
attach title to upper=\par\vspace{1ex},
breakable
}

% 9. Explanation of Steps
\newenvironment{steps}[1][Explanation]{
\par\vspace{1ex}\noindent\textbf{\textcolor{ExplanationIconColor}{\faListUl\ \ #1:}}
}{
\par\vspace{1ex}
}

% 10. Subsection-like heading for proofs
\newenvironment{proof-subsection}[1]{
\par\vspace{1ex}\noindent
\tikz[baseline=(A.base)]\node[fill=ProofSubBg, text=ProofSubText, rounded corners=2pt, inner xsep=4pt, inner ysep=2pt, font=\sffamily\bfseries, drop shadow] (A) {#1};
\vspace{0.5ex}\par
}{
\par\vspace{1ex}
}

% 11. Dynamic Proof Wrappers
\makeatletter

\newcommand{\qedbadge}{%
  \tikz[baseline=(Q.base)]\node[draw=ProofTint, fill=ProofQEDBg, line width=1pt, rounded corners=2pt, inner xsep=4pt, inner ysep=2pt, font=\sffamily\bfseries\small, text=ProofQEDText, drop shadow={shadow xshift=0.5ex, shadow yshift=-0.5ex, fill=ProofTint, opacity=0.4}] (Q) {Q.E.D.};%
}

% SHORT PROOF 
\renewenvironment{proof}[1][Proof]{%
    \par\addvspace{2ex plus 0.5ex minus 0.2ex}%
    \insideprooftrue 
    \noindent\textcolor{ProofTint}{\faPenNib}\hspace{0.5em}%
    {\normalfont\sffamily\bfseries\color{ProofTint}#1}%
    \par\nobreak\vspace{1ex plus .2ex}\vspace{-\parskip}%
    \@afterindentfalse\@afterheading
}{%
    \ifmmode
      \mathqed
    \else
      \ifhmode
        \unskip\nobreak\hfill\penalty50\hskip2em\hbox{}\nobreak\hfill\qedbadge\par\vspace{2ex}%
      \else
        \par\vspace{-\parskip}\noindent\hfill\qedbadge\par\vspace{2ex}%
      \fi
    \fi
}

% LONG PROOF 
\newenvironment{proofWrapper}[1][Proof]{%
    \par\addvspace{3.5ex plus 1ex minus .2ex}%
    \insideprooftrue 
    \noindent\makebox[0pt][r]{\textcolor{ProofTint}{\faPenNib}\hspace{10pt}}%
    {\normalfont\sffamily\Large\bfseries\color{ProofTint}#1}%
    \par\nobreak\vspace{1.5ex plus .2ex}\vspace{-\parskip}%
    \@afterindentfalse\@afterheading
}{%
    \ifmmode
      \mathqed
    \else
      \ifhmode
        \unskip\nobreak\hfill\penalty50\hskip2em\hbox{}\nobreak\hfill\qedbadge\par\vspace{2ex}%
      \else
        \par\vspace{-\parskip}\noindent\hfill\qedbadge\par\vspace{2ex}%
      \fi
    \fi
}
\makeatother

% 13. Lecture Break
\newtcolorbox{pause}[1][Lecture Break]{
enhanced,
frame hidden,
colback=white,
opacityback=0,
code={
  \ifinsideproof
    \tcbset{borderline north={0.5pt}{0pt}{EnvBreakProof, dashed}, borderline south={0.5pt}{0pt}{EnvBreakProof, dashed}}
  \else
    \tcbset{borderline north={0.5pt}{0pt}{EnvBreak, dashed}, borderline south={0.5pt}{0pt}{EnvBreak, dashed}}
  \fi
},
title={\sffamily\bfseries\centering\textcolor{\ifinsideproof EnvBreakProof\else EnvBreak\fi}{\faCoffee\ \ #1}},
attach title to upper=\par\vspace{1ex},
fontupper=\centering\itshape\color{\ifinsideproof EnvBreakProof\else EnvBreak\fi},
breakable
}

% ==========================================
% INVISIBLE LAYER & FALLBACK ENVIRONMENTS
% ==========================================
\NewEnviron{hidden}{}  % all AI scratchpads: plans, checkpoints, sketches

% Data Integrity Fallbacks
\newcommand{\NewFallbackBox}[1]{
\newtcolorbox{#1}[1][Fallback]{
enhanced,
colback=EnvAIBg,
dynamic shadow,
code={
  \ifinsideproof
    \tcbset{colframe=EnvAIFallbackProof}
  \else
    \tcbset{colframe=EnvAI}
  \fi
},
fonttitle=\sffamily\bfseries,
coltitle=white,
title=\faExclamationTriangle\ \ ##1,
breakable
}
}
\NewFallbackBox{ai-logical-gap}
\NewFallbackBox{ai-generation-loop-fallback}
\NewFallbackBox{ai-raw-ocr-fallback}
\NewFallbackBox{ai-phonetic-dump}
\NewFallbackBox{ai-modality-conflict}
\NewFallbackBox{ai-off-camera-state}
\NewFallbackBox{ai-async-board-update}
\NewFallbackBox{ai-kinetic-emphasis}

% Custom Lecture Chapter Command with Optional Semester/Week Number
\newif\ifrestorelecturesubcounters
\newcounter{currentweek}
\setcounter{currentweek}{0}
\newcounter{lecturecount}
\setcounter{lecturecount}{0}
\renewcommand{\theHchapter}{\thechapter.\arabic{lecturecount}}

\NewDocumentCommand{\lecturechapter}{ o m m m m }{%
  \restorelecturesubcountersfalse
  \stepcounter{lecturecount}%
  \IfValueT{#1}{%
    \ifblank{#1}{}{%
      \ifnum\numexpr#1\relax>\value{currentweek}%
        \setcounter{currentweek}{#1}%
        \phantomsection
        \addcontentsline{toc}{part}{Week #1}%
      \fi
      \ifnum\numexpr#1\relax=\value{chapter}%
        \restorelecturesubcounterstrue
        \edef\TempSec{\the\value{section}}%
        \edef\TempThm{\the\value{theorem}}%
        \edef\TempEq{\the\value{equation}}%
        \edef\TempFig{\the\value{figure}}%
        \edef\TempTab{\the\value{table}}%
      \fi
      \setcounter{chapter}{\numexpr#1-1\relax}%
    }%
  }%
  \chapter[\texorpdfstring{#2, #3\quad---\quad #5}{#2, #3 --- #5}]{{\Large\color{ChapSubtitleText}#2 #4} \texorpdfstring{\\[1ex]}{--- } {\LARGE #5}}%
  \ifrestorelecturesubcounters
    \setcounter{section}{\TempSec}%
    \setcounter{theorem}{\TempThm}%
    \setcounter{equation}{\TempEq}%
    \setcounter{figure}{\TempFig}%
    \setcounter{table}{\TempTab}%
  \fi
  \markboth{\chaptername\ \thechapter\ \ ---\ \ #2 #4\ \ ---\ \ #5}{}%
}

\titlespacing*{\chapter}{0pt}{-20pt}{18pt plus 4pt minus 2pt}

\titleformat{\chapter}[display]
{\normalfont\sffamily\Large\bfseries\color{ChapTitleText}}
{\raggedleft\large\color{ChapNumberText}\MakeUppercase{\chaptertitlename}\ \thechapter}
{1ex}
{\color{ChapLine}\rule{\textwidth}{1.5pt}\endgraf\nopagebreak\vspace{2ex}\raggedleft} 
[\endgraf\nopagebreak\vspace{1ex}\color{ChapLine}\rule{\textwidth}{1.5pt}]

% Extended Section Bars 
\newcommand{\sectionbar}{%
\tikz[baseline=0pt]{
\fill[\ifinsideproof SecBar1Proof\else SecBar1\fi] (-2.5pt, 0) rectangle (1.0pt, 0.72em);
\fill[\ifinsideproof SecBar2Proof\else SecBar2\fi] (2.0pt, 0) rectangle (3.5pt, 0.72em);
\fill[\ifinsideproof SecBar3Proof\else SecBar3\fi] (4.8pt, 0) rectangle (5.3pt, 0.72em);
\fill[\ifinsideproof SecBar4Proof\else SecBar4\fi] (6.6pt, 0) rectangle (8.1pt, 0.72em);
}%
\hspace{10pt}%
}

\newcommand{\subsectionbar}{%
\tikz[baseline=0pt]{
\fill[\ifinsideproof SubSecBar1Proof\else SubSecBar1\fi] (-2.5pt, 0) rectangle (0.3pt, 0.72em);
\shade[left color=\ifinsideproof SubSecBar1Proof\else SubSecBar1\fi, right color=\ifinsideproof SubSecBar2Proof\else SubSecBar2\fi] (1.3pt, 0) rectangle (2.0pt, 0.72em);
\fill[\ifinsideproof SubSecBar2Proof\else SubSecBar2\fi] (3.3pt, 0) rectangle (4.8pt, 0.72em);
}%
\hspace{10pt}%
}

\newcommand{\subsubsectionbar}{%
\tikz[baseline=0pt]{
\fill[\ifinsideproof SubSubSecBar1Proof\else SubSubSecBar1\fi] (-2.5pt, 0) rectangle (-1.0pt, 0.72em);
\fill[\ifinsideproof SubSubSecBar2Proof\else SubSubSecBar2\fi] (0,0) rectangle (2.2pt, 0.72em);
}%
\hspace{10pt}%
}

\titleformat{\section}{\normalfont\sffamily\Large\bfseries\color{\ifinsideproof SecTitleColorProof\else SecTitleColor\fi}}{}{0pt}{\makebox[0pt][r]{\sectionbar}}
\titleformat{\subsection}{\normalfont\sffamily\large\bfseries\color{\ifinsideproof SubSecTitleColorProof\else SubSecTitleColor\fi}}{}{0pt}{\makebox[0pt][r]{\subsectionbar}}
\titleformat{\subsubsection}{\normalfont\sffamily\normalsize\bfseries\color{\ifinsideproof SubSubSecTitleColorProof\else SubSubSecTitleColor\fi}}{}{0pt}{\makebox[0pt][r]{\subsubsectionbar}}

\titlespacing*{\section}{0pt}{3.5ex plus 1ex minus .2ex}{1.5ex plus .2ex}
\titlespacing*{\subsection}{0pt}{3ex plus 1ex minus .2ex}{1.2ex plus .2ex}
\titlespacing*{\subsubsection}{0pt}{2.5ex plus 1ex minus .2ex}{1ex plus .2ex}


% =====================================================================
%  BACKWARD COMPATIBILITY (ACTIVE)
%
%  Lets .tex files written against the OLD vocabulary compile against this
%  preamble, so existing lecture output does not have to be migrated.
%  Delete this block once every .tex file uses the new names.
%
%  These are \let aliases rather than re-declared environments on purpose:
%  the alias inherits the target's argument signature and default title
%  exactly, so nothing here can drift out of sync with the definitions
%  above. Re-stating "[1][Erklaerung]" by hand would silently change the
%  rendered default the moment the real definition changed.
% =====================================================================
\makeatletter
\newcommand{\AliasEnv}[2]{%
  \expandafter\let\csname #1\expandafter\endcsname\csname #2\endcsname
  \expandafter\let\csname end#1\expandafter\endcsname\csname end#2\endcsname
}
\makeatother

\AliasEnv{spoken-clean}{speech}
\AliasEnv{math-stroke}{content}
\AliasEnv{student-interaction}{student}
\AliasEnv{explanation-of-steps}{steps}
\AliasEnv{didactic-insight}{insight}
\AliasEnv{meta-note}{note}
\AliasEnv{nice-box}{overview}
\AliasEnv{color-box}{highlight}
\AliasEnv{lecture-break}{pause}
\AliasEnv{ai-note}{aiNote}

%  CAVEAT -- the one thing aliasing cannot paper over.
%  The swap gave `proof` to the SHORT form (it is the common case, 63 uses
%  against 19). Old files use `proof` for the LONG form and `short-proof`
%  for the short one, so `proof` means opposite things in the two
%  vocabularies and no single preamble can serve both.
%  `short-proof` is aliased to the new `proof`, which is correct. Old
%  `\begin{proof}` therefore renders with the short-proof heading --
%  slightly smaller, no hanging icon. It still compiles, still reads
%  "Beweis", still emits Q.E.D. Cosmetic only, and it disappears as soon
%  as the file is migrated.
\AliasEnv{short-proof}{proof}

%  Scratchpads: all of these swallow their body, so an empty definition is
%  exactly equivalent to aliasing them to `hidden`.
\NewEnviron{ai-tikz-planner-invisible-content}{}
\NewEnviron{ai-proof-skeleton-invisible-content}{}
\NewEnviron{ai-global-state-checkpoint-invisible-content}{}
\NewEnviron{ai-retroactive-patch}[1][]{}
\NewEnviron{ai-global-state-checkpoint}{}

%  NOT merged into `hidden`, deliberately: the fallback boxes below are
%  \newtcolorbox, not \NewEnviron -- they RENDER, as red warning panels. They
%  exist to surface OCR failures, off-camera state and modality conflicts for a
%  human to see. Making them invisible would silently swallow exactly the
%  content they were created to raise. They keep their own names and stay
%  visible: ai-logical-gap, ai-generation-loop-fallback, ai-raw-ocr-fallback,
%  ai-phonetic-dump, ai-modality-conflict, ai-off-camera-state,
%  ai-async-board-update, ai-kinetic-emphasis.


\begin{document}

\begin{titlepage}
\thispagestyle{empty}
\begin{tikzpicture}[remember picture, overlay]
  % Dezent geometrische Akzentformen oben rechts
  \fill[ThemeSecondary!12, rounded corners=18pt] ($(current page.north east) + (-6.5cm,-2.5cm)$) rectangle ($(current page.north east) + (-1.2cm,-7.8cm)$);
  \fill[ThemeWeekNumber!20, rounded corners=12pt] ($(current page.north east) + (-4.8cm,-4.2cm)$) rectangle ($(current page.north east) + (-0.6cm,-9.0cm)$);
\end{tikzpicture}

\vspace*{4.5cm}
\noindent\begin{minipage}{\textwidth}
    {\sffamily\fontsize{40pt}{48pt}\selectfont\bfseries\color{ThemePrimary!95!black} Algebra I \par}
    \vspace{1.2cm}
    {\color{ThemeWeekNumber}\rule{\linewidth}{2.5pt}\par}
    \vspace{0.2cm}
    {\color{ThemeAccent1}\rule{0.6\linewidth}{1.5pt}\par}
    \vspace{2.0cm}
    {\sffamily\fontsize{18pt}{24pt}\selectfont\bfseries\color{TextEmphColor} Lecture Notes \& Transcription \par}
    \vspace{0.6cm}
    {\sffamily\fontsize{15pt}{22pt}\selectfont\color{TextBoldColor} Course: \textbf{Algebra I} \par}
    \vspace{0.3cm}
    {\sffamily\fontsize{15pt}{22pt}\selectfont\color{TextBoldColor} Lecturer: \textbf{Prof. Rahul Pandharipande} \par}
    \vspace{0.3cm}
    {\sffamily\fontsize{14pt}{20pt}\selectfont\color{OliveGreen} Department of Mathematics -- ETH Z{\"u}rich \par}
    \vspace{2.0cm}
    {\sffamily\fontsize{14pt}{20pt}\selectfont\bfseries\color{ThemeWeekNumber}\faGraduationCap\quad Fall Semester 2026 \par}
\end{minipage}
\vfill
\noindent\begin{minipage}{\textwidth}
    \begin{tcolorbox}[
        arc=6pt,
        boxrule=0.7pt,
        colback=MidnightBlue!3!white,
        colframe=MidnightBlue!30!white,
        left=10pt, right=10pt, top=8pt, bottom=8pt
    ]
    {\sffamily\bfseries\small\color{MidnightBlue!90!black}\faInfoCircle\quad Notice \& Transcription Disclaimer \par}
    \vspace{0.25cm}
    \sffamily\small\color{black!75}
    \begin{itemize}[leftmargin=1.5em, itemsep=3pt, topsep=0pt, parsep=0pt]
        \item[\textcolor{MidnightBlue!80!black}{\small\faRobot}] \textbf{AI-Generated Notes:} This document was automatically generated using AI/LLM systems from the lecture recordings and board content.
        \item[\textcolor{MidnightBlue!80!black}{\small\faUserShield}] \textbf{Unofficial Document:} This is \emph{not} official course material of ETH Z{\"u}rich or the lecturers, and has \emph{not} been manually corrected or proofread by the staff (unreviewed draft).
        \item[\textcolor{MidnightBlue!80!black}{\small\faExclamationTriangle}] \textbf{Disclaimer:} Artificial intelligence can make errors (e.g., in notation, indices, or logical derivations). Use at your own risk --- without warranty of any kind regarding accuracy, completeness, or relevance to exams.
    \end{itemize}
    \end{tcolorbox}
\end{minipage}
\vspace*{0.5cm}
\end{titlepage}

\tableofcontents
\hypersetup{linkcolor=LinkColor}
\cleardoublepage

% actual content:

% week1:
% ==========================================
% LatexRefinement Step Output: step4-2026-09-16-wednesday-week1-algebra-1-offset-final.tex
% Model: gemini-3.8-flash (Refined & Integrated)
% Primary Language: English
% ==========================================

\lecturechapter[1]{Wednesday}{16 Sep}{16 September 2026}{Gaussian Integers, Ring Axioms, Ring Homomorphisms, and Ideals}

\begin{overview}[{Course Content Overview:\\[0.25ex] Gaussian Integers, Ring Axioms, Homomorphisms \& Ideals}]
In this lecture, we bridge arithmetic and abstract algebra by exploring the following core themes:
\begin{itemize}
    \item \textbf{Gaussian Integers and Primes:} Comparing factorization and units in $\mathbb{Z}$ and $\mathbb{Z}[i]$, and establishing how the primality of an ordinary integer prime in $\mathbb{Z}[i]$ connects directly to Fermat's two-square theorem (an odd prime $p$ is a sum of two squares $a^2 + b^2$ if and only if $p \equiv 1 \pmod 4$).
    \item \textbf{Axioms and Elementary Properties of Rings:} Formalizing the ring axioms $(R, +, \cdot, 0, 1)$, deriving fundamental arithmetic identities such as $0 \cdot a = 0$ and $(-1) \cdot a = -a$, analyzing the degenerate zero ring ($1 = 0$), and defining scalar multiples by integers.
    \item \textbf{Ring Homomorphisms:} Defining structure-preserving maps between rings, explaining why identity preservation ($f(1) = 1$) must be explicitly required while $f(0) = 0$ is automatic, and exploring canonical examples including standard subring inclusions, modular reduction $\mathbb{Z} \to \mathbb{Z}/n\mathbb{Z}$, and the initial evaluation map $\mathbb{Z} \to R$.
    \item \textbf{Ideals and Field Structure:} Introducing the kernel $\ker(f)$, defining ideals $I \subset R$ via additive closure and multiplicative absorption, proving that every kernel is an ideal, and classifying the ideals of a field (establishing that fields possess only $\{0\}$ and $F$).
    \item \textbf{Outlook:} Connecting the existence of bases in infinite-dimensional vector spaces to Zorn's lemma and the Axiom of Choice.
\end{itemize}
\end{overview}

\section{Gaussian Integers and Primes}
\subsection{Units and Primes in \texorpdfstring{$\mathbb{Z}$}{Z} and \texorpdfstring{$\mathbb{Z}[i]$}{Z[i]}}

\begin{speech}[00:00:00 - 00:00:53]
Okay, so we'll try to start. Maybe you could sit down.

On the practical aspects, there was some progress: the website is more or less correct now, I think, so you can look for the information on the website.

The second issue is that people asked, and now the lectures will be recorded. You can see yourself --- or maybe you can only see me, but anyway --- the lectures will be recorded, so you can find them later on some website somewhere.

At the beginning, I said that the website is more or less in reasonable form; in particular, the assignment for September 25th is there. Some of the technical parts we're still working on, so we'll get back to you on that, but...
\end{speech}

\begin{note}[Administrative Opening]
The lecturer makes brief administrative announcements regarding the course website, the availability of the first homework assignment due on September 25th, and lecture video recordings.
\end{note}

\begin{speech}[00:00:53 - 00:02:48]
If you spent last night thinking about the Gaussian integers \inlinemetanote{writes $\mathbb{Z}[i]$ on the blackboard}, then... well, maybe for the people who didn't spend the whole night thinking about the Gaussian integers, I'll say a few things here.

As I said in the previous lecture, we discussed $\mathbb{Z}$ and its factorization properties. And I said you can discuss the same thing for the Gaussian integers. It will be part of the theory of the course, but it's nice to think about the example.

The first thing is: what are... for $\mathbb{Z}$ \inlinemetanote{writes $\mathbb{Z}$}, there were the units, that's $\pm 1$ \inlinemetanote{writes units: $\pm 1$}. And those confused our multiplication always, because you could always put in as many $-1$'s and $+1$'s as you want.

So, what's the analogue here? What are the units, what are the invertible elements in here? Well, of course, there's still $\pm 1$ \inlinemetanote{writes $\pm 1$}, but there's also $\pm i$ \inlinemetanote{writes $\pm i$}. These are also invertible, because $i^2 = -1$, and then you can multiply it by $1$ (i.e., $(-i) \cdot i = 1$).

Is that clear that these four objects play in some sense the role of the two there? Is that clear?

So, these are the units. It means if you write down any multiplication, like any unique factorization into primes, before you had to worry about $\pm 1$; now you have to worry about more, okay?

In particular, the notion of prime here was that $p$ is prime \inlinemetanote{writes $p$ is prime} if the factors are $\pm 1$ and $\pm p$ \inlinemetanote{writes if the factors are $\pm 1$ and $\pm p$}. That's on the $\mathbb{Z}$ side. Here, if we have some element --- maybe I call it $q$ is prime \inlinemetanote{writes $q$ is prime} --- then we can consider factorizations here; but now there's a lot more: there's four choices, there could be any one of these. $\pm 1$ and $\pm i$ are factors, and also $\pm q$ and $\pm i q$ \inlinemetanote{writes $\pm 1, \pm i, \pm q, \pm i q$}. These are all factors. That's the analogous statement.

This is all some kind of nonsense. It's just saying that, you know, this is more confusion with these extra $i$'s.
\end{speech}

\begin{content}[Units and Primes in the Gaussian Integers]
\begin{definition}[Gaussian Integers]\label[definition]{def:gaussian-integers}
The ring of \newterm{Gaussian integers}, denoted by $\mathbb{Z}[i]$, is the subring of the complex numbers $\mathbb{C}$ given by:
\[
\mathbb{Z}[i] := \{ a + bi \in \mathbb{C} \mid a, b \in \mathbb{Z} \}, \quad \text{where } i^2 = -1.
\]
\end{definition}

\begin{proposition}[Units in $\mathbb{Z}$ and {$\mathbb{Z}[i]$}]\label[proposition]{prop:gaussian-units}
An element $u$ in a ring is called a \newterm{unit} if it possesses a multiplicative inverse in that ring.
\begin{itemize}
    \item \textbf{Units in $\mathbb{Z}$:} The group of units in $\mathbb{Z}$ consists of exactly two elements:
    \[
    \mathbb{Z}^\times = \{1, -1\} = \{\pm 1\}.
    \]
    \item \textbf{Units in {$\mathbb{Z}[i]$}:} The group of units in $\mathbb{Z}[i]$ consists of exactly four elements:
    \[
    \mathbb{Z}[i]^\times = \{1, -1, i, -i\} = \{\pm 1, \pm i\}.
    \]
\end{itemize}
\end{proposition}

\begin{definition}[Prime Elements and Divisors]\label[definition]{def:primes-associates}
In an integral domain (\cref{def:integral-domain}, Week~2), a non-zero element that is not a unit is called prime (irreducible) if its only divisors are units or units multiplied by the element itself:
\begin{itemize}
    \item \textbf{Primes in $\mathbb{Z}$:} An integer $p \in \mathbb{Z} \setminus \{0, \pm 1\}$ is prime if its only factors are:
    \[
    \pm 1 \quad \text{and} \quad \pm p.
    \]
    \item \textbf{Primes in {$\mathbb{Z}[i]$}:} An element $q \in \mathbb{Z}[i] \setminus \{0, \pm 1, \pm i\}$ is prime if its only factors are the four units and their associated multiples:
    \[
    \pm 1, \quad \pm i, \quad \pm q, \quad \pm i q.
    \]
\end{itemize}
\end{definition}

\begin{steps}[Algebraic Intuition: The Concept of Associates]
In a unique factorization domain (such as $\mathbb{Z}$ or $\mathbb{Z}[i]$), factorizations into irreducible elements are unique up to multiplication by invertible elements (\newterm{units}) and order. Two elements $x, y \in R$ are called \newterm{associates} if $x = u \cdot y$ for some unit $u \in R^\times$.
\begin{itemize}
    \item \textbf{Associates in $\mathbb{Z}$:} The only units are $\pm 1$, meaning every prime $p$ has two associates: $p$ and $-p$.
    \item \textbf{Associates in {$\mathbb{Z}[i]$}:} Since $i \cdot (-i) = 1$ (i.e., $(-i) \cdot i = 1$), the imaginary units $i$ and $-i$ are also units. Consequently, every Gaussian integer $q$ has four associates:
    \[
    q, \quad -q, \quad iq, \quad -iq.
    \]
\end{itemize}
\end{steps}

\begin{steps}[Verification of Invertibility via the Algebraic Norm]
The elements $\pm 1, \pm i$ are clearly invertible in $\mathbb{Z}[i]$ since $(\pm 1)^2 = 1$ and $i \cdot (-i) = -i^2 = 1$.

To confirm that these are the \emph{only} units in $\mathbb{Z}[i]$, we use the multiplicative field norm $N \colon \mathbb{Z}[i] \to \mathbb{Z}_{\ge 0}$ defined by:
\[
N(a + bi) := (a + bi)\overline{(a + bi)} = a^2 + b^2.
\]
If $\alpha \in \mathbb{Z}[i]$ is a unit, there exists $\beta \in \mathbb{Z}[i]$ such that $\alpha \beta = 1$. Taking the norm on both sides gives:
\[
N(\alpha) N(\beta) = N(\alpha \beta) = N(1) = 1.
\]
Since $N(\alpha), N(\beta) \in \mathbb{Z}_{\ge 0}$, we must have $N(\alpha) = a^2 + b^2 = 1$. The only integer solutions to $a^2 + b^2 = 1$ are $(a, b) \in \{(\pm 1, 0), (0, \pm 1)\}$, corresponding precisely to $\pm 1$ and $\pm i$. That $\mathbb{Z}[i]$ has no further units is formalized within the general theory of Euclidean domains in Week~3 (\cref{prop:units-gaussian-integers}).
\end{steps}
\end{content}

\subsection{When Does a Rational Prime Remain Prime in \texorpdfstring{$\mathbb{Z}[i]$}{Z[i]}?}

\begin{speech}[00:02:48 - 00:04:08]
So, let's do something that's not nonsense. A math question you can ask, which is a very nice question, is: if I take a prime, an actual prime in the integers \inlinemetanote{writes $p \in \mathbb{Z}$}, and I view that prime, that $p$, as an element of the Gaussian integers \inlinemetanote{writes and view that $p$ as an element of $\mathbb{Z}[i]$}, which it is --- because remember, the Gaussian integers are the set of integers plus $i$ times an integer \inlinemetanote{writes $\mathbb{Z}[i] = \{n + im \mid n, m \in \mathbb{Z}\}$}... So, the integers naturally sit inside the Gaussian integers \inlinemetanote{writes $\mathbb{Z} \subset \mathbb{Z}[i]$} by just an integer going here. Is that clear?

If I take a prime, an honest prime in the integers --- so this would be, you know, like $2, 3, 5, 7, 11, \dots$ \inlinemetanote{writes $p = 2, 3, 5, 7, 11, \dots$} and so on --- take one of them, and I consider the Gaussian integers: does it stay prime? Or maybe say: When does $p$ stay prime? \inlinemetanote{writes When does $p$ stay prime?}

Meaning: it's born prime in the integers; when does it stay prime when I consider it as a Gaussian integer? That's a real question. And we had an example in the class last time.
\end{speech}

\begin{content}[Embedding of $\mathbb{Z}$ and the Prime Divisibility Question]
The ring of integers embeds naturally into the Gaussian integers via the inclusion map:
\[
\mathbb{Z} \hookrightarrow \mathbb{Z}[i], \quad n \mapsto n + 0i.
\]

\begin{question}[When Does an Integer Prime Stay Prime in {$\mathbb{Z}[i]$}?]\label[question]{q:prime-permanence}
Let $p \in \mathbb{Z}$ be an integer prime ($p \in \{2, 3, 5, 7, 11, 13, \dots\}$). When viewed as an element of $\mathbb{Z}[i]$:
\[
\text{When does } p \text{ remain prime (irreducible) in } \mathbb{Z}[i]\text{?}
\]
\end{question}

\begin{steps}[Why Primes Can Factor in Larger Rings]
An element that has no non-trivial divisors in $\mathbb{Z}$ might acquire non-trivial divisors in $\mathbb{Z}[i]$ because $\mathbb{Z}[i]$ contains new elements with non-zero imaginary parts. For example, $5$ cannot be factored in $\mathbb{Z}$, but in $\mathbb{Z}[i]$ it factors as $5 = (2 + i)(2 - i)$, where neither factor is a unit.
\end{steps}

\par\vspace{1ex}\noindent
To investigate whether $p$ factors non-trivially in $\mathbb{Z}[i]$, suppose there exists a non-unit divisor:
\begin{equation}\label{eq:gaussian-divisor}
(a + bi) \mid p \quad \text{in } \mathbb{Z}[i] \quad (a, b \in \mathbb{Z}).
\end{equation}
\end{content}

\subsection{Complex Conjugation and the Sum of Two Squares}

\begin{speech}[00:04:08 - 00:05:18]
But now we try to think about it mathematically. We have $p$, and we have to consider whether it's divisible in the Gaussian integers. So, we think, okay, maybe it's divisible. How do I do that? Well, maybe that means $a + bi$ divides $p$ \inlinemetanote{writes $(a+bi) \mid p$} in the Gaussian integers. Let's try to think what that means. Now, if $a + bi$ divides $p$, that means $a + bi$ times something equals $p$ --- that's the definition of divisibility. If that happened, I could just complex conjugate the whole thing. So, if this happens --- you have to think about this --- then $a - bi$ also divides $p$ \inlinemetanote{writes $(a-bi) \mid p$}.

Do you believe that statement? This has to do with complex conjugation, and it depends on what your feeling is for the complex numbers. If you don't understand the statement, you should think about it, because what does this mean? This means $p$ is equal to something something something times $a + bi$ \inlinemetanote{writes $p = \dots \cdot (a+bi)$}. That's what divisibility means. And I just take the complex conjugate of everything. But the complex conjugate of $p$ is $p$, and all the something something something gets complex con--- I don't care, but at the end, this one becomes that \inlinemetanote{points to $a-bi$ and writes $p = \overline{\dots} \cdot (a-bi)$}. So, that means that that also happens, right?
\end{speech}

\begin{content}[Complex Conjugation of Divisors]
\begin{proposition}[Conjugate Divisor Property]\label[proposition]{prop:conjugate-divisor}
Let $p \in \mathbb{Z}$. If $(a + bi) \mid p$ in $\mathbb{Z}[i]$, then:
\begin{equation}\label{eq:div-conj}
(a - bi) \mid p \quad \text{in } \mathbb{Z}[i].
\end{equation}
\end{proposition}

\begin{proof}
By definition of divisibility in $\mathbb{Z}[i]$, $(a + bi) \mid p$ means there exists an element $z \in \mathbb{Z}[i]$ such that:
\[
p = z \cdot (a + bi).
\]
Applying the complex conjugation automorphism to both sides yields:
\[
\overline{p} = \overline{z \cdot (a + bi)} = \overline{z} \cdot \overline{(a + bi)}.
\]
Since $p \in \mathbb{Z}$ is real, $\overline{p} = p$. Furthermore, $\overline{a + bi} = a - bi$, and $\overline{z} \in \mathbb{Z}[i]$. Thus,
\[
p = \overline{z} \cdot (a - bi),
\]
which proves that $(a - bi) \mid p$ in $\mathbb{Z}[i]$.
\end{proof}
\end{content}

\begin{speech}[00:05:18 - 00:06:08]
Then you might think --- and this is the kind of thing we will make precise in this class --- that actually then both of them divide \inlinemetanote{writes $(a+bi)(a-bi) \mid p$}. That already requires some statement about unique factorization. And this is now a nice integer: $a^2 + b^2$ divides $p$ \inlinemetanote{writes $a^2 + b^2 \mid p$}. If it's an integer, then $p$ has to be of the form $a^2 + b^2$ \inlinemetanote{writes $p = a^2 + b^2$}. So, it has to be a prime that's a sum of two squares.

This argument will show that if you want $p$ to stay prime... so, if $p$ factors, then it's a sum of two squares. And this now connects to this world of trying to express primes as a sum of two squares, and there's facts about this.
\end{speech}

\begin{content}[Reduction to the Sum of Two Squares]
Multiplying the two conjugate factors together yields an ordinary integer:
\[
(a + bi)(a - bi) = a^2 + b^2 \in \mathbb{Z}.
\]
Using the field norm $N(x + yi) = x^2 + y^2 = (x + yi)(x - yi)$, which is multiplicative, the equation $p = (a + bi)z$ implies:
\[
N(p) = p^2 = N(a + bi) \cdot N(z) = (a^2 + b^2) \cdot N(z).
\]
Since $p$ is a prime in $\mathbb{Z}$, the integer $a^2 + b^2$ dividing $p^2$ is $1$, $p$ or $p^2$. But $a^2 + b^2 = 1$ would make $a + bi$ a unit (as $(a + bi)(a - bi) = 1$), and $a^2 + b^2 = p^2$ would force $N(z) = 1$, making $z$ a unit and $a + bi = p z^{-1}$ an associate of $p$. So if $a + bi$ is neither a unit nor an associate of $p$, then:
\begin{equation}\label{eq:sum-of-squares-condition}
a^2 + b^2 = p.
\end{equation}

\begin{steps}[Rigorous Derivation via the Multiplicative Norm]
The blackboard derivation outlines the necessity of $p$ being a sum of two squares. Rigorously, this follows from the multiplicative norm:
\begin{enumerate}[label=\textbf{(\arabic*)}]
    \item \textbf{Norm Equation:} If $p = z \cdot (a + bi)$ with $z = c + di$, taking norms yields:
    \[
    p^2 = N(p) = N(z) N(a + bi) = (c^2 + d^2)(a^2 + b^2).
    \]
    \item \textbf{Non-Trivial Factors:} Because $a + bi$ is not a unit, $N(a + bi) = a^2 + b^2 > 1$. Because $z$ is not a unit, $N(z) = c^2 + d^2 > 1$.
    \item \textbf{Prime Factorization in $\mathbb{Z}$:} Since $p$ is prime in $\mathbb{Z}$, the only divisors of $p^2$ in $\mathbb{Z}_{>0}$ are $1, p, p^2$. Thus, we must have:
    \[
    a^2 + b^2 = p \quad \text{and} \quad c^2 + d^2 = p.
    \]
\end{enumerate}
\end{steps}

\begin{steps}[Connection Between Gaussian Primes and Sums of Squares]
This algebraic deduction establishes a profound bridge:
\begin{itemize}
    \item If an integer prime $p$ factors non-trivially in $\mathbb{Z}[i]$ as $p = (a + bi)(a - bi)$, then $p$ must be expressible as the sum of two integer squares: $p = a^2 + b^2$.
    \item By contrapositive, if $p$ cannot be expressed as the sum of two squares, it cannot factor non-trivially into non-unit Gaussian integers, and hence $p$ must remain prime in $\mathbb{Z}[i]$.
    \item Conversely, $p = a^2 + b^2$ gives the factorization $p = (a + bi)(a - bi)$ with $N(a \pm bi) = p > 1$, so neither factor is a unit. Hence, $p$ stays prime in $\mathbb{Z}[i]$ \emph{if and only if} $p$ is not a sum of two squares.
\end{itemize}
\end{steps}
\end{content}

\begin{aiNote}[Intermediate Steps Added by Transcriber]
The lecturer argues that both $a + bi$ and $a - bi$ divide $p$, so their product $a^2 + b^2$ divides $p$, and remarks that this step \qt{already requires some statement about unique factorization}. The norm computation above reaches \eqref{eq:sum-of-squares-condition} without that statement: it only uses $N(zw) = N(z) \, N(w)$ and the primality of $p$ in $\mathbb{Z}$.
\end{aiNote}

\subsection{Fermat's Theorem on Sums of Two Squares}

\begin{speech}[00:06:08 - 00:08:12]
And now, this connects to this world of trying to express primes as a sum of two squares, and there are facts about this. The simple fact is that if $p$ is congruent to $3 \pmod 4$ \inlinemetanote{writes If $p \equiv 3 \pmod 4$}, this is impossible \inlinemetanote{writes this is impossible}. You can never write a prime that is $3 \pmod 4$ as a sum of two squares. That's easy to prove; I think it's part of the exercise, I'm not sure.

But the one that's tricky is: if $p$ is congruent to $1 \pmod 4$ \inlinemetanote{writes If $p \equiv 1 \pmod 4$}, then you can \inlinemetanote{writes then you can}. You can solve \inlinemetanote{writes you can solve} $p = a^2 + b^2$ \inlinemetanote{writes $p = a^2 + b^2$} in the integers \inlinemetanote{writes in the integers}. This is a theorem that's due to Fermat \inlinemetanote{writes Fermat}. And this is one of the problems on the problem set.

So, the answer to this question finally is that if you have a prime, and you want to know whether it stays prime in the Gaussian integers, the whole thing that matters is whether $p$ is equal to $1$ or $3 \pmod 4$. It's kind of an interesting question.

And as I said, we will return to these matters in the class, but it already shows us that the very naive idea of thinking about how things factor leads to interesting mathematics, and in this case it goes in a number-theoretic direction. That's this question about whether a prime can be expressed as a sum of two squares --- a really basic question.

We had an example yesterday, I remind you; the answer yesterday had to do with $5 = 2^2 + 1^2$ \inlinemetanote{writes $5 = 2^2 + 1^2$}. So, $5$ is a prime in the integers, $5$ is $1 \pmod 4$, and you can express it as a sum of two squares. The thing that's tricky on the board here is this statement: that if $p$ is equal to $1 \pmod 4$, then you can solve this equation, and that is the exercise on the homework, on the first homework set. There's a way to go through that --- a very clever way to go through that, not the way that Fermat used.
\end{speech}

\begin{content}[Fermat's Two-Square Theorem and Gaussian Primes]
\begin{theorem}[Fermat's Two-Square Theorem]\label[theorem]{thm:fermat-two-squares}
Let $p$ be an odd integer prime ($p > 2$).
\begin{enumerate}[label=\textbf{(\roman*)}]
    \item If $p \equiv 3 \pmod 4$, then $p$ cannot be written as the sum of two squares ($p \neq a^2 + b^2$). Consequently, $p$ remains prime in $\mathbb{Z}[i]$.
    \item If $p \equiv 1 \pmod 4$, then $p$ can be written uniquely (up to sign and order) as the sum of two squares:
    \[
    p = a^2 + b^2 = (a + bi)(a - bi) \quad (a, b \in \mathbb{Z}).
    \]
    Consequently, $p$ factors in $\mathbb{Z}[i]$ and does not remain prime.
\end{enumerate}
\end{theorem}

\begin{example}[{Decomposition of $5$ vs. Primality of $3$ in $\mathbb{Z}[i]$}]\label[example]{ex:primes-5-and-3}
\begin{itemize}
    \item \textbf{The prime $5$:} Since $5 \equiv 1 \pmod 4$, we have:
    \[
    5 = 2^2 + 1^2 = (2 + i)(2 - i).
    \]
    Thus, $5$ is composite in $\mathbb{Z}[i]$.
    \item \textbf{The prime $3$:} Since $3 \equiv 3 \pmod 4$, and squares modulo $4$ can only be $0$ or $1$ ($0^2 \equiv 0$, $1^2 \equiv 1$, $2^2 \equiv 0$, $3^2 \equiv 1$), the sum of two squares can only be $0, 1$, or $2 \pmod 4$. Hence, $a^2 + b^2 \equiv 3 \pmod 4$ has no integer solutions, and $3$ remains prime in $\mathbb{Z}[i]$.
\end{itemize}
\end{example}

\begin{steps}[Classification of Integer Primes in {$\mathbb{Z}[i]$}]
Every integer prime $p \in \mathbb{Z}$ falls into one of three categories in $\mathbb{Z}[i]$:
\begin{itemize}
    \item \textbf{Ramified Prime ($p = 2$):} $2 = 1^2 + 1^2 = (1 + i)(1 - i) = -i(1 + i)^2$.
    \item \textbf{Split Primes ($p \equiv 1 \pmod 4$):} $p = a^2 + b^2 = (a + bi)(a - bi)$ splits into two distinct non-associate Gaussian primes (e.g., $5, 13, 17, 29$).
    \item \textbf{Inert Primes ($p \equiv 3 \pmod 4$):} $p$ cannot be written as $a^2 + b^2$, so $p$ remains prime in $\mathbb{Z}[i]$ (e.g., $3, 7, 11, 19, 23$).
\end{itemize}
\end{steps}
\end{content}

\begin{aiNote}[Beyond the Lecture: The Prime $2$]
% [PERSISTENT, ADDITION]
The even prime is not covered by \cref{thm:fermat-two-squares}: $2 = 1^2 + 1^2 = (1 + i)(1 - i)$, and since $1 - i = -i(1 + i)$, in fact $2 = -i \, (1 + i)^2$, so $2$ does not stay prime in $\mathbb{Z}[i]$.
\end{aiNote}

\begin{note}[Board Transition]
The lecturer erases the center and right chalkboards with a sponge to prepare for reviewing the abstract definition and axioms of a ring.
\end{note}

\begin{hidden}
% timestamp: 00:08:15
% topic: Gaussian integers, units, primes, and Fermat's two-square theorem
% board_state: def:gaussian-integers, prop:conjugate-divisor, thm:fermat-two-squares
% next_goal: Formal definition of rings, axioms, and elementary arithmetic properties
% open_loops: none
\end{hidden}

\section{The Axioms and Basic Properties of Rings}
\subsection{Definition of a Ring}

\begin{speech}[00:08:15 - 00:09:45]
All right. So, then the main direction of the lecture last time was about rings, which we revisit now.

So, we go back to rings and the foundational material. I had defined a ring: $R$ with multiplication, $0$, and $1$ \inlinemetanote{writes $(R, +, \cdot, 0, 1)$}. This is a ring \inlinemetanote{writes is a ring}. And the axioms, if we rewrite them now differently: there's the axioms for $+$, and this axiom for $+$, we can say just that $(R, +, 0)$ is an abelian group \inlinemetanote{writes $+ \implies (R, +, 0) \text{ abelian group}$}. This is just the same axioms written differently; those are just the group axioms. Is that clear?

Then there's the axiom for multiplication \inlinemetanote{writes $\cdot \implies$}, and these were --- now I don't have a more efficient way to say this --- it's associative \inlinemetanote{writes associative}, and then there's the axiom of the unit \inlinemetanote{writes $1 \cdot a = a \cdot 1 = a$}. And then there's the axioms together \inlinemetanote{writes $+,\cdot \implies$}, and these are two times the distributive law \inlinemetanote{writes $\times 2 \text{ Dist}$}, which I don't write, okay?

This is just the definition which we've seen and you've seen before. And the only thing I've done now is abbreviated the $+$ as saying that if I look at the additive structure, it's an abelian group, because that's what the axioms were.
\end{speech}

\begin{content}[Axiomatic Definition of a Ring]
\begin{definition}[Ring]\label[definition]{def:ring}
A \newterm{ring} is a tuple $(R, +, \cdot, 0, 1)$ consisting of a set $R$, two binary operations $+$ (addition) and $\cdot$ (multiplication), and two distinguished elements $0, 1 \in R$ satisfying the following axioms:
\begin{enumerate}[label=\textbf{(\arabic*)}]
    \item \textbf{Additive Abelian Group Axioms:}
    $(R, +, 0)$ is an abelian group:
    \begin{itemize}
        \item \textbf{Associativity:} $\forall a, b, c \in R \colon (a + b) + c = a + (b + c)$.
        \item \textbf{Identity:} $\forall a \in R \colon a + 0 = 0 + a = a$.
        \item \textbf{Inverse:} $\forall a \in R \; \exists (-a) \in R \colon a + (-a) = (-a) + a = 0$.
        \item \textbf{Commutativity:} $\forall a, b \in R \colon a + b = b + a$.
    \end{itemize}

    \item \textbf{Multiplicative Monoid Axioms:}
    Multiplication is associative and possesses a unit element $1$:
    \begin{itemize}
        \item \textbf{Associativity:} $\forall a, b, c \in R \colon (a \cdot b) \cdot c = a \cdot (b \cdot c)$.
        \item \textbf{Identity (Unit):} $\forall a \in R \colon 1 \cdot a = a \cdot 1 = a$.
    \end{itemize}

    \item \textbf{Distributivity Axioms:}
    Multiplication distributes over addition from both the left and the right:
    \begin{itemize}
        \item \textbf{Left Distributivity:} $\forall a, b, c \in R \colon a \cdot (b + c) = (a \cdot b) + (a \cdot c)$.
        \item \textbf{Right Distributivity:} $\forall a, b, c \in R \colon (a + b) \cdot c = (a \cdot c) + (b \cdot c)$.
    \end{itemize}
\end{enumerate}
\end{definition}

\begin{steps}[Structural Synthesis of the Ring Axioms]
Summarizing the additive structure compactly as an abelian group $(R, +, 0)$ packages four standard axioms (associativity, commutativity, neutral element, and inverses) into a single standard algebraic concept. The distributive laws act as the essential bridge coupling the additive abelian group structure with the multiplicative monoid structure.

Two words that are easy to mix up: \emph{abelian} refers to the addition ($a + b = b + a$ is part of the axioms), whereas \emph{commutative ring} refers to the multiplication. The definition above does not require $a \cdot b = b \cdot a$; the standing convention of this course adds it (\cref{rem:ring-convention}). Nor does it require $1 \neq 0$: the case $1 = 0$ is allowed and treated in \cref{prop:zero-ring}.
\end{steps}
\end{content}

\subsection{Elementary Property: \texorpdfstring{$0 \cdot a = 0$}{0 * a = 0}}

\begin{speech}[00:09:45 - 00:11:32]
So, then there's a bunch of things to say at the beginning: some basic properties \inlinemetanote{writes Basic Properties}. And these are all in the books, but in these axioms, there are lots of things that you're familiar with from the integers and arithmetic that are not in the axioms. Like, one question is: what happens if I take $0 \times a$ \inlinemetanote{writes $0 \cdot a$}? Because $0$ is an object that lives on the additive side, and I'm multiplying it, so it's on the multiplicative side. You can wonder what this is, and you can imagine what the answer is going to be.

To prove it --- if we're going to formally prove it --- you, of course, have to use the distributive law. The distributive law is the only thing that relates the addition to the multiplication. So, we're going to try to prove this, and to start, we know we have to use the distributive law somewhere. One way to start is we can write this as $(0 + 0) \times a$ \inlinemetanote{writes $= (0+0) \cdot a$}. We have not used the distributive law yet; this uses the fact that $0 + 0 = 0$, which is the neutral part of the group law. And now we use the distributive law \inlinemetanote{writes $= (0 \cdot a) + (0 \cdot a)$}: $(0 \cdot a) + (0 \cdot a)$.

Okay. And now we're happy, because we have this equation, and then we go back to using the additive properties, which is the left cancellation --- the additive inverse, okay? So, now we can cancel \inlinemetanote{cancels $0 \cdot a$ on both sides with chalk}, and we happily see that $0 = 0 \times a$ \inlinemetanote{writes $0 = 0 \cdot a$ and draws a chalk box around it}, which is, of course, the basic law of arithmetic. It's not an axiom, but it follows immediately from the axioms, which is why it's not an axiom. You could actually just put it as an axiom, it wouldn't hurt anybody, but... Okay.
\end{speech}

\begin{content}[Derivation of Multiplication by Zero]
\begin{proposition}[Multiplication by Zero]\label[proposition]{prop:mult-by-zero}
In any ring $R$, for every element $a \in R$:
\begin{equation}\label{eq:mult-by-zero}
0 \cdot a = 0 \quad \text{and} \quad a \cdot 0 = 0.
\end{equation}
\end{proposition}

\begin{proof}
Because $0$ is the additive identity of $(R, +, 0)$, we have $0 + 0 = 0$. Substituting this into $0 \cdot a$:
\begin{align}
0 \cdot a &= (0 + 0) \cdot a \nonumber \\
&= (0 \cdot a) + (0 \cdot a) && \text{(by right distributivity)}. \label{eq:zero-expansion}
\end{align}
Because $(R, +, 0)$ is an abelian group, the element $0 \cdot a \in R$ possesses a unique additive inverse $-(0 \cdot a) \in R$. Adding $-(0 \cdot a)$ to both sides of \eqref{eq:zero-expansion}:
\begin{align*}
-(0 \cdot a) + (0 \cdot a) &= -(0 \cdot a) + \bigl((0 \cdot a) + (0 \cdot a)\bigr) \\
0 &= \bigl(-(0 \cdot a) + (0 \cdot a)\bigr) + (0 \cdot a) && \text{(by associativity)} \\
0 &= 0 + (0 \cdot a) \\
0 &= 0 \cdot a.
\end{align*}
By the exact same argument using left distributivity, $a \cdot 0 = a \cdot (0 + 0) = a \cdot 0 + a \cdot 0$, which upon adding $-(a \cdot 0)$ yields $a \cdot 0 = 0$.
\end{proof}

\begin{steps}[Why Distributivity is Indispensable]
The zero element $0$ is defined exclusively through addition ($a + 0 = a$). To deduce its behavior under multiplication ($0 \cdot a = 0$), one must invoke the distributive law, as it provides the only axiomatic connection linking addition to multiplication.
\end{steps}
\end{content}

\subsection{The Zero Ring (\texorpdfstring{$1 = 0$}{1 = 0})}

\begin{speech}[00:11:32 - 00:12:43]
In particular, if $1 = 0$ \inlinemetanote{writes If $1 = 0$, then} --- because I said for a ring we allow this extremely unpleasant boundary case where $1$ is equal to $0$, the degenerate case... So, if that happens --- if we have a ring where $1 = 0$ --- then we have $a$, for any $a$, is equal to $a \times 1$ \inlinemetanote{writes $a = a \cdot 1$}, which is the multiplication axiom. And then, if we use $1 = 0$, this is equal to $a \times 0$ \inlinemetanote{writes $= a \cdot 0$} using $1 = 0$, and then this is $0$ \inlinemetanote{writes $= 0$}.

So, it's a rather degenerate case: if you have a ring where $1 = 0$, then all elements are equal to $0$. It's not a super exciting ring; there's only one element. If I say this formally: if you have $0 = 1$, this implies your ring as a set consists of just one element \inlinemetanote{writes If $0 = 1 \implies R = \{0\}$}. Anyway, it's a degenerate case, it's unpleasant to think about, but we carry it around because it fits well in certain statements of theorems.
\end{speech}

\begin{content}[The Degenerate Zero Ring]
\begin{proposition}[Characterization of the Zero Ring]\label[proposition]{prop:zero-ring}
Let $R$ be a ring. If $1 = 0$ in $R$, then $R$ is the trivial ring containing only a single element:
\[
R = \{0\}.
\]
\end{proposition}

\begin{proof}
Assume $1 = 0$. Let $a \in R$ be an arbitrary element. Using the unit axiom ($a = a \cdot 1$), the hypothesis $1 = 0$, and \cref{prop:mult-by-zero}:
\[
a = a \cdot 1 = a \cdot 0 = 0.
\]
Because every element $a \in R$ must equal $0$, the set $R$ contains only the zero element:
\[
R = \{0\}.
\]
\end{proof}

\begin{remark}[The Zero Ring in Abstract Algebra]
The trivial ring $R = \{0\}$ is often called the \newterm{zero ring}. While fields strictly require $1 \neq 0$, the general definition of a ring intentionally admits $1 = 0$ so that the category of rings has a terminal object (the zero ring) and quotient rings $R/I$ remain rings even when $I = R$. Definitions that must exclude it say so explicitly, as for integral domains (\cref{def:integral-domain}, Week~2).
\end{remark}
\end{content}

\begin{hidden}
% timestamp: 00:12:50
% topic: Basic properties: 0*a = 0 and the zero ring 1 = 0
% board_state: def:ring, prop:mult-by-zero, prop:zero-ring
% next_goal: Additive inverses -a, multiples by integers, and the canonical map Z -> R
% open_loops: none
\end{hidden}

\subsection{Additive Inverses and the Element \texorpdfstring{$-a$}{-a}}

\begin{speech}[00:12:43 - 00:15:05]
Other things that are kind of convention: if you take $-1$... $1$ is in the ring because it's part of the data \inlinemetanote{writes $1 \in R$}. Because of the additive structure, the group axioms say that there exists an additive inverse. So, by convention, $-1$ is by definition the additive inverse of $1$ \inlinemetanote{writes $-1 \text{ by def the add. inv. of } 1$}, and you can prove that it's unique. Maybe I leave that as an exercise. Given the ring, there's $1$, and there's the additive inverse. Additive inverses are unique, you can prove that. So, that's $-1$.

And then $-a$, for any element, is defined to be $(-1) \times a$ \inlinemetanote{writes $-a := (-1) \cdot a$}; this is not a surprising definition.

And then the little lemma --- maybe it doesn't even get to be called a lemma --- is that $-a$ is the additive inverse of $a$ \inlinemetanote{writes Lemma: $-a \text{ is the additive inverse of } a.$}. That's very convenient; it would be really unpleasant if that weren't the case.

The proof of this is: we should start with $0 = (1 + (-1)) \cdot a$ \inlinemetanote{writes Proof: $0 = (1 + (-1)) \cdot a$}... that's $0$. And we've already proven that that's that (i.e., $0 \cdot a = 0$, \cref{prop:mult-by-zero}), okay? I've used two steps here: here on this side, this thing is $0$, and we've already proved $0 \times a = 0$, so... okay? And then I use the distributive law and distribute this side as $a$, and on this side as plus $(-1) \cdot a$ \inlinemetanote{writes $= a + (-1) \cdot a$}, and that exhibits this $(-1) \times a$ as the additive inverse of $a$.

So, there's a number of things --- I think we spent enough time on this --- but the way the axioms work, they work so that it fits with your intuition about how multiplication works.

Are you happy with this? You can also look in the book; there'll be a couple of pages of playing around with these. It's not deep, but it's good to be familiar with it.
\end{speech}

\begin{content}[Additive Inverses and Multiplication by \texorpdfstring{$-1$}{-1}]
\begin{definition}[The Element $-1$ and the Notation $-a$]\label[definition]{def:negative-elements}
Let $R$ be a ring.
\begin{enumerate}[label=\textbf{(\roman*)}]
    \item We denote by $-1 \in R$ the unique additive inverse (\cref{prop:uniqueness-inverse}) of the multiplicative identity $1 \in R$, so that:
    \[
    1 + (-1) = 0.
    \]
    \item For any element $a \in R$, we define:
    \begin{equation}\label{eq:def-negative-a}
    -a := (-1) \cdot a.
    \end{equation}
\end{enumerate}
\end{definition}

\begin{lemma}[$-a$ is the Additive Inverse of $a$]\label[lemma]{lem:additive-inverse-product}
For any element $a \in R$, the element $-a = (-1) \cdot a$ is the additive inverse of $a$, that is:
\[
a + (-a) = 0.
\]
\end{lemma}

\begin{proof}
Using \cref{prop:mult-by-zero} ($0 \cdot a = 0$), \cref{def:negative-elements} ($1 + (-1) = 0$), and distributivity:
\begin{align*}
0 &= 0 \cdot a \\
&= \bigl(1 + (-1)\bigr) \cdot a \\
&= (1 \cdot a) + \bigl((-1) \cdot a\bigr) && \text{(by distributivity)} \\
&= a + (-a) && \text{(since $1 \cdot a = a$ and $(-1) \cdot a = -a$)}.
\end{align*}
Since $a + \bigl((-1) \cdot a\bigr) = 0$, the element $(-1) \cdot a$ satisfies the defining property of the additive inverse of $a$. By uniqueness of inverses in $(R, +, 0)$, this confirms that $-a = (-1) \cdot a$ is indeed the unique additive inverse of $a$.
\end{proof}
\end{content}

\subsection{Integer Multiples and the Canonical Map \texorpdfstring{$\mathbb{Z} \to R$}{Z -> R}}

\begin{speech}[00:15:05 - 00:17:10]
Another convention about a ring is: I said that if I have $a$ as an element of the ring \inlinemetanote{writes $a \in R$}, $-a$ \inlinemetanote{writes $-a$} we just write as $-a$, and by definition $-a$ is just $(-1) \times a$ \inlinemetanote{writes $-a = -1 \cdot a$}, and it's a canonical element of the ring.

There's something else: we can also look at $a + a$ \inlinemetanote{writes $a + a$}, and this is defined as $2a$ \inlinemetanote{writes $= 2a$}. You can take $a + a + a$ \inlinemetanote{writes $a + a + a$} $n$ times \inlinemetanote{writes underbrace with $n$ times}; this is what this term $na$ is defined to be \inlinemetanote{writes $= na$}.

So, in any ring, if I have an element $a$ in a ring \inlinemetanote{writes $a \in R$}, I can take $17$ times that element \inlinemetanote{writes $17a \in R$}. That is an element of the ring, and I just mean you add $a$ to itself $17$ times. It's well-defined, just using the ring structure. You have no objections?

And I can also do $-17$ times \inlinemetanote{writes $-17a \in R$}, because I just, you know, either add $-a$ to itself $17$ times, or I add $a$ to itself $17$ times and take the minus, and they're the same. All right?

In particular, there's a map from $\mathbb{Z}$ to the ring \inlinemetanote{writes $\mathbb{Z} \to R$}. And this map requires no choices. This map sends $0$ to $0$ \inlinemetanote{writes $0 \mapsto 0$}, sends $1$ to $1$ \inlinemetanote{writes $1 \mapsto 1$}, sends $2$ to $1 + 1$ \inlinemetanote{writes $2 \mapsto 1 + 1$}, sends $3$ to $1 + 1 + 1$ \dots whatever \inlinemetanote{writes $3 \mapsto 1 + 1 + 1$}, okay? And it sends $-1$ to the additive inverse \inlinemetanote{writes $-1 \mapsto -1$}.

So, if you give me a ring, I know how to map $\mathbb{Z}$ to that ring. There's nothing going on with this map: I just map $1$ to $1$, and then everything else I just add up.

This is not just any map. It turns out this is what's called a ring homomorphism. So, we have to define this.
\end{speech}

\begin{content}[Scalar Multiplication by Integers and the Canonical Map $\mathbb{Z} \to R$]
\begin{definition}[Multiples of Ring Elements by Integers]\label[definition]{def:integer-multiples}
Let $R$ be a ring and $a \in R$. For any integer $n \in \mathbb{Z}$, the element $n \cdot a \in R$ is defined inductively:
\[
n \cdot a := \begin{cases}
0_R & \text{if } n = 0, \\
\underbrace{a + a + \dots + a}_{n \text{ times}} & \text{if } n > 0, \\
\underbrace{(-a) + (-a) + \dots + (-a)}_{|n| \text{ times}} & \text{if } n < 0.
\end{cases}
\]
In particular, for any specific integer such as $n = 17$ or $n = -17$:
\[
17a \in R \quad \text{and} \quad -17a = -(17a) \in R.
\]
\end{definition}

\begin{remark}[Two Ways to Form $-17 \cdot a$]\label[remark]{rem:negative-multiples}
The equality $n \cdot (-a) = -(n \cdot a)$ holds for $n > 0$: adding $-a$ to itself $n$ times gives the same element as adding $a$ to itself $n$ times and taking the additive inverse, because $(n \cdot a) + (n \cdot (-a)) = n \cdot \bigl(a + (-a)\bigr) = 0$ by commutativity and associativity of $+$. This is what the lecturer means by \qt{they're the same}.
\end{remark}

\begin{steps}[A Harmless Clash of Notation]
In $n \cdot a$ the factor $n$ is an integer, not an element of $R$, so this dot is not the ring multiplication. The two readings never conflict: $0 \cdot a = 0_R$ holds for the integer $0$ by definition and for the ring element $0$ by \cref{prop:mult-by-zero}, and in general $n \cdot a = (n \cdot 1_R) \cdot a$ by the distributive law.
\end{steps}

\begin{definition}[Canonical Map from $\mathbb{Z}$ into a Ring]\label[definition]{def:canonical-map-z}
For every ring $R$, there exists a canonical map $\phi \colon \mathbb{Z} \to R$ defined by evaluating integer multiples on the multiplicative unit $1_R$:
\[
\phi(n) := n \cdot 1_R.
\]
Explicitly, the map sends:
\begin{align*}
0 &\mapsto 0_R, \\
1 &\mapsto 1_R, \\
2 &\mapsto 1_R + 1_R, \\
3 &\mapsto 1_R + 1_R + 1_R, \\
-1 &\mapsto -1_R, \quad \dots
\end{align*}
\end{definition}

\begin{steps}[Universal Initial Property of the Integers $\mathbb{Z}$]
This map $\phi \colon \mathbb{Z} \to R$ requires absolutely no choices: it is entirely determined by the ring structure of $R$. As we will see next (\cref{ex:canonical-hom-z-to-r}), this map is not merely a set-theoretic function, but a ring homomorphism. Going beyond the lecture: it is even the unique ring homomorphism from $\mathbb{Z}$ to $R$ (any such $g$ has $g(1) = 1_R$, and additivity forces $g(n) = n \cdot 1_R$), establishing that $\mathbb{Z}$ is the \newterm{initial object} in the category of rings with unit (categories: \cref{def:category-full}, Week~2).
\end{steps}
\end{content}

\begin{hidden}
% timestamp: 00:17:10
% topic: Integer multiples and the canonical map Z -> R
% board_state: def:integer-multiples, def:canonical-map-z
% next_goal: Definition of ring homomorphisms and standing convention on commutativity
% open_loops: none
\end{hidden}

\section{Ring Homomorphisms}
\subsection{Standing Convention: Commutativity and Units}

\begin{speech}[00:17:10 - 00:18:44]
Let $R_1$ and $R_2$ be two rings \inlinemetanote{writes let $R_1$ and $R_2$ be two rings}.

Now, maybe I'll say a little bit more about these rings. For this class, as I said, if I consider totally abstract rings --- like totally arbitrary rings here \inlinemetanote{points to ring axioms on the upper left board} --- I have not assumed the multiplication is commutative. For a lot of things we don't, and sometimes we do, but when I say \qt{rings} in this class from now on, I will mean a ring that has commutative multiplication \inlinemetanote{writes Ring: commutative mult.} and always has a unit: $1$ \inlinemetanote{writes unit $1$}.

If you look in the literature about what the definition of a ring is, of course a ring does not require the multiplication to be commutative, and that's reasonable. If you get off the plane in Tokyo and you meet a mathematician and you say, \qt{I'm working on a ring,} that person will not know whether it's commutative or not; you have to tell them. And also, you can define rings without units. So, the fancy words for this would be a \emph{commutative ring with unit}. That's just telling the person that that's the version of ring you're signing up for. You can study anything you want, but in this class, we will be considering for the most part commutative rings with unit. I just wanted to say that.
\end{speech}

\begin{content}[Convention: Commutative Rings with Unit]
\begin{remark}[Standing Convention for Rings in this Course]\label[remark]{rem:ring-convention}
In the general mathematical literature, rings are not necessarily required to have commutative multiplication, nor are they always required to possess a multiplicative identity $1$ (sometimes called \newterm{rngs} or pseudo-rings).

\par\vspace{1ex}\noindent
\textbf{Standing Convention:} In this course, unless explicitly stated otherwise, the term \newterm{ring} will always designate a \textbf{commutative ring with unit}: a ring with unit $1$ in the sense of \cref{def:ring} such that $a \cdot b = b \cdot a$ for all $a, b \in R$. The zero ring ($1 = 0$, \cref{prop:zero-ring}) is still allowed.
\end{remark}
\end{content}

\begin{aiNote}[Terminology Across the Literature: Rings vs. Commutative Rings]
% [PERSISTENT, ADDITION]
Mathematical literature reflects varying conventions:
\begin{itemize}
    \item \textbf{Commutative Algebra \& Algebraic Geometry (Atiyah--Macdonald, Eisenbud):} A \qt{ring} implicitly means a commutative ring with multiplicative identity $1$.
    \item \textbf{General Algebra (Bourbaki, Lang, Artin):} A \qt{ring} possesses a multiplicative identity $1$, but multiplication is not required to be commutative.
    \item \textbf{Functional Analysis \& Operator Algebras:} Rings often lack a multiplicative identity $1$ (sometimes called \qt{rngs} or pseudo-rings).
\end{itemize}
In this course, Prof. Pandharipande adheres to the commutative algebra convention: unless explicitly stated otherwise, all rings are commutative and unital. The zero ring, where $1 = 0$, is allowed (\cref{prop:zero-ring}).
\end{aiNote}

\begin{aiNote}[Beyond the Lecture: What Changes Without Commutativity]
% [PERSISTENT, ADDITION]
Typical non-commutative rings are matrix rings $M_n(K)$ for $n \ge 2$ and endomorphism rings. Because this course assumes commutativity, the following distinctions never come up in the later lectures:
\begin{itemize}
    \item \textbf{Ideals (\cref{def:ideal}):} left ideals ($r \cdot x \in I$), right ideals ($x \cdot r \in I$) and two-sided ideals ($r \cdot x \cdot s \in I$) coincide in a commutative ring. In general they differ, and kernels of ring homomorphisms are always two-sided ideals.
    \item \textbf{Divisibility (\cref{def:divisibility-ring}, Week~2):} in general one distinguishes left divisibility ($b = a \cdot c$) from right divisibility ($b = c' \cdot a$).
    \item \textbf{Units (\cref{def:unit-ring}, Week~2):} in general a unit must have a two-sided inverse ($u w = w u = 1$). The units still form a group, with $(u_1 u_2)^{-1} = u_2^{-1} u_1^{-1}$, but it can be non-abelian, e.g., $U(M_n(K)) = \mathrm{GL}_n(K)$ for $n \ge 2$.
    \item \textbf{Integral domains (\cref{def:integral-domain}, Week~2):} the term presupposes commutativity. A non-commutative ring without zero divisors is just called a \emph{domain}, e.g., the real quaternions $\mathbb{H}$.
\end{itemize}
\end{aiNote}

\subsection{Definition of a Ring Homomorphism}

\begin{speech}[00:18:44 - 00:20:03]
Okay, so, the definition of a homomorphism: a homomorphism \inlinemetanote{writes A homomorphism} $f$ from $R_1$ to $R_2$ \inlinemetanote{writes $f \colon R_1 \to R_2$} --- maybe I put here: a ring homomorphism \inlinemetanote{writes Ring in front of homomorphism} --- is defined to satisfy some properties.

The first is that $f(a + b) = f(a) + f(b)$ \inlinemetanote{writes (i) $f(a+b) = f(a) + f(b)$}. This is saying that this map respects the linear structure, the additive structure of the two rings. So, this should remind you from linear algebra class: there's the idea of linear homomorphisms of vector spaces, and it has this linear condition --- that's this linear condition.

The second condition is that it respects the multiplication, so you can imagine what's going to be written here \inlinemetanote{writes (ii) $f(ab) = f(a)f(b)$}, okay? And the third condition is that it respects the identity: $f(1) = 1$ \inlinemetanote{writes (iii) $f(1) = 1$}.

Those are the three conditions. That's the definition of a homomorphism of rings.
\end{speech}

\begin{content}[Definition of a Ring Homomorphism]
\begin{definition}[Ring Homomorphism]\label[definition]{def:ring-homomorphism}
Let $R_1$ and $R_2$ be rings. A mapping $f \colon R_1 \to R_2$ is called a \newterm{ring homomorphism} if it satisfies the following three conditions for all $a, b \in R_1$:
\begin{enumerate}[label=\textbf{(\roman*)}]
    \item \textbf{Additive Compatibility:}
    \begin{equation}\label{eq:hom-add}
    f(a + b) = f(a) + f(b).
    \end{equation}
    \item \textbf{Multiplicative Compatibility:}
    \begin{equation}\label{eq:hom-mult}
    f(a \cdot b) = f(a) \cdot f(b).
    \end{equation}
    \item \textbf{Identity Preservation:}
    \begin{equation}\label{eq:hom-unit}
    f(1_{R_1}) = 1_{R_2}.
    \end{equation}
\end{enumerate}
\end{definition}

\begin{steps}[Comparison to Linear Algebra and Group Theory]
Condition \textbf{(i)} asserts that $f$ is a group homomorphism between the additive groups $(R_1, +)$ and $(R_2, +)$, mirroring the linearity condition $T(u + v) = T(u) + T(v)$ from linear algebra. Condition \textbf{(ii)} ensures compatibility with the multiplicative monoid structure. Condition \textbf{(iii)} explicitly guarantees that the multiplicative unit of the domain maps to the multiplicative unit of the codomain.
\end{steps}
\end{content}

\subsection{Preservation of the Additive Identity: \texorpdfstring{$f(0) = 0$}{f(0) = 0}}

\begin{speech}[00:20:03 - 00:22:25]
There are some puzzling aspects about this definition. You could say, \qt{Well, why don't I say it respects the zero?} That's a very reasonable question. So, you could say: what about $f(0)$ \inlinemetanote{writes What about $f(0)$?}?

Well, we can write $f(0 + 0)$ \inlinemetanote{writes $f(0+0)$}. I want to use property \textbf{(i)}: $f(0 + 0) = f(0) + f(0)$ \inlinemetanote{writes $= f(0) + f(0)$}, right? That's just applying condition \textbf{(i)} of the homomorphism \inlinemetanote{writes condition (i) of homomorphism}. But now, $0 + 0$ is $0$, we know that, because $0$ is the neutral element for addition. So, I can rewrite this as $f(0) = f(0) + f(0)$ \inlinemetanote{writes $f(0) = f(0) + f(0)$}, okay?

And then I can use left cancellation by taking the inverse. I've said that several times, and I hope you know what I mean by that: if I take such an equation... Let me just do it once formally, because I said it several times. I have one $f(0)$ here and two $f(0)$'s here, and I want to get rid of the $f(0)$ on both sides. So, what I do --- and this is what I mean by left cancellation --- is that there exists some $b \in R_2$ \inlinemetanote{writes $\exists b \in R_2$} such that $b + f(0) = 0$ \inlinemetanote{writes such that $b + f(0) = 0$}. That's the group property on $R_2$.

And then I just add $b$ to both sides of this. So, I get $b + f(0) = b + f(0) + f(0)$ \inlinemetanote{writes $b + f(0) = (b + f(0)) + f(0)$}. I use this cancellation, I use the additive inverse of $f(0)$ on both sides, and then, since that equals $0$, this is $0$ \inlinemetanote{writes $0$ under $b+f(0)$}, that's $0$ \inlinemetanote{writes $0$ under $(b+f(0))$}, and finally I conclude that $f(0) = 0$ \inlinemetanote{writes $0 = f(0)$ and boxes $f(0) = 0$}.

So, it is true that $f(0) = 0$ (i.e., $f(0_{R_1}) = 0_{R_2}$), but I don't have to put it here. It's just traditional not to. You could put it here, it wouldn't make any difference.
\end{speech}

\begin{content}[Automatic Preservation of the Additive Identity]
\begin{proposition}[Preservation of Zero]\label[proposition]{prop:homomorphism-preserves-zero}
Let $f \colon R_1 \to R_2$ be a ring homomorphism. Then $f$ automatically preserves the additive identity:
\begin{equation}\label{eq:hom-preserves-zero}
f(0_{R_1}) = 0_{R_2}.
\end{equation}
\end{proposition}

\begin{proof}
Because $0_{R_1}$ is the additive identity of $R_1$, we have $0_{R_1} + 0_{R_1} = 0_{R_1}$. Applying property \textbf{(i)} of \cref{def:ring-homomorphism}:
\begin{equation}\label{eq:fzero-double}
f(0_{R_1}) = f(0_{R_1} + 0_{R_1}) = f(0_{R_1}) + f(0_{R_1}).
\end{equation}
Since $(R_2, +, 0_{R_2})$ is an abelian group, the element $f(0_{R_1}) \in R_2$ possesses an additive inverse $b \in R_2$ such that $b + f(0_{R_1}) = 0_{R_2}$. Adding $b$ to both sides of \eqref{eq:fzero-double}:
\begin{align*}
b + f(0_{R_1}) &= b + \bigl( f(0_{R_1}) + f(0_{R_1}) \bigr) \\
0_{R_2} &= \bigl( b + f(0_{R_1}) \bigr) + f(0_{R_1}) && \text{(by associativity in $R_2$)} \\
0_{R_2} &= 0_{R_2} + f(0_{R_1}) \\
0_{R_2} &= f(0_{R_1}).
\end{align*}
Thus, $f(0_{R_1}) = 0_{R_2}$.
\end{proof}
\end{content}

\subsection{Why Identity Preservation Must Be Axiomatized}

\begin{speech}[00:22:25 - 00:24:16]
Then maybe the question is: \qt{Well, then why do I put $f(1) = 1$?} That's a reasonable question. They're treated not symmetrically.

And the answer to that... \inlinemetanote{erases the cancellation calculation for $f(0)$ on the board}

Yeah, it has to do with the fact that I cannot... I mean, maybe this is what you're hinting at, but I don't have this cancellation property in multiplication. I have this cancellation --- the existence of an additive inverse --- that is an axiom of the ring, so I can use it. But there is not an existence of a multiplicative inverse, so I cannot just follow these steps. So, I have to assume it by hand.

This is about the definitions, and as I said, the strange thing about mathematics is that when it's taught, the definitions are first, but somehow when it's understood, the definitions are last. So, it's somehow a constant difficulty. \inlinemetanote{erases more board space and begins writing}

So, the point is: we don't have \inlinemetanote{writes We don't have} multiplicative inverses \inlinemetanote{writes mult. inverses} --- we don't always have multiplicative inverses \inlinemetanote{inserts always} --- in $R_1$ \inlinemetanote{writes in $R_1$}, and so, we can't \inlinemetanote{writes so we can't} use this argument \inlinemetanote{writes use this argument} --- the one that's in the corner there (i.e., the cancellation argument for $f(0) = 0$, \cref{prop:homomorphism-preserves-zero}) \inlinemetanote{points at the proof of $f(0)=0$ in the corner of the board}.

So, to fix that, we just say at the beginning: $f(1) = 1$. We don't leave it to chance. You could hope, maybe, that even though we couldn't use this argument, maybe you could find some different argument. You could hope that, but it's not true.
\end{speech}

\begin{speech}[00:24:16 - 00:25:08]
For example, I could take a bad example \inlinemetanote{writes A bad example}: I could take $R$ and map it to... sorry, I could take $R_1 \to R_2$ \inlinemetanote{writes $R_1 \to R_2$} and map it so everything goes to zero. Say they're both healthy rings --- in fact, they could both be $\mathbb{Z}$ \inlinemetanote{writes $\mathbb{Z}$ above $R_1$ and $\mathbb{Z}$ above $R_2$}. I could map every element to zero \inlinemetanote{writes $a \mapsto 0$}.

If you map every element to zero, you certainly are not mapping $1$ to $1$. But you won't get a contradiction with every element being zero; all the other properties will be satisfied.

So, it's not a good idea to think about this, but the fact of the matter is: we need to just assume it. \inlinemetanote{erases the bad example} Okay.
\end{speech}

\begin{content}[The Asymmetry Between Addition and Multiplication]
\textbf{Board State Note:}
\begin{center}
\qt{We don't have mult.\ inverses in $R_2$, so we can't use this argument.}
\end{center}

\begin{proposition}[Multiplicative Idempotence of \texorpdfstring{$f(1)$}{f(1)}]\label[proposition]{prop:f1-idempotent}
Let $f \colon R_1 \to R_2$ satisfy conditions \textbf{(i)} and \textbf{(ii)} of a ring homomorphism. Because $1_{R_1} \cdot 1_{R_1} = 1_{R_1}$, multiplicativity implies:
\begin{equation}\label{eq:f1-idempotent}
f(1_{R_1}) = f(1_{R_1} \cdot 1_{R_1}) = f(1_{R_1}) \cdot f(1_{R_1}) = f(1_{R_1})^2.
\end{equation}
Thus, $f(1_{R_1})$ is an \newterm{idempotent} element in $R_2$. However, because elements in a general ring do not necessarily possess multiplicative inverses, we cannot cancel $f(1_{R_1})$ to deduce that $f(1_{R_1}) = 1_{R_2}$.
\end{proposition}

\begin{steps}[Why $f(1) = 1$ Does Not Follow from $f(ab) = f(a)f(b)$]
One might wonder why $f(1) = 1$ is required as an explicit axiom in \cref{def:ring-homomorphism}, whereas $f(0) = 0$ is a provable theorem.
\begin{itemize}
    \item \textbf{Additive Case:} In the additive group $(R_2, +)$, every element has an additive inverse. Hence, from $f(0) = f(0) + f(0)$, we can subtract $f(0)$ to obtain $f(0) = 0$.
    \item \textbf{Multiplicative Case:} Using multiplicativity, we know that $f(1) = f(1 \cdot 1) = f(1) \cdot f(1)$. However, in a general ring $R_2$, $f(1)$ does not necessarily possess a multiplicative inverse! While $x^2 = x$ in a division ring (or field) implies $x \in \{0, 1\}$, in a general ring there can be non-trivial idempotents ($e^2 = e$ with $e \neq 0, 1$), or $f(1)$ could simply equal $0$.
\end{itemize}
No cleverer argument can help either: the zero map in \cref{ex:zero-map-non-unital} satisfies the first two conditions but not $f(1) = 1$.
\end{steps}
\end{content}

\begin{content}[Counterexample: The Non-Unital Zero Map]
\begin{example}[The Trivial Zero Map Fails Condition \textbf{(iii)}]\label[example]{ex:zero-map-non-unital}
Let $R_1 = R_2 = \mathbb{Z}$ (or any non-zero rings). Consider the zero map:
\[
f \colon R_1 \to R_2, \quad f(a) := 0 \quad \forall a \in R_1.
\]
This map satisfies conditions \textbf{(i)} and \textbf{(ii)}:
\begin{itemize}
    \item \textbf{Addition:} $f(a + b) = 0 = 0 + 0 = f(a) + f(b)$,
    \item \textbf{Multiplication:} $f(a \cdot b) = 0 = 0 \cdot 0 = f(a) \cdot f(b)$.
\end{itemize}
However, $f(1_{R_1}) = 0 \neq 1_{R_2}$. Without demanding condition \textbf{(iii)} ($f(1) = 1$), the zero map would be considered a ring homomorphism between non-zero rings.
\end{example}
\end{content}

\subsection{Examples of Ring Homomorphisms: The Canonical Map \texorpdfstring{$\mathbb{Z} \to R$}{Z -> R}}

\begin{speech}[00:25:08 - 00:27:20]
Then there's lots of examples \inlinemetanote{writes Examples} of these ring homomorphisms, and it's the main way we compare rings --- by homomorphisms.

For example, the thing I said here (i.e., the canonical map of \cref{def:canonical-map-z}): $\mathbb{Z}$ goes to $R$ \inlinemetanote{writes $\mathbb{Z} \xrightarrow{\phi} R$}, this basic map $\phi$. This is a ring homomorphism \inlinemetanote{writes This is a ring hom.}. Not only does $\mathbb{Z}$ map to $R$ by just seeing where $1$ goes, it's actually a homomorphism. And to check it's a homomorphism, you have to do something.

I'm not going to do this for you, you have to do it, because it's some kind of internal tautology about addition... Where does $2$ go? We said $2$ goes to $1 + 1$ \inlinemetanote{writes $2 \mapsto 1+1$}. Where does $3$ go? $3$ goes to $1 + 1 + 1$ \inlinemetanote{writes $3 \mapsto 1+1+1$}. So, then we'd have to check that if I add these two, $5$ goes to five $1$'s, which looks pretty good \inlinemetanote{writes $2+3 = 5 \mapsto (1+1) + (1+1+1)$}.

And then you'd also have to check: where does $6$ go? $6$ is $2 \times 3$, and this goes to $(1+1) \times (1+1+1)$ \inlinemetanote{writes $6 = 2 \cdot 3 \mapsto (1+1)(1+1+1)$}. And now you'd use the distributive law a lot of times. If you use the distributive law a lot of times, you'll get this to be $1 + 1 + 1 + 1 + 1 + 1$ \inlinemetanote{writes $= 1+1+1+1+1+1$}, which is the definition of where $6$ goes. Okay?

So, if you do this, you'll see that there's essentially nothing going on except understanding that the distributive law is making this a homomorphism.

This kind of mathematics that's on this kind of foundational, definitional level, you pretty much have to do for yourself to understand, you know, because it has a certain kind of logical... And I hope this is clear. But there is some content --- there's a very small amount of content: it's about how this definition of $2$ as $1+1$ interacts with the distributive law. I think most people find this not so controversial, but if you want to understand it, you have to go through the mental effort of checking the things yourself. On the board, I could convince you of anything, probably. Okay.
\end{speech}

\begin{content}[The Canonical Homomorphism \texorpdfstring{$\mathbb{Z} \to R$}{Z -> R}]
\begin{example}[Canonical Homomorphism from the Integers]\label[example]{ex:canonical-hom-z-to-r}
For any ring $(R, +, \cdot, 0_R, 1_R)$, the canonical evaluation map:
\[
\phi \colon \mathbb{Z} \to R, \quad \phi(n) := n \cdot 1_R
\]
is a ring homomorphism.
\end{example}

\textbf{Verification of Homomorphism Axioms:}
\begin{itemize}
    \item \textbf{Preservation of Unit:} $\phi(1) = 1 \cdot 1_R = 1_R$.
    \item \textbf{Additive Compatibility (e.g., $2 + 3 = 5$):}
    \begin{align*}
    \phi(2) + \phi(3) &= (1_R + 1_R) + (1_R + 1_R + 1_R) \\
    &= 1_R + 1_R + 1_R + 1_R + 1_R = \phi(5) = \phi(2 + 3).
    \end{align*}
    \item \textbf{Multiplicative Compatibility (e.g., $2 \cdot 3 = 6$):}
    Using repeated applications of the distributive laws:
    \begin{align*}
    \phi(2) \cdot \phi(3) &= (1_R + 1_R)(1_R + 1_R + 1_R) \\
    &= 1_R(1_R + 1_R + 1_R) + 1_R(1_R + 1_R + 1_R) \\
    &= (1_R + 1_R + 1_R) + (1_R + 1_R + 1_R) \\
    &= 1_R + 1_R + 1_R + 1_R + 1_R + 1_R = \phi(6) = \phi(2 \cdot 3).
    \end{align*}
\end{itemize}

\begin{steps}[The General Check the Lecturer Leaves to You]
The two computations above are instances of two identities, valid for all $m, n \in \mathbb{Z}$:
\begin{itemize}
    \item \textbf{Additivity:} $(m + n) \cdot 1_R = m \cdot 1_R + n \cdot 1_R$. For $m, n \ge 0$ both sides are a sum of $m + n$ copies of $1_R$ (associativity of $+$); the cases with negative $m$ or $n$ follow by adding inverses on both sides (\cref{def:integer-multiples}, \cref{rem:negative-multiples}).
    \item \textbf{Multiplicativity:} $(m \cdot 1_R)(n \cdot 1_R) = (mn) \cdot 1_R$. Expanding the left side with the distributive laws gives $m n$ products $1_R \cdot 1_R = 1_R$ (with the sign of $mn$), exactly as in the computation for $2 \cdot 3 = 6$.
\end{itemize}
Together with $\phi(1) = 1_R$, this proves that $\phi$ is a ring homomorphism.
\end{steps}

\begin{steps}[Universal Initial Property of $\mathbb{Z}$]
In category theory, this property establishes that $\mathbb{Z}$ is the \newterm{initial object} in the category of rings (\cref{ex:cat-rings}): for every ring $R$, there exists a \emph{unique} ring homomorphism from $\mathbb{Z}$ to $R$.
\end{steps}
\end{content}

\begin{hidden}
% timestamp: 00:27:15
% topic: Homomorphism definition verification, canonical map Z -> R
% board_state: ex:canonical-hom-z-to-r, prop:homomorphism-preserves-zero
% next_goal: Subring inclusion homomorphisms, polynomial rings, residue class rings Z/nZ
% open_loops: none
\end{hidden}

\subsection{Inclusion Homomorphisms and Subrings}

\begin{speech}[00:27:20 - 00:28:24]
That's a very simple example of a homomorphism, but in fact there's lots of... all the ones we've seen so far: if you take $\mathbb{Z}$ inside $\mathbb{Q}$ \inlinemetanote{writes $\mathbb{Z} \subset \mathbb{Q}$}, if you view this as a map --- just the embedding map, if you view an integer as a rational number --- then this gives me a homomorphism. And that's a homomorphism inside $\mathbb{R}$ \inlinemetanote{writes $\subset \mathbb{R}$}, and this is a homomorphism inside $\mathbb{C}$ \inlinemetanote{writes $\subset \mathbb{C}$}. So, these inclusions are all homomorphisms. And I can also include this inside the ring of polynomials \inlinemetanote{writes $\mathbb{Z}[x] \subset \mathbb{Q}[x] \subset \mathbb{R}[x] \subset \mathbb{C}[x]$} --- polynomials in one variable \inlinemetanote{writes in one variable}.

These are many of the rings that we will study, but if you start at the bottom, $\mathbb{Z}$, they just include, and all these inclusions are homomorphisms. Because they all respect... you know, whenever I include, I'm not changing the definition of addition or multiplication for the smaller object, and $1$ is always going to $1$.

So, this is an example. Let's try some different kinds of examples... \inlinemetanote{draws a vertical line dividing the board}
\end{speech}

\begin{content}[Inclusion Homomorphisms of Standard Rings]
\begin{example}[Canonical Inclusions]\label[example]{ex:inclusions-standard-rings}
The canonical inclusion embeddings of the standard number systems:
\begin{equation}\label{eq:number-inclusions}
\mathbb{Z} \hookrightarrow \mathbb{Q} \hookrightarrow \mathbb{R} \hookrightarrow \mathbb{C}
\end{equation}
and their single-variable polynomial rings:
\begin{equation}\label{eq:poly-inclusions}
\mathbb{Z}[x] \hookrightarrow \mathbb{Q}[x] \hookrightarrow \mathbb{R}[x] \hookrightarrow \mathbb{C}[x]
\end{equation}
are all injective ring homomorphisms. For any subring $S \subseteq R$ (defined formally in Week~2, \cref{def:subring}), the inclusion map $\iota \colon S \hookrightarrow R$ defined by $\iota(s) = s$ is a ring homomorphism because $S$ inherits the operations and the unit $1_S = 1_R$ directly from $R$.
\end{example}
\end{content}

\subsection{The Ring of Residue Classes \texorpdfstring{$\mathbb{Z}/n\mathbb{Z}$}{Z/nZ} and the Projection Homomorphism}

\begin{speech}[00:28:24 - 00:30:48]
One example is that you've studied this ring $\mathbb{Z}/n\mathbb{Z}$ \inlinemetanote{writes $\mathbb{Z}/n\mathbb{Z}$}... Well, if we haven't, let me at least say what it is, because this also plays a role in the class. So, this ring here, this set here... \inlinemetanote{writes $(\mathbb{Z}/n\mathbb{Z}, +, \cdot, \bar{0}, \bar{1})$}. This is a standard name; probably in other courses it's also given other names. But the words are \qt{$\mathbb{Z} \bmod n\mathbb{Z}$}, and this is the set of residue classes $\bmod n$ \inlinemetanote{writes $= \{\text{set of residue classes mod } n\}$}. If you want to be explicit, this is equal to the set $\bar{0}, \bar{1}, \bar{2}, \dots, \overline{n-1}$ \inlinemetanote{writes $= \{\bar{0}, \bar{1}, \bar{2}, \dots, \overline{n-1}\}$}, because $n$ is zero. Is this notation reasonable notation? It's kind of semi-standard mathematical notation. There's exactly $n$ elements of this; it's just the residue classes $\bmod n$.

And you can add residue classes, you can multiply them, there's the zero, and there's the one. So, if I take this whole data together, this exactly is a commutative ring with unit \inlinemetanote{writes Ring!}. So, it's a ring for us. It's a ring of a slightly different nature than the other ones because it has finitely many elements. It's an interesting ring to think about. In some sense, one of the simplest and most confusing ones is this one \inlinemetanote{writes $\mathbb{Z}/2\mathbb{Z}$}: this is a ring with two elements \inlinemetanote{writes $= \{\bar{0}, \bar{1}\}$}, zero and one. That's kind of a fun ring, but it's a full-fledged ring. You know, a ring is a ring no matter how small.

And we get here: $\mathbb{Z}$ goes to $\mathbb{Z}/n\mathbb{Z}$ \inlinemetanote{writes $\mathbb{Z} \to \mathbb{Z}/n\mathbb{Z}$}; this is also a homomorphism \inlinemetanote{writes $f$ over the arrow}. It's a ring homomorphism where I take $f(a)$ --- for any integer $a$ --- to $\bar{a}$ \inlinemetanote{writes $f(a) = \bar{a}$}. Okay? So, that's another example.
\end{speech}

\begin{content}[Residue Class Rings and the Reduction Modulo $n$]
\begin{definition}[Ring of Residue Classes Modulo $n$]\label[definition]{def:residue-class-ring}
For any integer $n \ge 1$, the ring of \newterm{residue classes modulo $n$} (denoted $\mathbb{Z}/n\mathbb{Z}$ or $\mathbb{Z}_n$) is the set:
\begin{equation}\label{eq:zn-set}
\mathbb{Z}/n\mathbb{Z} := \{ \bar{0}, \bar{1}, \bar{2}, \dots, \overline{n-1} \},
\end{equation}
equipped with modular addition and multiplication:
\[
\bar{a} + \bar{b} := \overline{a + b}, \quad \bar{a} \cdot \bar{b} := \overline{a \cdot b},
\]
with additive identity $\bar{0}$ and multiplicative identity $\bar{1}$. Here $\bar{a}$ denotes the residue class of $a \in \mathbb{Z}$, and $\bar{a} = \bar{b}$ if and only if $n \mid (a - b)$.
\end{definition}

\begin{steps}[Why the Operations Are Well-Defined]
The formulas use representatives $a, b$ of the classes, so they must not depend on that choice. If $\bar{a} = \bar{a}'$ and $\bar{b} = \bar{b}'$, say $a' = a + kn$ and $b' = b + \ell n$, then
\[
a' + b' = (a + b) + (k + \ell) n, \qquad a' b' = ab + (a \ell + k b + k \ell n) n,
\]
so $\overline{a' + b'} = \overline{a + b}$ and $\overline{a' b'} = \overline{ab}$. The ring axioms are then inherited from $\mathbb{Z}$, and so is commutativity. For $n = 1$ one gets the zero ring $\{\bar{0}\}$.
\end{steps}

\begin{example}[The Two-Element Field $\mathbb{Z}/2\mathbb{Z}$]\label[example]{ex:z2-ring}
For $n = 2$, $\mathbb{Z}/2\mathbb{Z} = \{\bar{0}, \bar{1}\}$ is the smallest non-trivial commutative ring with unit ($1 \neq 0$). The addition and multiplication tables are given by:
\begin{center}
\begin{tabular}{c|cc}
$+$ & $\bar{0}$ & $\bar{1}$ \\ \hline
$\bar{0}$ & $\bar{0}$ & $\bar{1}$ \\
$\bar{1}$ & $\bar{1}$ & $\bar{0}$
\end{tabular}
\qquad\qquad
\begin{tabular}{c|cc}
$\cdot$ & $\bar{0}$ & $\bar{1}$ \\ \hline
$\bar{0}$ & $\bar{0}$ & $\bar{0}$ \\
$\bar{1}$ & $\bar{0}$ & $\bar{1}$
\end{tabular}
\end{center}
This is a non-trivial commutative ring with unit despite containing only two elements.
\end{example}

\begin{proposition}[Canonical Projection Homomorphism]\label[proposition]{prop:canonical-projection}
The reduction map modulo $n$ (written $f$ on the board):
\begin{equation}\label{eq:canonical-projection}
\pi \colon \mathbb{Z} \to \mathbb{Z}/n\mathbb{Z}, \quad a \mapsto \bar{a} := a \bmod n
\end{equation}
is a surjective ring homomorphism.
\end{proposition}

\begin{proof}
For any $a, b \in \mathbb{Z}$:
\begin{enumerate}[label=\textbf{(\roman*)}]
    \item \textbf{Additivity:} $\pi(a + b) = \overline{a + b} = \bar{a} + \bar{b} = \pi(a) + \pi(b)$.
    \item \textbf{Multiplicativity:} $\pi(a \cdot b) = \overline{a \cdot b} = \bar{a} \cdot \bar{b} = \pi(a) \cdot \pi(b)$.
    \item \textbf{Unit Preservation:} $\pi(1) = \bar{1}$.
\end{enumerate}
Surjectivity holds because every residue class $\bar{k} \in \mathbb{Z}/n\mathbb{Z}$ is the image of the integer $k \in \mathbb{Z}$.
\end{proof}
\end{content}

\section{The Kernel of a Ring Homomorphism and Ideals}
\subsection{Definition of the Kernel}

\begin{speech}[00:30:49 - 00:32:06]
There are lots of interesting examples; I'm not going to get to them all today. But I want to say that \inlinemetanote{pushes up the lower board} for homomorphisms, like for vector spaces, whenever you have one, you can consider the kernel and the image. \inlinemetanote{erases the board that came down: the definition of a ring homomorphism and the proof of $f(0)=0$} And it turns out, in a maybe surprising way, that for ring theory, the notion of the kernel turns out to be much more important than the notion of the image.

In linear algebra, if you have a homomorphism between two vector spaces --- did you use the word \qt{homomorphism}? You could also say \qt{linear transformation}; both those words are used --- if you have a linear transformation between two vector spaces, there's the notion of its kernel (those are the elements that go to zero) and the notion of its image. And they're both subvector spaces. In linear algebra, these things are kind of on similar footing. But in ring theory, when you have a homomorphism, there's again the kernel and the image. So, we will start with the kernel first; probably next time we'll say a little bit about the image.

So, this is an interesting object. Here we have: let...
\end{speech}

\begin{speech}[00:32:06 - 00:32:56]
Let $f$ from $R_1$ to $R_2$ \inlinemetanote{writes Let $f \colon R_1 \to R_2$} be a homomorphism of rings \inlinemetanote{writes be a hom. of rings}. Then we define the kernel of $f$ \inlinemetanote{writes Then we define the kernel of $f$}: $\ker(f)$ \inlinemetanote{writes $\ker(f) =$} --- it's those elements in the first ring, because the kernel lives in the first ring \inlinemetanote{writes $\{a \in R_1 \mid$}, such that $f(a)$ is equal to zero \inlinemetanote{writes $f(a) = 0\}$}. Okay? I hope this definition is clear: it's everybody that goes to zero.

In particular, you could say: what is the kernel of this? \inlinemetanote{writes $\ker(f) \subset \mathbb{Z} \xrightarrow{f} \mathbb{Z}/n\mathbb{Z}$} What is the kernel here?
\end{speech}

\begin{student}[Student Answer]
Multiples of $n$.
\end{student}

\begin{speech}[00:32:57 - 00:33:17]
Yeah. The kernel here is exactly all multiples of $n$, and this is often denoted $n\mathbb{Z}$ \inlinemetanote{writes $n\mathbb{Z}$ under $\ker(f)$}. This is equal to the set of multiples of $n$ \inlinemetanote{writes \qt{mults of $n$}}. That's the kernel.
\end{speech}

\begin{content}[Definition of the Kernel of a Ring Homomorphism]
\begin{definition}[Kernel of a Ring Homomorphism]\label[definition]{def:kernel-ring-hom}
Let $R_1, R_2$ be rings and let $f \colon R_1 \to R_2$ be a ring homomorphism. The \newterm{kernel} of $f$, denoted by $\ker(f)$, is the pre-image of the zero element $0_{R_2}$:
\begin{equation}\label{eq:kernel-def}
\ker(f) := \{ a \in R_1 \mid f(a) = 0_{R_2} \} \subseteq R_1.
\end{equation}
\end{definition}

\begin{example}[Kernel of the Modular Reduction Homomorphism]\label[example]{ex:kernel-of-projection}
For the canonical projection $\pi \colon \mathbb{Z} \to \mathbb{Z}/n\mathbb{Z}$ given by $\pi(a) = \bar{a}$, an integer $a$ maps to $\bar{0}$ if and only if $a$ is a multiple of $n$. Hence:
\begin{equation}\label{eq:kernel-zn}
\ker(\pi) = n\mathbb{Z} := \{ k \cdot n \mid k \in \mathbb{Z} \}.
\end{equation}
\end{example}

\begin{steps}[The Kernel as a Group Kernel]
Since $f$ respects addition and $f(0_{R_1}) = 0_{R_2}$ (\cref{prop:homomorphism-preserves-zero}), $\ker(f)$ is the kernel of the group homomorphism $f \colon (R_1, +, 0) \to (R_2, +, 0)$ in the sense of Week~2 (\cref{def:kernel}). The multiplication is what makes the ring kernel special, see \cref{prop:properties-of-kernel}.
\end{steps}

\begin{insight}[Kernels versus Images in Ring Theory]
In linear algebra, the kernel and image of a linear map $T \colon V \to W$ enjoy a symmetric status: both are vector subspaces, connected by the rank-nullity theorem $\dim V = \dim \ker(T) + \dim \operatorname{im}(T)$.

In ring theory, however, their algebraic roles diverge profoundly:
\begin{itemize}
    \item The image $\operatorname{im}(f) \subseteq R_2$ is merely a \textbf{subring} of the codomain.
    \item The kernel $\ker(f) \subseteq R_1$ is not merely a subring, but an \textbf{ideal} of the domain: it is closed under addition and absorbs multiplication from the entire ring ($r \cdot a \in \ker(f)$ for all $r \in R_1$ and $a \in \ker(f)$).
\end{itemize}
This absorption property makes the kernel the precise algebraic structure needed to construct quotient rings $R_1/\ker(f)$ and establish the First Isomorphism Theorem for rings ($R_1/\ker(f) \cong \operatorname{im}(f)$).
\end{insight}
\end{content}

\subsection{Properties of the Kernel}

\begin{speech}[00:33:17 - 00:34:49]
So, the kernel has very specific properties \inlinemetanote{writes Properties of the ker(f)}. The first property is that it's closed under addition: if we have $a$ and $b$ in the kernel of $f$ \inlinemetanote{writes bullet: $a, b \in \ker(f)$}, this implies $a + b$ is in the kernel \inlinemetanote{writes $\implies (a+b) \in \ker(f)$}. This is the linear algebra property. All of these things are easy to prove. To check whether $a + b$ is in the kernel, I have to calculate $f(a + b)$ \inlinemetanote{writes $f(a+b) =$}. Using the definition of homomorphism, this is $f(a) + f(b)$ \inlinemetanote{writes $= f(a) + f(b)$}. Then, using the fact that $a$ and $b$ are both in the kernel, well, this one is zero \inlinemetanote{writes $0$ under $f(a)$} and that one is zero \inlinemetanote{writes $0$ under $f(b)$}, so the whole thing is zero \inlinemetanote{writes $= 0$} --- and so, it's in the kernel, right? I should have said even earlier: we already know that zero is in the kernel \inlinemetanote{writes bullet: $0 \in \ker(f)$}.

If $a$ is in the kernel \inlinemetanote{writes bullet: If $a \in \ker(f)$} and $r$ is any element of the ring \inlinemetanote{writes $, r \in R$}, then $r \times a$ is in the kernel \inlinemetanote{writes $\implies ra \in \ker(f)$}.
\end{speech}

\begin{speech}[00:34:49 - 00:36:11]
That's a wider property. These are the linear properties --- we'll say more about them in a moment --- but this property is about the multiplication. It says that if I take an element of the kernel, I can then go pick any element of the ring and multiply it, and it'll be in the kernel. So, it's a more interesting property. You can see that here: if you have an element of the kernel, it's $n$ times something. If I multiply it by any integer, it's still $n$ times something.

This is not hard to prove either. We have to calculate $f(r \times a)$ \inlinemetanote{writes $f(ra) =$}, and then, of course, we use the multiplicative property \inlinemetanote{writes $= f(r)f(a)$}. And now, this is the point: in the multiplicative property, in order to deduce that this product is zero, we only need one of the factors to be zero. For the addition, we needed both, but for the multiplication, we only need to know that one of the factors is zero, because that'll kill the product. And that's already there because $a$ is in the kernel. So, this is zero \inlinemetanote{writes $0$ under $f(a)$}, and so, this is zero \inlinemetanote{writes $= 0$}, okay?

And so, in particular, if $a$ is in the kernel, $-a$ is in the kernel \inlinemetanote{writes bullet: $a \in \ker(f) \implies -a \in \ker(f)$}. That's the additive inverse, because for $-a$, $r$ is equal to $-1$. So, those are all properties enjoyed by the kernel...
\end{speech}

\begin{content}[Algebraic Properties of the Kernel]
\begin{proposition}[Axiomatic Properties of the Kernel]\label[proposition]{prop:properties-of-kernel}
Let $f \colon R_1 \to R_2$ be a ring homomorphism. The kernel $\ker(f) \subseteq R_1$ satisfies the following properties:
\begin{enumerate}[label=\textbf{(\arabic*)}]
    \item \textbf{Containment of Zero:} $0_{R_1} \in \ker(f)$.
    \item \textbf{Closure under Addition:} $\forall a, b \in \ker(f) \implies a + b \in \ker(f)$.
    \item \textbf{Closure under Additive Inverses:} $\forall a \in \ker(f) \implies -a \in \ker(f)$.
    \item \textbf{Multiplicative Absorption (Ideal Property):} $\forall a \in \ker(f), \forall r \in R_1 \implies r \cdot a \in \ker(f)$.
\end{enumerate}
\end{proposition}

\begin{proof}
We verify each property directly using the axioms of a ring homomorphism:
\begin{itemize}
    \item \textbf{Property (1):} By \cref{prop:homomorphism-preserves-zero}, $f(0_{R_1}) = 0_{R_2}$, so $0_{R_1} \in \ker(f)$.
    \item \textbf{Property (2):} If $a, b \in \ker(f)$, then $f(a) = 0$ and $f(b) = 0$. By additivity:
    \[
    f(a + b) = f(a) + f(b) = 0 + 0 = 0 \implies a + b \in \ker(f).
    \]
    \item \textbf{Property (4):} If $a \in \ker(f)$ and $r \in R_1$, then $f(a) = 0_{R_2}$. By multiplicativity and \cref{prop:mult-by-zero}:
    \[
    f(r \cdot a) = f(r) \cdot f(a) = f(r) \cdot 0_{R_2} = 0_{R_2} \implies r \cdot a \in \ker(f).
    \]
    \item \textbf{Property (3):} Setting $r = -1_{R_1}$ in Property (4) and using \cref{lem:additive-inverse-product}:
    \[
    -a = (-1_{R_1}) \cdot a \in \ker(f).
    \]
\end{itemize}
\end{proof}

\begin{steps}[Why Multiplicative Absorption Works: Additive vs. Multiplicative Absorption]
In the additive closure proof (Property (2)), both $f(a)$ and $f(b)$ had to be zero because the sum of two elements is zero only when one is the additive inverse of the other ($0 + 0 = 0$). In contrast, for Property (4), multiplication by zero in a ring is absorbing: as soon as one factor in a product is zero, the entire product vanishes ($f(r) \cdot 0_{R_2} = 0_{R_2}$), regardless of whether $f(r)$ is zero or not. Thus, multiplying an element of $\ker(f)$ by \emph{any} element $r \in R_1$ from the ambient ring keeps the result inside $\ker(f)$. This strong absorption property distinguishes ring kernels from mere subrings (\cref{def:subring}, Week~2).
\end{steps}
\end{content}

\subsection{Ideals in a Ring}

\begin{speech}[00:36:11 - 00:38:07]
Those properties can be said in a more concise way, which is a more natural way to say it, and it comes with a definition: the definition of an ideal. \inlinemetanote{pushes up the board and erases the lower board} This is the definition of an ideal $I$ in $R$ \inlinemetanote{writes Def: Ideal $I \subset R$}.

The properties are: $0$ must be in $I$ \inlinemetanote{writes $\bullet \ 0 \in I$}. And the second way to say things is that if I take $I$, and I take the $+$ from $R$, and I take the $0$, this is an abelian group \inlinemetanote{writes $\bullet \ (I, +, 0) \text{ abelian group}$}. Remember, whenever I write $R$, I always mean $R$ with the two operations and this data \inlinemetanote{writes $I \subset R = (R, +, \cdot, 0, 1)$}. So, to be an ideal, you have to have zero, and then when I consider this subset with $+$ and $0$, it should be an abelian group. This means it's closed under addition, and it's closed under additive inverses, okay?

And the second thing is that for every $a$ in $I$ and for every $r$ in $R$, $ra$ remains in $I$ \inlinemetanote{writes $\bullet \ \forall a \in I, \forall r \in R \implies ra \in I$}. Okay? That's the ideal condition.

This is an extremely important definition for the class. It tells you various subsets of $R$ are ideals, and they have this linear property, which is kind of like linear algebra, but then they also have this multiplicative property, which is not like linear algebra. And the role of an ideal will play a little bit like the role of a subspace in linear algebra.
\end{speech}

\begin{content}[Definition of an Ideal]
\begin{definition}[Ideal]\label[definition]{def:ideal}
Let $R = (R, +, \cdot, 0, 1)$ be a commutative ring with unit. A subset $I \subseteq R$ is called an \newterm{ideal} of $R$ if it satisfies:
\begin{enumerate}[label=\textbf{(\arabic*)}]
    \item \textbf{Additive Subgroup:} $(I, +, 0)$ is an abelian subgroup of $(R, +, 0)$:
    \begin{itemize}
        \item $0 \in I$,
        \item $\forall a, b \in I \implies a + b \in I$,
        \item $\forall a \in I \implies -a \in I$.
    \end{itemize}
    \item \textbf{Multiplicative Absorption (Ideal Condition):}
    \begin{equation}\label{eq:ideal-absorption}
    \forall a \in I, \ \forall r \in R \implies r \cdot a \in I.
    \end{equation}
\end{enumerate}
\end{definition}

\begin{center}
\begin{tikzpicture}[scale=1.3, >=Stealth]
    % \begin{hidden}
    % - Canvas: scale=1.3; R_1 ellipse at (0,0), radii 2.4 x 1.8; I = ker(f) ellipse at (0.2,-0.2), radii 1.25 x 0.95;
    %   R_2 ellipse at (6,0), radii 1.8 x 1.5 (left edge x = 4.2)
    % - Order: (1) R_1; (2) I inside R_1; (3) R_2; (4) points a, ra in I and r in R_1 \ I; (5) absorption arrow a -> ra;
    %   (6) arrow f between the ovals; (7) dashed arrow from I to 0_{R_2}
    % - Checks: r = (-1.65,-0.75) lies in R_1 (0.47 + 0.17 < 1) but not in I ((1.85/1.25)^2 > 1);
    %   ra = (-0.35,-0.65) lies in I (0.19 + 0.22 < 1); a = (0.75,-0.55) lies in I (0.19 + 0.14 < 1)
    % - Labels: a below right, ra below, r below; "multiply by r" above the arc a -> ra, clear of the I label at y = 0.35
    % - Collisions: none; the kernel arrow leaves I at its right edge (x = 1.45) and runs below f
    % \end{hidden}
    % Ring R_1 (domain)
    \draw[thick, MidnightBlue, fill=MidnightBlue!5] (0,0) ellipse (2.4cm and 1.8cm);
    \node[MidnightBlue, font=\bfseries] at (-1.55, 1.2) {$R_1$};

    % Ideal I = ker(f)
    \draw[thick, Fuchsia, fill=Fuchsia!15] (0.2,-0.2) ellipse (1.25cm and 0.95cm);
    \node[Fuchsia, font=\bfseries] at (0.2, 0.35) {$I = \ker(f)$};

    % Points a and ra inside I, r outside I
    \coordinate (A) at (0.75, -0.55);
    \coordinate (RA) at (-0.35, -0.65);
    \coordinate (Relem) at (-1.65, -0.75);
    \fill[Fuchsia] (A) circle (1.8pt) node[below right, font=\footnotesize] {$a$};
    \fill[Fuchsia] (RA) circle (1.8pt) node[below, font=\footnotesize] {$ra$};
    \fill[MidnightBlue] (Relem) circle (1.8pt) node[below, font=\footnotesize] {$r \in R_1$};

    % Absorption: multiplying a by r stays inside I
    \draw[->, dashed, thick, Fuchsia] (A) to[bend right=40] node[midway, above, font=\scriptsize] {multiply by $r$} (RA);

    % Codomain R_2 with its zero
    \draw[thick, ForestGreen, fill=ForestGreen!5] (6,0) ellipse (1.8cm and 1.5cm);
    \node[ForestGreen, font=\bfseries] at (6, 1.0) {$R_2$};
    \coordinate (Zero) at (6, -0.3);
    \fill[BrickRed] (Zero) circle (2.2pt) node[right, font=\bfseries] {$0_{R_2}$};

    % The homomorphism and the image of the kernel
    \draw[->, very thick, RoyalBlue] (2.5, 0.55) to[bend left=20] node[midway, above, font=\small\bfseries] {$f$} (4.25, 0.55);
    \draw[->, thick, dashed, Fuchsia] (1.45, -0.3) to[bend right=12] node[pos=0.42, below, font=\footnotesize] {$f(I) = \{0_{R_2}\}$} (Zero);
\end{tikzpicture}
\end{center}

\begin{steps}[Ideals as the Ring-Theoretic Analogue of Normal Subgroups and Subspaces]
In linear algebra, kernels of linear maps are vector subspaces ($T(u + v) = T(u) + T(v)$ and $T(c v) = c T(v)$). A vector subspace $W \subseteq V$ is closed under vector addition ($w_1 + w_2 \in W$) and scalar multiplication by field scalars ($\lambda \cdot w \in W$ for $\lambda \in k$).

In group theory, kernels of group homomorphisms are normal subgroups ($g N g^{-1} \subseteq N$).

In ring theory, kernels of ring homomorphisms are \newterm{ideals}: they are closed under addition ($a + b \in I$) and absorb multiplication by \emph{any} element from the entire ambient ring $R$ ($r \cdot a \in I$ for all $r \in R$).
\end{steps}
\end{content}

\begin{aiNote}[Notation: Ideals in the Lecture and in Textbooks]
% [PERSISTENT]
Many textbooks write $I \trianglelefteq R$ (or $I \lhd R$) for \qt{$I$ is an ideal of $R$}. The lecture does not use this symbol: on the board, the lecturer writes $I \subset R$ and says in words that $I$ is an ideal (\qt{Def: Ideal $I \subset R$}). These notes follow the board. His $\subset$ allows equality, like $\subseteq$; he explains this convention when he defines maximal ideals (\cref{def:maximal-ideal}).
\end{aiNote}

\subsection{The Kernel as an Ideal}

\begin{speech}[00:38:07 - 00:39:13]
The proposition we've proven with this is:

Proposition \inlinemetanote{writes Prop:}: The kernel of a homomorphism $f \colon R_1 \to R_2$ \inlinemetanote{writes The ker of a hom $f \colon R_1 \to R_2$} is an ideal \inlinemetanote{writes is an ideal} in $R_1$ \inlinemetanote{writes in $R_1$}.

That's what we've proven on the board up there (\cref{prop:properties-of-kernel}): with the kernel, you always get an ideal. And rings can have... we will do some examples next time, I'm running out of time here, but for $\mathbb{Z}$, if I take inside here $n\mathbb{Z}$ \inlinemetanote{writes $\mathbb{Z} \supset n\mathbb{Z} = I$}, that's an ideal (it is the kernel of $\pi$, \cref{ex:kernel-of-projection}). Rings can have lots of ideals.

What about a field? I'll just say one more thing about this. Is that clear?

We have lots of examples of rings and lots of examples of ideals, but suppose $F$ is a field \inlinemetanote{writes $F = \text{field}$}, $(F, +, \cdot, 0, 1)$ \inlinemetanote{writes $(F, +, \cdot, 0, 1)$}. What are its ideals? \inlinemetanote{writes What are the ideals?}
\end{speech}

\begin{content}[The Kernel is an Ideal]
\begin{proposition}[Kernel as an Ideal]\label[proposition]{prop:kernel-is-ideal}
Let $f \colon R_1 \to R_2$ be a ring homomorphism. Then $\ker(f)$ is an ideal of $R_1$:
\[
\ker(f) \subset R_1 \text{ is an ideal.}
\]
\end{proposition}

\begin{proof}
This is the synthesis of \cref{prop:properties-of-kernel}:
\begin{enumerate}[label=\textbf{(\arabic*)}]
    \item $(\ker(f), +, 0)$ is an abelian subgroup of $(R_1, +, 0)$ because $0 \in \ker(f)$, $a + b \in \ker(f)$ for all $a, b \in \ker(f)$, and $-a \in \ker(f)$ for all $a \in \ker(f)$.
    \item Multiplicative absorption holds because for all $a \in \ker(f)$ and all $r \in R_1$:
    \[
    f(r \cdot a) = f(r) \cdot f(a) = f(r) \cdot 0 = 0 \implies r \cdot a \in \ker(f).
    \]
\end{enumerate}
Hence, $\ker(f)$ satisfies all axioms of an ideal.
\end{proof}
\end{content}

\subsection{Ideals in a Field}

\begin{speech}[00:39:13 - 00:40:05]
What are its ideals? What are the ideals of a field? Because a field is a ring, so it has the notion of ideals. Every abstract concept of a ring, we can apply to a field, because a field is a particular type of ring: it's a ring where we have multiplicative inverses for non-zero elements. So, what could be an ideal here? What are the ideals?

I'll start with one simple thing: if I just take $I = \{0\}$ \inlinemetanote{writes $I = \{0\}$}. Then zero is in $I$. It is an abelian subgroup just by itself, because $0 + 0 = 0$, and everything... you know, zero is sufficient unto itself \inlinemetanote{laughter}, one might say. And it has this multiplicative property, because if I multiply by anybody, it stays zero. So, zero is always an ideal.

So, what else? Yeah? \inlinemetanote{points at student}
\end{speech}

\begin{student}[Student Answer]
The whole field?
\end{student}

\begin{speech}[00:40:07 - 00:40:21]
The whole field is always an ideal: ideal equals the whole field \inlinemetanote{writes $I = F$}. And, of course, there can be no doubt about that, because if you take the whole thing, it's an abelian group and you're not going to get out of it by multiplying.

Can there be any more? Why not?
\end{speech}

\begin{student}[Student Answer]
Because as soon as you have another element, then you have the inverse...
\end{student}

\begin{speech}[00:40:26 - 00:40:28]
Multiplicative inverse.
\end{speech}

\begin{student}[Student Answer, continued]
...multiplicative inverse, so you have one...
\end{student}

\begin{speech}[00:40:30 - 00:40:37]
One, and then you have everything.

Okay, so, this argument is correct. So, the claim is: for a field, these are the only ideals \inlinemetanote{writes \qt{only ideals!} next to $I = \{0\}$ and $I = F$}.
\end{speech}

\begin{speech}[00:40:37 - 00:42:20]
Which is an important thing to always keep in mind. And the reason is the argument he said: suppose that I have an ideal, and suppose this ideal is not this silly zero ideal \inlinemetanote{writes $I \neq (0)$}. That means there exists somebody in the ideal that's not zero \inlinemetanote{writes $\exists \ a \in I \text{ s.t. } a \neq 0$}. That somebody, because it's a field, has a multiplicative inverse \inlinemetanote{writes $\exists \ a^{-1}$}. There exists an $a^{-1}$, so that $a^{-1} \cdot a$ is also in the ideal \inlinemetanote{writes $a^{-1} a \in I$}, because you're allowed to take anybody and multiply --- this is anybody, and this is somebody in the ideal. So, if we start with $a$ in the ideal, I can multiply by anybody; in particular, I can multiply by its inverse. But then this is one \inlinemetanote{writes $1 \in I$ above $a^{-1}a$, joined by $=$}!

Once I have one in the ideal, then I can easily get anybody in the ideal, because I can just take any $r \times 1$. Okay?

It's an elementary argument, and it's very good to understand. So, fields have very few ideals --- extremely few. But rings can have quite a number of ideals. As I said here, this is the simplest example to see: the ideals in the integers come from multiples of any integer. And ideals in... yeah, we will see some examples of it. But was this argument clear? Sometimes I do these things very quickly, it's just, you know...

The standard thing to say here is that you can view fields as sellouts, because they've lost all their ideals. \inlinemetanote{laughter}
\end{speech}

\begin{content}[Classification of Ideals in a Field]
\begin{theorem}[Ideals in a Field]\label[theorem]{thm:ideals-in-a-field}
Let $F$ be a field. The only ideals of $F$ are the trivial zero ideal $\{0\}$ and the unit ideal $F$:
\begin{equation}\label{eq:field-ideals}
I \subset F \text{ an ideal} \implies I = \{0\} \quad \text{or} \quad I = F.
\end{equation}
\end{theorem}

\begin{proof}
Let $I \subseteq F$ be an ideal of $F$. Clearly, $I = \{0\}$ and $I = F$ are ideals.

Now, suppose $I \neq \{0\}$. Then there exists a non-zero element:
\[
a \in I \quad \text{with } a \neq 0.
\]
Since $F$ is a field, every non-zero element has a multiplicative inverse $a^{-1} \in F$. By the multiplicative absorption property of ideals (\cref{def:ideal}), choosing $r = a^{-1} \in F$:
\[
a^{-1} \cdot a \in I \implies 1_F \in I.
\]
Once the unit $1_F$ lies in $I$, applying absorption once more with any arbitrary element $r \in F$:
\[
r = r \cdot 1_F \in I \quad \forall r \in F.
\]
Therefore, $F \subseteq I$, which forces $I = F$.
\end{proof}

\begin{steps}[Implications for Field Homomorphisms]
A direct consequence of this theorem is that every ring homomorphism $f \colon F \to R$ from a field $F$ to a non-zero ring $R$ is automatically \textbf{injective}:
\[
\ker(f) \subset F \text{ an ideal} \implies \ker(f) = \{0\} \quad (\text{since } f(1_F) = 1_R \neq 0_R \implies \ker(f) \neq F).
\]
Because $\ker(f) = \{0\}$, field homomorphisms into non-zero rings are always embeddings ($f(a) = f(b)$ gives $f(a - b) = 0$, so $a = b$). This is why the extension $g \colon \mathbb{Q} \to F$ in Week~3 is automatically injective (\cref{thm:universal-property-q}).
\end{steps}
\end{content}

\begin{insight}[Fields as \qt{Sellouts}]
The lecturer's humorous pedagogical quip --- that \qt{fields are sellouts because they have lost all their ideals} --- highlights an essential contrast between field theory and ring theory:
\begin{itemize}
    \item In a field $F$, the abundance of units ($F^\times = F \setminus \{0\}$) destroys all non-trivial ideals, as any non-zero element $a \in I$ immediately pulls the identity into the ideal: $1 = a^{-1} \cdot a \in I$ (with $r = a^{-1} \in F$ and $a \in I$), forcing $I = F$.
    \item In general rings (such as $\mathbb{Z}$ or polynomial rings), the scarcity of units allows rich, complex lattices of proper non-zero ideals to flourish (e.g., the ideals $n\mathbb{Z} \subset \mathbb{Z}$, which by \cref{thm:ideals-of-z-are-principal} are all the ideals of $\mathbb{Z}$), which form the foundational building blocks for quotient rings $R/I$ and algebraic geometry.
\end{itemize}
\end{insight}

\begin{aiNote}[Beyond the Lecture: Homomorphisms out of a Field]
% [PERSISTENT, ADDITION]
Every ring homomorphism $f \colon F \to R$ from a field $F$ into a non-zero ring $R$ is injective: $\ker(f)$ is an ideal of $F$, so by \cref{thm:ideals-in-a-field}, either $\ker(f) = \{0\}$ or $\ker(f) = F$. Since $f(1_F) = 1_R \neq 0_R$, $1_F \notin \ker(f)$, which forces $\ker(f) = \{0\}$.
\end{aiNote}

\section{Outlook: Zorn's Lemma and Vector Space Bases}
\subsection{Existence of Bases in Infinite-Dimensional Vector Spaces}

\begin{speech}[00:42:20 - 00:43:29]
For this week, I did some material that was billed to the second week, and I didn't cover Zorn, so I'll say two words about Zorn and then I'll say a little bit more next week.

The first question is about linear algebra, which is this statement which \inlinemetanote{erases the lower left board and prepares it} probably you've seen in the first year, which is a strange statement, but it's about... So, first of all, I will spend some time also next week on Zorn's lemma, so you can read about that in the book \inlinemetanote{writes Zorn's Lemma}.

And it's related to the statement:

Let $V$ be a vector space over a field $k$ \inlinemetanote{writes Let $V$ be a vector space over a field $k$}.

Does $V$ have a basis \inlinemetanote{writes Does $V$ have a basis?}?
\end{speech}

\begin{content}[Zorn's Lemma and the Existence of Vector Space Bases]
\begin{question}[Existence of a Basis for Vector Spaces]\label[question]{q:vector-space-basis}
Let $V$ be a vector space over a field $k$. Does $V$ always have a basis?
\end{question}

\begin{steps}[Finite-Dimensional versus Infinite-Dimensional Vector Spaces]
\begin{itemize}
    \item \textbf{Finite-Dimensional Spaces:} In standard first-year linear algebra, if $V$ is spanned by finitely many vectors, an explicit basis is easily constructed by extracting a maximal linearly independent subset.
    \item \textbf{Infinite-Dimensional Spaces:} When $V$ is infinite-dimensional (e.g., function spaces such as $C(\mathbb{R}, \mathbb{R})$ or $\mathbb{R}^\mathbb{N}$), constructing an explicit (Hamel) basis by hand is impossible. Proving that \emph{every} vector space possesses a basis requires non-constructive set-theoretic axioms---specifically, \newterm{Zorn's Lemma}, which is equivalent to the Axiom of Choice ($\text{AC}$).
\end{itemize}
\end{steps}

\begin{remark}[Infinite-Dimensional Vector Spaces and the Axiom of Choice]\label[remark]{rem:zorns-lemma-basis}
For finite-dimensional vector spaces, the existence of a basis is constructive via Gaussian elimination or the Steinitz exchange lemma. For arbitrary, infinite-dimensional vector spaces (such as function spaces $C(\mathbb{R}, \mathbb{R})$ or sequence spaces $\mathbb{R}^\mathbb{N}$), the assertion that \emph{every vector space possesses a basis} (known as a \newterm{Hamel basis}) cannot be proven constructively in Zermelo--Fraenkel set theory ($\text{ZF}$). It is logically equivalent to the \newterm{Axiom of Choice} ($\text{AC}$) and is standardly established by applying \newterm{Zorn's Lemma} to the poset of linearly independent subsets ordered by inclusion.
\end{remark}
\end{content}

\begin{speech}[00:43:29 - 00:44:01]
Did you consider this question in your first year?

Of course, if it's finite-dimensional, of course it has a basis. And this issue is about whether, if you have some massively infinite-dimensional field (i.e., vector space)... has a basis. And the traditional way to do this is an application of Zorn's lemma, which is related to the axiom of choice in logic --- not my favorite direction to think about, but we will say a little bit about it. But...

Okay, and I think that's it. All right, so... \inlinemetanote{The lecturer puts down the chalk, unclips the microphone, students applaud, and the lecture concludes.}
\end{speech}

\begin{overview}[Summary of Key Concepts]
This lecture established core bridges connecting concrete arithmetic with axiomatic ring theory:
\begin{itemize}
    \item \textbf{Gaussian Integers $\mathbb{Z}[i]$ and Primes:} The ring $\mathbb{Z}[i] = \{a + bi \mid a, b \in \mathbb{Z}\}$ possesses four units $\{\pm 1, \pm i\}$. An ordinary prime $p \in \mathbb{Z}$ factors in $\mathbb{Z}[i]$ if and only if $p = a^2 + b^2$, which occurs precisely when $p = 2$ or $p \equiv 1 \pmod 4$ (Fermat's two-square theorem). Integer primes $p \equiv 3 \pmod 4$ remain prime in $\mathbb{Z}[i]$.
    \item \textbf{Rings and Elementary Arithmetic:} A ring $(R, +, \cdot, 0, 1)$ consists of an additive abelian group $(R, +, 0)$ and a multiplicative monoid $(R, \cdot, 1)$ coupled by the distributive laws. Distributivity yields the universal identities $0 \cdot a = 0$ and $(-1) \cdot a = -a$. When $1 = 0$, the ring collapses to the degenerate zero ring $R = \{0\}$.
    \item \textbf{Canonical Evaluation Map $\mathbb{Z} \to R$:} Defining integer multiples $n \cdot a$ yields the unique canonical ring homomorphism $\phi \colon \mathbb{Z} \to R$ with $\phi(n) = n \cdot 1_R$, establishing $\mathbb{Z}$ as the initial object in the category of unital rings.
    \item \textbf{Ring Homomorphisms:} A ring homomorphism preserves addition, multiplication, and the multiplicative identity ($f(1) = 1$). While $f(0) = 0$ is automatic via additive group cancellation, $f(1) = 1$ must be explicitly axiomatized because elements in arbitrary rings generally lack multiplicative inverses.
    \item \textbf{Ideals and Kernels:} An ideal $I \subset R$ is an additive subgroup that multiplicatively absorbs the entire ambient ring ($r \cdot a \in I$ for all $a \in I, r \in R$). The kernel $\ker(f) = f^{-1}(\{0\})$ of every ring homomorphism is an ideal.
    \item \textbf{Ideals in Fields:} Because every non-zero element in a field $F$ is invertible, any non-zero ideal contains $1_F$ and equals $F$. Thus, fields have only the trivial ideals $\{0\}$ and $F$, ensuring that every homomorphism from a field into a non-zero ring is injective.
    \item \textbf{Outlook on Zorn's Lemma:} The fundamental problem of whether every infinite-dimensional vector space possesses a basis relies on non-constructive set theory via Zorn's Lemma.
\end{itemize}
\end{overview}


% week2:
% ==========================================
% LatexRefinement Step Output: step4-2026-09-22-tuesday-week2-algebra-1-offset-final.tex
% Model: gemini-3.6-flash
% Temperature: 0.35
% TopP: 0.9
% TopK: 10
% MaxOutputTokens: 65535
% ThinkingBudget: 4096
% ThinkingLevel: HIGH
% Prompt Tokens: 70’067
% Candidates Tokens: 37’364 (inkl. Thinking Tokens)
% Cached Tokens: 0
% Processed on: 2026-10-03 21:06:35
% ==========================================

% Primary Language: English
% End of the video: 01:30:33

\lecturechapter[2]{Tuesday}{22 Sep}{22 September 2026}{Groups, Rings, Ideals, Subrings, and Introduction to Categories}

\begin{overview}[Course Content Overview]
This lecture reviews the fundamental algebraic structures introduced in Week 1---groups, rings, and fields---and explores their intrinsic properties, including the uniqueness of neutral elements and inverses. We formalize structure-preserving mappings (\newterm{group and ring homomorphisms/isomorphisms}), examine classical isomorphism classes over $\mathbb{R}$, and investigate substructures: \newterm{subgroups}, \newterm{ideals}, and \newterm{subrings} (including the classification of subrings and ideals of $\mathbb{Z}$, and ideals in smooth function spaces). Finally, we introduce the foundational language of \newterm{category theory}, defining categories, objects, morphisms, categorical isomorphisms, and \newterm{functors} with concrete algebraic examples.
\end{overview}

\section{Organizational Remarks}

\begin{speech}[00:00:00 - 00:01:48]
Okay, so we start, and there has been some progress on the website. There should be instructions there. Also, first of all, you should be able to register for the problem sessions now. 

The second thing is that there should be instructions on the website about how to upload your homework. 

And the third thing is about the quiz, and I think there's a link there, too. We'll see if it works on Friday. So, the rule for the quiz will be that you can take it anytime from $6$ in the morning on Friday to midnight. And if this is not enough freedom, we can also discuss it. It's not so...

So, another course hack is that if you look at this Metaphor page \inlinemetanote{writes \qt{Metaphor 2018} on the blackboard} in $2018$ --- it was the last time I taught this course ---, that Metaphor page is there. This course is going to follow roughly the same path. In fact, the topics and homework --- I mean, there's maybe some slight switching order of lectures ---, but you can find all the information there. In particular, all the problem sets that Lieke will eventually upload, you can already look at them now if you want to get ahead. 

And also, they're all the answers to all the problem sets. You can also look at them if you want, but in the modern world it's easy to find answers to these questions, so I don't see there's any reason to hide them from you. So, you can find everything there.

So, the other thing is that if you want to ask me a question, I'm happy to talk about math with anybody. You have to find me. Normally, I have an office, but I was kicked out because of this \qt{Fenstersanierung}. And it's going to last forever. So, I have a secret office on this floor. If you find me, you can talk to me. You just have to find it. It's actually pretty hard to find. But if you do, you can knock on the door. Sometimes I'm busy, sometimes I'm not, but I'm happy to talk to anyone. But you have to find the office yourself.
\end{speech}

\begin{note}[Organizational Information: Course Website and Metaphor Archive]
The lecturer outlines administrative details for the semester:
\begin{itemize}
    \item \textbf{Problem Sessions \& Homework:} Registration for exercise sessions is open on the course website, where instructions for online homework submission are posted.
    \item \textbf{Weekly Quizzes:} Friday quizzes can be completed asynchronously between 06:00 and 24:00 (midnight).
    \item \textbf{Course Archive (Metaphor 2018):} Lecture topics, problem sets, and solution sets from the 2018 iteration of the course are publicly accessible for advance study.
    \item \textbf{Office Hours:} Due to ongoing building renovations (\qt{Fenstersanierung}), the lecturer occupies a temporary office on the floor.
\end{itemize}
\end{note}

\section{Review of Groups, Rings, and Fields}

\subsection{Group Axioms and Uniqueness of Identities}

\begin{speech}[00:01:48 - 00:03:22]
Okay, so last week we discussed kind of carefully these three objects which are the centerpiece for this course. 

$G$ is a group \inlinemetanote{writes $(G, \cdot, e)$ on the blackboard}, and last week I used the star ($\star$) for the product. No one does that. That was purely for didactic purposes, and probably pointless didactic purposes. The standard for a group is dot ($\cdot$), but as I said, you can name this anything you want. 

So, this is the standard data of the group, and the group --- I don't repeat everything, but has the three axioms: associativity, existence of the neutral element, and existence of inverses.

So, now this is written now in slightly more abstract language than the explicit equations. And there's many things I didn't say, because they are all, in some sense, part of the definition. Like I didn't say explicitly... you could ask yourself: this neutral element, could there be some other element that also has the neutrality property? So, these are things you should think about yourself. I can't teach you everything; it's not possible. 

Could there be another $\hat{e} \in G$ that has the neutrality property? And the answer is no, because then you could look at $e \cdot \hat{e}$. If I use the fact that $e$ is the neutral element, this equals $\hat{e}$, and if $\hat{e}$ is a pretend neutral element, then this equals $e$. Okay? Is that clear?
\end{speech}

\begin{content}[Definition of a Group and Uniqueness of the Identity]
\begin{definition}[Group]\label[definition]{def:group}
A \newterm{group} is a triple $(G, \cdot, e)$, consisting of a non-empty set $G$, a binary operation $\cdot \colon G \times G \to G$, and a distinguished element $e \in G$ (called the \newterm{identity} or \newterm{neutral element}), satisfying the following three group axioms:
\begin{enumerate}[label=\textbf{(\arabic*)}]
    \item \textbf{Associativity:} For all $a, b, c \in G$,
    \[
    (a \cdot b) \cdot c = a \cdot (b \cdot c).
    \]
    \item \textbf{Existence of a Neutral Element:} There exists an element $e \in G$ such that for all $a \in G$,
    \[
    a \cdot e = e \cdot a = a.
    \]
    \item \textbf{Existence of Inverses:} For every $a \in G$, there exists an element $b \in G$ such that
    \[
    a \cdot b = b \cdot a = e.
    \]
\end{enumerate}
\end{definition}

\begin{proposition}[Uniqueness of the Neutral Element]\label[proposition]{prop:uniqueness-neutral-element}
Let $(G, \cdot)$ be a group. The neutral element $e \in G$ is unique.
\end{proposition}

\begin{proof}
Suppose $e, \hat{e} \in G$ are both neutral elements in $G$. Consider the product $e \cdot \hat{e}$:
\begin{itemize}
    \item Since $e$ is a neutral element, $e \cdot \hat{e} = \hat{e}$.
    \item Since $\hat{e}$ is a neutral element, $e \cdot \hat{e} = e$.
\end{itemize}
Equating both expressions yields:
\[
\hat{e} = e \cdot \hat{e} = e.
\]
Thus, the neutral element is unique.
\end{proof}
\end{content}

\subsection{Uniqueness of Inverses}

\begin{speech}[00:03:22 - 00:05:04]
And then other questions you could have: if $x$ is an element of the group, then this existence of inverse says that there exists a $y$ such that $x \cdot y$ equals the neutral element. And you could say: could there be another $\hat{y}$ where $x \cdot \hat{y}$ is equal to $e$? 

And the answer is no, again. No! And how do you prove this? Well, you look at this equation, and you multiply both sides by the actual --- the first one:
\[
y \cdot (x \cdot \hat{y}) = y \cdot e = y.
\]
And then this thing is equal by the original property --- so the neutral element property is both sides ---, so this thing is equal to $e$, so this is equal to $\hat{y}$. Okay?

Is that clear? So, all of this is embedded in some sense in the definitions. You could define a group by saying that there exists a neutral element and it's unique, there exist inverses and they're unique. You don't have to; you could do that. That's making the definition bigger. 

You can also try to make the definition smaller. There's some kind of strange direction where people think that, well, there's three axioms. We'd like a group to only have two axioms. And some guy in the '60s found some crazy way to write it with just one axiom, but it's just logically equivalent. You can go look that up if you want. I don't remember; it's some complicated axioms, you can imagine.
\end{speech}

\begin{content}[Uniqueness of Inverses in a Group]
\begin{proposition}[Uniqueness of Inverses]\label[proposition]{prop:uniqueness-inverse}
Let $(G, \cdot, e)$ be a group. For each element $x \in G$, the inverse element $y \in G$ satisfying $x \cdot y = y \cdot x = e$ is unique.
\end{proposition}

\begin{proof}
Let $x \in G$. Suppose $y, \hat{y} \in G$ are both inverses of $x$, so that $x \cdot y = y \cdot x = e$ and $x \cdot \hat{y} = \hat{y} \cdot x = e$.

Multiplying the equation $x \cdot \hat{y} = e$ on the left by $y$ gives:
\[
y \cdot (x \cdot \hat{y}) = y \cdot e = y.
\]
Using the associativity of the group operation on the left-hand side:
\[
(y \cdot x) \cdot \hat{y} = y.
\]
Since $y$ is an inverse of $x$, we have $y \cdot x = e$. Substituting this in:
\[
e \cdot \hat{y} = y \implies \hat{y} = y.
\]
Therefore, the inverse of $x$ is unique. We denote this unique inverse by $x^{-1}$.
\end{proof}

\begin{steps}[Why Redundant Axioms Are Omitted]
The standard definition of a group requires only the \emph{existence} of an identity and of inverses. As demonstrated above, their \emph{uniqueness} is not an independent axiom but an immediate logical consequence of associativity and the existence axioms.
\end{steps}
\end{content}

\subsection{Rings and Fields}

\begin{speech}[00:05:04 - 00:06:55]
Okay, so that's the group, and the other structure we had was this ring, which I don't repeat everything again. But the way to think about this is that if I look at just the additive part, this is an abelian group, and then I have some axioms for multiplication, and then two times the distributive laws. So, those were the axioms for the ring, and we did some examples.

So, certainly one part of this class is really internalizing these axioms. And there's some kind of idea about mathematics which is some kind of asymptotic understanding. Which is not just that you understand it, but like the really asymptotic understanding that you try to hope for, which you never really achieve, is that not only do you understand it, but you cannot imagine the state of not understanding it. That's the thing you should really aim for. And you know, you can only do that kind of asymptotically. You can never really achieve that, but with these... So, this is just at the axiomatic level.
\end{speech}

\begin{speech}[00:06:55 - 00:08:02]
And then there was a notion of a field, and the field is somehow... the group is by itself an algebraic object. The field is a particular case of a ring, meaning a field is a ring with some additional structure, and the additional structure is that non-zero elements have multiplicative inverses. 

So, you can think about this: in the whole world of rings, some of them are fields. But rings and groups are separate, although they are related by the fact that the additive part of the ring is an abelian group.

Okay, so those were the structures, and in algebra, after you have structures, there's the notion of the maps that respect the structures, which are the homomorphisms. And you've seen this before, and we discussed a little bit last week, but I want to put them on equal footing here.
\end{speech}

\begin{content}[Review of Rings and Fields]
\begin{definition}[Ring]\label[definition]{def:ring}
A \newterm{ring} is a $5$-tuple $(R, +, \cdot, 0, 1)$ consisting of a set $R$, two binary operations $+$ (addition) and $\cdot$ (multiplication), and two distinguished elements $0, 1 \in R$ ($0 \neq 1$) such that:
\begin{enumerate}[label=\textbf{(\arabic*)}]
    \item $(R, +, 0)$ is an \newterm{abelian group} (i.e., addition is associative, commutative, has identity $0$, and every element has an additive inverse $-a$).
    \item Multiplication $\cdot$ is associative with identity $1$:
    \[
    (a \cdot b) \cdot c = a \cdot (b \cdot c) \quad \text{and} \quad a \cdot 1 = 1 \cdot a = a \quad \forall a, b, c \in R.
    \]
    \item Multiplication distributes over addition from both the left and the right:
    \begin{align*}
    a \cdot (b + c) &= (a \cdot b) + (a \cdot c), \\
    (a + b) \cdot c &= (a \cdot c) + (b \cdot c) \quad \forall a, b, c \in R.
    \end{align*}
\end{enumerate}
\end{definition}

\begin{definition}[Field]\label[definition]{def:field}
A \newterm{field} is a commutative ring $(F, +, \cdot, 0, 1)$ in which every non-zero element possesses a multiplicative inverse:
\[
\forall x \in F \setminus \{0\}, \ \exists x^{-1} \in F \text{ such that } x \cdot x^{-1} = 1.
\]
Equivalently, $(F \setminus \{0\}, \cdot, 1)$ is an abelian group under multiplication.
\end{definition}
\end{content}

\begin{insight}[The Hierarchy of Algebraic Structures]
Algebra organizes mathematical systems by layering structural axioms:
\begin{itemize}
    \item \textbf{Groups $(G, \cdot, e)$:} A single operation with associativity, identity, and inverses.
    \item \textbf{Rings $(R, +, \cdot, 0, 1)$:} Two operations where $(R, +, 0)$ is an abelian group, multiplication is associative with $1$, and the operations are tied together via the distributive laws.
    \item \textbf{Fields $(F, +, \cdot, 0, 1)$:} Commutative rings where $(F \setminus \{0\}, \cdot, 1)$ is also an abelian group.
\end{itemize}
\end{insight}

\begin{aiNote}[Standing Convention: Commutativity of Rings]
As established in Week 1 (\cref{rem:ring-convention}), unless explicitly stated otherwise, the term \textbf{ring} in this course always refers to a \textbf{commutative ring with unity} ($a \cdot b = b \cdot a$ for all $a, b \in R$). Non-commutative rings (such as matrix algebras $M_n(K)$ or endomorphism rings) are treated as specialized exceptions.
\end{aiNote}

\section{Group Homomorphisms and Isomorphisms}

\subsection{Definition and Properties of Homomorphisms}

\begin{speech}[00:08:02 - 00:09:40]
So, if you have two groups, then we want to consider maps that respect the group structure. So, this is called $f$ is a group homomorphism if:
\[
f(a \cdot b) = f(a) \cdot f(b) \quad \text{for all } a, b \in G.
\]

Okay, so last week I did something complicated where we have $G$ is $(G, \star)$ and then $\hat{G}$ is $(\hat{G}, \hat{\star})$ to be completely clear about the products being different. I'm never going to do that again. No one ever does that.

So, here when you read this, this multiplication is the operation in the first one, and that's the operation in the second one. Is that clear? They're both just called dot ($\cdot$). And in fact, many people just don't write it at all. It'd be kind of normal to look in a book and there'd just be nothing, because that's the standard way you deal with algebra. You don't write $x \cdot y$, you just write $xy$.

Okay, so that's the definition of a group homomorphism. It respects the structure. It says you can multiply here and then go over by $f$, or you can go over separately and multiply. That's what's respecting the structure.
\end{speech}

\begin{speech}[00:09:40 - 00:10:40]
Then you could say: what about the neutral element? Why am I not saying anything about the neutral element? I could. I could add $f(e) = \hat{e}$, or if I want to be clear, $\hat{e}$, the $e$ on that side. 

I could add that, but I don't have to, because if I put $e$ here, if I just use this property, I put $f(e \cdot e)$ --- there's only one thing you can do always --- so I just stick $e$ in here. You agree?

And then I use the neutral property on this side: we get $f(e) = f(e) \cdot f(e)$.

And then I can use the cancellation property of the group: $f(e)$ has some inverse. There exists some $y \in \hat{G}$ such that $y \cdot f(e) = \hat{e}$, right? This is the cancellation property. We've used it many times, and at some point I will stop explaining it. 

So, we multiply by $y$ on both sides: we get $y \cdot f(e) = y \cdot (f(e) \cdot f(e))$. So, then this is $\hat{e} = (y \cdot f(e)) \cdot f(e) = \hat{e} \cdot f(e) = f(e)$. So, $f(e) = \hat{e}$. Okay?

So, I don't have to, but you could put it in there.
\end{speech}

\begin{content}[Group Homomorphisms and Preservation of Identities]
\begin{definition}[Group Homomorphism]\label[definition]{def:group-homomorphism}
Let $(G, \cdot, e)$ and $(\hat{G}, \cdot, \hat{e})$ be groups. A function $f \colon G \to \hat{G}$ is called a \newterm{group homomorphism} (or simply \newterm{homomorphism}) if it respects the group operation, that is:
\begin{equation}\label{eq:group-homomorphism}
f(a \cdot b) = f(a) \cdot f(b) \quad \forall a, b \in G.
\end{equation}
\end{definition}

\begin{proposition}[Preservation of the Neutral Element]\label[proposition]{prop:homomorphism-neutral-element}
Let $f \colon G \to \hat{G}$ be a group homomorphism. Then $f$ automatically maps the identity of $G$ to the identity of $\hat{G}$:
\begin{equation}\label{eq:hom-preserves-identity}
f(e) = \hat{e}.
\end{equation}
\end{proposition}

\begin{proof}
Since $e \cdot e = e$ in $G$, applying the homomorphism property yields:
\[
f(e) = f(e \cdot e) = f(e) \cdot f(e).
\]
Since $f(e) \in \hat{G}$ belongs to a group, it possesses an inverse $y \in \hat{G}$ such that $y \cdot f(e) = \hat{e}$. Multiplying both sides of the equation on the left by $y$:
\begin{align}
y \cdot f(e) &= y \cdot \bigl(f(e) \cdot f(e)\bigr) \nonumber \\
\hat{e} &= \bigl(y \cdot f(e)\bigr) \cdot f(e) && \text{(by Associativity)} \nonumber \\
\hat{e} &= \hat{e} \cdot f(e) \nonumber \\
\hat{e} &= f(e). \label{eq:f-e-equals-ehat}
\end{align}
Thus, $f(e) = \hat{e}$.
\end{proof}
\end{content}

\subsection{Group Isomorphisms}

\begin{speech}[00:10:40 - 00:11:47]
So, that's for the group. The group is the algebraic structure, and if you have two groups, you compare them with a homomorphism. It's a map that respects the structure.

And this is following the pattern of this kind of analysis in linear algebra also. There you have two vector spaces, and the map is a linear transformation. That's the notion of the map that respects the vector space structure on both sides.

And then there's some standard things to say: we say $f$ is a group isomorphism if $f$ is a group homomorphism and $f$ is bijective. Okay? So, this should not surprise you. You've seen these terms in linear algebra, right?

Which language did you learn linear algebra in?
\end{speech}

\begin{student}[Student Answer]
English.
\end{student}

\begin{speech}[00:11:49 - 00:13:00]
English, okay, well, that's no problem then.

And there's another way to say this: $f$ is a group isomorphism if and only if there exists an inverse, there exists a group homomorphism in the other direction. There exists a group homomorphism $g \colon \hat{G} \to G$ such that if I do $f$ first and then $g$, I get the identity on the first group, and if I do $g$ first and then $f$, I get the identity on the second group. These are equivalent. Okay?

So, if you have the bottom, it's certainly bijective, because you have the forward and backwards map. And if you have the top, since it's bijective there exists an inverse function, and you have to show that that's a homomorphism, which I leave to you to show. Okay? So, I hope this is clear.
\end{speech}

\begin{content}[Definition and Characterization of Group Isomorphisms]
\begin{definition}[Group Isomorphism]\label[definition]{def:group-isomorphism}
Let $G$ and $\hat{G}$ be groups. A map $f \colon G \to \hat{G}$ is a \newterm{group isomorphism} if:
\begin{enumerate}[label=\textbf{(\arabic*)}]
    \item $f$ is a group homomorphism, and
    \item $f$ is bijective.
\end{enumerate}
If such an isomorphism exists, we say that $G$ and $\hat{G}$ are \newterm{isomorphic} and write $G \cong \hat{G}$.
\end{definition}

\begin{proposition}[Alternative Invertibility Characterization]\label[proposition]{prop:isomorphism-inverse-characterization}
A map $f \colon G \to \hat{G}$ is a group isomorphism if and only if there exists a group homomorphism $g \colon \hat{G} \to G$ such that:
\begin{equation}\label{eq:isomorphism-inverses}
g \circ f = \operatorname{id}_G \quad \text{and} \quad f \circ g = \operatorname{id}_{\hat{G}}.
\end{equation}
\end{proposition}

\begin{proof}
\textbf{Direction ($\Leftarrow$):} If such a homomorphism $g$ exists, the two-sided inverse equations $g \circ f = \operatorname{id}_G$ and $f \circ g = \operatorname{id}_{\hat{G}}$ imply by standard set theory that $f$ is a bijection. Since $f$ is also given as a homomorphism, $f$ is an isomorphism.

\textbf{Direction ($\Rightarrow$):} If $f$ is an isomorphism, it is a bijective homomorphism, so a set-theoretic inverse function $g = f^{-1} \colon \hat{G} \to G$ exists satisfying \eqref{eq:isomorphism-inverses}. It remains only to verify that $g$ is a homomorphism. For any $\hat{a}, \hat{b} \in \hat{G}$, let $a = g(\hat{a})$ and $b = g(\hat{b})$, so that $f(a) = \hat{a}$ and $f(b) = \hat{b}$. Then:
\[
f(a \cdot b) = f(a) \cdot f(b) = \hat{a} \cdot \hat{b}.
\]
Applying $g = f^{-1}$ to both sides gives:
\[
g(\hat{a} \cdot \hat{b}) = a \cdot b = g(\hat{a}) \cdot g(\hat{b}).
\]
Thus, $g$ is a group homomorphism.
\end{proof}
\end{content}

\section{Examples of Groups and Isomorphisms over \texorpdfstring{$\mathbb{R}$}{R}}

\subsection{Three Classical Groups of Real Numbers}

\begin{speech}[00:13:00 - 00:15:26]
All right, so maybe an example. 

If I take the real numbers, one group is the real numbers with addition and zero: $(\mathbb{R}, +, 0)$. This is a rather basic group: real numbers, zero, you can add them.

Another group which you maybe don't think about as much is you could take the non-zero real numbers, multiplication, and one: $(\mathbb{R}^*, \cdot, 1)$. So, here I just remind you --- this is pretty standard ---, when you have a field and you put a star ($^*$), this usually means you take the field and you throw out the zero:
\[
\mathbb{R}^* := \mathbb{R} \setminus \{0\}.
\]
Is that clear? So, this is the non-zero real numbers, and I put multiplication, one. That is also a group. Are you happy with that? Because you can invert non-zero elements.

And there's another one here: there's kind of two ways to be non-zero: you can be positive, or you can be negative. So, I could take the positive ones: $(\mathbb{R}_{>0}^*, \cdot, 1)$. That's still a group. Is that clear?

If I try to then take the negative ones, I kind of get in trouble because if I multiply negative ones I get positive, and the one is not there. So, anyway, there's these three groups.

So, which of these are isomorphic? Or which ones are not isomorphic? Do you know this? Let me write the question: Which are isomorphic? Which means: for which ones can I find an isomorphism?
\end{speech}

\begin{content}[Three Standard Groups over $\mathbb{R}$]
We consider the following three algebraic groups defined on subsets of $\mathbb{R}$:
\begin{enumerate}[label=\textbf{(\arabic*)}]
    \item \textbf{The Additive Group of Real Numbers:}
    \[
    (\mathbb{R}, +, 0)
    \]
    where $+$ is standard addition, $0$ is the neutral element, and the inverse of $x$ is $-x$.
    \item \textbf{The Multiplicative Group of Non-Zero Real Numbers:}
    \[
    (\mathbb{R}^*, \cdot, 1) \quad \text{where } \mathbb{R}^* := \mathbb{R} \setminus \{0\},
    \]
    equipped with standard multiplication, identity $1$, and inverse $x^{-1} = \frac{1}{x}$.
    \item \textbf{The Multiplicative Group of Positive Real Numbers:}
    \[
    (\mathbb{R}_{>0}^*, \cdot, 1) \quad \text{where } \mathbb{R}_{>0}^* := \{x \in \mathbb{R} \mid x > 0\},
    \]
    which forms a group under multiplication since the product and inverse of positive numbers remain strictly positive.
\end{enumerate}
\end{content}

\subsection{Non-Isomorphism between \texorpdfstring{$(\mathbb{R}^*, \cdot, 1)$}{(R*, *, 1)} and \texorpdfstring{$(\mathbb{R}_{>0}^*, \cdot, 1)$}{(R>0*, *, 1)}}

\begin{speech}[00:15:26 - 00:16:21]
So, let's start with these --- this is maybe the more confusing one. If I take the non-zero real numbers, multiplication, $1$, and I take the positive ones, $(\mathbb{R}_{>0}^*, \cdot, 1)$, could these be isomorphic?

It's kind of strange, because if you draw the picture, one of them is kind of both sides, and the other one is just one side. So, it seems kind of unlikely, although, you know, there's still bijections that do this. There are set-theoretic bijections, because both sets are just uncountable sets; they have the same cardinality.

What do you think?

Yeah?
\end{speech}

\begin{student}[Student Answer]
No, because we run into a problem for $-1$, because if we had if $f(-1)$...
\end{student}

\begin{speech}[00:16:30 - 00:16:33]
Which way is $f$ going?
\end{speech}

\begin{student}[Student Answer continued]
So, $f$ goes from $\mathbb{R}^*$ to $\mathbb{R}_{>0}^*$.
\end{student}

\begin{speech}[00:16:37 - 00:16:39]
Yeah.
\end{speech}

\begin{student}[Student Answer continued]
And then $f(1)$ has to be $1$.
\end{student}

\begin{speech}[00:16:42 - 00:16:45]
Has to be $1$, that's the...
\end{speech}

\begin{student}[Student Answer continued]
Yeah, that's $1$, but that's also $f(-1)$ squared.
\end{student}

\begin{speech}[00:16:48 - 00:16:50]
Yeah, so $f(-1)$ equals what? That's a good...
\end{speech}

\begin{student}[Student Answer continued]
$f(-1) \cdot f(-1)$ will be $1$, so the only positive number such that times itself is $1$ will be $1$, so $f(-1)$ has to be $1$. So, $f(1) = f(-1)$, and then it's not bijective.
\end{student}

\begin{speech}[00:17:09 - 00:18:00]
Yeah, this is all correct. So, the simplest way to say it is that $f(-1)$ has to be a square root of $1$, okay, by this law. And there's only one square root of $1$ in the positive, and that's already taken by $1$. So, it violates the bijectivity.

Do you understand this statement? The point is that $1$ has two square roots in here. There are two elements when you multiply them by themselves you get $1$: $+1$ and $-1$. And those two elements have to go to two different elements here, because it's a bijection. If there were a bijection, they would have to go to two different elements. But here there's only one element whose square is $1$. So, it's impossible that there is a bijection. Okay? 

So, this is a solid argument. This is the standard argument.
\end{speech}

\begin{content}[Non-Isomorphism of \texorpdfstring{$(\mathbb{R}^*, \cdot, 1)$}{(R*, *, 1)} and \texorpdfstring{$(\mathbb{R}_{>0}^*, \cdot, 1)$}{(R>0*, *, 1)}]
\begin{proposition}\label[proposition]{prop:rstar-not-isom-rpos}
The groups $(\mathbb{R}^*, \cdot, 1)$ and $(\mathbb{R}_{>0}^*, \cdot, 1)$ are not isomorphic:
\[
(\mathbb{R}^*, \cdot, 1) \not\cong (\mathbb{R}_{>0}^*, \cdot, 1).
\]
\end{proposition}

\begin{proof}
Suppose for contradiction that there exists an isomorphism $f \colon \mathbb{R}^* \to \mathbb{R}_{>0}^*$.

By \cref{prop:homomorphism-neutral-element}, any homomorphism must preserve the identity:
\begin{equation}\label{eq:f-of-one}
f(1) = 1.
\end{equation}
Now consider the element $-1 \in \mathbb{R}^*$. Since $(-1) \cdot (-1) = 1$, the homomorphism property implies:
\[
f(-1)^2 = f(-1) \cdot f(-1) = f\bigl((-1) \cdot (-1)\bigr) = f(1) = 1.
\]
Thus, $f(-1)$ must be a real root of the equation $t^2 = 1$ in the codomain $\mathbb{R}_{>0}^*$.
In the set of positive real numbers $\mathbb{R}_{>0}^*$, the only solution to $t^2 = 1$ is $t = 1$. Therefore:
\begin{equation}\label{eq:f-of-minus-one}
f(-1) = 1.
\end{equation}
Comparing \eqref{eq:f-of-one} and \eqref{eq:f-of-minus-one}, we obtain $f(1) = f(-1) = 1$. Since $1 \neq -1$ in $\mathbb{R}^*$, this contradicts the injectivity of $f$. Consequently, no such isomorphism can exist.
\end{proof}

\begin{steps}[Algebraic Invariants: Torsion Elements]
This proof illustrates the power of \newterm{algebraic invariants}: an algebraic property preserved by all isomorphisms. In $(\mathbb{R}^*, \cdot, 1)$, the element $-1$ is an element of order $2$ (an element $x \neq e$ such that $x^2 = e$, called a \newterm{torsion element}). However, $(\mathbb{R}_{>0}^*, \cdot, 1)$ contains no elements of order $2$. Since an isomorphism must map elements of order $k$ to elements of order $k$, the two groups cannot be isomorphic despite having the same uncountable cardinality ($2^{\aleph_0}$) as sets.
\end{steps}
\end{content}

\subsection{Isomorphism between \texorpdfstring{$(\mathbb{R}, +, 0)$}{(R, +, 0)} and \texorpdfstring{$(\mathbb{R}_{>0}^*, \cdot, 1)$}{(R>0*, *, 1)}}

\begin{speech}[00:18:00 - 00:18:22]
And what about this one? $(\mathbb{R}, +, 0)$ and $(\mathbb{R}_{>0}^*, \cdot, 1)$?

If I make any mistakes on the board, you have to tell me, all right?

So, what about this? Could these be isomorphic?

Yeah?
\end{speech}

\begin{student}[Student Answer]
I think it could, we can take $f$ from $\mathbb{R}$ to $\mathbb{R}_{>0}^*$ by the function $e$ to the power of $x$.
\end{student}

\begin{speech}[00:18:32 - 00:19:30]
Yeah, yeah. This is like the most famous function in mathematics, and this is its purpose, actually. 

This is by $\exp$, $\exp(x)$, so $f(x) = \exp(x)$. So, you take a number, and you get a positive number. If you take $0$, you get $1$. And the main property of $\exp$ is that it intertwines the addition here with the multiplication there.

So, this is actually the most famous function in mathematics, basically, and this is its role. It gives you an isomorphism, a group isomorphism, between the additive structure and the multiplicative structure.

So, that's a kind of weird thing. It's saying that actually from the point of view of the group, the group doesn't see the difference between the additive structure and the multiplicative structure of the reals.
\end{speech}

\begin{content}[The Exponential Isomorphism]
\begin{proposition}\label[proposition]{prop:r-additive-isom-rpos-mult}
The additive group of real numbers $(\mathbb{R}, +, 0)$ is isomorphic to the multiplicative group of positive real numbers $(\mathbb{R}_{>0}^*, \cdot, 1)$:
\[
(\mathbb{R}, +, 0) \cong (\mathbb{R}_{>0}^*, \cdot, 1).
\]
\end{proposition}

\begin{proof}
Define the map $f \colon \mathbb{R} \to \mathbb{R}_{>0}^*$ by the exponential function:
\begin{equation}\label{eq:exp-map}
f(x) := \exp(x) = e^x.
\end{equation}
We verify the two isomorphism conditions:
\begin{enumerate}[label=\textbf{(\arabic*)}]
    \item \textbf{Homomorphism Property:} For all $x, y \in \mathbb{R}$, by the functional equation of the exponential function:
    \[
    f(x + y) = e^{x+y} = e^x \cdot e^y = f(x) \cdot f(y).
    \]
    Thus, $f$ transforms addition in $(\mathbb{R}, +)$ into multiplication in $(\mathbb{R}_{>0}^*, \cdot)$.
    \item \textbf{Bijectivity:} The exponential function $\exp \colon \mathbb{R} \to \mathbb{R}_{>0}^*$ is strictly monotonically increasing and continuous, with an explicit two-sided inverse given by the natural logarithm:
    \[
    g \colon \mathbb{R}_{>0}^* \to \mathbb{R}, \quad g(y) := \ln(y).
    \]
    For all $x \in \mathbb{R}$, $g(f(x)) = \ln(e^x) = x = \operatorname{id}_{\mathbb{R}}(x)$, and for all $y \in \mathbb{R}_{>0}^*$, $f(g(y)) = e^{\ln(y)} = y = \operatorname{id}_{\mathbb{R}_{>0}^*}(y)$.
\end{enumerate}
Therefore, $f$ is a group isomorphism.
\end{proof}
\end{content}

\begin{insight}[The Fundamental Role of the Exponential Function]
From an algebraic perspective, the exponential function $x \mapsto e^x$ is not merely an analytic power series or a solution to $y' = y$; its primary algebraic purpose is to establish an exact equivalence (\newterm{group isomorphism}) between the continuous additive world $(\mathbb{R}, +)$ and the continuous positive multiplicative world $(\mathbb{R}_{>0}^*, \cdot)$.
\end{insight}

\section{The Kernel and Subgroups}

\subsection{Definition and Properties of the Kernel}

\begin{speech}[00:19:30 - 00:21:10]
There's a notion of a subgroup. We did this for the ring; I'm going to repeat it all for the ring in one second. So, if you have this map, $f \colon G \to \hat{G}$ --- and I'm not going to keep saying group homomorphism. The standard is to say homomorphism, which is abbreviated $\operatorname{hom}$, and you're supposed to understand, since I'm talking about groups, it's a group homomorphism.

All right, so if $f$ is a group homomorphism, then from the spirit of linear algebra, which is the algebra you're supposed to be most familiar with, if you have a linear transformation between two vector spaces, the thing you immediately study is the kernel and the image of that linear transformation. So, we can do that also with groups.

So, then we can consider the kernel and the image, but let's start with the kernel. So, what properties does this kernel enjoy?

The kernel of a group homomorphism is defined to be those elements in $G$ such that when I apply $f$, I get the neutral element on the second side. In the linear algebra language, it's things that went to zero, so now it's the elements that go to the neutral element. Is that clear?

Well, the neutral element itself is in here. That's one property, because the homomorphism takes the neutral element here to the neutral element there. So, that's there, right?
\end{speech}

\begin{speech}[00:21:10 - 00:22:15]
Second property: suppose $a$ and $b$ are in the kernel of $f$, then what can we say? We can consider what happens to $f(ab)$. 

There's only one thing we can do here. These proofs at this level are strong proofs, meaning that there's only one thing you can do at every stage, and you just have to do it. 

So, this is $f(a) \cdot f(b)$, and now the assumption is they're both in the kernel, so these are both $\hat{e} \cdot \hat{e}$, and then we use the identity property, so that's $\hat{e}$. 

Okay, so what have we learned here? That says that if $a$ and $b$ are in the kernel, we have proven that $a \cdot b$ is in the kernel. Okay?
\end{speech}

\begin{content}[Definition and Closure Properties of the Kernel]
\begin{definition}[Kernel of a Group Homomorphism]\label[definition]{def:kernel}
Let $f \colon G \to \hat{G}$ be a group homomorphism between groups $(G, \cdot, e)$ and $(\hat{G}, \cdot, \hat{e})$. The \newterm{kernel} of $f$, denoted $\ker(f)$, is the pre-image of the identity element $\hat{e} \in \hat{G}$:
\begin{equation}\label{eq:kernel-def}
\ker(f) := \bigl\{ x \in G \;\big|\; f(x) = \hat{e} \bigr\} \subseteq G.
\end{equation}
\end{definition}

\begin{proposition}[Properties of the Kernel: Identity and Product Closure]\label[proposition]{prop:kernel-properties-part1}
Let $f \colon G \to \hat{G}$ be a group homomorphism. Then $\ker(f)$ satisfies:
\begin{enumerate}[label=\textbf{(\arabic*)}]
    \item \textbf{Contains the Identity:} $e \in \ker(f)$.
    \item \textbf{Closed under the Group Product:} If $a, b \in \ker(f)$, then $a \cdot b \in \ker(f)$.
\end{enumerate}
\end{proposition}

\begin{proof}
\textbf{Property (1):} By \cref{prop:homomorphism-neutral-element}, $f(e) = \hat{e}$. Thus, $e \in \ker(f)$ by definition.

\textbf{Property (2):} Let $a, b \in \ker(f)$, which means $f(a) = \hat{e}$ and $f(b) = \hat{e}$. Using the homomorphism property of $f$:
\[
f(a \cdot b) = f(a) \cdot f(b) = \hat{e} \cdot \hat{e} = \hat{e}.
\]
Therefore, $a \cdot b \in \ker(f)$.
\end{proof}
\end{content}

\subsection{Closure under Inverses and the Definition of Subgroups}

\begin{speech}[00:22:15 - 00:24:05]
So, we're looking at properties of the kernel, and we have discovered that the neutral element's in the kernel, and the kernel is closed under the group law, the product.

We might be curious now whether the inverses are there also.

So, that's the third property of the kernel: if $x$ is in the kernel of $f$, then the inverse... 

I should make this comment that very often the inverse for the group law is denoted $x^{-1}$. That's just the name of it. The axioms say there exists an inverse. This is very often the name that's given to it. This kind of presupposes that the product you're calling dot ($\cdot$). If the product you're calling is plus ($+$), then you usually use the minus sign ($-x$) here. But this is just notation. You could call it anything you want. You could call it Fred, or an elephant, anything. We're calling it this, okay?

Then that's the inverse is also in the kernel. Okay?
\end{speech}

\begin{speech}[00:24:05 - 00:26:00]
And how are you going to prove this? You have to start out by taking $f(x \cdot x^{-1})$. That's the one thing you know. 

On the one hand, this is $f(e)$, and that's $\hat{e}$. On the other hand, this is $f(x) \cdot f(x^{-1})$, and this thing is assumed to be in the kernel, so this is equal to $\hat{e} \cdot f(x^{-1})$.

So, we get that the neutral element is equal to neutral element times that ($\hat{e} = \hat{e} \cdot f(x^{-1})$), so that must be the neutral element. Okay? We could go one more if we want to be really explicit:
\[
f(x^{-1}) = \hat{e} \cdot f(x^{-1}) = \hat{e}.
\]
So, now we have the full chain: $f(x^{-1}) = \hat{e}$.

So, we have these three properties: that the kernel enjoys these three properties. It contains the neutral element of $G$, the source group, it's closed under the group product, and it's closed under inverses.

Definition: if $H$ is a subset of $G$, then $H$ is a subgroup if $e \in H$, $H$ is closed under the product, and $H$ is closed under inverses. Meaning if anybody's in $H$, then its multiplicative inverse is also in $H$, okay?
\end{speech}

\begin{content}[Closure of the Kernel under Inverses and the Subgroup Axioms]
\begin{proposition}[Closure of the Kernel under Inverses]\label[proposition]{prop:kernel-closure-inverses}
Let $f \colon G \to \hat{G}$ be a group homomorphism. If $x \in \ker(f)$, then $x^{-1} \in \ker(f)$.
\end{proposition}

\begin{proof}
Let $x \in \ker(f)$, so that $f(x) = \hat{e}$. Consider the image of the identity element $e = x \cdot x^{-1}$:
\begin{align*}
\hat{e} &= f(e) && \text{(by \cref{prop:homomorphism-neutral-element})} \\
&= f(x \cdot x^{-1}) \\
&= f(x) \cdot f(x^{-1}) && \text{(homomorphism property)} \\
&= \hat{e} \cdot f(x^{-1}) && \text{(since } x \in \ker(f)\text{)} \\
&= f(x^{-1}) && \text{(neutral element property in } \hat{G}\text{)}.
\end{align*}
Thus, $f(x^{-1}) = \hat{e}$, which demonstrates that $x^{-1} \in \ker(f)$.
\end{proof}

\begin{definition}[Subgroup]\label[definition]{def:subgroup}
Let $(G, \cdot, e)$ be a group, and let $H \subseteq G$ be a non-empty subset. $H$ is called a \newterm{subgroup} of $G$ (denoted $H \le G$) if it satisfies the following three conditions:
\begin{enumerate}[label=\textbf{(\arabic*)}]
    \item \textbf{Identity Element:} $e \in H$.
    \item \textbf{Closure under Multiplication:} For all $a, b \in H$, $a \cdot b \in H$.
    \item \textbf{Closure under Inverses:} For all $x \in H$, $x^{-1} \in H$.
\end{enumerate}
\end{definition}

\begin{steps}[Inherited Group Structure]
If $H \subseteq G$ satisfies these three axioms, the restriction of the group operation $\cdot|_{H \times H}$ defines a binary operation on $H$, associativity is inherited from $G$, and $e \in H$ acts as the identity with inverses in $H$. Consequently, $(H, \cdot, e)$ is itself a group in its own right.
\end{steps}
\end{content}

\subsection{The Kernel and the Image as Subgroups}

\begin{speech}[00:26:00 - 00:26:50]
So, in particular, this implies that if I take $H$ and I take the group law, which makes sense on $H$ because $H$ is closed under it, and I take the zero, which also makes sense because zero's in $H$, that this is a group.

Which is a very nice thing. A subgroup is a group. It'd be terrible if it weren't, right?

So, that's the definition for a subgroup. And now we've proven a little proposition: The kernel of $f$ is a subgroup of $G$.

Is that clear?
\end{speech}

\begin{content}[Subgroup Property of the Kernel]
\begin{proposition}[Kernel as a Subgroup]\label[proposition]{prop:kernel-is-subgroup}
Let $f \colon G \to \hat{G}$ be a group homomorphism. Then the kernel $\ker(f)$ is a subgroup of the domain group $G$:
\begin{equation}\label{eq:kernel-subgroup}
\ker(f) \le G.
\end{equation}
\end{proposition}

\begin{proof}
By \cref{prop:homomorphism-neutral-element}, $e \in \ker(f)$. By \cref{prop:kernel-properties-part1}, $\ker(f)$ is closed under multiplication. By \cref{prop:kernel-closure-inverses}, $\ker(f)$ is closed under taking inverses. Hence, by \cref{def:subgroup}, $\ker(f) \le G$.
\end{proof}
\end{content}

\subsection{The Image of a Group Homomorphism}

\begin{speech}[00:26:50 - 00:28:30]
Okay, so that's the investigation of the kernel and the notion of a subgroup. But there's also the image. And so proposition --- and I leave this for you, you can prove it ---: The image of $f$ is a subgroup of $\hat{G}$.

So, how do you prove this? You just check that the image contains the neutral element... it has to because the neutral element from the first group goes to the neutral element of the second group, so the image has to contain the neutral element. And then you have to check that if we have two elements in the image, then if I apply the group law for $\hat{G}$, they stay in the image, and the inverse. The proofs are all just the same. You just follow your nose. There's only one thing you can do, and you can do it.

All right. So, this is a nice situation for groups, that the understanding of the group is by these very simple axioms. There's lots of examples; in fact, part of this class is the examples for the groups, which will come later. And whenever you have a homomorphism, there's the kernel and the image.

So, this one I leave to you. You can look in the books too, but it's not worth looking in the books. You should read the books, I take it back. But the problem with reading such a simple thing in the book is that you have to learn what notation the author is using and all that. It's much simpler just to write the proof of this yourself.
\end{speech}

\begin{content}[The Image of a Group Homomorphism]
\begin{definition}[Image of a Homomorphism]\label[definition]{def:group-image}
Let $f \colon G \to \hat{G}$ be a group homomorphism. The \newterm{image} of $f$, denoted $\operatorname{Im}(f)$ (or $f(G)$), is the set:
\begin{equation}\label{eq:image-def}
\operatorname{Im}(f) := \bigl\{ \hat{y} \in \hat{G} \;\big|\; \exists x \in G \text{ such that } f(x) = \hat{y} \bigr\} \subseteq \hat{G}.
\end{equation}
\end{definition}

\begin{proposition}[Image as a Subgroup]\label[proposition]{prop:image-is-subgroup}
Let $f \colon G \to \hat{G}$ be a group homomorphism. Then the image $\operatorname{Im}(f)$ is a subgroup of the codomain group $\hat{G}$:
\begin{equation}\label{eq:image-subgroup}
\operatorname{Im}(f) \le \hat{G}.
\end{equation}
\end{proposition}

\begin{proof}
We verify the three subgroup conditions for $\operatorname{Im}(f) \subseteq \hat{G}$:
\begin{enumerate}[label=\textbf{(\arabic*)}]
    \item \textbf{Identity Element:} Since $f(e) = \hat{e}$ by \cref{prop:homomorphism-neutral-element}, $\hat{e} \in \operatorname{Im}(f)$.
    \item \textbf{Closure under Multiplication:} Let $\hat{y}_1, \hat{y}_2 \in \operatorname{Im}(f)$. Then there exist $x_1, x_2 \in G$ such that $f(x_1) = \hat{y}_1$ and $f(x_2) = \hat{y}_2$. Using the homomorphism property:
    \[
    \hat{y}_1 \cdot \hat{y}_2 = f(x_1) \cdot f(x_2) = f(x_1 \cdot x_2).
    \]
    Since $x_1 \cdot x_2 \in G$, we have $\hat{y}_1 \cdot \hat{y}_2 \in \operatorname{Im}(f)$.
    \item \textbf{Closure under Inverses:} Let $\hat{y} \in \operatorname{Im}(f)$ with $\hat{y} = f(x)$ for some $x \in G$. Since $x^{-1} \in G$, we compute:
    \[
    \hat{y} \cdot f(x^{-1}) = f(x) \cdot f(x^{-1}) = f(x \cdot x^{-1}) = f(e) = \hat{e}.
    \]
    By uniqueness of inverses in $\hat{G}$, $\hat{y}^{-1} = f(x^{-1}) \in \operatorname{Im}(f)$.
\end{enumerate}
Thus, $\operatorname{Im}(f) \le \hat{G}$.
\end{proof}
\end{content}

\begin{speech}[00:28:30 - 00:29:39]
But actually, I'm a big fan of books. It's one way to understand this kind of classical history of mathematics: read a lot of math books. And I know that somehow in the modern world, this is less popular, but in my generation we read a lot of math books.

And on the topic of books, there's the book for the course, which was picked a while ago, not by me even, and it has the advantage that it's freely available on the ETH website, so you can just click it and... Then there's some other references I gave which are, I think, a little bit better books. They're more advanced in perspective, they're written better.

But I think they're not freely available on the webpage, so it wasn't chosen for this class. But it's a good idea to read books. You'll learn something.
\end{speech}

\begin{note}[Pedagogical Aside: Recommended Literature]
The lecturer emphasizes the pedagogical value of reading classical mathematical texts and consulting textbook references, noting the freely available course textbook on the ETH website alongside advanced supplementary reading.
\end{note}

\section{Ring Homomorphisms, Ideals, and Subrings}

\subsection{Ring Homomorphisms and Isomorphisms}

\begin{speech}[00:29:39 - 00:31:49]
Okay, so for the ring, we had also defined what a homomorphism is: it's a ring homomorphism. If I have two rings, then we had defined $f$ is a ring hom, which, as I said, I will just say hom and you're supposed to understand they're rings. And then the axioms for a ring hom, if you remember, the first one is:
\[
f(a+b) = f(a) + f(b).
\]
And the interpretation of this: it says that the additive part of this is a group and the additive part of that is a group, and the ring hom should be a group hom for those parts. Is that clear? So, additively $f$ is a group hom for the plus part, okay?

Of course, we also have this, which respects the multiplication: $f(ab) = f(a)f(b)$.

And then we also have that the identity has to be respected: $f(1) = 1$.

Okay? So, those are the notions of the ring homomorphism. And again, a ring isomorphism is a bijective ring homomorphism.

And another way to think about it is if there's an isomorphism, then that means that the inverse is a ring homomorphism in the other direction, just as I said with groups.
\end{speech}

\begin{content}[Review of Ring Homomorphisms and Isomorphisms]
\begin{definition}[Ring Homomorphism]\label[definition]{def:ring-homomorphism}
Let $(R, +, \cdot, 0, 1)$ and $(\hat{R}, +, \cdot, \hat{0}, \hat{1})$ be rings. A map $f \colon R \to \hat{R}$ is a \newterm{ring homomorphism} if for all $a, b \in R$:
\begin{enumerate}[label=\textbf{(\arabic*)}]
    \item \textbf{Additive Homomorphism:} $f(a + b) = f(a) + f(b)$ (i.e., $f \colon (R, +, 0) \to (\hat{R}, +, \hat{0})$ is a group homomorphism on the additive group).
    \item \textbf{Multiplicative Homomorphism:} $f(a \cdot b) = f(a) \cdot f(b)$.
    \item \textbf{Preservation of the Multiplicative Identity:} $f(1) = \hat{1}$.
\end{enumerate}
\end{definition}

\begin{definition}[Ring Isomorphism]\label[definition]{def:ring-isomorphism}
A ring homomorphism $f \colon R \to \hat{R}$ is called a \newterm{ring isomorphism} if $f$ is bijective. Equivalently, $f$ is an isomorphism if and only if there exists a ring homomorphism $g \colon \hat{R} \to R$ such that $g \circ f = \operatorname{id}_R$ and $f \circ g = \operatorname{id}_{\hat{R}}$.
\end{definition}
\end{content}

\subsection{The Kernel of a Ring Homomorphism and Ideals}

\begin{speech}[00:31:49 - 00:34:03]
Okay, so then we can do the parallel study: if $f \colon R \to \hat{R}$ is a ring homomorphism, there are two things I can study as before: the kernel of $f$ and the image of $f$. And we had studied the kernel last week: for the group, both the kernel and image were subgroups: the former was the subgroup of the source group and the latter was the subgroup of the target group. But for rings, the symmetry is broken, so one has to concentrate a bit here.

And that's also true for vector spaces. For vector spaces, you have a linear map, the kernel is the subspace of the first one, and the image is the subspace of the second one. For the ring, something different happens.

And for the kernel of $f$, we had proven last week that the kernel of $f$ is something called an ideal. An ideal $I$ of a ring satisfies the following condition: one condition on the additive structure is that $(I, +, 0)$ inside $(R, +, 0)$ is a subgroup. That's the additive structure.
\end{speech}

\begin{speech}[00:34:03 - 00:35:23]
So for an ideal, there's kind of two pieces. One is how it works additively, and it must be a subgroup, and it's not a surprise because we already said that $f$ on the additive structures is a group homomorphism, and we already know that the kernel from the point of view of the group homomorphism is an additive subgroup.

And then the second has to do with the multiplicative part. The kernel satisfies this property that if I take any element of the ring $r \in R$ and any element of the ideal $x \in I$, then $r \cdot x$ is in the ideal. So, these are the two conditions that define an ideal. And the proposition last week was: the kernel is an ideal. All right?

So, what about the image? Is the image an ideal? What is the image? The answer is no. Ideal is the notion for the kernel. The image is something else: the image is a subring.
\end{speech}

\begin{content}[Ideals in Rings and the Kernel as an Ideal]
\begin{definition}[Ideal]\label[definition]{def:ideal}
Let $(R, +, \cdot, 0, 1)$ be a commutative ring. A subset $I \subseteq R$ is called an \newterm{ideal} (denoted $I \trianglelefteq R$) if it satisfies:
\begin{enumerate}[label=\textbf{(\arabic*)}]
    \item \textbf{Additive Subgroup:} $(I, +, 0)$ is a subgroup of $(R, +, 0)$, i.e.,
    \[
    0 \in I, \quad \forall x, y \in I \implies x + y \in I, \quad \text{and} \quad \forall x \in I \implies -x \in I.
    \]
    \item \textbf{Multiplicative Absorption (Ideal Property):} For every element of the ring $r \in R$ and every element of the ideal $x \in I$, the product belongs to $I$:
    \[
    \forall r \in R, \ \forall x \in I \implies r \cdot x \in I.
    \]
\end{enumerate}
\end{definition}

\begin{aiNote}[Ideals and Commutativity]
Because $R$ is assumed commutative, left ideals ($r \cdot x \in I$), right ideals ($x \cdot r \in I$), and two-sided ideals ($r \cdot x \cdot s \in I$) coincide. In non-commutative ring theory, these concepts diverge, and kernels of ring homomorphisms are always two-sided ideals.
\end{aiNote}

\begin{proposition}[The Kernel of a Ring Homomorphism is an Ideal]\label[proposition]{prop:kernel-ring-hom-ideal}
Let $f \colon R \to \hat{R}$ be a ring homomorphism. Then the kernel $\ker(f) := \{x \in R \mid f(x) = \hat{0}\}$ is an ideal of $R$.
\end{proposition}

\begin{proof}
We verify the two ideal conditions for $\ker(f)$:
\begin{enumerate}[label=\textbf{(\arabic*)}]
    \item \textbf{Additive Subgroup:} Since $f \colon (R, +, 0) \to (\hat{R}, +, \hat{0})$ is a group homomorphism between abelian groups, $\ker(f)$ is an additive subgroup of $(R, +, 0)$ by \cref{prop:kernel-is-subgroup}.
    \item \textbf{Multiplicative Absorption:} Let $r \in R$ and $x \in \ker(f)$, so $f(x) = \hat{0}$. Using the multiplicative homomorphism property:
    \[
    f(r \cdot x) = f(r) \cdot f(x) = f(r) \cdot \hat{0} = \hat{0}.
    \]
    Thus, $r \cdot x \in \ker(f)$.
\end{enumerate}
Therefore, $\ker(f) \trianglelefteq R$.
\end{proof}
\end{content}

\begin{insight}[Broken Symmetry: Subrings vs. Ideals]
In group theory and linear algebra, substructures are uniform: the kernel and image of a homomorphism are both subgroups (or vector subspaces). In ring theory, this symmetry is broken:
\begin{itemize}
    \item The \textbf{kernel} $\ker(f)$ is an \textbf{ideal} (it absorbs multiplication from the entire ring, but does not contain $1$ unless $\hat{R} = \{0\}$).
    \item The \textbf{image} $\operatorname{Im}(f)$ is a \textbf{subring} (it contains the identity $1$, but does not absorb multiplication from outside the image).
\end{itemize}
\end{insight}

\subsection{Subrings and the Image of a Ring Homomorphism}

\begin{speech}[00:36:04 - 00:37:50]
So, what is a subring? Let $R$ be a ring. Whenever I say let $R$ be a ring, you should be thinking that I have in mind all the data: $(R, +, \cdot, 0, 1)$.

A subring $S$ of $R$ satisfies the following properties --- this is the definition. The first property is that $0$ and $1$ should be in the subring. The second is that $S$ is closed under the two operations: whenever I have elements of the subring, I can either add them or multiply them, and they stay in the subring, okay? And the third is that $(S, +, \cdot, 0, 1)$ is a ring. A subring inside here has $0$ and $1$, it's closed under the operations, and together it's a ring. This means it has the additive inverses and things like that, okay?
\end{speech}

\begin{speech}[00:37:50 - 00:40:21]
And so the proposition here is: if I have a ring homomorphism, then the image of $f$ in the second one is a subring.

The kernel is an ideal, the image is a subring. You have to check that the image satisfies these properties. Does the image have $0$? The answer is yes, because on the additive structure, it's a homomorphism of groups, and $0$ goes to $0$. Does the image have $1$? Well, that's part of the definition of a ring homomorphism, so it has $1$. Then you have to check whether it's closed under $+$ and $\cdot$, and you just check that by writing out the equations. And whether this is a ring, you have to check whether it has the additive inverses, but that's also clear.

Okay. So, it's good to have these ideas extremely clear in your head: you have the group, you have the ring, there's a notion of group hom, there's a notion of ring hom, and there's kernel and image in both cases. For the group, the kernel and the image are subgroups; for the ring, something different happens: the image is a subring, but the kernel is an ideal. And so now we should do some examples.
\end{speech}

\begin{content}[Definition of a Subring and the Image of a Ring Homomorphism]
\begin{definition}[Subring]\label[definition]{def:subring}
Let $(R, +, \cdot, 0, 1)$ be a commutative ring with unity. A subset $S \subseteq R$ is called a \newterm{subring} of $R$ (denoted $S \le R$) if:
\begin{enumerate}[label=\textbf{(\roman*)}]
    \item \textbf{Contains Neutral Elements:} $0 \in S$ and $1 \in S$.
    \item \textbf{Closure under Addition and Multiplication:} For all $a, b \in S$,
    \[
    a + b \in S \quad \text{and} \quad a \cdot b \in S.
    \]
    \item \textbf{Ring Axioms (Closure under Additive Inverses):} $(S, +, \cdot, 0, 1)$ is itself a commutative ring with the induced operations (in particular, $-a \in S$ for all $a \in S$).
\end{enumerate}
\end{definition}

\begin{proposition}[The Image of a Ring Homomorphism is a Subring]\label[proposition]{prop:image-ring-hom-subring}
Let $f \colon R \to \hat{R}$ be a ring homomorphism. Then the image $\operatorname{Im}(f) \subseteq \hat{R}$ is a subring of $\hat{R}$.
\end{proposition}

\begin{proof}
We check the subring conditions for $\operatorname{Im}(f)$:
\begin{enumerate}[label=\textbf{(\arabic*)}]
    \item \textbf{Neutral Elements:} Since $f \colon (R, +, 0) \to (\hat{R}, +, \hat{0})$ is a group homomorphism, $f(0) = \hat{0} \in \operatorname{Im}(f)$. By axiom \textbf{(3)} of ring homomorphisms, $f(1) = \hat{1} \in \operatorname{Im}(f)$.
    \item \textbf{Closure under Addition and Multiplication:} Let $\hat{a}, \hat{b} \in \operatorname{Im}(f)$ with $\hat{a} = f(a)$ and $\hat{b} = f(b)$ for $a, b \in R$. Then:
    \[
    \hat{a} + \hat{b} = f(a) + f(b) = f(a + b) \in \operatorname{Im}(f),
    \]
    \[
    \hat{a} \cdot \hat{b} = f(a) \cdot f(b) = f(a \cdot b) \in \operatorname{Im}(f).
    \]
    \item \textbf{Additive Inverses:} For $\hat{a} = f(a) \in \operatorname{Im}(f)$, $-\hat{a} = -f(a) = f(-a) \in \operatorname{Im}(f)$.
\end{enumerate}
Thus, $\operatorname{Im}(f)$ is a subring of $\hat{R}$.
\end{proof}
\end{content}

\section{Examples of Ideals and Subrings}

\subsection{The Even Integers \texorpdfstring{$2\mathbb{Z}$}{2Z} in \texorpdfstring{$\mathbb{Z}$}{Z}}

\begin{speech}[00:40:21 - 00:40:57]
Okay, let's try here: $2\mathbb{Z}$ inside $\mathbb{Z}$. So, this is a ring, $\mathbb{Z}$. I can take the even integers. 

So, what is this inside that? Is it an ideal? Is it a subring? 

What's the answer?
\end{speech}

\begin{student}[Student Answer]
I think it is an ideal.
\end{student}

\begin{speech}[00:41:00 - 00:41:35]
It's an ideal. So, claim: this is an ideal. 

Why is it an ideal? We have to check that additively it's a subgroup: if you add two even numbers you get an even number, $0$ is an even number, if you have an even number the minus of it is an even number, so you're totally solid on the subgroup. Then if you take any other element of the integers and multiply it by an even number, you stay even. So, that looks pretty good. Answer is yes. 

Okay, what about a subring? Is it a subring?
\end{speech}

\begin{student}[Student Answer]
No, because we don't have the $1$.
\end{student}

\begin{speech}[00:41:38 - 00:41:48]
Yeah, so it's not a subring. Ideal: yes. Subring: no. 

Okay, so let's try something else.
\end{speech}

\begin{content}[The Even Integers: Ideal vs. Subring]
\begin{example}[The Even Integers $2\mathbb{Z} \subset \mathbb{Z}$]\label[example]{ex:even-integers-ideal}
Let $R = \mathbb{Z}$ and consider the subset of even integers $2\mathbb{Z} := \{ 2k \mid k \in \mathbb{Z} \}$.
\begin{itemize}
    \item \textbf{$2\mathbb{Z}$ is an Ideal of $\mathbb{Z}$ (\textcolor{ForestGreen}{YES}):}
    \begin{enumerate}[label=\textbf{(\arabic*)}]
        \item $(2\mathbb{Z}, +, 0)$ is an additive subgroup of $\mathbb{Z}$: $0 = 2 \cdot 0 \in 2\mathbb{Z}$, the sum of two even integers $2k + 2\ell = 2(k + \ell)$ is even, and $-(2k) = 2(-k) \in 2\mathbb{Z}$.
        \item Multiplicative absorption holds: For any $r \in \mathbb{Z}$ and any $x = 2k \in 2\mathbb{Z}$,
        \[
        r \cdot x = r \cdot (2k) = 2(r \cdot k) \in 2\mathbb{Z}.
        \]
    \end{enumerate}
    \item \textbf{$2\mathbb{Z}$ is NOT a Subring of $\mathbb{Z}$ (\textcolor{BrickRed}{NO}):}
    The multiplicative identity $1 \notin 2\mathbb{Z}$ since $1$ is odd.
\end{itemize}
\end{example}
\end{content}

\subsection{Subrings between \texorpdfstring{$\mathbb{Z}$}{Z} and \texorpdfstring{$\mathbb{Q}$}{Q}}

\begin{speech}[00:41:48 - 00:42:37]
Okay, this is maybe a trickier question. 

Certainly $\mathbb{Z}$ is a subring of $\mathbb{Q}$. I hope that's not controversial. You accept that $\mathbb{Z}$ is a subring of $\mathbb{Q}$? 

Can you find a subring that's in between $\mathbb{Z}$ and $\mathbb{Q}$? Does there exist a subring in between $\mathbb{Z} \subsetneq S \subsetneq \mathbb{Q}$? 

How about that? Can you do that? You have $3$ minutes.
\end{speech}

\begin{student}[Student Answer]
So we look at pairs of integers and denominator is just powers of two?
\end{student}

\begin{speech}[00:42:44 - 00:44:10]
Okay, that's good. So, he is saying that we let $S$ be the set of rational numbers of this form $\frac{a}{2^n}$, where $a$ is an integer and $n$ is also an integer. Is this exactly what you're saying? 

Okay, so that's a proposal. One thing that's good about this proposal is this set $S$ is certainly more than $\mathbb{Z}$ because it has like $3/2$ in it, and it's certainly less than $\mathbb{Q}$ because it doesn't have $2/3$ in it. So, that part is solid. 

Then the claim is you have to check that this satisfies all these properties. What are the properties about having $0$ and $1$? Basically the answer is yes, it's going to satisfy, because you need to know how to add fractions, which I don't teach you. I've never taught anyone how to add fractions; I'm not going to teach you. But when you think about the rule about adding fractions, you keep the $2^n$ in the bottom, that's the main thing here. When you add fractions, the denominator doesn't give you things that you didn't put in. 

And so the answer is yes. This is true. 

So, that's an example. That's a very important one. Of course, $2$ is not important here; you could have put here $17$, it'd be fine. That's another example. And in fact, being prime is not that important either.
\end{speech}

\begin{content}[Subrings of the Rational Numbers]
\begin{example}[{Dyadic Rationals $\mathbb{Z}[\frac{1}{2}] \subset \mathbb{Q}$}]\label[example]{ex:dyadic-rationals-subring}
Consider the problem of finding an intermediate subring strictly between $\mathbb{Z}$ and $\mathbb{Q}$:
\[
\mathbb{Z} \subsetneq S \subsetneq \mathbb{Q}.
\]
Define the set of \newterm{dyadic rationals} $S = \mathbb{Z}[\frac{1}{2}]$ by:
\begin{equation}\label{eq:dyadic-rationals-def}
S := \left\{ \frac{a}{2^n} \;\middle|\; a \in \mathbb{Z}, \ n \in \mathbb{N}_{\ge 0} \right\} \subset \mathbb{Q}.
\end{equation}
\begin{itemize}
    \item \textbf{Proper Inclusions:} $S \supsetneq \mathbb{Z}$ since $\frac{1}{2} = \frac{1}{2^1} \in S \setminus \mathbb{Z}$, and $S \subsetneq \mathbb{Q}$ since $\frac{1}{3} \notin S$ (powers of $2$ cannot produce a denominator divisible by $3$).
    \item \textbf{Subring Verification:}
    \begin{enumerate}[label=\textbf{(\roman*)}]
        \item $0 = \frac{0}{2^0} \in S$ and $1 = \frac{1}{2^0} \in S$.
        \item Closed under addition: For $\frac{a}{2^n}, \frac{b}{2^m} \in S$ with $n \ge m$,
        \[
        \frac{a}{2^n} + \frac{b}{2^m} = \frac{a + b \cdot 2^{n-m}}{2^n} \in S.
        \]
        \item Closed under multiplication: $\frac{a}{2^n} \cdot \frac{b}{2^m} = \frac{a \cdot b}{2^{n+m}} \in S$.
        \item Closed under additive inverses: $-\left(\frac{a}{2^n}\right) = \frac{-a}{2^n} \in S$.
    \end{enumerate}
\end{itemize}
Thus, $\mathbb{Z}[\frac{1}{2}]$ is a subring of $\mathbb{Q}$. More generally, for any $m \in \mathbb{Z}_{\ge 2}$ (e.g., $m = 17$), the localized ring $\mathbb{Z}[\frac{1}{m}] := \{ \frac{a}{m^n} \mid a \in \mathbb{Z}, n \ge 0 \}$ forms a proper intermediate subring $\mathbb{Z} \subsetneq \mathbb{Z}[\frac{1}{m}] \subsetneq \mathbb{Q}$.
\end{example}
\end{content}

\subsection{Polynomial Subrings: \texorpdfstring{$\mathbb{Z}[x^2] \subset \mathbb{Z}[x]$}{Z[x^2] in Z[x]}}

\begin{speech}[00:44:10 - 00:45:05]
Okay, so another example. 

Do we get a bell? I don't remember. Do I get to go to the bell, or do I have to stop? 

Okay, let's try one more example before the bell, hopefully. 

If I look at polynomials here, these are polynomials in the variable $x$, that's a ring: $\mathbb{Z}[x]$. Inside it are polynomials in the variable $x^2$: $\mathbb{Z}[x^2]$. Every polynomial in the variable $x^2$ is a polynomial in $x$, so this is a subset: $\mathbb{Z}[x^2] \subset \mathbb{Z}[x]$. I hope that's clear. 

So, is this an ideal, or a subring, or both, or neither, or what?

You ran out of time. Okay, we take a break, you can think about it.
\end{speech}

\begin{pause}[Mid-Lecture Break]
The lecturer announces a short pause for the students to consider whether $\mathbb{Z}[x^2] \subset \mathbb{Z}[x]$ is an ideal or a subring. The lecture resumes at timestamp 00:45:06.
\end{pause}

\begin{speech}[00:45:06 - 00:45:48]
Okay, so maybe we start, or is it a minute? 

I don't know whether we get a bell, I don't remember, but anyway. So, what's the answer to this question?

Yeah?
\end{speech}

\begin{student}[Student Answer]
I'd say it's a subring but not an ideal.
\end{student}

\begin{speech}[00:45:48 - 00:46:22]
Yeah, so subring: yes. Ideal: no. So, for the subring, again, you just have to check that everything's there. 

\inlinemetanote{The bell rings}
Oh, we do get a bell. 

So, why is it not an ideal? Yeah?
\end{speech}

\begin{student}[Student Answer]
If you take $x^2$, and multiply by $x$, then you get $x^3$, which is not in the set.
\end{student}

\begin{speech}[00:46:22 - 00:46:37]
So, $x^2$ is in this, and if it were an ideal, then I should be able to take any $r$ in the ring and multiply it by this $x^2$. So, I take any $r$ in the ring, multiply by $x^2$, it should stay in here. That's if it were an ideal. 

So, I could take $r$ to be $x$, and then this would be odd, and it wouldn't be in there, okay? So, you can check it yourself. 

These are not supposed to be hard questions; it's just about whether you have absorbed the definitions. Which, as I said, is not an easy task.
\end{speech}

\begin{content}[Polynomial Subrings with Even Powers]
\begin{example}[{Even Degree Polynomials $\mathbb{Z}[x^2] \subset \mathbb{Z}[x]$}]\label[example]{ex:even-polynomials-subring}
Let $R = \mathbb{Z}[x]$ be the polynomial ring over $\mathbb{Z}$, and consider the subset $\mathbb{Z}[x^2] \subset \mathbb{Z}[x]$ of polynomials containing only even powers of $x$:
\begin{equation}\label{eq:even-polynomials-def}
\mathbb{Z}[x^2] := \left\{ \sum_{k=0}^m a_k (x^2)^k = a_0 + a_1 x^2 + a_2 x^4 + \dots + a_m x^{2m} \;\middle|\; a_k \in \mathbb{Z} \right\}.
\end{equation}
\begin{itemize}
    \item \textbf{$\mathbb{Z}[x^2]$ is a Subring of $\mathbb{Z}[x]$ (\textcolor{ForestGreen}{YES}):}
    $0, 1 \in \mathbb{Z}[x^2]$, the sum and difference of even polynomials are even, and the product of even powers yields even powers ($x^{2k} \cdot x^{2\ell} = x^{2(k+\ell)}$). Thus, $\mathbb{Z}[x^2] \le \mathbb{Z}[x]$.
    \item \textbf{$\mathbb{Z}[x^2]$ is NOT an Ideal of $\mathbb{Z}[x]$ (\textcolor{BrickRed}{NO}):}
    Taking $r = x \in \mathbb{Z}[x]$ and $p(x) = x^2 \in \mathbb{Z}[x^2]$, their product is:
    \[
    r \cdot p(x) = x \cdot x^2 = x^3 \notin \mathbb{Z}[x^2]
    \]
    since $x^3$ has an odd degree. Thus, $\mathbb{Z}[x^2]$ fails the multiplicative absorption property.
\end{itemize}
\end{example}
\end{content}

\subsection{Subrings of the Integers \texorpdfstring{$\mathbb{Z}$}{Z}}

% [TIMESTAMP-SUSPECT: inspect manually — original: 00:46:37 - 00:47:30]
\begin{speech}[00:46:37 - 00:47:30]
Okay, so let's go to the integers. What are the subrings? 

We've seen the rational numbers have lots of interesting subrings, because we found some: there's the integers, and you could invert $2$, invert $3$, $17$. But what are the subrings of $\mathbb{Z}$?

Yeah?
\end{speech}

\begin{student}[Student Answer]
Only $\mathbb{Z}$.
\end{student}

\begin{speech}[00:46:54 - 00:47:30]
Yeah. There's no real interesting subrings: only $\mathbb{Z}$ in $\mathbb{Z}$. That's the only subring. 

The whole ring is always a subring, just by definition. And this is the only one. 

And why is that? Because the subring would have to contain $1$. And once it contains $1$, it contains every positive number. But the subring also is a subgroup, so it has to contain the negative ones too. 

Okay? So, there's nothing you can do.
\end{speech}

\begin{content}[Classification of the Subrings of the Integers]
\begin{proposition}[Subrings of $\mathbb{Z}$]\label[proposition]{prop:subrings-of-z}
The only subring of $(\mathbb{Z}, +, \cdot, 0, 1)$ is $\mathbb{Z}$ itself:
\begin{equation}\label{eq:subrings-of-z}
S \le \mathbb{Z} \implies S = \mathbb{Z}.
\end{equation}
\end{proposition}

\begin{proof}
Let $S \le \mathbb{Z}$ be a subring.
\begin{enumerate}[label=\textbf{(\arabic*)}]
    \item By definition of a subring, $1 \in S$.
    \item Since $S$ is closed under addition, by induction $n = \underbrace{1 + 1 + \dots + 1}_{n \text{ times}} \in S$ for all positive integers $n \in \mathbb{N}_{>0}$.
    \item By definition, $0 \in S$.
    \item Since $(S, +, 0)$ is an additive subgroup, $S$ is closed under additive inverses: $-n \in S$ for all $n \in \mathbb{N}_{>0}$.
\end{enumerate}
Combining these, $S$ contains all integers $\mathbb{Z} \subseteq S$. Since $S \subseteq \mathbb{Z}$ by assumption, $S = \mathbb{Z}$.
\end{proof}

\begin{steps}[Contrast: Ideals vs. Subrings in $\mathbb{Z}$]
This result highlights the sharp contrast between subrings and ideals in the ring of integers $\mathbb{Z}$:
\begin{itemize}
    \item \textbf{Subrings of $\mathbb{Z}$:} There is only \textbf{one} subring, namely $\mathbb{Z}$ itself (because a subring must contain $1$, which generates all of $\mathbb{Z}$ additively).
    \item \textbf{Ideals of $\mathbb{Z}$:} There are \textbf{infinitely many} distinct ideals of $\mathbb{Z}$, each of the form $n\mathbb{Z} = \langle n \rangle$ for $n \in \mathbb{N}_{\ge 0}$ (because an ideal does not need to contain $1$).
\end{itemize}
\end{steps}
\end{content}

\subsection{Classification of Ideals of \texorpdfstring{$\mathbb{Z}$}{Z}}

\begin{speech}[00:47:30 - 00:48:25]
So, what are the ideals? What are the ideals of $\mathbb{Z}$? 

What would you guess the answer to this question is?

Yeah?
\end{speech}

\begin{student}[Student Answer]
The set of multiples of an integer.
\end{student}

\begin{speech}[00:48:25 - 00:49:35]
Yeah, so you pick some integer $n$, and you look at the multiples of $n$. 

That's the standard notation: $n \cdot \mathbb{Z}$. So, if $n$ is $2$, it's the even numbers; if $n$ is $3$, it's the numbers divisible by $3$, and so on. 

Comment 1: this is certainly an ideal. It's just because the way we checked carefully that $2$, the even ones, were an ideal, it's the same argument, okay?

Comment 2: it's not sensitive to whether you put $+n$ or $-n$ here, because inside $\mathbb{Z}$ is a $-1$. So, if you put $+7$ or $-7$, it's the same set. Is that also clear?

So, typically when you write $n\mathbb{Z}$, you should take $n \ge 0$. You don't have to, but that's the convention. In the literature, you'll never see minus... it's just a weird thing to do. There's no reason...
\end{speech}

\begin{content}[Principal Ideals of the Integers $\mathbb{Z}$]
\begin{definition}[Principal Ideals of $\mathbb{Z}$]\label[definition]{def:principal-ideals-z}
For each integer $n \in \mathbb{Z}$, the set of integer multiples of $n$:
\begin{equation}\label{eq:n-z-def}
n\mathbb{Z} := \{ n \cdot k \mid k \in \mathbb{Z} \}
\end{equation}
is an ideal of $\mathbb{Z}$.
\end{definition}

\begin{remark}[Conventions on $n\mathbb{Z}$]\label[remark]{rem:nz-conventions}
\begin{itemize}
    \item \textbf{Ideals of $\mathbb{Z}$:} For every $n \in \mathbb{Z}$, $(n\mathbb{Z}, +, 0)$ is an additive subgroup of $\mathbb{Z}$, and for any $r \in \mathbb{Z}$ and $x = nk \in n\mathbb{Z}$, $r \cdot x = n(rk) \in n\mathbb{Z}$.
    \item \textbf{Sign Invariance:} Since $-1 \in \mathbb{Z}$ is a unit, $(-n)\mathbb{Z} = n\mathbb{Z}$. By convention, we restrict the generator index to $n \in \mathbb{N}_{\ge 0} = \{0, 1, 2, \dots\}$.
\end{itemize}
\end{remark}
\end{content}

\subsection{Historical Origin of Ideals and David Hilbert}

\begin{speech}[00:49:35 - 00:50:40]
So, we're going to prove these are all the ideals. But this is an important example because the word ideal, where does it come from? The idea is that it's an idealized number, idealized integer. 

For the integers, the ideals are the integers, but when you think about other rings, the ideals are some kind of idea of a number. This was in some sense the origin of this word ideal here: ideal number. This terminology goes back to about $1900$, and there were lots of people around, but certainly one person who played a role was the great David Hilbert, who was a mathematician in Göttingen. And a piece of math culture: does anyone know what the epitaph on the grave of David Hilbert is?

Not one person?

Yeah?
\end{speech}

\begin{student}[Student Question]
What is an epitaph?
\end{student}

\begin{speech}[00:50:40 - 00:51:45]
Ah, epitaph! So, you have a gravestone \inlinemetanote{draws a sketch of an arched gravestone on the blackboard}, and then what's written here? There's the person's name, and then there's like a little motto or slogan. That's the epitaph. 

Okay, so this is a real piece of math culture, what's the epitaph of David Hilbert's grave. This is the kind of thing you learn to be a mathematician; it's part of being in the culture. Yeah, do you know it?
\end{speech}

\begin{student}[Student Answer]
Wir müssen wissen, wir werden wissen.
\end{student}

\begin{speech}[00:51:45 - 00:52:30]
Yeah: \qt{Wir müssen wissen, wir werden wissen.} Two lines. It's in Göttingen. 

\inlinemetanote{The lecturer writes \qt{Motto / Slogan} inside the drawn gravestone on the board, and then erases it}

I have to teach you some culture, too, you know. 

Okay, so we have to prove something.
\end{speech}

\begin{content}[Historical Context: Ideal Numbers and David Hilbert]
\begin{center}
\begin{tikzpicture}[scale=1.1]
    % Gravestone body
    \draw[thick, fill=gray!10] (-1.5, 0) -- (-1.5, 2.0) arc (180:0:1.5) -- (1.5, 0) -- cycle;
    \draw[thick] (-1.8, 0) -- (1.8, 0);
    
    % Inscription
    \node[font=\bfseries] at (0, 2.3) {David Hilbert};
    \node[font=\footnotesize, text=gray!70] at (0, 1.9) {1862 -- 1943};
    \draw[thin, gray!50] (-1.0, 1.6) -- (1.0, 1.6);
    \node[font=\itshape\small] at (0, 1.1) {\qt{Wir müssen wissen,}};
    \node[font=\itshape\small] at (0, 0.6) {\qt{wir werden wissen.}};
\end{tikzpicture}
\end{center}

\begin{steps}[Origin of the Term \qt{Ideal}]
The concept of an \newterm{ideal} originated in algebraic number theory with Ernst Kummer (1840s) as \qt{ideal complex numbers} to restore unique prime factorization in rings of algebraic integers such as $\mathbb{Z}[\zeta_p]$, and was later formalized by Richard Dedekind and David Hilbert around 1900. In $\mathbb{Z}$, every ideal is represented by a concrete integer ($I = n\mathbb{Z}$); in more general rings, ideals generalize the behavior of numbers.
\end{steps}
\end{content}

\subsection{Proof: Every Ideal in \texorpdfstring{$\mathbb{Z}$}{Z} is Principal}

\begin{content}
\begin{theorem}[Ideals of $\mathbb{Z}$ are Principal]\label[theorem]{thm:ideals-of-z-are-principal}
Every ideal $I$ of the ring of integers $\mathbb{Z}$ is of the form:
\begin{equation}\label{eq:ideal-is-nz}
I = n\mathbb{Z} = \langle n \rangle
\end{equation}
for a uniquely determined non-negative integer $n \in \mathbb{N}_{\ge 0}$.
\end{theorem}
\end{content}

\begin{proofWrapper}[Proof that Every Ideal in $\mathbb{Z}$ is of the Form $n\mathbb{Z}$]

\begin{speech}[00:52:30 - 00:54:10]
What do we have to prove? 

Let $I \subseteq \mathbb{Z}$ be an ideal. 

I want to prove it's $n\mathbb{Z}$. What about math what I like is truth. Proof is good, but truth is better. And this stuff is true, but anyway, we should prove it; it's a math course. 

So, if you want to prove it: \textbf{Case (1)} is that $I$ could be just the zero set ($\{0\}$). And this means it's $0 \cdot \mathbb{Z}$. It's a degenerate case. That's possible. And that's an ideal: $0$ is an ideal. 

So, \textbf{Case (2)} is that $I$ has some non-zero element.
\end{speech}

\begin{content}[Proof: Case 1 and Identification of Positive Elements]
\textbf{Case 1 ($I = \{0\}$):} \\
If $I = \{0\}$, then $I = 0\mathbb{Z}$ with $n = 0$. This settles the degenerate case.

\textbf{Case 2 ($I \neq \{0\}$):} \\
If $I$ contains a non-zero element $x \in I$ with $x \neq 0$:
\begin{itemize}
    \item If $x > 0$, then $I$ contains a positive integer.
    \item If $x < 0$, since $I$ is an ideal and $-1 \in \mathbb{Z}$, $-x = (-1) \cdot x \in I$, where $-x > 0$.
\end{itemize}
In either case, the set of positive elements $I \cap \mathbb{N}_{>0}$ is non-empty.
\end{content}

\begin{speech}[00:54:10 - 00:55:40]
Because $I$ has to have $0$ because $I$ is an additive subgroup. If you're not in this case, you have to have some non-zero element. 

Let's call that element $x$. $x$ is in $I$. 

Now, $x$ is an integer. 

So, we can multiply by $\pm 1$ and get a positive integer. One of these is a positive integer. 

So, if it has a non-zero element, then $I$ has a positive integer. 

And if it has a positive integer, what am I going to say now? 

Yeah?
\end{speech}

\begin{student}[Student Answer]
It has a minimum.
\end{student}

\begin{speech}[00:55:40 - 00:57:00]
Yeah, has a least positive integer. That's the next line. 

So, $I$ has a least positive integer. 

Basically every property of integers at some point you have to use this principle. So, once it has a least positive integer, let's call that one $n$. So, $n$ is equal to the least positive integer of $I$. 

So, now we have to show that $I$ is actually $n\mathbb{Z}$. That's what we want to show, all right? 

You look at $y \in I$. There's only one thing you can do now, which is the division algorithm. The division algorithm says $y$ is equal to some multiple of $n$ plus a remainder term:
\[
y = q \cdot n + r.
\]
That is the division, long division.
\end{speech}

\begin{content}[Application of the Well-Ordering Principle and Division Algorithm]
\textbf{Selection of the Minimal Generator $n$:} \\
By the \newterm{Well-Ordering Principle} of $\mathbb{N}$, the non-empty subset $I \cap \mathbb{N}_{>0} \subseteq \mathbb{N}$ contains a unique smallest element:
\begin{equation}\label{eq:n-least-positive}
n := \min(I \cap \mathbb{N}_{>0}).
\end{equation}

\textbf{Claim:} $I = n\mathbb{Z}$.
\begin{itemize}
    \item \textbf{Inclusion $n\mathbb{Z} \subseteq I$:} Since $n \in I$ and $I$ is an ideal, for any $k \in \mathbb{Z}$, $k \cdot n \in I$. Thus, $n\mathbb{Z} \subseteq I$.
    \item \textbf{Inclusion $I \subseteq n\mathbb{Z}$:} Let $y \in I$ be an arbitrary element. By the \newterm{Division Algorithm} (division with remainder in $\mathbb{Z}$), there exist unique integers $q, r \in \mathbb{Z}$ such that:
    \begin{equation}\label{eq:division-algorithm}
    y = q \cdot n + r \quad \text{with} \quad 0 \le r < n.
    \end{equation}
\end{itemize}
\end{content}

\begin{speech}[00:57:00 - 00:58:30]
And the important part of this remainder term is $0 \le r < n$. 

Suppose $y$ is not in $n\mathbb{Z}$, that means that it's not divisible, so there is a remainder term. 

But this remainder term is equal to $y - qn$, and so the remainder term is also in $I$, because $y$ is in $I$, $n$ is in $I$ so $qn$ is in $I$, and since $I$ is an additive subgroup, this is also in $I$. 

It can't be zero because otherwise it would divide, so the long division says $r$ is between $0$ and $n$. 

But then $r$ is a lesser positive integer, which is a contradiction. Okay, that's the proof.
\end{speech}

\begin{speech}[00:58:30 - 00:59:10]
So, we use what you have to use for integers: you use this fact that non-empty subsets have a least one. Once you have that least one, you know what you're trying to prove, and then you use the long division algorithm and use the remainder. This is what you always do with these kind of elementary properties of integers. 

So, that proves this claim. 

And that's the most basic properties of the ideals in the integers.
\end{speech}

\begin{content}[Conclusion of the Proof: Minimality Contradiction]
\begin{proof}[Continuation of Proof of \cref{thm:ideals-of-z-are-principal}]
Rearranging equation \eqref{eq:division-algorithm} for the remainder $r$:
\[
r = y - q \cdot n.
\]
Since $y \in I$, $n \in I$, and $I$ is an ideal, $q \cdot n \in I$. Since $(I, +, 0)$ is an additive subgroup:
\[
r = y - q \cdot n \in I.
\]
We now analyze the remainder $r$:
\begin{itemize}
    \item If $r > 0$, then $r \in I \cap \mathbb{N}_{>0}$ with $r < n$. But this directly contradicts the defining minimality of $n = \min(I \cap \mathbb{N}_{>0})$ from \eqref{eq:n-least-positive}.
    \item Therefore, we must have $r = 0$.
\end{itemize}
Setting $r = 0$ in \eqref{eq:division-algorithm} gives $y = q \cdot n \in n\mathbb{Z}$. Since $y \in I$ was arbitrary, $I \subseteq n\mathbb{Z}$.
Combining both inclusions yields $I = n\mathbb{Z}$.
\end{proof}
\end{content}

\end{proofWrapper}

\begin{insight}[Principal Ideal Domains (PIDs)]
A ring in which every ideal is generated by a single element (i.e., $I = \langle a \rangle = aR$) is called a \newterm{Principal Ideal Domain} (\newterm{PID}). \Cref{thm:ideals-of-z-are-principal} proves that the ring of integers $\mathbb{Z}$ is a PID.
\end{insight}

\subsection{Rings of Smooth Functions and Compactly Supported Ideals}

\begin{speech}[00:59:10 - 01:00:50]
For any ring $R$, we've seen before that $\{0\}$ is an ideal, and the whole ring $R$ is an ideal. 

And we haven't discussed fields at all today, but we discussed last time that if you have a field, these are the only two ideals. 

But the integers have a lot of ideals: all these $n\mathbb{Z}$ are ideals. 

Okay, so let's throw in another example. One object you certainly know in mathematics: if I take the real numbers, you can take $C^\infty$ functions. 

You don't have to take $C^\infty$, you could take $17$ times differentiable functions, but we'll just go $\infty$. These are infinitely differentiable functions. 

So an element here is some function $f \colon \mathbb{R} \to \mathbb{R}$, which I can differentiate as many times as I want, so it has a nice graph; I can graph it like this \inlinemetanote{sketches a wavy function curve on the board}. 

And this beautiful object is a ring. 

Because why? You can add functions, you can multiply functions, $0$ is just the constant function $0$, and $1$ is the constant function $1$. 

Is that clear? It's a ring.
\end{speech}

\begin{content}[The Ring of Smooth Functions \texorpdfstring{$C^\infty(\mathbb{R})$}{C^infty(R)}]
\begin{definition}[Ring of Smooth Functions]\label[definition]{def:smooth-functions-ring}
The set of infinitely differentiable functions on the real line:
\begin{equation}\label{eq:c-infty-def}
C^\infty(\mathbb{R}) := \bigl\{ f \colon \mathbb{R} \to \mathbb{R} \;\big|\; f \text{ is smooth} \bigr\}
\end{equation}
forms a commutative ring $(C^\infty(\mathbb{R}), +, \cdot, 0, 1)$ under pointwise operations:
\begin{align*}
(f + g)(x) &:= f(x) + g(x), \\
(f \cdot g)(x) &:= f(x) \cdot g(x),
\end{align*}
with additive identity $0(x) \equiv 0$ and multiplicative identity $1(x) \equiv 1$.
\end{definition}

\begin{center}
\begin{tikzpicture}[scale=1.2]
    \draw[->, thick, gray!70] (-3.0, 0) -- (3.5, 0) node[right, text=black] {$x$};
    \draw[->, thick, gray!70] (0, -1.2) -- (0, 2.0) node[above, text=black] {$y$};
    
    \draw[ultra thick, MidnightBlue] (-2.8, -0.5) 
        .. controls (-2.0, 1.8) and (-1.0, 0.2) .. (0.0, 1.4)
        .. controls (1.0, 2.2) and (2.0, -1.0) .. (3.0, 0.8)
        node[right] {$f \in C^\infty(\mathbb{R})$};
\end{tikzpicture}
\end{center}
\end{content}

\begin{speech}[01:00:50 - 01:02:05]
\inlinemetanote{A student asks a question about the board text}
Here? Which board? Infinitely differentiable function. 

Yeah, my handwriting's not good: function. And it's not because I have this wrist injury \inlinemetanote{laughs}; I just have bad handwriting. 

And I don't even know how I got this wrist injury. It's not like I'm an MMA fighter or anything like that. It just started hurting at one point. I thought it was going to be a career-ending injury, but it turns out that if I wear this thing, it's fine.
\end{speech}

\begin{speech}[01:02:05 - 01:03:45]
In this class, it's not primarily about such rings of differentiable functions --- that's some other class ---, but I can show you something interesting in it. 

There's an interesting ideal inside $C^\infty(\mathbb{R})$, and these are functions with compact support: $C_c^\infty(\mathbb{R})$. 

It's the set of functions with compact support. 

A general function can do whatever it wants, totally unconstrained. But a function with compact support means that its entire life happens somewhere compact. So it has to be zero here, and then suddenly it comes to life, and then it goes to zero. So it's like this kind of hat function that looks like that, okay? That's a function of compact support.
\end{speech}

\begin{speech}[01:03:45 - 01:05:10]
So, is this an ideal? 

To be an ideal, it has to be the case that zero's there, and very fortunately zero has compact support. What does not have compact support is the constant one function; that does not have compact support. 

So, this whatever it is, it's surely not a subring, because it does not have the constant one function. 

If you add compact support things, they have compact support, that's true. And the interesting thing is if I take a function of compact support and I multiply it by any function, it still has compact support, because once this function is zero, it just cannot be brought back to life, no matter how much you multiply. 

So, this is an example.
\end{speech}

\begin{content}[Compactly Supported Smooth Functions as an Ideal]
\begin{definition}[Functions with Compact Support]\label[definition]{def:compactly-supported-functions}
The set of smooth functions with \newterm{compact support} on $\mathbb{R}$, denoted $C_c^\infty(\mathbb{R})$, is:
\begin{equation}\label{eq:cc-infty-def}
C_c^\infty(\mathbb{R}) := \bigl\{ f \in C^\infty(\mathbb{R}) \;\big|\; \operatorname{supp}(f) \text{ is compact} \bigr\},
\end{equation}
where $\operatorname{supp}(f) := \overline{\{x \in \mathbb{R} \mid f(x) \neq 0\}}$ is bounded (i.e., $f(x) = 0$ outside some finite interval $[-M, M]$).
\end{definition}

\begin{center}
\begin{tikzpicture}[scale=1.2]
    \draw[->, thick, gray!70] (-3.5, 0) -- (3.5, 0) node[right, text=black] {$x$};
    \draw[->, thick, gray!70] (0, -0.3) -- (0, 1.8) node[above, text=black] {$y$};
    
    % Flat tails
    \draw[ultra thick, BrickRed] (-3.2, 0) -- (-1.5, 0);
    \draw[ultra thick, BrickRed] (1.5, 0) -- (3.2, 0);
    
    % Smooth bump (hat function)
    \draw[ultra thick, BrickRed] (-1.5, 0) 
        .. controls (-1.0, 0.0) and (-0.8, 1.3) .. (-0.3, 1.4)
        .. controls (0.0, 1.5) and (0.3, 1.1) .. (0.6, 1.3)
        .. controls (0.9, 1.5) and (1.0, 0.0) .. (1.5, 0);
        
    \node[BrickRed, above] at (0, 1.5) {$f \in C_c^\infty(\mathbb{R})$};
    \draw[dashed, ForestGreen] (-1.5, 0) -- (-1.5, -0.3) node[below] {$-M$};
    \draw[dashed, ForestGreen] (1.5, 0) -- (1.5, -0.3) node[below] {$M$};
    \node[ForestGreen, font=\footnotesize] at (0, -0.3) {$\operatorname{supp}(f) \subseteq [-M, M]$};
\end{tikzpicture}
\end{center}

\begin{proposition}\label[proposition]{prop:cc-infty-is-ideal}
$C_c^\infty(\mathbb{R})$ is an ideal of $C^\infty(\mathbb{R})$, but it is not a subring.
\end{proposition}

\begin{proof}
\begin{enumerate}[label=\textbf{(\arabic*)}]
    \item \textbf{Additive Subgroup:} The zero function $0 \in C_c^\infty(\mathbb{R})$ has $\operatorname{supp}(0) = \emptyset$. For $f, g \in C_c^\infty(\mathbb{R})$, $\operatorname{supp}(f + g) \subseteq \operatorname{supp}(f) \cup \operatorname{supp}(g)$, which is compact as the union of two compact sets.
    \item \textbf{Multiplicative Absorption (Ideal Property):} For any arbitrary $g \in C^\infty(\mathbb{R})$ and $f \in C_c^\infty(\mathbb{R})$,
    \[
    \operatorname{supp}(g \cdot f) \subseteq \operatorname{supp}(f).
    \]
    Since $\operatorname{supp}(f)$ is compact and closed subsets of compact sets in $\mathbb{R}$ are compact, $g \cdot f \in C_c^\infty(\mathbb{R})$.
    \item \textbf{Not a Subring:} The constant unit function $1(x) \equiv 1$ has $\operatorname{supp}(1) = \mathbb{R}$, which is not compact. Since $1 \notin C_c^\infty(\mathbb{R})$, $C_c^\infty(\mathbb{R})$ is not a subring.
\end{enumerate}
\end{proof}
\end{content}

\begin{speech}[01:05:10 - 01:06:40]
Maybe I ask you a question: is there some ideal $J$ that's bigger? 

Can you find an ideal that's bigger than this, and not the whole ring, of course, because that's always an ideal that's bigger? 

The traditional vision is there's an algebra direction, an analysis direction, a geometry direction. But actually, it doesn't really make sense to separate them all. It's all just one way of looking at numbers. 

Can you think of a way to make this ideal bigger?
\end{speech}

\begin{speech}[01:06:40 - 01:07:55]
Well, we used the fact that it's zero on both sides here, but it's not necessary, right? Can't we just say functions that eventually go to zero on the right? Those are more. And it satisfies the same thing. 

So here you could say $J$ is functions that become zero for all sufficiently large $x$. 

Those are functions that look like: they do whatever they want, and then eventually they become zero. These functions are certainly inside there, but there are more of them. 

And these things are also an ideal, because if you multiply by something that's zero here, it's going to stay zero.
\end{speech}

\begin{content}[A Strictly Larger Ideal: Functions Vanishing at Positive Infinity]
\begin{example}[Functions Vanishing for Large $x$]\label[example]{ex:ideal-vanishing-at-infinity}
Consider the set of smooth functions that vanish identically for all sufficiently large positive $x$:
\begin{equation}\label{eq:j-ideal-def}
J := \bigl\{ f \in C^\infty(\mathbb{R}) \;\big|\; \exists M \in \mathbb{R} \text{ such that } f(x) = 0 \text{ for all } x \ge M \bigr\}.
\end{equation}
\begin{itemize}
    \item \textbf{Proper Inclusions:} $C_c^\infty(\mathbb{R}) \subsetneq J \subsetneq C^\infty(\mathbb{R})$.
    Any compactly supported function vanishes for large $x$, so $C_c^\infty(\mathbb{R}) \subseteq J$. However, a function that is non-zero on all of $(-\infty, 0)$ and vanishes on $[1, \infty)$ belongs to $J \setminus C_c^\infty(\mathbb{R})$. Furthermore, $1 \notin J$, so $J \neq C^\infty(\mathbb{R})$.
    \item \textbf{Ideal Verification:} If $f \in J$ vanishes on $[M, \infty)$ and $g \in C^\infty(\mathbb{R})$, then $(g \cdot f)(x) = g(x) \cdot 0 = 0$ for all $x \ge M$, so $g \cdot f \in J$. Thus, $J$ is an ideal.
\end{itemize}

\begin{center}
\begin{tikzpicture}[scale=1.2]
    \draw[->, thick, gray!70] (-3.5, 0) -- (3.5, 0) node[right, text=black] {$x$};
    \draw[->, thick, gray!70] (0, -0.3) -- (0, 1.8) node[above, text=black] {$y$};
    
    % Non-zero on left, decaying to 0 at M=1.0
    \draw[ultra thick, BurntOrange] (-3.2, 0.8) 
        .. controls (-2.5, 1.6) and (-1.8, -0.2) .. (-1.0, 1.2)
        .. controls (-0.5, 1.8) and (0.2, 1.2) .. (1.0, 0);
    \draw[ultra thick, BurntOrange] (1.0, 0) -- (3.2, 0);
    
    \draw[dashed, gray] (1.0, 0) -- (1.0, -0.3) node[below] {$M$};
    \node[BurntOrange, above right] at (-1.0, 1.3) {$f \in J \setminus C_c^\infty(\mathbb{R})$};
    \node[gray!80, font=\footnotesize] at (2.2, -0.3) {$f(x) = 0 \ \forall x \ge M$};
\end{tikzpicture}
\end{center}
\end{example}
\end{content}

\subsection{The Restriction Homomorphism}

\begin{speech}[01:07:55 - 01:09:12]
If I look at $C^\infty$ functions of $\mathbb{R}$, there's a ring homomorphism to $C^\infty$ functions of positive reals ($\mathbb{R}_{>0}$):
\[
C^\infty(\mathbb{R}) \xrightarrow{\operatorname{res}} C^\infty(\mathbb{R}_{>0}).
\]
This is a good example of a homomorphism. 

And this is called restriction, $\operatorname{res}$. 

Meaning if I have a $C^\infty$ function here, it's a function from all of $\mathbb{R}$ to $\mathbb{R}$. If I have one, I can just restrict it:
\[
f \colon \mathbb{R} \to \mathbb{R} \implies f|_{\mathbb{R}_{>0}} \colon \mathbb{R}_{>0} \to \mathbb{R}.
\]
Is that clear? Every element here, I can just restrict, and that gives me a map. First, it's a map on sets, and then you have to show it respects the zero, it respects the function that's constant $1$, and addition and multiplication behave well under restriction. So, this is a legitimate ring homomorphism. 

So for function spaces, restriction is a very useful way of making all these ring homomorphisms. 

Now your analysis question: is this surjective? 

That's properly a question in your analysis class. You can think about that.
\end{speech}

\begin{content}[The Restriction Homomorphism on Function Spaces]
\begin{definition}[Restriction Homomorphism]\label[definition]{def:restriction-homomorphism}
Let $U \subset \mathbb{R}$ be an open subset (e.g., $U = \mathbb{R}_{>0} = (0, \infty)$). The \newterm{restriction mapping}:
\begin{equation}\label{eq:restriction-map-def}
\operatorname{res} \colon C^\infty(\mathbb{R}) \to C^\infty(U), \quad f \mapsto f|_U
\end{equation}
is a ring homomorphism satisfying:
\begin{itemize}
    \item $\operatorname{res}(f + g) = (f + g)|_U = f|_U + g|_U = \operatorname{res}(f) + \operatorname{res}(g)$,
    \item $\operatorname{res}(f \cdot g) = (f \cdot g)|_U = f|_U \cdot g|_U = \operatorname{res}(f) \cdot \operatorname{res}(g)$,
    \item $\operatorname{res}(1_{\mathbb{R}}) = 1_U$.
\end{itemize}
\end{definition}
\end{content}

\section{Introduction to Category Theory}

\subsection{Categories, Objects, and Morphisms}

\begin{speech}[01:09:12 - 01:10:48]
There are two foundational topics which I've avoided discussing, but I think I'll discuss now. And it has to do with category theory and Zorn's lemma.

So, you had a foundations, this Grundstrukturen class. Did you discuss categories there?

You can say no.

Okay, so category. We discuss category. 

A category consists of objects and homs.

Homs are sometimes called morphisms, or arrows. The languages differ across different authors.

So, the first thing to know about a category: it has objects and it has homs.

More specifically, given two objects $A, B$ in the category, there is a set of homs, and this set has a name:
\[
\operatorname{Hom}(A, B).
\]
\end{speech}

\begin{content}[Foundational Data of a Category]
\begin{definition}[Category: Objects and Hom-Sets]\label[definition]{def:category-foundations}
A \newterm{category} $\mathcal{C}$ consists of:
\begin{enumerate}[label=\textbf{(\arabic*)}]
    \item A collection of \newterm{objects}, denoted $\operatorname{Obj}(\mathcal{C})$. We write $A \in \operatorname{Obj}(\mathcal{C})$ to indicate that $A$ is an object of $\mathcal{C}$.
    \item For every pair of objects $A, B \in \operatorname{Obj}(\mathcal{C})$, a set of \newterm{morphisms} (also called \newterm{homs} or \newterm{arrows}), denoted:
    \begin{equation}\label{eq:hom-set}
    \operatorname{Hom}_{\mathcal{C}}(A, B) \quad (\text{or } \operatorname{Mor}_{\mathcal{C}}(A, B)).
    \end{equation}
    An element $f \in \operatorname{Hom}_{\mathcal{C}}(A, B)$ is represented diagrammatically as an arrow $f \colon A \to B$ or $A \xrightarrow{f} B$, where $A$ is the \newterm{domain} (source) and $B$ is the \newterm{codomain} (target).
\end{enumerate}
\end{definition}
\end{content}

\subsection{Examples of Categories: Groups, Rings, and Vector Spaces}

\begin{speech}[01:10:48 - 01:12:06]
For example, the category of groups.

The category of groups has objects, and the objects are groups.

An object is a group $(G, \cdot, e)$. That's the objects. In the world of all groups, every possible group is in here. 

And what are the homs going to be? The homs between two particular groups are the set of group homomorphisms from $G$ to $\hat{G}$.

Another example is the category of rings. There the objects are rings with the full data, and the homs are ring homomorphisms.

From last year you could consider the category of vector spaces over the complex numbers: the objects are vector spaces over $\mathbb{C}$, and the homs are linear maps.

To be a category, I have to tell you what the objects are, given two objects I have to tell you what the homs are, and the resulting structure has to have some properties.
\end{speech}

\begin{speech}[01:12:06 - 01:13:00]
Did you figure out this question about the restriction homomorphism?

\inlinemetanote{A student in the lecture hall answers}
\end{speech}

\begin{student}[Student Answer]
I think $\sin(1/x)$...
\end{student}

\begin{speech}[01:13:00 - 01:14:00]
Yeah, so the answer is no, it's not surjective.

A picture is worth a thousand words. 

The picture is an element here that's not there: $\sin(1/x)$ oscillating near $x = 0$. \inlinemetanote{draws a sketch of $\sin(1/x)$ oscillating near $x = 0$ on the blackboard}

You should recognize that from your analysis class. You know this function, right?
\end{speech}

\begin{content}[Non-Surjectivity of Restriction and the Topologist's Sine Curve]
\begin{proposition}[Non-Surjectivity of Restriction]\label[proposition]{prop:restriction-not-surjective}
The restriction ring homomorphism:
\[
\operatorname{res} \colon C^\infty(\mathbb{R}) \to C^\infty(\mathbb{R}_{>0}), \quad f \mapsto f|_{\mathbb{R}_{>0}}
\]
is \textbf{not} surjective.
\end{proposition}

\begin{proof}
To show that $\operatorname{res}$ is not surjective, it suffices to exhibit a smooth function $g \in C^\infty(\mathbb{R}_{>0})$ that cannot be extended to any continuous (and hence no smooth) function $\tilde{g} \colon \mathbb{R} \to \mathbb{R}$.

Consider the oscillating function:
\begin{equation}\label{eq:sine-one-over-x}
g(x) := \sin\left(\frac{1}{x}\right) \quad \text{for } x > 0.
\end{equation}
Since the composition of smooth functions on $(0, \infty)$ is smooth, $g \in C^\infty(\mathbb{R}_{>0})$. However, as $x \to 0^+$, the argument $\frac{1}{x} \to \infty$, and $g(x)$ oscillates infinitely often between $-1$ and $+1$ on every neighborhood $(0, \varepsilon)$. Consequently, the limit $\lim_{x \to 0^+} \sin(1/x)$ does not exist. Thus, $g$ cannot be extended continuously to $x = 0$, proving that $g \notin \operatorname{Im}(\operatorname{res})$.
\end{proof}

\begin{center}
\begin{tikzpicture}[scale=1.3]
    \draw[->, thick, gray!70] (-0.5, 0) -- (4.0, 0) node[right, text=black] {$x$};
    \draw[->, thick, gray!70] (0, -1.4) -- (0, 1.4) node[above, text=black] {$y$};
    
    % Oscillating function sin(1/x)
    \draw[domain=0.08:3.5, samples=300, smooth, thick, BrickRed] 
        plot (\x, {sin((1.0 / \x) * 180 / 3.14159265)});
        
    % Dense oscillation region indicator near 0
    \draw[thick, BrickRed, opacity=0.7] (0.02, -1.0) -- (0.02, 1.0);
    \draw[thick, BrickRed, opacity=0.5] (0.04, -1.0) -- (0.04, 1.0);
    \draw[thick, BrickRed, opacity=0.3] (0.06, -1.0) -- (0.06, 1.0);
    
    \node[BrickRed, above right] at (1.5, 0.7) {$g(x) = \sin\left(\frac{1}{x}\right)$};
    \node[gray!80, font=\footnotesize, below left] at (0, 0) {$0$};
    \draw[dashed, gray!60] (0, 1) -- (3.5, 1) node[right, font=\footnotesize] {$+1$};
    \draw[dashed, gray!60] (0, -1) -- (3.5, -1) node[right, font=\footnotesize] {$-1$};
\end{tikzpicture}
\end{center}
\end{content}

\begin{content}[Standard Algebraic Categories]
\begin{example}[The Category of Groups $\mathbf{Grp}$]\label[example]{ex:cat-groups}
The \newterm{category of groups}, denoted $\mathbf{Grp}$, is defined by:
\begin{itemize}
    \item \textbf{Objects:} $\operatorname{Obj}(\mathbf{Grp}) = \bigl\{ (G, \cdot, e) \;\big|\; G \text{ is a group} \bigr\}$.
    \item \textbf{Morphisms:} For any two groups $G, \hat{G} \in \operatorname{Obj}(\mathbf{Grp})$,
    \begin{equation}\label{eq:grp-homs}
    \operatorname{Hom}_{\mathbf{Grp}}(G, \hat{G}) := \bigl\{ f \colon G \to \hat{G} \;\big|\; f \text{ is a group homomorphism} \bigr\}.
    \end{equation}
\end{itemize}
\end{example}

\begin{example}[The Category of Rings $\mathbf{Ring}$]\label[example]{ex:cat-rings}
The \newterm{category of rings}, denoted $\mathbf{Ring}$, is defined by:
\begin{itemize}
    \item \textbf{Objects:} $\operatorname{Obj}(\mathbf{Ring}) = \bigl\{ (R, +, \cdot, 0, 1) \;\big|\; R \text{ is a ring} \bigr\}$.
    \item \textbf{Morphisms:} For any two rings $R, \hat{R} \in \operatorname{Obj}(\mathbf{Ring})$,
    \begin{equation}\label{eq:ring-homs}
    \operatorname{Hom}_{\mathbf{Ring}}(R, \hat{R}) := \bigl\{ f \colon R \to \hat{R} \;\big|\; f \text{ is a ring homomorphism (with } f(1) = \hat{1}\text{)} \bigr\}.
    \end{equation}
\end{itemize}
\end{example}

\begin{example}[The Category of Vector Spaces $\mathbf{Vect}_{\mathbb{C}}$]\label[example]{ex:cat-vect}
The \newterm{category of complex vector spaces} $\mathbf{Vect}_{\mathbb{C}}$ is defined by:
\begin{itemize}
    \item \textbf{Objects:} Vector spaces $V$ over $\mathbb{C}$.
    \item \textbf{Morphisms:} Linear maps $T \colon V \to W$, so $\operatorname{Hom}_{\mathbf{Vect}_{\mathbb{C}}}(V, W) = \mathcal{L}(V, W)$.
\end{itemize}
\end{example}
\end{content}

\subsection{Axioms of a Category: Composition and Identity}

\begin{speech}[01:14:00 - 01:15:45]
Okay, so the properties that the homs have to satisfy for the category are very simple properties.

The first property is that $\operatorname{Hom}(A, B)$ is a set.

The second is that there's composition.

Composition means that if I have three objects, $A, B$, and $C$, then here I have the set of homs from $A$ to $B$, and homs from $B$ to $C$, and composition means there has to be a rule that tells you how, if you have a hom from $A$ to $B$ and a hom from $B$ to $C$, this gives me a hom from $A$ to $C$.

If I'm telling you a category, it's up to me to tell you how to compose homs. And in all examples, you just compose them as functions.

And inside $\operatorname{Hom}(A, A)$ there's a distinguished element, called the identity ($\operatorname{id}_A$). That's part of the data.
\end{speech}

\begin{speech}[01:15:45 - 01:17:26]
And there's one more condition, which is how the identity relates to composition. 

If you take the identity in $A$, and you take any element $f \in \operatorname{Hom}(A, B)$, then if I do the identity first and then I do $f$, what should this be equal to?

Yeah, $f$! That's the identity axiom:
\[
f \circ \operatorname{id}_A = f.
\]
And also the other way: if you do $f$ first and then $\operatorname{id}_B$, then this should just be $f$:
\[
\operatorname{id}_B \circ f = f.
\]

The data of this category is to capture what happens when we think about these algebraic objects abstractly. We have objects, we have homs, the homs compose, there's an identity hom, and they satisfy these identity properties and associativity.
\end{speech}

\begin{content}[The Axioms of a Category]
\begin{definition}[Category (Full Formal Definition)]\label[definition]{def:category-full}
A \newterm{category} $\mathcal{C}$ consists of the following data:
\begin{enumerate}[label=\textbf{(\arabic*)}]
    \item A collection $\operatorname{Obj}(\mathcal{C})$ of \newterm{objects}.
    \item For any two objects $A, B \in \operatorname{Obj}(\mathcal{C})$, a set $\operatorname{Hom}_{\mathcal{C}}(A, B)$ of \newterm{morphisms}.
    \item For any three objects $A, B, C \in \operatorname{Obj}(\mathcal{C})$, a \newterm{composition law}:
    \begin{equation}\label{eq:composition-law}
    \circ \colon \operatorname{Hom}_{\mathcal{C}}(B, C) \times \operatorname{Hom}_{\mathcal{C}}(A, B) \to \operatorname{Hom}_{\mathcal{C}}(A, C), \quad (g, f) \mapsto g \circ f.
    \end{equation}
    \item For every object $A \in \operatorname{Obj}(\mathcal{C})$, a distinguished \newterm{identity morphism}:
    \begin{equation}\label{eq:identity-morphism}
    \operatorname{id}_A \in \operatorname{Hom}_{\mathcal{C}}(A, A).
    \end{equation}
\end{enumerate}
These data are required to satisfy the following category axioms:
\begin{itemize}
    \item \textbf{Associativity of Composition:} For all $f \in \operatorname{Hom}_{\mathcal{C}}(A, B)$, $g \in \operatorname{Hom}_{\mathcal{C}}(B, C)$, and $h \in \operatorname{Hom}_{\mathcal{C}}(C, D)$,
    \begin{equation}\label{eq:category-associativity}
    (h \circ g) \circ f = h \circ (g \circ f).
    \end{equation}
    \item \textbf{Identity Axioms (Unit Laws):} For any morphism $f \in \operatorname{Hom}_{\mathcal{C}}(A, B)$, composing with the identity morphism from either the left or the right acts as the identity:
    \begin{equation}\label{eq:category-identity-laws}
    f \circ \operatorname{id}_A = f \quad \text{and} \quad \operatorname{id}_B \circ f = f.
    \end{equation}
\end{itemize}
\end{definition}

\begin{center}
\begin{tikzpicture}[scale=1.3]
    \node (A) at (0,0) [circle, draw, fill=MidnightBlue!10, thick, font=\bfseries] {$A$};
    \node (B) at (3.5,0) [circle, draw, fill=MidnightBlue!10, thick, font=\bfseries] {$B$};
    \node (C) at (1.75,2.2) [circle, draw, fill=MidnightBlue!10, thick, font=\bfseries] {$C$};
    
    \draw[->, thick, RoyalBlue] (A) -- node[below] {$f$} (B);
    \draw[->, thick, RoyalBlue] (B) -- node[above right] {$g$} (C);
    \draw[->, thick, ForestGreen] (A) -- node[above left] {$g \circ f$} (C);
    
    % Identity loops
    \draw[->, thick, BurntOrange] (A) to[out=150, in=210, looseness=4] node[left] {$\operatorname{id}_A$} (A);
    \draw[->, thick, BurntOrange] (B) to[out=30, in=-30, looseness=4] node[right] {$\operatorname{id}_B$} (B);
\end{tikzpicture}
\end{center}
\end{content}

\subsection{Category-Theoretic Definition of Isomorphism}

\begin{speech}[01:17:26 - 01:19:58]
If you have a category, built into the category is a notion of isomorphic objects.

Suppose $\mathbf{Cat}$ is a category. I want to define the notion of isomorphism purely in categorical terms.

Is there a purely arrow-theoretic notion of isomorphism? Category-theoretic?

And the answer is yes.
\end{speech}

% [TIMESTAMP-SUSPECT: inspect manually — original: 01:19:58 - 01:22:08]
\begin{speech}[01:19:58 - 01:22:08]
So, I have this category $\mathbf{Cat}$, and I have two objects: $A$ and $B$. If I have $f \in \operatorname{Hom}(A, B)$, then $f$ is an isomorphism if...

What do you normally need for an isomorphism? You need a map the other way. 

If there exists a $g \in \operatorname{Hom}(B, A)$ going in the other direction such that what?
\end{speech}

\begin{student}[Student Answer]
The two compositions give the identity.
\end{student}

\begin{speech}[01:21:22 - 01:22:08]
Yeah, there exists a $g$ in the other direction such that:
\[
f \circ g = \operatorname{id}_B \quad \text{and} \quad g \circ f = \operatorname{id}_A.
\]
That's the definition. That's a purely category-theoretic definition of isomorphism.

Once I tell you the category of groups or rings, I tell you the homs, the identity, and this composition property, that already tells you what the notion of isomorphism is.
\end{speech}

\begin{content}[Categorical Isomorphisms]
\begin{definition}[Isomorphism in a Category]\label[definition]{def:categorical-isomorphism}
Let $\mathcal{C}$ be a category, and let $A, B \in \operatorname{Obj}(\mathcal{C})$. A morphism $f \in \operatorname{Hom}_{\mathcal{C}}(A, B)$ is called an \newterm{isomorphism} if there exists a morphism $g \in \operatorname{Hom}_{\mathcal{C}}(B, A)$ such that:
\begin{equation}\label{eq:categorical-isom-def}
g \circ f = \operatorname{id}_A \quad \text{and} \quad f \circ g = \operatorname{id}_B.
\end{equation}
The morphism $g$ is uniquely determined and is called the \newterm{inverse morphism} of $f$, denoted $g = f^{-1}$.

Two objects $A$ and $B$ are said to be \newterm{isomorphic} in $\mathcal{C}$ (written $A \cong B$) if there exists an isomorphism $f \in \operatorname{Hom}_{\mathcal{C}}(A, B)$.
\end{definition}

\begin{center}
\begin{tikzpicture}[scale=1.5]
    \node (A) at (0,0) [circle, draw, fill=MidnightBlue!10, thick, font=\bfseries] {$A$};
    \node (B) at (3.5,0) [circle, draw, fill=MidnightBlue!10, thick, font=\bfseries] {$B$};
    
    \draw[->, thick, RoyalBlue] (A) to[bend left=30] node[above] {$f$} (B);
    \draw[->, thick, BurntOrange] (B) to[bend left=30] node[below] {$g = f^{-1}$} (A);
    
    \draw[->, thick, gray!60] (A) to[out=150, in=210, looseness=3.5] node[left] {$\operatorname{id}_A = g \circ f$} (A);
    \draw[->, thick, gray!60] (B) to[out=30, in=-30, looseness=3.5] node[right] {$\operatorname{id}_B = f \circ g$} (B);
\end{tikzpicture}
\end{center}

\begin{steps}[Arrow-Theoretic Universality of Isomorphisms]
Notice that \cref{def:categorical-isomorphism} makes \textbf{no reference to internal set elements, injectivity, or surjectivity}. Bijectivity is replaced purely by the existence of a two-sided structural inverse arrow. In concrete categories:
\begin{itemize}
    \item In $\mathbf{Grp}$, categorical isomorphisms are precisely bijective group homomorphisms.
    \item In $\mathbf{Ring}$, categorical isomorphisms are bijective ring homomorphisms.
    \item In $\mathbf{Vect}_K$, categorical isomorphisms are bijective linear transformations (invertible matrices).
\end{itemize}
\end{steps}
\end{content}

\begin{insight}[The Power of Category-Theoretic Abstraction]
Category theory unifies mathematics by defining concepts purely in terms of external interactions (\newterm{morphisms}) between systems rather than internal set elements. The definition of an isomorphism $g \circ f = \operatorname{id}_A$ and $f \circ g = \operatorname{id}_B$ is universally identical across group theory, ring theory, linear algebra, topology, and algebraic geometry.
\end{insight}

\section{Functors Between Categories}

\subsection{Definition of a Covariant Functor}

\begin{speech}[01:22:08 - 01:23:55]
There's a notion of functor.

A functor is a notion of a map from one category to another. Functors compare two different categories:
\[
F \colon \mathbf{Cat} \to \widehat{\mathbf{Cat}}.
\]

And now what's the data of the functor? 

For every object $A$ in the first category, I have to give you an object $F(A)$ in the second category. That's part of the data of the functor.
\end{speech}

\begin{speech}[01:23:55 - 01:25:28]
That's not enough: whenever I have an arrow $A \xrightarrow{f} B$ on the $\mathbf{Cat}$ side, I also have to tell you how this hom goes to a hom between these on the other side:
\[
F(f) \colon F(A) \to F(B).
\]

And it has to respect all of the structure, which means this rule must take the identity to the identities:
\[
F(\operatorname{id}_A) = \operatorname{id}_{F(A)},
\]
and it must respect composition:
\[
F(g \circ f) = F(g) \circ F(f).
\]
That's the notion of a functor.
\end{speech}

\begin{content}[Definition of a Covariant Functor]
\begin{definition}[Covariant Functor]\label[definition]{def:functor}
Let $\mathcal{C}$ and $\mathcal{D}$ be categories. A \newterm{(covariant) functor} $F \colon \mathcal{C} \to \mathcal{D}$ consists of:
\begin{enumerate}[label=\textbf{(\arabic*)}]
    \item An \textbf{Object Mapping}: Assigning to each object $A \in \operatorname{Obj}(\mathcal{C})$ an object $F(A) \in \operatorname{Obj}(\mathcal{D})$.
    \item A \textbf{Morphism Mapping}: Assigning to each morphism $f \in \operatorname{Hom}_{\mathcal{C}}(A, B)$ a morphism:
    \begin{equation}\label{eq:functor-morphism-map}
    F(f) \in \operatorname{Hom}_{\mathcal{D}}\bigl(F(A), F(B)\bigr).
    \end{equation}
\end{enumerate}
Such that the following two functoriality axioms are satisfied:
\begin{itemize}
    \item \textbf{Preservation of Identities:} For every object $A \in \operatorname{Obj}(\mathcal{C})$,
    \begin{equation}\label{eq:functor-identity}
    F(\operatorname{id}_A) = \operatorname{id}_{F(A)}.
    \end{equation}
    \item \textbf{Preservation of Composition:} For all composable morphisms $A \xrightarrow{f} B \xrightarrow{g} C$ in $\mathcal{C}$,
    \begin{equation}\label{eq:functor-composition}
    F(g \circ f) = F(g) \circ F(f).
    \end{equation}
\end{itemize}
\end{definition}

\begin{center}
\begin{tikzpicture}[scale=1.3]
    % Category C Box
    \draw[thick, draw=MidnightBlue!40, fill=MidnightBlue!5, rounded corners] (-0.8,-0.8) rectangle (2.8,1.8);
    \node[MidnightBlue, font=\bfseries] at (-0.3, 1.5) {$\mathcal{C}$};
    \node (A) at (0,0.5) [circle, draw, fill=MidnightBlue!10, thick] {$A$};
    \node (B) at (2.0,0.5) [circle, draw, fill=MidnightBlue!10, thick] {$B$};
    \draw[->, thick, RoyalBlue] (A) -- node[above] {$f$} (B);
    
    % Category D Box
    \draw[thick, draw=ForestGreen!40, fill=ForestGreen!5, rounded corners] (4.7,-0.8) rectangle (8.8,1.8);
    \node[ForestGreen, font=\bfseries] at (5.2, 1.5) {$\mathcal{D}$};
    \node (FA) at (5.5,0.5) [circle, draw, fill=ForestGreen!10, thick] {$F(A)$};
    \node (FB) at (8.0,0.5) [circle, draw, fill=ForestGreen!10, thick] {$F(B)$};
    \draw[->, thick, ForestGreen] (FA) -- node[above] {$F(f)$} (FB);
    
    % Functor Arrow
    \draw[->, ultra thick, BurntOrange] (3.1, 0.5) -- node[above] {$F$} (4.4, 0.5);
\end{tikzpicture}
\end{center}
\end{content}

\subsection{Examples of Functors: The Forgetful Functor and the Double Dual Functor}

\begin{speech}[01:25:28 - 01:28:46]
What's an example of a functor? There's one that's been staring us in the face the whole time: if I take the category of rings, there's a way to get to the category of groups. It's just taking the additive structure of the ring. That's a rule to go from rings to groups, and that is a functor.

The first example of a functor is: take the category of rings $\mathbf{Ring}$ and the category of groups $\mathbf{Grp}$. There's a functor $F \colon \mathbf{Ring} \to \mathbf{Grp}$ defined by: if I take a ring $(R, +, \cdot, 0, 1)$, I associate to it the additive group $(R, +, 0)$.

And then if I have a ring homomorphism on the ring side, it also gives you an additive group homomorphism on the group side. That's the forgetful functor.
\end{speech}

\begin{content}[The Additive Forgetful Functor from Rings to Groups]
\begin{example}[The Additive Functor $\mathbf{Ring} \to \mathbf{Grp}$]\label[example]{ex:additive-forgetful-functor}
There is a canonical \newterm{forgetful functor} $F \colon \mathbf{Ring} \to \mathbf{Grp}$ mapping rings to their underlying additive abelian groups:
\begin{itemize}
    \item \textbf{On Objects:} For each ring $(R, +, \cdot, 0, 1) \in \operatorname{Obj}(\mathbf{Ring})$,
    \begin{equation}\label{eq:forgetful-obj}
    F\bigl((R, +, \cdot, 0, 1)\bigr) := (R, +, 0) \in \operatorname{Obj}(\mathbf{Grp}).
    \end{equation}
    \item \textbf{On Morphisms:} For any ring homomorphism $f \colon R \to \hat{R}$,
    \begin{equation}\label{eq:forgetful-mor}
    F(f) := f \colon (R, +, 0) \to (\hat{R}, +, \hat{0}).
    \end{equation}
    Since a ring homomorphism satisfies $f(a + b) = f(a) + f(b)$ and $f(0) = \hat{0}$, $F(f)$ is automatically a group homomorphism.
    \item \textbf{Functoriality:} $F(\operatorname{id}_R) = \operatorname{id}_{(R, +)}$ and $F(g \circ f) = g \circ f = F(g) \circ F(f)$ hold trivially.
\end{itemize}
\end{example}
\end{content}

\begin{speech}[01:28:46 - 01:30:08]
A more interesting example that you learned: if I take the category of complex vector spaces $\mathbf{Vect}_{\mathbb{C}}$, there's a very nice functor here to itself. 

This functor takes $V$ to its double dual $(V^*)^*$. 

If you take the dual $V^*$, it reverses the arrows (contravariant). But the double dual $(V^*)^*$ doesn't switch the arrows; it goes forward (covariant). Every vector space, you can associate its double dual.

And if you have a linear map $V \xrightarrow{T} W$, then the dual gives $W^* \xrightarrow{T^*} V^*$, and doing it again gives $(V^*)^* \xrightarrow{(T^*)^*} (W^*)^*$. 

So this functor associates $T \mapsto (T^*)^*$.
\end{speech}

% [TIMESTAMP-SUSPECT: inspect manually — original: 01:30:08 - 01:30:33]
\begin{speech}[01:30:08 - 01:30:33]
\inlinemetanote{A student asks a question about vector space dimensions}
\end{speech}

\begin{student}[Student Question]
Is the dimension of vector space only finite or can it also be infinite?
\end{student}

\begin{speech}[01:30:08 - 01:30:33]
I didn't say anything. You could, if you want, say you want the category of finite-dimensional vector spaces. But the way I said it, it is anything.

For finite-dimensional vector spaces it happens the double dual is isomorphic. I'm just saying there's a functor giving a vector space its double dual.

Usually in linear algebra you spend all your time dualing and double dualing, right? 

Okay, good. That's it. \inlinemetanote{students applaud / lecturer concludes the lecture}
\end{speech}

\begin{content}[The Double Dual Endofunctor on Vector Spaces]
\begin{example}[The Double Dual Functor on $\mathbf{Vect}_{\mathbb{C}}$]\label[example]{ex:double-dual-functor}
Let $\mathbf{Vect}_{\mathbb{C}}$ be the category of complex vector spaces and $\mathbb{C}$-linear maps. The \newterm{double dual operation} defines an \newterm{endofunctor} (a functor from a category to itself):
\[
(-)^{**} \colon \mathbf{Vect}_{\mathbb{C}} \to \mathbf{Vect}_{\mathbb{C}}
\]
defined by:
\begin{enumerate}[label=\textbf{(\arabic*)}]
    \item \textbf{On Objects:} For each vector space $V$,
    \begin{equation}\label{eq:double-dual-obj}
    V \mapsto V^{**} := (V^*)^* = \operatorname{Hom}_{\mathbb{C}}(V^*, \mathbb{C}).
    \end{equation}
    \item \textbf{On Morphisms:} For any linear map $T \colon V \to W$, the single dual map $T^* \colon W^* \to V^*$ reverses arrows (acting \newterm{contravariantly}: $T^*(\psi) = \psi \circ T$). Taking the dual a second time restores the original arrow direction, yielding the \newterm{double dual map}:
    \begin{equation}\label{eq:double-dual-mor}
    T^{**} := (T^*)^* \colon V^{**} \to W^{**}, \quad T^{**}(\xi) := \xi \circ T^*.
    \end{equation}
\end{enumerate}
\end{example}

\begin{center}
\begin{tikzpicture}[scale=1.4]
    \node (V) at (0, 2.0) [font=\bfseries] {$V$};
    \node (W) at (3.5, 2.0) [font=\bfseries] {$W$};
    \draw[->, thick, RoyalBlue] (V) -- node[above] {$T$} (W);
    
    \node (Wstar) at (3.5, 1.0) [font=\bfseries] {$W^*$};
    \node (Vstar) at (0, 1.0) [font=\bfseries] {$V^*$};
    \draw[->, thick, BrickRed] (Wstar) -- node[above] {$T^* \text{ (Dual map / Contravariant)}$} (Vstar);
    
    \node (Vstarstar) at (0, 0) [font=\bfseries] {$V^{**}$};
    \node (Wstarstar) at (3.5, 0) [font=\bfseries] {$W^{**}$};
    \draw[->, thick, ForestGreen] (Vstarstar) -- node[above] {$T^{**} = (T^*)^* \text{ (Double dual / Covariant)}$} (Wstarstar);
    
    \draw[|->, dashed, gray!70] (1.75, 1.8) -- (1.75, 1.2);
    \draw[|->, dashed, gray!70] (1.75, 0.8) -- (1.75, 0.2);
\end{tikzpicture}
\end{center}

\begin{steps}[Functoriality of Double Duality]
The composition of two arrow-reversing (contravariant) operations produces an arrow-preserving (covariant) functor:
\begin{itemize}
    \item \textbf{Preservation of Identities:} $(\operatorname{id}_V)^{**} = (\operatorname{id}_{V^*})^* = \operatorname{id}_{V^{**}}$.
    \item \textbf{Preservation of Composition:} For $U \xrightarrow{S} V \xrightarrow{T} W$, the dual satisfies $(T \circ S)^* = S^* \circ T^*$. Taking the dual again yields:
    \[
    (T \circ S)^{**} = (S^* \circ T^*)^* = (T^*)^* \circ (S^*)^* = T^{**} \circ S^{**}.
    \]
\end{itemize}
This holds for all vector spaces regardless of dimension. When $\dim(V) < \infty$, the canonical evaluation map $\operatorname{ev} \colon V \to V^{**}$, $v \mapsto (\psi \mapsto \psi(v))$ is an isomorphism, making $V \cong V^{**}$ a natural isomorphism of functors.
\end{steps}
\end{content}

\begin{overview}[Summary of Key Concepts]
This lecture established foundational concepts across modern abstract algebra and category theory:
\begin{itemize}
    \item \textbf{Groups, Rings, and Fields:} Reviewed core definitions, proving the uniqueness of identity elements and inverses. Explored the exponential isomorphism $(\mathbb{R}, +) \cong (\mathbb{R}_{>0}^*, \cdot)$ and proved $(\mathbb{R}^*, \cdot) \not\cong (\mathbb{R}_{>0}^*, \cdot)$ via torsion invariants.
    \item \textbf{Structure-Preserving Maps:} Formalized group and ring homomorphisms/isomorphisms. Analyzed how symmetry breaks between rings and groups: for group homomorphisms, kernels and images are subgroups; for ring homomorphisms, the kernel is an \emph{ideal} while the image is a \emph{subring}.
    \item \textbf{Substructures and Ideals:} Investigated subrings and ideals of $\mathbb{Z}$ (proving that $\mathbb{Z}$ is a Principal Ideal Domain where every ideal is $n\mathbb{Z}$) and smooth function rings $C^\infty(\mathbb{R})$ (where compactly supported functions $C_c^\infty(\mathbb{R})$ form an ideal).
    \item \textbf{Category Theory:} Introduced categories $\mathcal{C}$ (objects, hom-sets, composition, identity laws) and arrows-only definitions of categorical isomorphisms. Defined covariant functors $F \colon \mathcal{C} \to \mathcal{D}$ with canonical examples: the additive forgetful functor $\mathbf{Ring} \to \mathbf{Grp}$ and the double dual endofunctor $(-)^{**} \colon \mathbf{Vect}_{\mathbb{C}} \to \mathbf{Vect}_{\mathbb{C}}$.
\end{itemize}
\end{overview}
% ==========================================
% LatexRefinement Step Output: step4-2026-09-23-wednesday-week2-algebra-1-offset-final.tex
% Model: gemini-3.6-flash
% Temperature: 0.35
% TopP: 0.9
% TopK: 10
% MaxOutputTokens: 65535
% ThinkingBudget: 4096
% ThinkingLevel: HIGH
% Prompt Tokens: 52’384
% Candidates Tokens: 20’975 (inkl. Thinking Tokens)
% Cached Tokens: 0
% Processed on: 2026-10-03 21:34:13
% ==========================================

% Primary Language: English
% End of the video: 00:42:33

\lecturechapter[2]{Wednesday}{23 Sep}{23 September 2026}{Category Theory Basics, Integral Domains, and Groups of Units}
% Old chapter title: Category Theory Basics \& Integral Domains

\begin{overview}[Lecture Overview \& Context]
This lecture concludes the foundational discussion of category theory before returning to core ring theory and modular arithmetic:
\begin{itemize}
    \item \textbf{Category Theory Basics:} Abstract definitions of objects, morphism sets (hom-sets), composition laws, identity morphisms, isomorphisms, monomorphisms (left-cancellation / injectivity), and epimorphisms (right-cancellation / surjectivity) along with categorical duality.
    \item \textbf{Integral Domains:} Zero divisors, formal definition of an integral domain (including the mandatory exclusion of the zero ring), and the proof that $\mathbb{Z}/m\mathbb{Z}$ is an integral domain if and only if $m$ is prime.
    \item \textbf{Divisibility and Units:} Divisibility in arbitrary commutative rings, units, and precise linguistic distinctions (\qt{the units}, \qt{the unit}, \qt{a unit}).
    \item \textbf{Groups of Units \& Modular Isomorphisms:} Proof that the set of units $U(R)$ forms a multiplicative group, characterization of $U(\mathbb{Z}/m\mathbb{Z})$ via coprimality, and explicit proof that $U(\mathbb{Z}/7\mathbb{Z}) \cong (\mathbb{Z}/6\mathbb{Z}, +)$ as a cyclic group generated by $\bar{3}$.
\end{itemize}
\end{overview}

\section{Administrative Remarks}

\begin{speech}[00:00:00 - 00:01:26]
Okay, so we start. So, I wanted to make this comment: there were several questions about the course. The first one is if you want to watch the lectures, you have to be enrolled. So, if you're not enrolled, it will stop you. And there's not much I can do about that, or if there is, I don't know what to do, but the idea is you should enroll in the course if you want to watch the lectures.

Other people had some difficulty uploading the problem sets. Maybe you have to be enrolled in the course for that, too --- I don't know ---, but many people were successful, so if you have difficulty, find someone who's successful, because I don't know how to solve the problem in another way.

Okay, a question that's more about the actual course was that every week there are these suggested pages to read in the book, and the question is: how do they relate to the course? That's a good question.

And the answer is that, roughly speaking, I do that so you can find the parallel material on which I'm lecturing in the book. And I don't always follow exactly the order there, I don't follow exactly the notation, and I don't follow exactly the text of the definitions. But I do something that's logically equivalent, so it can be helpful. As far as matters like the exam, I will only test you on material that has appeared in the lecture.

Okay? So, this is my best statement about this so far. Okay.
\end{speech}

\section{Category Theory: Abstract Definitions}

\subsection{Definition and Axioms of a Category}

\begin{speech}[00:01:26 - 00:02:50]
So, the last topic yesterday was categories, and I wanted to say a few things.

When I had this discussion about categories yesterday, it was mixed in with all the examples, so I want to just very quickly give the pure abstract definition and two small points. But this course is not about category theory. It's a topic here because it's natural to introduce when you talk about these algebraic objects, so, anyway, the course is not about category theory.

So, a category has two types of data: it has objects, as we said, and it has homs. So, $A$ can be an object, and for homs, if you pick any two objects, you get a hom space, sometimes called morphisms.

And the abstract properties of this hom space --- I said them, but I will repeat ---: first, that it's a set, that's a technical property. And that's to avoid some contradictions in foundations, which I'm not going to discuss now.

The second is that there's a composition law.

And the third is that there are identity elements --- existence of identity.

So... \inlinemetanote{writes on board} Okay, that's the data. So, is there any question about this data?
\end{speech}

\begin{content}[Constituents of a Category: Data]
\recap{Categories were introduced the day before, mixed in with examples (\cref{def:category-foundations}, \cref{def:category-full}). The lecturer now repeats the abstract definition on its own.}

\begin{definition}[Category Data]\label[definition]{def:category-data-w2wed}
A \newterm{category} $\mathcal{C}$ consists of the following mathematical data:
\begin{itemize}
    \item \textbf{Objects:} A class of objects, denoted $\operatorname{Obj}(\mathcal{C})$. We write $A \in \operatorname{Obj}(\mathcal{C})$ (or simply $A \in \mathcal{C}$) to denote that $A$ is an object of $\mathcal{C}$.
    \item \textbf{Morphism Sets (Homs):} For any pair of objects $A, B \in \operatorname{Obj}(\mathcal{C})$, a set of morphisms from $A$ to $B$, denoted:
    \[
    \operatorname{Hom}(A, B) \quad \text{or} \quad \operatorname{Hom}_{\mathcal{C}}(A, B).
    \]
    An element $f \in \operatorname{Hom}(A, B)$ is represented diagrammatically as an arrow $f \colon A \to B$.
    \item \textbf{Composition Law:} For any three objects $A, B, C \in \operatorname{Obj}(\mathcal{C})$, a binary composition mapping:
    \begin{align*}
        \circ \colon \operatorname{Hom}(B, C) \times \operatorname{Hom}(A, B) &\to \operatorname{Hom}(A, C), \\
        (g, f) &\mapsto g \circ f.
    \end{align*}
    \item \textbf{Identities:} For each object $A \in \operatorname{Obj}(\mathcal{C})$, a designated morphism called the identity morphism on $A$:
    \[
    \operatorname{Id}_A \in \operatorname{Hom}(A, A) \quad (\text{or } 1_A \in \operatorname{Hom}(A, A)).
    \]
\end{itemize}
\end{definition}

\begin{steps}[Foundational Remark on Hom-Sets]
The requirement that $\operatorname{Hom}(A, B)$ forms a legitimate set (rather than a proper class) is a standard set-theoretic convention defining a \emph{locally small category}. This prevents Russell-style self-referential paradoxes when indexing collections of arrows.
\end{steps}
\end{content}

\begin{speech}[00:02:50 - 00:03:45]
And then there's the axioms. The first one is that the composition law is associative, which means, you know, if you have three things, you can compose them in the two different ways. I hope I don't have to say more than that.

And then the second is the axioms about the identity, which is roughly: if I do this with $f$, this is $f$, and the other way around also. The identity doesn't change $f$. So, we discussed this last time.
\end{speech}

\begin{content}[Associativity and Identity Axioms]
\begin{definition}[Category Axioms]\label[definition]{ax:category-axioms-w2wed}
The data of a category $\mathcal{C}$ must satisfy the following two fundamental axioms:
\begin{enumerate}[label=\textbf{(\arabic*)}]
    \item \textbf{Associativity of Composition:} For all objects $A, B, C, D \in \operatorname{Obj}(\mathcal{C})$ and all morphisms $f \in \operatorname{Hom}(A, B)$, $g \in \operatorname{Hom}(B, C)$, and $h \in \operatorname{Hom}(C, D)$, composition is associative:
    \begin{equation}\label{eq:cat-assoc}
    (h \circ g) \circ f = h \circ (g \circ f).
    \end{equation}
    \item \textbf{Identity Axioms (Neutrality):} For any morphism $f \in \operatorname{Hom}(A, B)$, composition with the respective identity morphisms leaves $f$ unchanged:
    \begin{align}
        f \circ \operatorname{Id}_A &= f, \label{eq:cat-id-right} \\
        \operatorname{Id}_B \circ f &= f. \label{eq:cat-id-left}
    \end{align}
\end{enumerate}
\end{definition}

{\centering
\begin{tikzpicture}[scale=1.4, >=Stealth, auto]
    \node (A) at (0, 0) {$A$};
    \node (B) at (2, 1.3) {$B$};
    \node (C) at (4, 0) {$C$};

    \draw[->, thick] (A) to node {$f$} (B);
    \draw[->, thick] (B) to node {$g$} (C);
    \draw[->, thick, MidnightBlue] (A) to node[below] {$g \circ f$} (C);

    \path[->, thick] (A) edge[out=150, in=210, looseness=4] node[left] {$\operatorname{Id}_A$} (A);
    \path[->, thick] (B) edge[out=60, in=120, looseness=4] node[above] {$\operatorname{Id}_B$} (B);
    \path[->, thick] (C) edge[out=330, in=30, looseness=4] node[right] {$\operatorname{Id}_C$} (C);
\end{tikzpicture}\par}
\end{content}

\subsection{Categorical Morphisms: Isomorphisms, Monomorphisms, and Epimorphisms}

\begin{speech}[00:03:45 - 00:05:35]
So, that's the abstract definition, so there's really not that much here. And we discussed this last time, and we had lots of examples.

And so you could say: well, why do we want to do this? And as I said, there's a setting where many concepts that we talk about when we talk about groups, rings, fields, vector spaces --- they're categorical concepts, meaning that they only depend on this structure. So, it simplifies thinking to consider them all together.

So, I already told you about isomorphism. The categorical isomorphism is that if I have this arrow from $A$ to $B$ --- this picture means this: $f \in \operatorname{Hom}(A, B)$, okay? That's why in category theory, you can say this, or you can say that, and you mean the same thing --- that this $f$ is an isomorphism --- that's a definition --- if and only if there exists $g$ the other way, $g \colon B \to A$, such that the two compositions are identities.

And if you want to say which identity, you have to be careful. So, $g$ starts with $B$, so if you do $g$ first, then $f$ back, this should be the identity in $B$, and the other is the identity in $A$.

So, that's the abstract definition of isomorphism. And in all of the categories we discussed, that coincides with the notion of isomorphism in the category that you already know, like, for example, isomorphism of vector spaces.
\end{speech}

\begin{content}[Categorical Isomorphism]
\recap{The categorical isomorphism was defined the day before (\cref{def:categorical-isomorphism}); the lecturer restates it.}

\begin{definition}[Isomorphism]\label[definition]{def:categorical-isomorphism-w2wed}
Let $\mathcal{C}$ be a category and let $f \in \operatorname{Hom}(A, B)$ (written $f \colon A \to B$). The morphism $f$ is an \newterm{isomorphism} if and only if there exists a morphism $g \in \operatorname{Hom}(B, A)$ (written $g \colon B \to A$) such that:
\begin{align}
    f \circ g &= \operatorname{Id}_B, \label{eq:isom-right-inverse} \\
    g \circ f &= \operatorname{Id}_A. \label{eq:isom-left-inverse}
\end{align}
If such a morphism $g$ exists, it is uniquely determined, called the \newterm{inverse} of $f$, and denoted $f^{-1} := g$. Two objects $A$ and $B$ are said to be \newterm{isomorphic} (written $A \cong B$) if there exists an isomorphism between them.
\end{definition}

{\centering
\begin{tikzpicture}[scale=1.5, >=Stealth, auto]
    \node (A) at (0,0) {$A$};
    \node (B) at (3,0) {$B$};

    \draw[->, thick, RoyalBlue] (A) to[bend left=25] node {$f$} (B);
    \draw[->, thick, BurntOrange] (B) to[bend left=25] node {$g = f^{-1}$} (A);

    % caption below the label of the lower arc (arc dips to y = -0.33, its label to about y = -0.55)
    \node at (1.5, -1.15) {\footnotesize $f \circ g = \operatorname{Id}_B \quad \text{and} \quad g \circ f = \operatorname{Id}_A$};
\end{tikzpicture}\par}

\begin{steps}[Consistency across Concrete Categories]
In concrete categories where objects are sets endowed with structure and morphisms are structure-preserving maps (homomorphisms):
\begin{itemize}
    \item In $\mathbf{Set}$, isomorphisms are precisely bijective functions.
    \item In $\mathbf{Grp}$, $\mathbf{Ring}$, and $\mathbf{Vect}_K$, isomorphisms are bijective homomorphisms (for groups: \cref{prop:isomorphism-inverse-characterization}).
\end{itemize}
The power of the categorical definition is that it never inspects elements $x \in A$; it characterizes invertibility purely through diagrammatic composition.
\end{steps}
\end{content}

\begin{hidden}
% timestamp: 00:05:35
% topic: Abstract Definition of Category, Axioms, and Isomorphism
% board_state: def:category-data-w2wed, ax:category-axioms-w2wed, def:categorical-isomorphism-w2wed
% next_goal: Introduce monomorphisms and epimorphisms as categorical analogues of injectivity and surjectivity
% open_loops: none
\end{hidden}

\begin{speech}[00:05:35 - 00:07:20]
There are two other notions, and that's the last thing I want to say. And this is the notion of monomorphism.

So, $f$ from $A$ to $B$ is a monomorphism --- that's a long word, that's a definition --- if for all $g_1$ arrows from another object, $X$, to $A$, and another one, $g_2$ --- so we take any two ---, then the composition $f \circ g_1 = f \circ g_2$ implies $g_1 = g_2$. I hope I wrote that clearly here, okay?

So, it's a little bit complicated to say. If I have an arrow in the category, it may or may not be a monomorphism. I'm defining a monomorphism. And the condition of being a monomorphism means that whenever I take an object $X$ and map it two ways to $A$, the composition with $f$ being equal implies they were equal to start with. Of course, if they're equal to start with, the composition is equal, but it's the other implication. That's the definition of a monomorphism.

And, you know, you have to think about this and unwind it. But once you have this definition, you can imagine: what is a monomorphism in all of those categories you know? And what do you think it is?
\end{speech}

\begin{speech}[00:07:20 - 00:08:15]
So, I leave this for you to think about, but roughly speaking, monomorphism is being injective.

Just imagine they're sets. If this is injective, and you take any two maps, you can compare them after you do $f$. So, you have to think about it, but this is the definition.

And once you take this definition, you can say: okay, I can take the category of sets and functions, and what is a monomorphism there? And the answer is: injective. That's the \qt{mono} in monomorphism.

So, this is a categorical formulation of being injective.

I don't know if you find this appealing or not, but that's what this discussion is about. Because the nice thing about this formulation is that it applies to all examples of categories. Every single one we discussed, and all the other ones you can think about, there's a notion of monomorphism.
\end{speech}

\begin{content}[Left-Cancellable Morphisms: Monomorphisms]
\begin{definition}[Monomorphism]\label[definition]{def:monomorphism}
Let $\mathcal{C}$ be a category. A morphism $f \in \operatorname{Hom}(A, B)$ (or $f \colon A \to B$) is called a \newterm{monomorphism} (often abbreviated as a \emph{monic morphism} or \emph{mono}) if it is left-cancellable: for every object $X \in \operatorname{Obj}(\mathcal{C})$ and all pairs of parallel morphisms $g_1, g_2 \in \operatorname{Hom}(X, A)$,
\begin{equation}\label{eq:monomorphism-cond}
f \circ g_1 = f \circ g_2 \implies g_1 = g_2.
\end{equation}
\end{definition}

{\centering
\begin{tikzpicture}[scale=1.5, >=Stealth, auto]
    \node (X) at (0,0) {$X$};
    \node (A) at (2.2,0) {$A$};
    \node (B) at (4.4,0) {$B$};

    \draw[->, thick] (X) to[bend left=25] node {$g_1$} (A);
    \draw[->, thick] (X) to[bend right=25] node[below] {$g_2$} (A);
    \draw[->, thick, ForestGreen] (A) to node {$f$} (B);

    % caption below the g_2 label (lower arc dips to y = -0.24, its label to about y = -0.5)
    \node at (2.2, -1.0) {\footnotesize $f \circ g_1 = f \circ g_2 \implies g_1 = g_2 \quad (\text{Categorical Injective})$};
\end{tikzpicture}\par}

\begin{steps}[Why Monomorphisms Characterize Injectivity]
In the category of sets ($\mathbf{Set}$), suppose $f \colon A \to B$ is a monomorphism. To test whether $f$ is injective on points $a_1, a_2 \in A$, let $X = \{*\}$ be a one-element set and define two maps $g_1, g_2 \colon \{*\} \to A$ by $g_1(*) = a_1$ and $g_2(*) = a_2$.
Then $(f \circ g_1)(*) = f(a_1)$ and $(f \circ g_2)(*) = f(a_2)$. Thus:
\[
f(a_1) = f(a_2) \implies f \circ g_1 = f \circ g_2 \implies g_1 = g_2 \implies a_1 = a_2.
\]
Hence, left-cancellation with respect to arbitrary test objects $X$ precisely recovers element-wise injectivity without mentioning elements.
\end{steps}
\end{content}

\begin{speech}[00:08:15 - 00:09:20]
And this one has a sibling, or a partner --- I don't know their relationship ---, which is... whenever you talk about injective, you can also talk about surjective, and the categorical term there is epimorphism.

So, you can define epimorphism. So, if you have $f \colon A \to B$, it is an epimorphism if --- it's a definition --- for all...

And this is now about this kind of feeling of mathematics: how do you think I'm going to define the opposite notion of injective? Like here we had injective by $X$ mapping to $A$. So, what do you think I should do here?

You should guess. Guessing is like one of the most important tools of a mathematician.

Yeah?
\end{speech}

\begin{student}[Student Answer]
Well, we do the same, except we also switch like we compose\dots
\end{student}

\begin{speech}[00:09:25 - 00:10:35]
This is a long answer. Here we had $X$ mapping to $A$, so the opposite of that is mapping from $B$. It's like the other side.

So, for every $B \xrightarrow[g_2]{g_1} X$, then... So, this is the opposite of that, by just flipping everything to the other side. So, for every one of these, then we can consider $f$...

In that one it was the $g$'s first, then $f$; now it's $f$ first, then the $g$'s: $g_1 \circ f = g_2 \circ f \implies g_1 = g_2$.

So, that's the categorical notion of epimorphism, and it's somehow dual to monomorphism. If you like this world, it almost makes epimorphism and monomorphism the same definition up to this flipping. And, you know, thinking about things categorically is good for such understanding.

If you look at this definition, and you look at it for the examples we know, this is surjectivity.

Okay, so that's the end of this discussion about categories. If you like it, you can go read more. As I said, this course is not a course on categories. There's a bunch of other things that I would say here, but I'm not going to, because, yeah, maybe I will at the end if I have time.
\end{speech}

\begin{content}[Right-Cancellable Morphisms: Epimorphisms]
\begin{definition}[Epimorphism]\label[definition]{def:epimorphism}
Let $\mathcal{C}$ be a category. A morphism $f \in \operatorname{Hom}(A, B)$ (or $f \colon A \to B$) is called an \newterm{epimorphism} (often abbreviated as an \emph{epic morphism} or \emph{epi}) if it is right-cancellable: for every object $X \in \operatorname{Obj}(\mathcal{C})$ and all pairs of parallel morphisms $g_1, g_2 \in \operatorname{Hom}(B, X)$,
\begin{equation}\label{eq:epimorphism-cond}
g_1 \circ f = g_2 \circ f \implies g_1 = g_2.
\end{equation}
\end{definition}

{\centering
\begin{tikzpicture}[scale=1.5, >=Stealth, auto]
    \node (A) at (0,0) {$A$};
    \node (B) at (2.2,0) {$B$};
    \node (X) at (4.4,0) {$X$};

    \draw[->, thick, BrickRed] (A) to node {$f$} (B);
    \draw[->, thick] (B) to[bend left=25] node {$g_1$} (X);
    \draw[->, thick] (B) to[bend right=25] node[below] {$g_2$} (X);

    % caption below the g_2 label (lower arc dips to y = -0.24, its label to about y = -0.5)
    \node at (2.2, -1.0) {\footnotesize $g_1 \circ f = g_2 \circ f \implies g_1 = g_2 \quad (\text{Categorical Surjective})$};
\end{tikzpicture}\par}
\end{content}

\begin{aiNote}[Epimorphisms Need Not Be Surjective]
The lecturer says that for the examples we know, epimorphism means surjectivity. This holds in $\mathbf{Set}$, $\mathbf{Grp}$ and $\mathbf{Vect}_K$, but not in $\mathbf{Ring}$: the inclusion $\iota \colon \mathbb{Z} \hookrightarrow \mathbb{Q}$ is an epimorphism without being surjective. Indeed, if $g_1, g_2 \colon \mathbb{Q} \to R$ are ring homomorphisms with $g_1 \circ \iota = g_2 \circ \iota$, then for $b \neq 0$ the element $g_k(b)$ is invertible with inverse $g_k(\frac{1}{b})$, so $g_1(\frac{a}{b}) = g_1(a) \, g_1(b)^{-1} = g_2(a) \, g_2(b)^{-1} = g_2(\frac{a}{b})$. For monomorphisms the slogan is safe: in $\mathbf{Set}$, $\mathbf{Grp}$, $\mathbf{Ring}$ and $\mathbf{Vect}_K$ they are exactly the injective homomorphisms.
\end{aiNote}

\begin{insight}[Categorical Duality: Reversing Arrows]
Monomorphisms and epimorphisms represent our first concrete encounter with \newterm{categorical duality}. If we construct the opposite category $\mathcal{C}^{\text{op}}$ by reversing all morphism arrows ($f^{\text{op}} \colon B \to A$ for every $f \colon A \to B$), then:
\[
f \text{ is a monomorphism in } \mathcal{C} \iff f^{\text{op}} \text{ is an epimorphism in } \mathcal{C}^{\text{op}}.
\]
Every theorem or statement about monomorphisms immediately yields an identical, dual statement about epimorphisms simply by reversing the order of all composition arrows.
\end{insight}

\section{Ring Theory: Zero Divisors and Integral Domains}

\subsection{Zero Divisors and the Definition of an Integral Domain}

\begin{speech}[00:10:35 - 00:12:35]
So, let us return to rings now. This was kind of a foundational side note.

So, we have this ring $(R, +, \cdot, 0, 1)$.

And the rings here, there's kind of a different flavor of rings: there are the rings that are like $\mathbb{Z}$, where when you have things that are non-zero and you multiply them, they stay non-zero. That's one of the principles of $\mathbb{Z}$. But it's not true for rings like $\mathbb{Z}/6\mathbb{Z}$, this ring.

I can take two elements here, like the residue class of $2$ times the residue class of $3$, neither of which is zero. So, this is not zero, and that is not zero. So, I take two elements that are legitimately not zero, but when I multiply these residue classes, this is $6$, and that is $0$.

So, that's a different type of multiplication. It's a multiplication where when you multiply non-zero things, you get zero, which is unexpected from the point of view of the integers.

And this is also related to the fact that there's no hope of having a multiplicative inverse if you're mapping things to zero... Okay, so there are these two differences, and there's a definition for this:
\end{speech}

\begin{content}[Zero Divisors in Modular Arithmetic]
Let $(R, +, \cdot, 0, 1)$ be a commutative ring with unity $1 \neq 0$.

\begin{example}[Multiplication in \texorpdfstring{$\mathbb{Z}/6\mathbb{Z}$}{Z/6Z}]\label[example]{ex:z6z-zero-divisors}
In the ring of residue classes modulo $6$, denoted $\mathbb{Z}/6\mathbb{Z}$ (or $\mathbb{Z}_6$):
\begin{align*}
    \bar{2} &\neq \bar{0}, \\
    \bar{3} &\neq \bar{0}.
\end{align*}
However, their product in the ring is:
\[
\bar{2} \cdot \bar{3} = \overline{2 \cdot 3} = \bar{6} = \bar{0} \in \mathbb{Z}/6\mathbb{Z}.
\]
\end{example}

\begin{steps}[Zero Divisors vs. Invertibility]
Non-zero elements whose product is zero are called \newterm{zero divisors}. If an element $a \in R$ is a zero divisor (with $a \cdot b = 0$ for some $b \neq 0$), it can never possess a multiplicative inverse $a^{-1}$:
\[
b = 1 \cdot b = (a^{-1} \cdot a) \cdot b = a^{-1} \cdot (a \cdot b) = a^{-1} \cdot 0 = 0,
\]
which contradicts $b \neq 0$. Hence, rings containing zero divisors fail to exhibit regular arithmetic cancellation.
\end{steps}
\end{content}

\begin{speech}[00:12:35 - 00:14:05]
A ring $R$ is an integral domain... So, it's a definition of this term \qt{integral domain}. And some authors would just say it's a domain, and also if you speak about rings, very often mathematicians will say it's a domain... if whenever I have $a$ and $b$, neither of which is zero --- so maybe I'll just be really explicit --- whenever I have $a \neq 0$ and $b \neq 0$, this implies the product is not equal to zero ($a \cdot b \neq 0$).

So, that's the definition. It's a really basic definition in ring theory, and it separates this type of ring from that type of ring. Is that clear?

Fields, for example, are always domains, because you have the multiplicative inverse, if you think about it, okay?

So, we can ask...
\end{speech}

\begin{content}[Definition of an Integral Domain (Initial Board State)]
\begin{definition}[Integral Domain / Domain]\label[definition]{def:integral-domain-prelim}
A commutative ring $R$ is called an \newterm{integral domain} (or simply a \newterm{domain}) if for all $a, b \in R$:
\begin{equation}\label{eq:integral-domain-cond}
a \neq 0 \quad \text{and} \quad b \neq 0 \implies a \cdot b \neq 0.
\end{equation}
Equivalently, taking the contrapositive:
\[
a \cdot b = 0 \implies a = 0 \quad \text{or} \quad b = 0.
\]
\end{definition}

\begin{example}[Fields are Domains]\label[example]{ex:fields-are-domains}
Every field $F$ is an integral domain: if $a \neq 0$ and $a \cdot b = 0$, then multiplying by the inverse $a^{-1} \in F$ immediately yields:
\[
b = 1 \cdot b = (a^{-1} a) b = a^{-1} (a b) = a^{-1} \cdot 0 = 0.
\]
\end{example}
\end{content}

\begin{speech}[00:14:05 - 00:14:42]
So, see, I did that already there: I called it a domain, not an integral domain, because it takes a lot of syllables to say \qt{integral domain}. Formally, this is defined as an integral domain; colloquially, almost everyone says a domain.

So, these terms, monomorphism and epimorphism, have you ever heard these terms before? In which context?
\end{speech}

\begin{student}[Student Question/Answer]
Linear algebra.
\end{student}

\begin{speech}[00:14:42 - 00:14:52]
Ah, okay. And they were just used synonymously with injection and surjection?
\end{speech}

\begin{student}[Student Answer]
We were introduced to them, but then we went back to using injection and surjection.
\end{student}

\begin{speech}[00:14:52 - 00:14:58]
Okay, yeah.
\end{speech}

\begin{student}[Student Answer]
So, they were like injective and surjective homomorphisms.
\end{student}

\begin{speech}[00:14:58 - 00:15:20]
Yes, exactly, yes. Injective and surjective homomorphisms. That's precisely what this categorical thing says, right? Because in the category of groups, the arrows are homomorphisms, and when a homomorphism is a monomorphism, it's an injective homomorphism.

In my experience with undergraduates, some really like the idea of categories and really like to think in that language, and others like the examples, so either way is okay.
\end{speech}

\subsection{Characterization of \texorpdfstring{$\mathbb{Z}/m\mathbb{Z}$}{Z/mZ} as an Integral Domain}

\begin{speech}[00:15:20 - 00:16:24]
All right, so here we come to the first proposition, which is something my guess is you are already familiar with, but we do it again. So, we take this ring $\mathbb{Z}/m\mathbb{Z}$, and the claim is that this $\mathbb{Z}/m\mathbb{Z}$ is a domain if and only if --- and what am I going to write here? Yeah?
\end{speech}

\begin{student}[Student Answer]
$m$ is prime.
\end{student}

\begin{speech}[00:16:04 - 00:16:24]
$m$ is prime. $m$ is prime.

Oh, yeah, I should say one thing about domains which is important, I didn't say, and it's kind of annoying: part of the definition... I actually should put that in.
\end{speech}

\begin{speech}[00:16:24 - 00:17:15]
It's a super annoying thing to talk about, but I talk about any ring here, so my definition has to deal with the zero ring, the one that, you know, is unpleasant to think about. And you could think about: what would this definition be for the zero ring? Because for all $a$ and $b$, there are no $a$'s and $b$'s that are not equal to zero. So, formally this statement in the formal deductive logical sense is true for the zero ring.

So, it is part of the definition: a ring $R$ which is not the zero ring ($R \neq \{0\}$) is a domain. Zero rings are not allowed to be domains, by fiat, in the definition.
\end{speech}

\begin{content}[Correction: Exclusion of the Zero Ring]
The lecturer corrects \cref{def:integral-domain-prelim} by adding condition (i):

\begin{definition}[Integral Domain (Formal Standard Definition)]\label[definition]{def:integral-domain}
A commutative ring $R$ is an \newterm{integral domain} (or \newterm{domain}) if:
\begin{enumerate}[label=\textbf{(\roman*)}]
    \item $R \not\cong \{0\}$ (that is, $1_R \neq 0_R$, excluding the zero ring; cf.\ \cref{prop:zero-ring}).
    \item For all $a, b \in R$:
    \[
    a \neq 0 \quad \text{and} \quad b \neq 0 \implies a \cdot b \neq 0.
    \]
\end{enumerate}
\end{definition}

\begin{steps}[Why the Zero Ring Satisfies the Conditional Vacuously]
In the zero ring $R = \{0\}$, the only element is $0$. Therefore, there exist no non-zero elements $a, b \in R$. The premise $(a \neq 0 \land b \neq 0)$ is vacuously false for all elements, making the logical implication $(a \neq 0 \land b \neq 0 \implies ab \neq 0)$ vacuously true! To avoid having the degenerate zero ring classified as a domain, mathematicians add $R \neq 0$ explicitly by fiat.
\end{steps}
\end{content}

\begin{content}[Proposition: When Residue Class Rings are Domains]
\begin{proposition}\label[proposition]{prop:z-mod-m-domain}
Let $m \in \mathbb{Z}_{>0}$ be a positive integer. The residue class ring $(\mathbb{Z}/m\mathbb{Z}, +, \cdot, \bar{0}, \bar{1})$ is an integral domain if and only if $m$ is a prime number:
\begin{equation}\label{eq:z-mod-m-equiv}
\mathbb{Z}/m\mathbb{Z} \text{ is an integral domain} \iff m \text{ is prime}.
\end{equation}
\end{proposition}
\end{content}

\begin{hidden}
% timestamp: 00:17:15
% topic: Proof of Proposition: Z/mZ is an integral domain iff m is prime
% board_state: def:integral-domain, prop:z-mod-m-domain
% next_goal: Direction 1 (m not prime implies not domain) and Direction 2 (m prime implies domain via Euclid's lemma)
% open_loops: none
\end{hidden}

\begin{proofWrapper}[Proof of Proposition \ref{prop:z-mod-m-domain}]
\begin{speech}[00:17:15 - 00:19:24]
Okay, $\mathbb{Z}/m\mathbb{Z}$ is a domain if and only if $m$ is prime. Okay, so let's try to prove this. So, proof.

This is like one of the issues I said: this issue about the zero ring is a little bit problematic everywhere in the statements, because for any statement you write, \qt{let $R$ be a ring}, you have to think --- if you want that statement to be literally true in the mathematical sense from the definitions in this class ---, you have to think about what that statement could mean for the zero ring and whether it's true. And as I said, I don't view it as a deep or really important thing to do, but in terms of bookkeeping and saying things correctly, you have to say it...

Okay, so let's try to prove this: $\mathbb{Z}/m\mathbb{Z}$ is a domain if and only if $m$ is prime.

So, \textbf{(1)} suppose $m$ is not prime.

Here we have $m$ as some positive integer ($m \in \mathbb{Z}_{>0}$). Suppose $m$ is not prime. Then this means that $m$ admits a factorization into $a \cdot b$, where $a$ is strictly greater than $1$ and strictly less than $m$, and $b$ is strictly like this ($1 < b < m$). That's what it means. If it's not prime, I can factor it, and if I can factor it, I break it into two smaller pieces.

And then if I look at this equation in residue classes, this equals that, which is equal to zero in $\mathbb{Z}/m\mathbb{Z}$; and this shows that that's not zero ($\bar{a} \neq \bar{0}$), and this shows that that's not zero ($\bar{b} \neq \bar{0}$).

Okay? So, this says that if $m$ is not prime, there is no way for it to be a domain, because if it's not prime, it factors, and those factors precisely give the violation of that property. All right?
\end{speech}

\begin{content}[Proof of Proposition: Direction 1 (Non-Prime Case)]
\textbf{Direction 1: \texorpdfstring{$m$}{m} Not Prime \texorpdfstring{$\implies \mathbb{Z}/m\mathbb{Z}$}{=> Z/mZ} Not a Domain.}

Suppose $m$ is not prime. If $m = 1$, then $\mathbb{Z}/1\mathbb{Z} = \{\bar{0}\}$ is the zero ring, which is not a domain by condition (i) of \cref{def:integral-domain} (see also \cref{rem:boundary-cases-mod-rings}). Otherwise $m > 1$ is composite and admits a proper non-trivial factorization in $\mathbb{Z}$:
\begin{align}
    m &= a \cdot b, \label{eq:composite-m-factor} \\
    1 &< a < m, \nonumber \\
    1 &< b < m. \nonumber
\end{align}
Passing to the residue classes in the quotient ring $\mathbb{Z}/m\mathbb{Z}$:
\begin{equation}\label{eq:prod-zero-divisors}
\bar{a} \cdot \bar{b} = \overline{a \cdot b} = \bar{m} = \bar{0} \in \mathbb{Z}/m\mathbb{Z}.
\end{equation}
However, because $1 < a < m$ and $1 < b < m$, neither $a$ nor $b$ is divisible by $m$. Therefore:
\begin{equation}\label{eq:nonzero-factors}
\bar{a} \neq \bar{0} \quad \text{and} \quad \bar{b} \neq \bar{0}.
\end{equation}
The existence of non-zero elements $\bar{a}, \bar{b} \in \mathbb{Z}/m\mathbb{Z}$ satisfying $\bar{a} \cdot \bar{b} = \bar{0}$ demonstrates that $\mathbb{Z}/m\mathbb{Z}$ possesses zero divisors, so $\mathbb{Z}/m\mathbb{Z}$ is not an integral domain.
\end{content}

\begin{speech}[00:19:24 - 00:21:50]
So, then the other thing is suppose it is prime. Suppose $m = p$ is prime.

Then, we let $a$ and $b$ be integers such that $\bar{a} \cdot \bar{b} = \bar{0}$ in $\mathbb{Z}/p\mathbb{Z}$. All right?

What are we going to do here? This means that if I translate this equation on residue classes to integers, this means that $a \cdot b$ is equal to some integer $d \cdot p$.

So, I am using this notation and arithmetic in $\mathbb{Z}/m\mathbb{Z}$ freely here, because I assume you've seen this before. Is this true? If it's not, you should tell me. All right. If you haven't seen it, you should study it.

So, if I take these two residue classes and their product is zero, that exactly means that the product is divisible by $p$. Is that clear?

And then this means that $p$ divides $a \cdot b$, because $a \cdot b = d \cdot p$.

And then there's the principle of Euclid we discussed in the first week that says either $p$ divides $a$ or $p$ divides $b$. It wasn't such a long time ago, it was just last week.

And this means if $p \mid a$, then that means $\bar{a} = \bar{0}$; and if $p \mid b$, that means $\bar{b} = \bar{0}$.

So, we've proven that if the product is zero in $\mathbb{Z}/p\mathbb{Z}$, the only way that happens is one of them is zero, or the other is zero. And that's logically equivalent to this thing: if they were both not zero, then the product is not zero. Or the logical equivalent, which is the contrapositive, is that if the product is zero, then one of the factors must be zero, which we've established. So, this proposition is proven.
\end{speech}

\begin{content}[Proof of Proposition: Direction 2 (Prime Case via Euclid's Lemma)]
\textbf{Direction 2: \texorpdfstring{$m = p$}{m=p} Prime \texorpdfstring{$\implies \mathbb{Z}/p\mathbb{Z}$}{=> Z/pZ} Is a Domain.}

Suppose $m = p$ is a prime number. Let $\bar{a}, \bar{b} \in \mathbb{Z}/p\mathbb{Z}$ and assume that their product vanishes:
\begin{equation}\label{eq:prime-prod-zero}
\bar{a} \cdot \bar{b} = \bar{0} \in \mathbb{Z}/p\mathbb{Z}.
\end{equation}
By definition of modular multiplication, this implies $\overline{a \cdot b} = \bar{0}$, meaning that $a \cdot b$ is congruent to $0 \pmod p$. Hence, there exists some integer $d \in \mathbb{Z}$ such that:
\begin{equation}\label{eq:divisible-by-p}
a \cdot b = d \cdot p \implies p \mid (a \cdot b).
\end{equation}
By Euclid's Lemma (which holds whenever $p$ is prime):
\begin{equation}\label{eq:euclids-lemma}
p \mid (a \cdot b) \implies p \mid a \quad \text{or} \quad p \mid b.
\end{equation}
Translating each case back into modular equivalence in $\mathbb{Z}/p\mathbb{Z}$:
\begin{itemize}
    \item If $p \mid a$, then $\bar{a} = \bar{0} \in \mathbb{Z}/p\mathbb{Z}$.
    \item If $p \mid b$, then $\bar{b} = \bar{0} \in \mathbb{Z}/p\mathbb{Z}$.
\end{itemize}
Thus, $\bar{a} \cdot \bar{b} = \bar{0}$ implies $\bar{a} = \bar{0}$ or $\bar{b} = \bar{0}$. This establishes that $\mathbb{Z}/p\mathbb{Z}$ has no zero divisors. Because $p \ge 2$, $\mathbb{Z}/p\mathbb{Z}$ contains at least two distinct elements ($\bar{0} \neq \bar{1}$), so $\mathbb{Z}/p\mathbb{Z}$ is not the zero ring. Therefore, $\mathbb{Z}/p\mathbb{Z}$ is an integral domain.
\end{content}

\begin{speech}[00:21:50 - 00:23:19]
So, you could say again in this proposition, when someone writes this, I here put $m$ to be a positive number ($m \in \mathbb{Z}_{>0}$), and now it matches carefully, but these boundary cases... you know, there's some aspect of this foundational part where you have to be careful yourself about boundary cases.

So, if I put $m = 1$, then I get the zero ring, but $1$ is not considered a prime number, so it matches. You could say: what happens if I put zero? It's not here; we don't think about primality of zero, so zero is left out of this, but if $m$ is zero, this is $\mathbb{Z}$, and it is a domain. Okay? All right.

\inlinemetanote{erases the blackboard}

Sometimes being too clear is actually more confusing in lectures. So, there's this kind of principle of lecturing where you lie about all the little details, and the lectures come out cleaner if you do that. And it's not clear didactically that that isn't a better way to do it.

There's a certain mathematician that takes pleasure in these edge cases, and you can tell that I'm not one of them.
\end{speech}
\end{proofWrapper}

\begin{content}[Pedagogical Analysis of Boundary Cases]
\begin{remark}[Boundary Cases: \texorpdfstring{$m = 1$}{m=1} and \texorpdfstring{$m = 0$}{m=0}]\label[remark]{rem:boundary-cases-mod-rings}
The proposition specifies $m \in \mathbb{Z}_{>0}$. Let us examine the boundary cases:
\begin{itemize}
    \item \textbf{Case $m = 1$:} The ideal is $1\mathbb{Z} = \mathbb{Z}$. The quotient ring is the zero ring:
    \[
    \mathbb{Z}/1\mathbb{Z} \cong \{0\},
    \]
    where $\bar{1} = \bar{0}$. Since the definition of an integral domain explicitly requires $R \neq \{0\}$, $\mathbb{Z}/1\mathbb{Z}$ is \emph{not} an integral domain. This perfectly matches the number-theoretic convention that $1$ is not a prime number.
    \item \textbf{Case $m = 0$:} The zero ideal is $0\mathbb{Z} = \{0\}$. The quotient ring is canonically isomorphic to the integers:
    \[
    \mathbb{Z}/0\mathbb{Z} \cong \mathbb{Z}.
    \]
    The integers $\mathbb{Z}$ are the prototypical integral domain. In modern commutative algebra, $\{0\} \subset \mathbb{Z}$ is a prime ideal precisely because $\mathbb{Z}/\{0\} \cong \mathbb{Z}$ is an integral domain!
\end{itemize}
\end{remark}
\end{content}

\section{Divisibility and Units in Commutative Rings}

\subsection{Divisibility in Arbitrary Rings}

\begin{speech}[00:23:19 - 00:24:29]
Okay, so something more fun is we get to talk about units, and this is actually pretty fun.

This has to do with every ring, so let's take a ring $R$.

And this is now related to last week; in fact, it was the point of the discussion last week. So, there we discussed the ring $\mathbb{Z}$, but now you can take any ring.

So, the definition here is that if we have elements $a$ and $b$ in the ring ($a, b \in R$), then we say --- this is the definition of the term ---, we say \qt{a divides b}. What this means is there exists some element of the ring, $c$ ($c \in R$), such that $a \cdot c = b$.

Okay, that's just the actual definition of division, and it's the same definition as with the integers.

So, this is a notion of division for any ring. Okay?
\end{speech}

\begin{content}[Divisibility Relations]
Let $R$ be a ring (assumed commutative with unity $1 \neq 0$).

\begin{definition}[Divisibility in a Ring]\label[definition]{def:divisibility-ring}
Let $a, b \in R$. We say that \newterm{$a$ divides $b$}, denoted:
\begin{equation}\label{eq:divisibility-notation}
a \mid b,
\end{equation}
if there exists an element $c \in R$ such that:
\begin{equation}\label{eq:divisibility-def}
a \cdot c = b.
\end{equation}
In this case, we also say that $b$ is a \newterm{multiple} of $a$, or that $a$ is a \newterm{divisor} of $b$.
\end{definition}

\begin{steps}[Connection with Principal Ideals]
In ideal-theoretic language, $a \mid b$ in a commutative ring $R$ is equivalent to the statement that the principal ideal generated by $b$ is contained in the principal ideal generated by $a$:
\[
a \mid b \iff b \in aR \iff (b) \subseteq (a).
\]
\end{steps}
\end{content}

\subsection{The Definition of a Unit}

\begin{speech}[00:24:29 - 00:26:22]
Now, $u$ in $R$ ($u \in R$) is a unit --- that's the definition of a unit.

The definition is if $u$ divides $1$ ($u \mid 1$). Oh, I should say that also: \qt{a divides b}, that's how you say it in English, and the symbol for that is $a \mid b$. Okay?

So, here this is now a definition: an element is a unit if $u$ divides $1$ ($u \mid 1$). And that is equivalent --- this is just equivalent --- to there existing a $w$ ($w \in R$) such that $u \cdot w = 1$. I never have to worry about the order here because I said that for the rings in this course, unless I tell you explicitly otherwise, we're always in commutative rings. Okay?

So, this is the definition, this is unwinding division, and if I want to unwind this using English --- so this is the definition of division, and this is just English ---, we say that if and only if $u$ has a multiplicative inverse in $R$.

Okay? So, this is a really fundamental definition: a unit in our ring. So, are you happy with this? That's the definition, this is just translation, and that's what you say: $u$ has a multiplicative inverse, $w$ is the multiplicative inverse. So, is there any problem with this?
\end{speech}

\begin{content}[Invertible Elements]
\begin{definition}[Unit in a Ring]\label[definition]{def:unit-ring}
Let $R$ be a commutative ring with unity $1$. An element $u \in R$ is called a \newterm{unit} if $u$ divides $1$:
\begin{equation}\label{eq:unit-div-1}
u \mid 1.
\end{equation}
By the definition of divisibility (\cref{def:divisibility-ring}), this is equivalent to:
\begin{equation}\label{eq:unit-mult-inv}
\exists w \in R \quad \text{such that} \quad u \cdot w = 1.
\end{equation}
In words, $u$ is a unit if and only if $u$ has a multiplicative inverse in $R$. The element $w$ is uniquely determined and denoted $u^{-1} := w$.
\end{definition}

\begin{steps}[Equivalence of Characterizations on the Blackboard]
The blackboard outlines three equivalent perspectives on units:
\begin{enumerate}[label=\textbf{(\arabic*)}]
    \item \textbf{Divisibility Perspective:} $u \mid 1$ (the element divides the multiplicative identity).
    \item \textbf{Algebraic Perspective (Unwinding Division):} $\exists w \in R$ such that $u \cdot w = 1$.
    \item \textbf{Verbal Perspective (English):} $u$ possesses a multiplicative inverse in $R$.
\end{enumerate}
Since all rings considered here are commutative, left-inverses and right-inverses coincide.
\end{steps}
\end{content}

\begin{hidden}
% timestamp: 00:26:22
% topic: Definition of divisibility and units in arbitrary commutative rings
% board_state: def:divisibility-ring, def:unit-ring
% next_goal: Linguistic distinction between "the unit", "the units", and "a unit", followed by the units of Z
% open_loops: none
\end{hidden}

\subsection{Linguistic Nuances: \qt{The Units} vs. \qt{The Unit} vs. \qt{A Unit}}

\begin{speech}[00:26:22 - 00:28:16]
Actually, people asked me some questions about language, and you know, English is a language I know, which means that I speak it without thinking. It was pointed out to me that the use of \qt{unit} is subtle, and I was doing it already in the first week; if you were sensitive to language, you'd have seen. So, for $\mathbb{Z}$ here, the units: what are the units for $\mathbb{Z}$?

In $\mathbb{Z}$, and it comes as no surprise, this is $+1$ and $-1$. Those are the only invertibles in it.

So, the units: if I say \qt{the units} here, that means all of them. The units.

If I say \qt{the unit} --- this is standard use of mathematical English ---, if I say \qt{the unit}, that's just $1$. If I have a ring and say \qt{let's consider the unit}, that means $1$.

You could say \qt{a unit}. If I say \qt{let's consider a ring, and let's consider a unit}, that could be either one; we don't know. Could be either. A unit --- this could be either.

So, that's actually standard usage. It's a little confusing, but once you get used to it, you know, it's just language. Is it clear? It was pointed out to me that I was doing that, and anyway, that's right. But if someone comes to you with a ring and says \qt{consider the unit}, that person means $1$.
\end{speech}

\begin{content}[Units of the Integers and Linguistic Distinctions]
\begin{example}[Units of the Ring of Integers \texorpdfstring{$\mathbb{Z}$}{Z}]\label[example]{ex:units-of-integers}
In the ring of integers $\mathbb{Z}$, the only elements possessing multiplicative inverses are $1$ and $-1$:
\begin{equation}\label{eq:units-of-integers}
\text{The units of } \mathbb{Z} = \{1, -1\}.
\end{equation}
\end{example}

\begin{remark}[Linguistic Distinction in Mathematical English]\label[remark]{rem:unit-terminology}
In mathematical English, the word \qt{unit} carries three distinct grammatical usages depending on the article:
\begin{itemize}
    \item \textbf{\qt{The units} (Plural):} Refers to the entire set of all invertible elements in the ring $R$, denoted $U(R)$ or $R^\times$. For $\mathbb{Z}$, the units are $\{1, -1\}$.
    \item \textbf{\qt{The unit} (Singular, Definite Article):} Refers specifically to the multiplicative identity element $1_R \in R$.
    \item \textbf{\qt{A unit} (Singular, Indefinite Article):} Refers to an arbitrary invertible element $u \in R$ (i.e., any element $u \in U(R)$). For $\mathbb{Z}$, \qt{a unit} could be either $1$ or $-1$.
\end{itemize}
\end{remark}
\end{content}

\section{The Group of Units and Modular Arithmetic}

\subsection{The Group of Units of a Ring}

\begin{speech}[00:28:16 - 00:29:11]
Okay, so we discussed the units of the Gaussian integers (Recall: $\pm 1, \pm i$, Week~1, \cref{def:gaussian-integers}), and now, since everything is defined, I recommend proving what the units of the Gaussian integers are. We wrote them last time.

But let's do something else now.

So, this is an exercise... \inlinemetanote{writes on board} Exercise: prove the units of the Gaussian integers are $\pm 1, \pm i$. You need an argument for that, and that argument comes from the norm in complex analysis, so prove it.

Okay, so observation...
\end{speech}

\begin{content}[Exercise: Units of the Gaussian Integers]
\begin{exercise}[Units of the Gaussian Integers]\label[exercise]{exer:units-gaussian-integers}
Prove that the units of the ring of Gaussian integers $\mathbb{Z}[i] = \{a + bi \mid a, b \in \mathbb{Z}\}$ are precisely:
\begin{equation}\label{eq:units-gaussian-integers}
U(\mathbb{Z}[i]) = \{\pm 1, \pm i\}.
\end{equation}
\end{exercise}

The lecturer writes out this proof in Week~3 (\cref{prop:units-gaussian-integers}).

\begin{steps}[Proof Strategy via the Multiplicative Norm]
Consider the field norm $N \colon \mathbb{Q}(i) \to \mathbb{Q}$ defined by $N(a + bi) = a^2 + b^2$. For all $z, w \in \mathbb{Z}[i]$, the norm is multiplicative:
\[
N(z \cdot w) = N(z) \cdot N(w).
\]
If $z \in \mathbb{Z}[i]$ is a unit, there exists $w \in \mathbb{Z}[i]$ with $z \cdot w = 1$. Taking norms yields:
\[
N(z) \cdot N(w) = N(1) = 1.
\]
Since $a, b \in \mathbb{Z}$, $N(z) = a^2 + b^2$ is a non-negative integer. The only positive integer divisor of $1$ in $\mathbb{Z}$ is $1$, so $a^2 + b^2 = 1$. The only integer solutions are:
\[
(a, b) \in \{(\pm 1, 0), (0, \pm 1)\} \implies z \in \{\pm 1, \pm i\}.
\]
\end{steps}
\end{content}

\begin{speech}[00:29:11 - 00:30:47]
...and it's a really important observation: if I take any ring $R$ --- let $R$ be a ring --- and let $U$ --- you could write $U(R)$ if you want --- be the set of units of $R$.

Then this observation is: if I take the units, I take multiplication --- that's multiplication in the ring ---, and I take the unit $1$, this is a group.

Okay? What does it say? If I take two units and multiply them, I get a unit.

And that's clear because, you know, if you have two elements here, if I multiply them, I have to show that's a unit. That means I have to show that if I multiply by some more stuff, I get $1$.

But this is a unit, so there's already its multiplicative inverse, and that one is a unit, so there's already its multiplicative inverse; so if I multiply these together, I get $1$, because I can flip the orders around.

So, this says the units are closed under multiplication.

And there's the neutral element, and then you have to show the existence of inverses, but a unit is a unit because it has an inverse.

So, that's the proof of this. This is a good thing to think about: this is always a group.
\end{speech}

\begin{content}[The Group of Units]
\begin{proposition}[Group of Units]\label[proposition]{prop:group-of-units}
Let $R$ be a ring with unity $1_R$, and let $U(R)$ (also denoted $U_R$ or $R^\times$) be the set of units of $R$:
\begin{equation}\label{eq:unit-set-def}
U(R) := \{u \in R \mid \exists w \in R \colon u \cdot w = 1_R\}.
\end{equation}
Then $(U(R), \cdot, 1_R)$ is a group under the multiplication of $R$.
\end{proposition}

\begin{proof}
We verify the group axioms for $(U(R), \cdot, 1_R)$:
\begin{itemize}
    \item \textbf{Closure under Multiplication:} Let $u_1, u_2 \in U(R)$. By definition, there exist $w_1, w_2 \in R$ such that $u_1 w_1 = 1_R$ and $u_2 w_2 = 1_R$. Since $R$ is commutative:
    \[
    (u_1 u_2)(w_2 w_1) = (u_1 w_1)(u_2 w_2) = 1_R \cdot 1_R = 1_R.
    \]
    Thus, $u_1 u_2$ has a two-sided multiplicative inverse in $R$, which shows $u_1 u_2 \in U(R)$.
    \item \textbf{Associativity:} Multiplicative associativity in $U(R)$ is directly inherited from the ring $R$.
    \item \textbf{Identity Element:} Since $1_R \cdot 1_R = 1_R$, the unity $1_R$ is its own inverse, so $1_R \in U(R)$. For all $u \in U(R)$, $1_R \cdot u = u$.
    \item \textbf{Existence of Inverses:} For any $u \in U(R)$, there exists $w \in R$ such that $u w = 1_R$. By symmetry of the relation $w u = 1_R$, $w$ is also invertible with inverse $u$, so $w = u^{-1} \in U(R)$.
\end{itemize}
Hence, $(U(R), \cdot, 1_R)$ is a group in the sense of \cref{def:group}. If $R$ is commutative, $U(R)$ is an abelian group.
\end{proof}
\end{content}

\subsection{Characterization of Units in \texorpdfstring{$\mathbb{Z}/m\mathbb{Z}$}{Z/mZ}}

\begin{speech}[00:30:47 - 00:31:44]
So, a very natural question, and one of the first things to think about regarding $\mathbb{Z}/m\mathbb{Z}$, if I take this ring $R$ \inlinemetanote{writes on board}...

And it was implicit there: I will always consider $m > 0$ when I do this, because there's no reason to consider negative ones, as that's just counting everything twice. And then there's the confusing case when $m = 0$, which is usually a special case, so I exclude it.

Usually we allow $m = 1$, which is the zero ring, but anyway, this is kind of standard. When I write $\mathbb{Z}/m\mathbb{Z}$, I always mean this, and typically I also mean $m > 1$, but anyway.

So, the question is: what are the units of $\mathbb{Z}/m\mathbb{Z}$?
\end{speech}

\begin{speech}[00:31:44 - 00:33:33]
So, let's say: suppose we have $\bar{x} \in \mathbb{Z}/m\mathbb{Z}$, and $\bar{x}$ is a unit.

This means that there exists $\bar{y} \in \mathbb{Z}/m\mathbb{Z}$ such that $\bar{x} \cdot \bar{y} = \bar{1}$ in $\mathbb{Z}/m\mathbb{Z}$.

I can write this, or sometimes you write $\bmod 1$, but anyway, that's what it means to be a unit --- I'm just repeating it, right?

And then I can translate: what does this mean in this ring? What does it mean for the integers $x$, $y$, and $1$? This is equivalent to $x \cdot y = 1 + d \cdot m$.

Do you agree with that? That's what it means, for some $d$. Okay?

And this, if you remember from the first week, means that $x$ and $m$ have $\operatorname{gcd} = 1$.

The only thing that could divide both $x$ and $m$, it would also have to divide $1$. So, that's why we say $\operatorname{gcd} = 1$.

In English, this condition is usually said as $x$ and $m$ are relatively prime.

Those are just words that express $\operatorname{gcd}(x, m) = 1$. Okay?
\end{speech}

\begin{hidden}
% timestamp: 00:33:33
% topic: Units of Z/mZ characterized via coprimality (gcd(x, m) = 1)
% board_state: prop:group-of-units, thm:units-z-mod-m
% next_goal: Formulate explicit theorem that U(Z/mZ) consists of coprime residue classes and is a group
% open_loops: none
\end{hidden}

\begin{speech}[00:33:33 - 00:35:02]
So, we have now discovered what the units of $\mathbb{Z}/m\mathbb{Z}$ are: they are all the elements which are relatively prime.

Because if it is a unit, it's the residue class of something relatively prime; if you have $\operatorname{gcd} = 1$, then you go backwards in this diagram.

So, I state that as a result.

So, if you write this as a result, then it says that the units of $\mathbb{Z}/m\mathbb{Z}$ consist of all elements $\bar{x}$ where $\operatorname{gcd}(x, m) = 1$. Did I write $\operatorname{gcd}$ here? Maybe $\operatorname{gcd}$. Okay? Is that clear?

And so the thing we've learned is that this units group, $U(\mathbb{Z}/m\mathbb{Z})$, is a group.

So, it says in particular, if you take $\mathbb{Z}/m\mathbb{Z}$, and you take the residue classes that are relatively prime, they form a group under multiplication.

So, let me do an example.
\end{speech}

\begin{content}[Units of the Quotient Ring \texorpdfstring{$\mathbb{Z}/m\mathbb{Z}$}{Z/mZ}]
\begin{theorem}[Units in $\mathbb{Z}/m\mathbb{Z}$]\label[theorem]{thm:units-z-mod-m}
Let $m \in \mathbb{Z}_{>0}$ be a positive integer. An element $\bar{x} \in \mathbb{Z}/m\mathbb{Z}$ is a unit if and only if $x$ and $m$ are coprime (this does not depend on the representative $x$, since $\operatorname{gcd}(x + km, m) = \operatorname{gcd}(x, m)$):
\begin{equation}\label{eq:units-z-mod-m}
U(\mathbb{Z}/m\mathbb{Z}) = (\mathbb{Z}/m\mathbb{Z})^\times = \left\{ \bar{x} \in \mathbb{Z}/m\mathbb{Z} \;\big|\; \operatorname{gcd}(x, m) = 1 \right\}.
\end{equation}
The order of this group is given by Euler's totient function $\phi(m) := |U(\mathbb{Z}/m\mathbb{Z})|$ (named by the lecturer in Week~3, \cref{def:eulers-totient}).
\end{theorem}

\begin{proof}
We establish the equivalence by unwinding the definition of divisibility:
\begin{itemize}
    \item \textbf{Forward Direction ($\implies$):} Suppose $\bar{x} \in U(\mathbb{Z}/m\mathbb{Z})$. Then there exists $\bar{y} \in \mathbb{Z}/m\mathbb{Z}$ such that:
    \[
    \bar{x} \cdot \bar{y} = \bar{1} \in \mathbb{Z}/m\mathbb{Z}.
    \]
    In terms of integer divisibility, this means $x y \equiv 1 \pmod m$. Hence, there exists an integer $d \in \mathbb{Z}$ such that:
    \[
    x y = 1 + d \cdot m \iff x y - d \cdot m = 1.
    \]
    Any common divisor $g \in \mathbb{Z}$ of $x$ and $m$ must divide the linear combination $x y - d m = 1$. Therefore, $g \mid 1$, which forces $\operatorname{gcd}(x, m) = 1$.
    \item \textbf{Reverse Direction ($\impliedby$):} Conversely, suppose $\operatorname{gcd}(x, m) = 1$. By Bézout's Identity, there exist integers $y, k \in \mathbb{Z}$ such that:
    \[
    x \cdot y + m \cdot k = 1 \implies x \cdot y \equiv 1 \pmod m.
    \]
    Projecting this relation into residue classes in $\mathbb{Z}/m\mathbb{Z}$, we obtain $\bar{x} \cdot \bar{y} = \bar{1}$. Thus, $\bar{x}$ is invertible, so $\bar{x} \in U(\mathbb{Z}/m\mathbb{Z})$.
\end{itemize}
By \cref{prop:group-of-units}, $(U(\mathbb{Z}/m\mathbb{Z}), \cdot, \bar{1})$ forms a finite abelian group.
\end{proof}
\end{content}

\subsection{Example: The Unit Group of \texorpdfstring{$\mathbb{Z}/7\mathbb{Z}$}{Z/7Z}}

\begin{speech}[00:35:02 - 00:36:09]
$\mathbb{Z}/7\mathbb{Z}$.

So, that's actually a prime one, but let me now do this example. So, the elements of this are $\bar{0}, \bar{1}, \bar{2}, \bar{3}, \bar{4}, \bar{5}, \bar{6}$, right?

And which of these residue classes come here? The classes of $x$ where $x$ and $7$ are relatively prime.

So, the question is: what are the units? Of course, it's some of these, but which of these residue classes correspond to $x$ and $7$ being relatively prime?

What would you say? As I said, if you don't know, guess.
\end{speech}

\begin{student}[Student Answer]
All of them.
\end{student}

\begin{speech}[00:36:09 - 00:36:14]
All of them? What about this one \inlinemetanote{points at $\bar{0}$}?
\end{speech}

\begin{student}[Student Answer]
Only $1$. $1$ to $6$ are the units.
\end{student}

\begin{speech}[00:36:14 - 00:36:50]
Exactly. So, the answer to this question is these: $\{\bar{1}, \bar{2}, \bar{3}, \bar{4}, \bar{5}, \bar{6}\}$.

And the reason is because, look, this was the answer: it's the residue classes of those numbers that are relatively prime. So, this is the residue class of those numbers that are relatively prime to $7$. And these are those residue classes.

If you're $0 \bmod 7$, then you're not relatively prime to $7$. Like an example of $0 \bmod 7$ would be $14$, and that's not relatively prime to $7$. So, it's these. Is that clear?
\end{speech}

\begin{content}[The Unit Group of \texorpdfstring{$\mathbb{Z}/7\mathbb{Z}$}{Z/7Z}]
\begin{example}[Units of $\mathbb{Z}/7\mathbb{Z}$]\label[example]{ex:units-z7z}
Consider the quotient ring $\mathbb{Z}/7\mathbb{Z}$. The elements of the ambient ring are:
\[
\mathbb{Z}/7\mathbb{Z} = \{\bar{0}, \bar{1}, \bar{2}, \bar{3}, \bar{4}, \bar{5}, \bar{6}\}.
\]
Since $7$ is a prime number, every non-zero residue class $\bar{x} \in \{\bar{1}, \bar{2}, \bar{3}, \bar{4}, \bar{5}, \bar{6}\}$ satisfies $\operatorname{gcd}(x, 7) = 1$. The zero class $\bar{0}$ is not coprime to $7$ (since $\operatorname{gcd}(0, 7) = 7 \neq 1$, or for example $14 \equiv 0 \pmod 7$ where $\operatorname{gcd}(14, 7) = 7$).
Consequently, the group of units is:
\begin{equation}\label{eq:units-z7z-set}
U(\mathbb{Z}/7\mathbb{Z}) = (\mathbb{Z}/7\mathbb{Z})^\times = \{\bar{1}, \bar{2}, \bar{3}, \bar{4}, \bar{5}, \bar{6}\}.
\end{equation}
The order of the unit group is $|U(\mathbb{Z}/7\mathbb{Z})| = \phi(7) = 7 - 1 = 6$.
\end{example}
\end{content}

\subsubsection{Powers of the Element \texorpdfstring{$a = \bar{3}$}{a = 3} and Cyclic Structure}

\begin{speech}[00:36:50 - 00:37:09]
Okay, and the claim is this is a group under multiplication. Well, anyway, this is a group under addition, but if you take the relatively prime ones, this is a group under multiplication.

So, we should be able to figure out what group it is. So, what's the neutral element of this group?
\end{speech}

\begin{student}[Student Answer]
One.
\end{student}

\begin{speech}[00:37:09 - 00:37:56]
One, it's going to be $\bar{1}$.

What is this group? It's a commutative group of order $6$. And so later we'll see that we know all commutative groups of order $6$. We'd like to know which one this is. It's isomorphic to some group, we know.

And so it matters what you pick. I propose we pick $3$. So, let me say we take this element in this group, $a = \bar{3}$.

And let's try to consider its powers. So, there's the neutral element, then I can take $a, a^2, a^3, a^4$. You should know the powers of $3$. Okay, let's see how many powers of $3$ you know.
\end{speech}

\begin{speech}[00:37:56 - 00:39:44]
So, what is this one? So I'm just going to do this: this is $\bar{1}$. So, what is $a$? Well, that's just that: $\bar{3}$.

What is this? That's the square of $a$, so that's the residue class of $3^2$, which is $9$. So, I should put $\bar{2}$ here, right? Are you happy with this?

And this is $a^3$: $3^3$ is $27$. For $27$, the residue $\bmod 7$ is $\bar{6}$, right?

$a^4$ is $3^4$, which is $81$, and the residue $\bmod 7$ of $81$ is $\bar{4}$, I think.

$3^5$ is $243$, and its residue is going to be $\bar{5}$.

What do you think this is going to be? Do you know what $3^6$ is? $3^6$ is $729$. And if I look at its residue class $\bmod 7$, well, $700$ is already here, and $28$ is divisible, so its residue class is $\bar{1}$. This is back to $1$.

So, this is understanding the group. You should know the powers of $3$; they're very useful, you know.

So, this is what happens when you multiply in this group: you take $a$, $a^2$, the next one $a^3$, the next one $a^4$, the next one $a^5$, and then when you go to the sixth, you go back to the identity.

What group do you know this happens in? Does this remind you of anything? Yeah?
\end{speech}

\begin{student}[Student Answer]
A cyclic group.
\end{student}

\begin{speech}[00:39:44 - 00:39:54]
It's $\mathbb{Z}/6\mathbb{Z}$. So, the proposition that we've kind of proven here...
\end{speech}

\begin{content}[Powers of the Generator \texorpdfstring{$a = \bar{3}$}{a = 3} in \texorpdfstring{$U(\mathbb{Z}/7\mathbb{Z})$}{U(Z/7Z)}]
To examine the internal structure of the abelian group $U(\mathbb{Z}/7\mathbb{Z})$ of order $6$, we choose the candidate generator $a := \bar{3} \in U(\mathbb{Z}/7\mathbb{Z})$ and compute its successive powers:

\begin{center}
\begin{tabular}{c|ccccccc}
\toprule
\textbf{Power} & $a^0$ & $a^1$ & $a^2$ & $a^3$ & $a^4$ & $a^5$ & $a^6$ \\
\midrule
\textbf{Symbolic} & $\bar{1}$ & $a$ & $a^2$ & $a^3$ & $a^4$ & $a^5$ & $a^6$ \\
\textbf{Value in $\mathbb{Z}$} & $1$ & $3$ & $9$ & $27$ & $81$ & $243$ & $729$ \\
\textbf{Residue modulo $7$} & $\bar{1}$ & $\bar{3}$ & $\bar{2}$ & $\bar{6}$ & $\bar{4}$ & $\bar{5}$ & $\bar{1}$ \\
\bottomrule
\end{tabular}
\end{center}

\begin{steps}[Modular Arithmetic Derivations for Powers of $3$]
\begin{align*}
    a^0 &= \bar{1}, \\
    a^1 &= \bar{3}, \\
    a^2 &= \overline{3^2} = \bar{9} = \bar{2} && (9 = 7 \cdot 1 + 2), \\
    a^3 &= \overline{3^3} = \overline{27} = \bar{6} && (27 = 7 \cdot 3 + 6), \\
    a^4 &= \overline{3^4} = \overline{81} = \bar{4} && (81 = 7 \cdot 11 + 4), \\
    a^5 &= \overline{3^5} = \overline{243} = \bar{5} && (243 = 7 \cdot 34 + 5), \\
    a^6 &= \overline{3^6} = \overline{729} = \bar{1} && (729 = 7 \cdot 104 + 1).
\end{align*}
Because the powers of $a = \bar{3}$ exhaust all $6$ elements of $U(\mathbb{Z}/7\mathbb{Z})$ before returning to the identity $\bar{1}$ at exponent $6$, the element $\bar{3}$ is a \newterm{generator} (or \newterm{primitive root} modulo $7$). Consequently, $U(\mathbb{Z}/7\mathbb{Z})$ is a cyclic group of order $6$.
\end{steps}
\end{content}

\begin{hidden}
% timestamp: 00:39:54
% topic: Cyclic generator of U(Z/7Z) and establishment of powers of 3 mod 7
% board_state: ex:units-z7z, generator-powers-table
% next_goal: Formalize the group isomorphism U(Z/7Z) \cong (Z/6Z, +, 0) via explicit discrete logarithm map
% open_loops: none
\end{hidden}

\subsubsection{The Group Isomorphism \texorpdfstring{$(U(\mathbb{Z}/7\mathbb{Z}), \cdot) \cong (\mathbb{Z}/6\mathbb{Z}, +)$}{(U(Z/7Z), *) = (Z/6Z, +)}}

\begin{speech}[00:39:54 - 00:41:24]
...we've proven what that calculation of powers of $3$... and I picked $3$. If you picked another one, it would have been slightly different, but that calculation proves a group isomorphism, which maybe I'll write here and we think about.

You should check my powers of $3$, and then after you take the power of $3$, you have to take the residue $\bmod 7$.

So, this great 20th century Hungarian mathematician (Paul Erd\H{o}s) said that if numbers aren't beautiful, I don't know what is. That's kind of the mathematical perspective on numbers.

So, we have here the units of $\mathbb{Z}/7\mathbb{Z}$, and then there's another group, which is $\mathbb{Z}/6\mathbb{Z}$ with $+$ and $\bar{0}$.

And I claim these are isomorphic. So, this is the units with multiplication and $\bar{1}$. This group is isomorphic to that group. And what is the isomorphism? It's written on the board.

The isomorphism --- this is a slightly painful way to write the isomorphism, but let's just be painful about it ---: so, here $\bar{1}$ is going to go to $\bar{0}$. And bar means something different: it means $\bmod 7$ here, and $\bmod 6$ here, right?
\end{speech}

\begin{speech}[00:41:24 - 00:42:33]
$a$, which is $3$ here, $a$ is $\bar{3}$. So, the nice way to do it is: $a$ goes to $\bar{1}$, $a^2$ goes to $\bar{2}$. I guess I could write this in a slightly better notation, but anyway, you get the pattern. There are not so many more: $a^3, a^4, a^5$. These go to $\bar{3}, \bar{4}, \bar{5}$.

So, this is the simple way to write it. And the reason this is true is because on this side when we multiply, we just add up the exponents, and on this side we add up those. And the important thing is the sixth one goes to zero, which is what we have here, okay?

And if you want to say this explicitly, you have to write these $a$'s in terms of the elements of the units, which I've done there. So, this is totally explicit.

Okay, that's all. \inlinemetanote{students applaud and pack up; video ends}
\end{speech}

\begin{content}[Group Isomorphism: \texorpdfstring{$U(\mathbb{Z}/7\mathbb{Z}) \cong (\mathbb{Z}/6\mathbb{Z}, +, \bar{0})$}{U(Z/7Z) isomorphic to (Z/6Z, +, 0)}]
\begin{proposition}[Cyclic Group Isomorphism]\label[proposition]{prop:z7z-cyclic-isomorphism}
The multiplicative group of units $(U(\mathbb{Z}/7\mathbb{Z}), \cdot, \bar{1})$ is isomorphic to the additive cyclic group $(\mathbb{Z}/6\mathbb{Z}, +, \bar{0})$:
\begin{equation}\label{eq:z7z-isom-group}
(U(\mathbb{Z}/7\mathbb{Z}), \cdot, \bar{1}) \cong (\mathbb{Z}/6\mathbb{Z}, +, \bar{0}).
\end{equation}
\end{proposition}

\begin{proof}
Let $a = \bar{3} \in U(\mathbb{Z}/7\mathbb{Z})$. We define the mapping $\Phi \colon U(\mathbb{Z}/7\mathbb{Z}) \to \mathbb{Z}/6\mathbb{Z}$ by sending powers of the generator $a$ to their integer exponents modulo $6$:
\begin{align*}
    \Phi \colon U(\mathbb{Z}/7\mathbb{Z}) &\to \mathbb{Z}/6\mathbb{Z}, \\
    \bar{1} = a^0 &\mapsto \bar{0}, \\
    \bar{3} = a^1 &\mapsto \bar{1}, \\
    \bar{2} = a^2 &\mapsto \bar{2}, \\
    \bar{6} = a^3 &\mapsto \bar{3}, \\
    \bar{4} = a^4 &\mapsto \bar{4}, \\
    \bar{5} = a^5 &\mapsto \bar{5}.
\end{align*}
\textbf{Homomorphism Property:} For any exponents $j, k \in \{0, 1, 2, 3, 4, 5\}$, multiplication in the unit group corresponds to addition of exponents:
\[
a^j \cdot a^k = a^{j+k} = a^{(j+k) \bmod 6},
\]
since $a^6 = \bar{1}$. Applying the map $\Phi$:
\[
\Phi(a^j \cdot a^k) = \overline{j + k} = \bar{j} + \bar{k} = \Phi(a^j) + \Phi(a^k) \in \mathbb{Z}/6\mathbb{Z}.
\]
Hence, $\Phi$ is a group homomorphism. By the table of powers, $a^0, \dots, a^5$ are exactly the six units, and $\Phi$ sends them to the six distinct classes $\bar{0}, \dots, \bar{5}$; so $\Phi$ is bijective, which establishes the isomorphism.
\end{proof}

\begin{steps}[Disambiguating Residue Class Notations Across Groups]
The bar notation $\bar{\cdot}$ operates over two distinct moduli in this isomorphism:
\begin{itemize}
    \item In the domain $(U(\mathbb{Z}/7\mathbb{Z}), \cdot)$, bars denote residue classes modulo $7$ under multiplication: $\bar{x} = x + 7\mathbb{Z}$.
    \item In the codomain $(\mathbb{Z}/6\mathbb{Z}, +)$, bars denote residue classes modulo $6$ under addition: $\bar{k} = k + 6\mathbb{Z}$.
\end{itemize}
The isomorphism $\Phi$ is the \newterm{discrete logarithm} with respect to the base $3$ modulo $7$.
\end{steps}
\end{content}

\begin{insight}[Paul Erd\H{o}s on the Beauty of Numbers]
The lecturer references the celebrated quote by the Hungarian mathematician Paul Erd\H{o}s (1913--1996): \qt{If numbers are not beautiful, I don't know what is.} The unexpected isomorphism between the multiplicative units of a prime field $\mathbb{F}_7^\times$ and the simple additive clock arithmetic of $\mathbb{Z}/6\mathbb{Z}$ is a classic embodiment of this aesthetic: multiplication modulo $7$ transforms into addition modulo $6$ via the powers of a single primitive root.
\end{insight}

\begin{overview}[Summary of Key Concepts]
\begin{itemize}
    \item \textbf{Category Theory:} A category $\mathcal{C}$ comprises objects $\operatorname{Obj}(\mathcal{C})$, hom-sets $\operatorname{Hom}(A, B)$, associative composition $\circ$, and identity morphisms $\operatorname{Id}_A$.
    \item \textbf{Categorical Morphisms:}
    \begin{itemize}
        \item \emph{Isomorphism:} Invertible morphism ($f \circ g = \operatorname{Id}_B$, $g \circ f = \operatorname{Id}_A$).
        \item \emph{Monomorphism:} Left-cancellable morphism ($f \circ g_1 = f \circ g_2 \implies g_1 = g_2$), generalizing injectivity.
        \item \emph{Epimorphism:} Right-cancellable morphism ($g_1 \circ f = g_2 \circ f \implies g_1 = g_2$), generalizing surjectivity.
    \end{itemize}
    \item \textbf{Integral Domains:} A commutative ring with $1 \neq 0$ having no zero divisors ($a \neq 0 \land b \neq 0 \implies ab \neq 0$). The quotient ring $\mathbb{Z}/m\mathbb{Z}$ is an integral domain if and only if $m$ is prime.
    \item \textbf{Units and Divisibility:} An element $u \in R$ is a unit if $u \mid 1$, i.e., $u$ has a multiplicative inverse. The units form a multiplicative group $U(R)$.
    \item \textbf{Units of Residue Class Rings:} $U(\mathbb{Z}/m\mathbb{Z}) = \{ \bar{x} \in \mathbb{Z}/m\mathbb{Z} \mid \operatorname{gcd}(x, m) = 1 \}$ with order $\phi(m)$. For prime $p = 7$, $U(\mathbb{Z}/7\mathbb{Z}) \cong (\mathbb{Z}/6\mathbb{Z}, +)$ generated by $\bar{3}$.
\end{itemize}
\end{overview}

% week3:
% ==========================================
% LatexRefinement Step Output: step4-2026-09-29-tuesday-week3-algebra-1-offset-final.tex
% Model: gemini-3.8-flash
% Temperature: 0.35
% TopP: 0.9
% TopK: 10
% MaxOutputTokens: 65535
% ThinkingBudget: 4096
% ThinkingLevel: HIGH
% Prompt Tokens: 72’032
% Candidates Tokens: 38’622 (inkl. Thinking Tokens)
% Cached Tokens: 0
% Processed on: 2026-10-04 13:15:25
% ==========================================

% Primary Language: English
% End of the video: 01:30:00

\lecturechapter[3]{Tuesday}{29 Sep}{29 September 2026}{Groups of Units, Fields of Fractions, and Polynomial Rings}
% Old chapter title: Rings, Groups of Units, and Fields of Fractions

\begin{overview}[Course Content Overview]
This lecture investigates the transition from commutative rings to fields and polynomials, focusing on the invertible elements of rings, the localization process that produces the minimal ambient field of an integral domain, and the algebraic construction of polynomial rings:
\begin{itemize}
    \item \textbf{The Group of Units $\operatorname{Unit}(R)$:} Structure of multiplicatively invertible elements, characterization of units in $\mathbb{Z}$ and the Gaussian integers $\mathbb{Z}[i]$ via the norm, and the group of units of residue class rings $(\mathbb{Z}/m\mathbb{Z})^*$.
    \item \textbf{Euler's Totient Function $\phi(m)$:} Definition as the order of $(\mathbb{Z}/m\mathbb{Z})^*$, with explicit computations for prime powers $\phi(p) = p-1$ and $\phi(p^2) = p(p-1)$.
    \item \textbf{Universal Property of $\mathbb{Q}$:} Rigorous formulation of $\mathbb{Q}$ as the smallest field containing $\mathbb{Z}$, and the unique extension of ring embeddings into fields.
    \item \textbf{Construction of the Field of Fractions $\operatorname{Frac}(R)$:} Generalization to arbitrary integral domains via equivalence classes of formal pairs $(x, y) \in R \times (R \setminus \{0\})$, verification of well-defined operations, and canonical embedding $i \colon R \hookrightarrow \operatorname{Frac}(R)$.
    \item \textbf{Concrete Quotient Fields:} Field of fractions of $\mathbb{Z}$, the Gaussian rationals $\mathbb{Q}(i) = \operatorname{Frac}(\mathbb{Z}[i])$, subrings of $\mathbb{Q}$ such as $\mathbb{Z}[\frac{1}{3}]$, and self-isomorphism for existing fields $\operatorname{Frac}(K) \cong K$.
    \item \textbf{Polynomial Rings $R[x]$:} Intuitive formulation via formal linear combinations, rigorous foundation using almost-zero sequences in $R^{\mathbb{N}_0}$, componentwise addition, and Cauchy product (convolution) multiplication.
\end{itemize}
\end{overview}

% --- TEIL 1 ---

\section{Review: Rings and the Group of Units}

\begin{speech}[00:00:00 - 00:00:32]
Okay, so we start. 

About the quiz: as I said, I'm very much pro-guessing in mathematics; you're never penalized for getting the wrong answer. Lika said that there was some implementation problem and some people were penalized for getting the right answer, and she's going to fix that, so. 

If you think the quiz was too easy, or too hard, or too medium, tell us. As I said, I've never gone through this quiz thing before. Okay, so then...
\end{speech}

\subsection{The Group of Units of a Ring}

\begin{speech}[00:00:32 - 00:01:51]
Last week we discussed rings, and associated to a ring there's an additive group. This is easy, because it's part of the definition, but there's another group which was also discussed in the last lecture (Recall: \cref{prop:group-of-units}), which is the group of units. 

And I remind you of this group... \inlinemetanote{writes \qt{the group of unit} $\operatorname{Unit}_R = \{ x \mid \exists y, xy = 1 \}$} Various people use different notation for describing it, but here I just say $\operatorname{Unit} R$ so there's no question: it's the group of units of the ring $R$. It's those elements such that there exists somebody else where the product is $1$, so it's those elements with a multiplicative inverse. And this is a group --- we discussed this last week. 

So, this group here is $\operatorname{Unit} R$, and it has, well, the units by this definition, multiplication, and $1$. \inlinemetanote{writes $\operatorname{Unit}_R = (\text{Units}, \cdot, 1)$}

Whenever you see a ring, you immediately have these two groups. And this one is, you know, often the more complicated one to think about, because it's not part of the definition in some sense.
\end{speech}

\begin{content}[The Group of Units]
\recap{Units and the group of units were introduced in Week~2 as $U(R)$ (\cref{def:unit-ring}, \cref{prop:group-of-units}); the lecturer now writes $\operatorname{Unit}(R)$.}

Let $(R, +, \cdot, 0, 1)$ be a ring with unity. Associated with $R$ are two natural canonical groups:
\begin{enumerate}[label=\textbf{(\arabic*)}]
    \item \textbf{The Additive Group:} $(R, +, 0)$, which forms an abelian group by the ring axioms.
    \item \textbf{The Group of Units:} The set of multiplicatively invertible elements in $R$:
    \begin{equation}\label{eq:unit-group-def-w3}
    \operatorname{Unit}(R) := \{ x \in R \mid \exists y \in R \text{ such that } xy = 1 \}.
    \end{equation}
\end{enumerate}

Under ring multiplication, $(\operatorname{Unit}(R), \cdot, 1)$ forms a group with identity element $1$.

\begin{steps}[Multiplicative Group Structure]
The group of units is closed under multiplication: if $x, z \in \operatorname{Unit}(R)$ with inverses $x^{-1}, z^{-1} \in R$, then $(xz)(z^{-1}x^{-1}) = 1$, so $xz \in \operatorname{Unit}(R)$. Associativity is inherited from the ring, and the multiplicative identity $1$ is invertible since $1 \cdot 1 = 1$.
\end{steps}
\end{content}

\subsection{Units in the Integers and Gaussian Integers}

\begin{speech}[00:01:51 - 00:03:03]
We've seen this --- I just review, but you know this --- if you take $\mathbb{Z}$, the units are $\pm 1$. And I said this many times, but maybe since I didn't actually ever give a proof: in $\mathbb{Z}[i]$, the units are $\pm 1, \pm i$ (Recall: \cref{exer:units-gaussian-integers}). 

If you're taking complex analysis, you should be able to prove this, but let's actually write a proof. I don't think I ever wrote the proof for this. 

So, proof: this is a proof that these are the only invertible elements in this ring. Is it clear what I'm trying to prove? 

Proof: let $a + bi \in \mathbb{Z}[i]$ be a unit. When I write this, of course that means that $a$ and $b$ are integers, okay? That's what being in there means. 

All right, so then there exists somebody else, $c + di$, also here, such that $(a+bi)(c+di) = 1$. And this is multiplication of complex numbers.
\end{speech}

\begin{speech}[00:03:03 - 00:03:58]
The whole idea of the proof is understanding multiplication of complex numbers. And if you've taken a complex numbers class, or you've met complex numbers, however you know it: when you multiply two complex numbers, their angles add and their lengths multiply, right? 

So, we just take the norm --- the lengths --- take the norms. The norm here is $\sqrt{a^2 + b^2}$; that's the length of this one. I hope that's not controversial, that's Euclid or something. This one is $\sqrt{c^2 + d^2}$, and their norms multiply, so this has to be just $1$. 

Okay? So, this is pretty restrictive: $a, b, c$, and $d$ are all integers.
\end{speech}

\begin{hidden}
% timestamp: 00:03:58
% topic: Determining units in the Gaussian integers Z[i] via norm multiplication
% board_state: eq:unit-group-def, units of Z, units of Z[i], norm equation
% next_goal: Deduce the 4 units (+-1, +-i) from the Diophantine equation a^2 + b^2 = 1 and draw the diagram
% open_loops: none
\end{hidden}

\begin{speech}[00:03:58 - 00:05:13]
So, we can square that equation: $(a^2+b^2)(c^2+d^2) = 1^2$, which is $1$. Now you have an equation where you have integers all multiplying to $1$, and you know the only way this can happen is if this integer itself is $\pm 1$, because you can't factor $1$ in the integers without having just $1 \cdot 1$ or $(-1) \cdot (-1)$. Is that clear? 

How could this be $1$? Well, $a^2$ for an integer is either $0$, or $1$, or something bigger. So, the only way this could be $\pm 1$ is if $a$ is equal to $0$ or $\pm 1$, and $b$ is equal to $0$ or $\pm 1$ --- those are the only cases. If these are both positive integers, it's going to be bigger than $1$. So, you only have to check these cases. Is that clear? 

And if you look at these cases, well, these are the answers. Okay? That's a complete proof; that's as complete a proof as you get in this class. I think this proof was $100$ percent complete. 

That's a nice fact, and if you draw them, they look like this \inlinemetanote{draws a sketch of the complex plane with four points on the axes}. Those are the four units in the Gaussian integers.
\end{speech}

\begin{content}[Classification of Units in the Gaussian Integers]
\recap{The units of $\mathbb{Z}$ were determined in Week~2 (\cref{ex:units-of-integers}), and those of $\mathbb{Z}[i]$ were stated in Week~1 and set as an exercise (\cref{exer:units-gaussian-integers}). The proof below settles that exercise.}

\begin{example}[Units of the Integers]\label[example]{ex:units-integers-w3}
For the ring of integers $\mathbb{Z}$:
\[
\operatorname{Unit}(\mathbb{Z}) = \mathbb{Z}^* = \{1, -1\}.
\]
\end{example}

\begin{proposition}[Units of the Gaussian Integers]\label[proposition]{prop:units-gaussian-integers}
The group of units of the Gaussian integers $\mathbb{Z}[i] = \{ a + bi \mid a, b \in \mathbb{Z} \}$ is:
\[
\operatorname{Unit}(\mathbb{Z}[i]) = \{1, -1, i, -i\}.
\]
\end{proposition}

\begin{proof}
Let $z = a + bi \in \mathbb{Z}[i]$ be a unit with $a, b \in \mathbb{Z}$. By definition, there exists $w = c + di \in \mathbb{Z}[i]$ ($c, d \in \mathbb{Z}$) such that:
\begin{equation}\label{eq:gaussian-unit-prod}
(a + bi)(c + di) = 1.
\end{equation}
Taking the complex modulus (Euclidean length) on both sides of \eqref{eq:gaussian-unit-prod}:
\[
|a + bi| \cdot |c + di| = |1| = 1 \implies \sqrt{a^2 + b^2} \cdot \sqrt{c^2 + d^2} = 1.
\]
Squaring this relation yields an equation in non-negative integers:
\begin{equation}\label{eq:norm-mult-equation}
(a^2 + b^2)(c^2 + d^2) = 1.
\end{equation}
Since $a^2 + b^2 \in \mathbb{Z}_{\ge 0}$ and divides $1$, we must have:
\[
a^2 + b^2 = 1.
\]
Because $a, b \in \mathbb{Z}$, the only integer solutions to $a^2 + b^2 = 1$ are:
\begin{align*}
a = \pm 1, \quad b = 0 &\implies z = \pm 1, \\
a = 0, \quad b = \pm 1 &\implies z = \pm i.
\end{align*}
Conversely, $1 \cdot 1 = 1$, $(-1)(-1) = 1$, and $i(-i) = 1$, so all four elements are invertible. Thus:
\[
\operatorname{Unit}(\mathbb{Z}[i]) = \{1, -1, i, -i\}.
\]
\end{proof}

\begin{center}
\begin{tikzpicture}[scale=1.5]
% \begin{hidden}
% - Canvas: scale=1.5, x in [-1.5, 1.5], y in [-1.5, 1.5]
% - Elements: coordinate axes, dashed unit circle, four unit points marked with labels
% - Coordinates: (1,0) for 1, (-1,0) for -1, (0,1) for i, (0,-1) for -i
% - Occlusion: axes in background, dashed circle in midground, dots and labels in foreground
% \end{hidden}
    % Axes
    \draw[->, thick, gray!70] (-1.5, 0) -- (1.5, 0) node[right, text=black] {$\operatorname{Re}$};
    \draw[->, thick, gray!70] (0, -1.5) -- (0, 1.5) node[above, text=black] {$\operatorname{Im}$};

    % Unit Circle (radius 1 coordinate unit)
    \draw[dashed, MidnightBlue!50, thick] (0,0) circle (1);

    % Four Units
    \fill[BrickRed] (1, 0) circle (2pt) node[below right, text=BrickRed, font=\bfseries] {$1$};
    \fill[BrickRed] (-1, 0) circle (2pt) node[below left, text=BrickRed, font=\bfseries] {$-1$};
    \fill[BrickRed] (0, 1) circle (2pt) node[above right, text=BrickRed, font=\bfseries] {$i$};
    \fill[BrickRed] (0, -1) circle (2pt) node[below right, text=BrickRed, font=\bfseries] {$-i$};
\end{tikzpicture}
\end{center}

\begin{steps}[Algebraic Meaning of the Norm]
The field norm $N(a+bi) = a^2 + b^2$ is multiplicative: $N(zw) = N(z)N(w)$. Since $N(z) \in \mathbb{Z}_{\ge 0}$ for any Gaussian integer, the norm of any invertible element must divide $N(1) = 1$, forcing $N(z) = 1$. The only lattice points on the circle of radius $1$ are the four points shown above.
\end{steps}
\end{content}

\subsection{Units in Residue Class Rings}

\begin{speech}[00:05:14 - 00:07:21]
We did a much more interesting example last time, which was also with proof, but since various people asked me about it afterwards, I want to do just another example. And that was this ring: $\mathbb{Z}/m\mathbb{Z}$. 

That's a nice ring, and so it has a group of units, $\operatorname{Unit}(\mathbb{Z}/m\mathbb{Z})$. This was a result I proved last time (Recall: \cref{thm:units-z-mod-m}): the units are equal to those classes such that this (i.e., $\gcd(x, m)$) is equal to $1$. And as I said, it's standard here to say $m \ge 2$, so we're not dealing with the zero ring, okay? 

So, that's the units, and these are the residue classes which are relatively prime. I did an example last time (Recall: $\mathbb{Z}/7\mathbb{Z}$, \cref{ex:units-z7z}). This is a very interesting thing to study; it's just not a trivial thing, in the sense that one might say that if I take the additive structure here, this is very simple to think about. It's just adding up integers mod $m$. In some sense, most people learn about this in one name or the other in school. 

But the units are more complicated, more interesting. And so I wanted to do another... there were various people who asked me some questions. So, let's try this: $\mathbb{Z}/10\mathbb{Z}$, just to say a little bit...

The standard notation for this is that: $(\mathbb{Z}/m\mathbb{Z})^*$. This is the usual notation for the units of $\mathbb{Z}/m\mathbb{Z}$. And also this notation is sometimes used in general: $R^*$ sometimes denotes the units, depending on the author. I've just been writing $\operatorname{Unit}$, so there's no question about what I mean, but in this case it's so common to use the star here that I'll also use the star. Is that clear?
\end{speech}

\begin{speech}[00:07:21 - 00:08:37]
How to think about this? Well, I'll just do an example, because as I said, there were various questions. How to think about this? This is all the residue classes that are relatively prime. So, you can just write them all first and cross them out: $0, 1, 2, 3, 4, 5$... maybe I should have picked something less than $10$... 

These are all the residue classes, and then we have to cancel all the ones that are not relatively prime. So, which ones? Is $0$ relatively prime? No, so we cross that out. $1$ is relatively prime; $1$ is always relatively prime. What about $2$? No. $3$? Yes, so we skip it. $4$? That has a $2$ in it, and $10$ has a $2$, so we cross that out. $5$? $10$ has a $5$ in it, so we cross that one out, cross that one out, and what's left? $4$, yeah. So, these are the four residue classes left. 

The group of units here consists of these residue classes, okay? And the claim is that this is a group under multiplication. That's what that claim is. So, for example, $3$ is in this group, so what is its inverse?
\end{speech}

\begin{student}[Student Answer]
$7$.
\end{student}

\begin{speech}[00:08:42 - 00:08:52]
Yeah, $7$, because $7 \times 3$ is $21$, and so on. So, anyway, you can play with this.
\end{speech}

\begin{hidden}
% timestamp: 00:08:45
% topic: Group of units in residue class rings Z/mZ and example Z/10Z
% board_state: eq:unit-group-def, def of (Z/mZ)^*, calculation of (Z/10Z)^* = {1, 3, 7, 9}
% next_goal: Introduce Euler's totient function phi(m) as the order of (Z/mZ)^*
% open_loops: none
\end{hidden}

\begin{content}[Units in \texorpdfstring{$\mathbb{Z}/m\mathbb{Z}$}{Z/mZ} and the Example \texorpdfstring{$\mathbb{Z}/10\mathbb{Z}$}{Z/10Z}]
\recap{This characterization was proved in Week~2 (\cref{thm:units-z-mod-m}); the lecturer restates it with the notation $(\mathbb{Z}/m\mathbb{Z})^*$.}

\begin{proposition}[Units in the Residue Class Ring]\label[proposition]{prop:units-mod-m-w3}
Let $m \in \mathbb{Z}_{\ge 2}$. The group of units of the residue class ring $\mathbb{Z}/m\mathbb{Z}$ (frequently denoted by $(\mathbb{Z}/m\mathbb{Z})^*$ or $R^*$ for general rings) consists of those residue classes whose representatives are coprime to $m$:
\begin{equation}\label{eq:units-mod-m}
(\mathbb{Z}/m\mathbb{Z})^* = \operatorname{Unit}(\mathbb{Z}/m\mathbb{Z}) = \{ \bar{x} \in \mathbb{Z}/m\mathbb{Z} \mid \gcd(x, m) = 1 \}.
\end{equation}
\end{proposition}

\begin{example}[Computing the Units of $\mathbb{Z}/10\mathbb{Z}$]\label[example]{ex:units-z10z}
Consider the ring $\mathbb{Z}/10\mathbb{Z} = \{ \bar{0}, \bar{1}, \bar{2}, \bar{3}, \bar{4}, \bar{5}, \bar{6}, \bar{7}, \bar{8}, \bar{9} \}$. To find the units, we eliminate all classes that share a common divisor with $10$ greater than $1$:
\begin{itemize}
    \item \textbf{Class $\bar{0}$:} $\gcd(0, 10) = 10 \ne 1$ (eliminate $\bar{0}$),
    \item \textbf{Even classes:} $\gcd(2, 10) = 2$, $\gcd(4, 10) = 2$, $\gcd(6, 10) = 2$, $\gcd(8, 10) = 2$ (eliminate $\bar{2}, \bar{4}, \bar{6}, \bar{8}$),
    \item \textbf{Class $\bar{5}$:} $\gcd(5, 10) = 5$ (eliminate $\bar{5}$).
\end{itemize}
The remaining coprime residue classes form the group of units:
\[
(\mathbb{Z}/10\mathbb{Z})^* = \{ \bar{1}, \bar{3}, \bar{7}, \bar{9} \}.
\]
This is an abelian group of order $4$ under multiplication modulo $10$. For instance, the inverse of $\bar{3}$ is $\bar{7}$, since:
\[
\bar{3} \cdot \bar{7} = \overline{21} = \bar{1} \in \mathbb{Z}/10\mathbb{Z}.
\]
\end{example}
\end{content}

\section{Euler's Totient Function}

\subsection{Definition and Elementary Calculations}

\begin{speech}[00:08:52 - 00:10:35]
This is a more interesting object, and in fact you could ask: what is the order? This is a basic question in elementary number theory: what is the order of this group? 

By order, this word means for groups the number of elements. So, when I say order of a group, if it's a finite group, I just mean the number of elements. What is the order of this? For example, sometimes this order is denoted just by straight lines, the bars. You should be familiar with this set-theoretic notation. 

These two bars usually denote order. Is that clear? That's kind of standard set-theoretic notation. So, first of all, let's do $10$: what is this? Well, the answer is $4$; you just have to compute these.
\end{speech}

\begin{speech}[00:10:35 - 00:12:12]
The order of $(\mathbb{Z}/m\mathbb{Z})^*$, the group of units, has a name. Do you know the name of this function? Yes, it's $\phi$ --- Euler's totient function. That's just the definition. That's the definition of this function: it's just the order. 

And you could ask: can you compute this? There are different ways to compute it, but it has to do with the factorization, and we will get to that later. So, you could say this is maybe a question for you now: try to compute Euler's totient function.

Euler is like the most famous Swiss person that ever lived. A few years ago it might have been Federer, but now he's retired, so it's Euler. And it turns out that Euler and Federer were both from Basel, so you can check it. 

You know, my first trip to Switzerland was when I was a kid in the eighties, and the Swiss still had Euler on the $10$ Swiss Franc note. I don't know if you've ever seen that; it's a beautiful note. You should go look it up. 

Okay, so let's do one example to compute this. Well, let's do two examples. Suppose you have a prime $p$, then what is this?
\end{speech}

\begin{student}[Student Answer]
$p - 1$.
\end{student}

\begin{speech}[00:12:13 - 00:12:15]
And why is that?
\end{speech}

\begin{student}[Student Answer]
Because any number less than $p$ is coprime to $p$.
\end{student}

\begin{speech}[00:12:22 - 00:13:48]
In other words, if you write this, you'll start... well, you'll write $p$ elements, and you'll cross out $0$, and everybody else is relatively prime. So, it's $p - 1$. 

And so what about this one: $\phi(p^2) = {?}$ Do you know how to do this? 

What you do --- I mean, of course once you do this, you can do the general case for powers of a prime --- you write all the elements (this is a really primitive way to do it, but let's just do it): all the elements up to $p^2 - 1$. That's what we did there, right? So, it starts out, and then you cross out all of the... 

Here you're going to cross out zero, so it's $p^2 - 1$, but you have to take away... you know that for these classes, of course you don't want zero, but we have to take away all the multiples of $p$, right?

And so now you have to concentrate really hard: how many multiples of $p$? Well, there's $p$ here, there's $2p$ here, and what's the last one? The last one you have to cross out is $p(p - 1)$; that's the last one. Is that clear? 

So, we have to take away $p - 1$, so the answer is $p^2 - p$. The answer to this is $p^2 - p$. 

Okay, that's the beginning; you can play with it a bit. All right?
\end{speech}

\begin{hidden}
% timestamp: 00:13:40
% topic: Euler's totient function and calculations for prime powers phi(p) and phi(p^2)
% board_state: def of phi(m) = |(Z/mZ)^*|, phi(p) = p - 1, phi(p^2) = p^2 - p
% next_goal: Transition to fields of fractions and the universal property of Q
% open_loops: none
\end{hidden}

\begin{content}[Euler's Totient Function and Prime Power Formulas]
\begin{definition}[Euler's Totient Function]\label[definition]{def:eulers-totient}
The \newterm{Euler totient function} (or Euler's $\phi$-function) $\phi \colon \mathbb{Z}_{\ge 1} \to \mathbb{Z}_{\ge 1}$ is defined as the order of the group of units of $\mathbb{Z}/m\mathbb{Z}$:
\begin{equation}\label{eq:totient-def}
\phi(m) := |(\mathbb{Z}/m\mathbb{Z})^*| = \#\{ x \in \{0, 1, \dots, m-1\} \mid \gcd(x, m) = 1 \}.
\end{equation}
\end{definition}

\begin{proposition}[Totient Values for Primes and Prime Squares]\label[proposition]{prop:phi-prime-powers}
Let $p$ be a prime number. Then:
\begin{align}
\phi(p) &= p - 1, \label{eq:phi-prime} \\
\phi(p^2) &= p^2 - p = p(p - 1). \label{eq:phi-prime-square}
\end{align}
\end{proposition}

\begin{proof}
We evaluate the number of coprime residue classes in each case:
\begin{itemize}
    \item \textbf{For a prime $p$:} The representatives are $\{0, 1, 2, \dots, p - 1\}$. The only representative sharing a non-trivial factor with $p$ is $0$. Thus, the remaining $p - 1$ elements are all coprime to $p$, giving $\phi(p) = p - 1$ (e.g., $\phi(7) = 6$, as in \cref{ex:units-z7z}).
    \item \textbf{For $p^2$:} The representatives are $\{0, 1, 2, \dots, p^2 - 1\}$. An element $x$ shares a non-trivial factor with $p^2$ if and only if $p \mid x$. The multiples of $p$ in this range are:
    \[
    0, \ p, \ 2p, \ \dots, \ (p - 1)p.
    \]
    There are exactly $p$ such multiples. Subtracting these from the total of $p^2$ elements yields:
    \[
    \phi(p^2) = p^2 - p = p(p - 1).
    \]
\end{itemize}
\end{proof}

\begin{steps}[Generalization to Prime Powers]
For any prime power $p^k$ with $k \ge 1$, an integer $x \in \{0, 1, \dots, p^k - 1\}$ shares a factor with $p^k$ if and only if $p \mid x$. The multiples of $p$ are $0, p, 2p, \dots, (p^{k-1}-1)p$, totaling $p^{k-1}$ elements. Hence, $\phi(p^k) = p^k - p^{k-1} = p^k (1 - 1/p)$.
\end{steps}
\end{content}

\section{Fields of Fractions}

\subsection{The Universal Property of the Rationals}

\begin{speech}[00:13:48 - 00:15:15]
The next topic in the course has to do with fields of fractions. This is a peculiar subject. 

It starts with this idea that $\mathbb{Z}$ is inside $\mathbb{Q}$. $\mathbb{Z}$ is a ring, it's an integral domain, and $\mathbb{Q}$ is a field, and $\mathbb{Z}$ sits inside $\mathbb{Q}$. But their relationship is not just that of any ring inside any field. In some sense, $\mathbb{Q}$ is the smallest field that contains $\mathbb{Z}$. 

This understanding will lead to the study of the field of fractions. So, let's try to first make precise the qualitative statement that... 

Here's the statement: \qt{$\mathbb{Q}$ is the smallest field which contains $\mathbb{Z}$}. That's a statement that probably you feel is true. It's not a $100$ percent mathematical statement, so let's try to make this precise. You don't contradict the meaning of this statement, right?
\end{speech}

\begin{speech}[00:15:15 - 00:17:24]
One thing that happens when one starts studying these fields of fractions and $\mathbb{Q}$ is that you get this strange feeling that maybe you never checked everything in $\mathbb{Q}$ works. Have you checked everything in $\mathbb{Q}$ works? It's kind of a strange feeling, because you learn about the rational numbers in school, but it's not like the definitions there are checked to work. You just accept them, which is pretty reasonable --- they are correct. 

What I mean by that is that $\mathbb{Q}$ is a much more complicated thing than $\mathbb{Z}$ conceptually, because $\mathbb{Z}$ is really great: there are these objects and they have names, and everybody has one name. That's what's great about $\mathbb{Z}$. 

The thing that's confusing about $\mathbb{Q}$ is that there are objects, but they have many names. You can have $5/7 = 60/84$, and that's kind of confusing: all these numbers have many, many different names. 

And when I say \qt{have you ever checked it works}, it means that whatever operations you do, you have to be able to do them with any name, and it has to give the same answer, so to speak. That's what I mean. 

So, the question is: have you ever checked that? It's a weird question. But kind of the ideal math question in first grade or second grade, when you're adding fractions, is that you should say: \qt{But why is it well-defined?} That's the question you should ask. And if you'd asked that question, then we wouldn't have to give this lecture today, actually.
\end{speech}

\begin{speech}[00:17:24 - 00:19:45]
Okay, so I will try to make precise the statement that $\mathbb{Q}$ is the smallest field which contains $\mathbb{Z}$. There are different ways to do it, but one way is to say it like this: 

Suppose I have $f \colon \mathbb{Z} \to F$, where $F$ is a field, so this big $F$ is a field, and this little $f$ is an injection of rings --- an injective homomorphism of rings. Okay, so I start with this. I take any possible field, like maybe a field that I might think is smaller than $\mathbb{Q}$, or whatever field --- take any field you want --- and I start with a ring homomorphism which is injective. That means that I view this as $\mathbb{Z}$ sitting inside $F$, okay? 

And then I claim... The typical way you write this is with a diagram: so this is $F$, and then I claim there exists a unique $g$ --- so this is $\mathbb{Z}$ sitting inside $\mathbb{Q}$, and there exists a unique $g$ --- such that $g$ is an injective homomorphism of fields, and the diagram commutes. That means I can take an integer here and go here, or I can consider it in $\mathbb{Q}$ and go this way, so that for all $x \in \mathbb{Z}$, $f(x) = g(x)$. 

We would say the diagram commutes because if you have $x$ here, I can get to here by going straight across, or I can go inside $\mathbb{Q}$ and come around; that would say that $f(x) = g(x)$.
\end{speech}

\begin{speech}[00:19:45 - 00:20:08]
Okay. I'll give a little proof of this, but this is one way --- this is now totally precise mathematics --- that makes that statement, which is somehow slightly qualitative, into a mathematical statement. 

Do you accept that? When I say $\mathbb{Q}$ is the smallest field which contains $\mathbb{Z}$, it means that if I have any other candidate field $F$, I can actually squeeze $\mathbb{Q}$ into it. And that means that $\mathbb{Q}$ is smaller than $F$. Does that make sense?
\end{speech}

\begin{hidden}
% timestamp: 00:20:00
% topic: Formulation of the Universal Property of Q as the field of fractions of Z
% board_state: Z \subset Q, f: Z -> F injection, g: Q -> F unique injection, commutative diagram
% next_goal: Construct g(x/y) = f(x)/f(y) and verify well-definedness
% open_loops: none
\end{hidden}

\begin{content}[Universal Property of the Rational Field]
\textbf{Qualitative Observation:} $\mathbb{Q}$ is the \qt{smallest} field containing the integral domain $\mathbb{Z}$.

\begin{theorem}[Universal Property of $\mathbb{Q}$]\label[theorem]{thm:universal-property-q}
Let $F$ be a field, and let $f \colon \mathbb{Z} \to F$ be an injective ring homomorphism (embedding of rings).
Then there exists a unique field homomorphism $g \colon \mathbb{Q} \to F$ such that the following diagram commutes:
\begin{center}
\begin{tikzpicture}[scale=1.5, >=Stealth]
% \begin{hidden}
% - Canvas: scale=1.5, triangle diagram with vertices Z at (0,0), Q at (0,1.5), F at (2.2,0)
% - Nodes: Z at (0,0), Q at (0,1.5), F at (2.2,0)
% - Arrows: Z -> Q canonical inclusion \iota (vertical up)
%   Z -> F map f (horizontal right)
%   Q -> F unique map g (dashed diagonal down-right)
% \end{hidden}
    \node (Z) at (0, 0) {$\mathbb{Z}$};
    \node (Q) at (0, 1.5) {$\mathbb{Q}$};
    \node (F) at (2.2, 0) {$F$};

    \draw[right hook->, thick] (Z) -- (Q) node[midway, left] {$\iota$};
    \draw[->, thick] (Z) -- (F) node[midway, below] {$f$};
    \draw[->, dashed, thick, BrickRed] (Q) -- (F) node[midway, above right] {$\exists! \, g$};
\end{tikzpicture}
\end{center}
that is, $g|_{\mathbb{Z}} = f$, meaning $g(x) = f(x)$ for all $x \in \mathbb{Z}$. Furthermore, $g$ is automatically injective.
\end{theorem}

Injectivity is automatic, because every ring homomorphism from a field to a non-zero ring is injective (see \cref{thm:ideals-in-a-field} and the steps after it, Week~1).

\begin{steps}[Universal Property as Categorical Minimality]
In category theory, this universal mapping property identifies $\mathbb{Q}$ as the initial object among all field extensions receiving $\mathbb{Z}$. Any field $F$ containing a copy of $\mathbb{Z}$ must necessarily contain a subfield isomorphic to $\mathbb{Q}$, showing that $\mathbb{Q}$ cannot be compressed or embedded into any strictly smaller field.
\end{steps}
\end{content}

\subsection{Construction and Well-Definedness of the Field Extension}

\begin{speech}[00:20:08 - 00:21:23]
Okay, but I haven't proven it, so let's try to prove it. I won't prove the whole thing, because it's, of course, kind of painful. 

Let us see: I have to prove this, which means I get to start with the data I suppose. That's what \qt{suppose} is about. \qt{Suppose} means I get to start with this data, which is here, all right? And then the whole claim is existence. So, to prove this, what I have to do is produce $g$. 

Proof, more or less: I have to construct $g$. So, I have to write down... rational numbers: $x$ over $y$, where $x$ and $y$ are integers, and $y \ne 0$. That's what a rational number is: $x$ over $y$, $x$ and $y$ are integers, and $y$ is not equal to zero. I hope you're happy with that; that's what a rational number is, okay? 

And I have to define this --- this is my responsibility --- or we define: how am I going to define this?
\end{speech}

\begin{student}[Student Answer]
$f(x) / f(y)$.
\end{student}

\begin{speech}[00:21:26 - 00:21:51]
Yeah. Formally, of course, we'll write it the way you said, but formally we write it like this: $f(x) \cdot f(y)^{-1}$, okay? Because $f(x)$ is in the field $F$, and $f(y)$ is in the field $F$, and I want to take its inverse. This is the kind of field language. 

Why am I allowed to take its inverse? You cannot take the inverse of everything. Yeah?
\end{speech}

\begin{student}[Student Answer]
Because $F$ is a field and we chose, um... $y \ne 0$.
\end{student}

\begin{speech}[00:21:58 - 00:23:12]
Yeah, but you have to check --- and this is the whole thing --- you have to check that this is legal so long as $f(y)$ is not zero. And why is $f(y)$ not zero? It's because $f$ is injective and $y \ne 0$. That's actually using this thing here, okay? So, this is legal. Is that clear?

And as was implicit in the answer, we will often write this as $f(x)/f(y)$, because it's charming to do that. In a field, if you see this, all that means is that the fellow downstairs can be put upstairs by multiplicative inverse. To write this, you're only legally allowed to write it if you know this is not zero. Just like normally, you can't divide by zero. So, when we write this, it means in particular this one's not zero, and you can put it on either side --- it doesn't matter, because our fields are commutative, okay?
\end{speech}

\begin{speech}[00:23:12 - 00:24:24]
Okay, so we managed to... but is this a definition? The answer, of course, is not! Because all you did was pick a particular name of this and give a definition. You have to check that all the names give the same definition. That's the business about the rational numbers which no one ever thinks about, and it's better not to think about it, but unfortunately today we have to think about it.
\inlinemetanote{The lecturer pauses to lower and erase the middle chalkboards}
\end{speech}

% --- TEIL 2 ---

\begin{content}[Candidate Definition of the Field Homomorphism \texorpdfstring{$g$}{g}]
\begin{definition}[Candidate Map for Field Embedding]\label[definition]{def:candidate-map-g}
Let $f \colon \mathbb{Z} \to F$ be an injective ring homomorphism into a field $F$. For a rational number represented by $\frac{x}{y} \in \mathbb{Q}$ with $x, y \in \mathbb{Z}$ and $y \ne 0$, we define:
\begin{equation}\label{eq:candidate-g-def}
g\left(\frac{x}{y}\right) := f(x) \cdot f(y)^{-1} = \frac{f(x)}{f(y)} \in F.
\end{equation}
\end{definition}

\begin{steps}[Legitimacy of Inversion in the Field $F$]
The inverse $f(y)^{-1}$ exists in $F$ if and only if $f(y) \ne 0_F$:
\begin{enumerate}[label=\textbf{(\arabic*)}]
    \setcounter{enumi}{0} \item \textbf{Non-Zero Denominator:} In the rational representation $\frac{x}{y}$, we require $y \ne 0$ in $\mathbb{Z}$.
    \setcounter{enumi}{1} \item \textbf{Injectivity of Ring Embedding:} Since $f \colon \mathbb{Z} \to F$ is an injective ring homomorphism, its kernel is trivial: $\ker(f) = \{0\}$.
    \setcounter{enumi}{2} \item \textbf{Non-Vanishing Image:} It follows immediately that $y \ne 0 \implies f(y) \ne f(0) = 0_F$ (\cref{prop:homomorphism-preserves-zero}). Since $F$ is a field, every non-zero element has a unique multiplicative inverse $f(y)^{-1} \in F$.
\end{enumerate}
\end{steps}
\end{content}

\subsection{Well-Definedness of the Field Extension}

\begin{speech}[00:24:25 - 00:26:11]
Okay, so what do we have to do? We aren't done with the definition, because we must further check that if $x_1/y_1$ is equal to $x_2/y_2$, then $g(x_1/y_1)$ equals $g(x_2/y_2)$. You understand we have to check this, because that's the kind of slipperiness of the rational numbers: they don't have a unique name. All right. But you accept we have to check this, otherwise whatever was there was nonsense.

So, we have to check, and how are we going to check it? We have $g(x_1/y_1)$, and this is defined to be $f(x_1)/f(y_1)$, and $g(x_2/y_2)$ is $f(x_2)/f(y_2)$. We want to figure out whether these things are equal. So, we need to know: do we have this equality?
\end{speech}

\begin{speech}[00:26:11 - 00:27:58]
Using the fact that these are invertible and we can multiply by them, this question is equivalent to what's called clearing the denominators. By clearing the denominators, this is equivalent to $f(x_1) f(y_2) = f(x_2) f(y_1)$ by just multiplying both sides by $f(y_1) f(y_2)$, okay? That's a legal move in fields. And since these are invertible, this implies that, and it implies the other way, because they're both non-zero. Everything here is invertible.

We're using all the time that when we write this, neither $y_1$ nor $y_2$ is zero, which implies $f(y_1)$ and $f(y_2)$ are not zero, okay?

And now, these things are all integers, and everything's upstairs --- there's no division. This is if and only if --- using the homomorphism property --- $f(x_1 y_2) = f(x_2 y_1)$, right? Because $f$ is a homomorphism, okay?

And this is if and only if $f(x_1 y_2 - x_2 y_1) = 0$, okay? And why would this be the case? Yeah?
\end{speech}

\begin{student}[Student Answer]
We could rearrange...
\end{student}

\begin{speech}[00:27:59 - 00:29:02]
Yeah, we have to say what this means. What does this mean? This goes to the heart of the issue about whether you learned the rational numbers correctly, and the answer is: this means that $x_1 y_2 = x_2 y_1$. It means nothing more and nothing less than that --- it just actually means that, okay?

So, if we start with this, we have this equation, which is $x_1 y_2 - x_2 y_1 = 0$ in the integers. That's what this means.

And very fortunately, that's what we need: we need this to be zero. \inlinemetanote{puts a checkmark on the board} So, this is true. And so everything is true, and this is well-defined, all right?

That means I have defined a map. But there are a lot of things to check, which of course I'm not going to check. You have to do it, so to speak, on your own time.
\end{speech}

\begin{hidden}
% timestamp: 00:28:40
% topic: Well-definedness verification of the canonical embedding g: Q -> F
% board_state: candidate-g-def, well-definedness chain: g(x1/y1) = g(x2/y2) <=> f(x1 y2 - x2 y1) = 0
% next_goal: Discuss field homomorphism properties (preserving addition and multiplication)
% open_loops: none
\end{hidden}

\begin{content}[Well-Definedness of the Extension Map \texorpdfstring{$g$}{g}]
\begin{claim}[Well-Definedness of $g$]\label[claim]{clm:g-well-defined}
The map $g \colon \mathbb{Q} \to F$ defined in \eqref{eq:candidate-g-def} is well-defined; that is, if $\frac{x_1}{y_1} = \frac{x_2}{y_2}$ in $\mathbb{Q}$, then $g\left(\frac{x_1}{y_1}\right) = g\left(\frac{x_2}{y_2}\right)$ in $F$.
\end{claim}

\begin{proof}
Let $\frac{x_1}{y_1}, \frac{x_2}{y_2} \in \mathbb{Q}$ with $y_1, y_2 \in \mathbb{Z} \setminus \{0\}$ such that $\frac{x_1}{y_1} = \frac{x_2}{y_2}$.

\textbf{1. Meaning of Equality in $\mathbb{Q}$:}
By definition of the rational numbers as fractions, equality $\frac{x_1}{y_1} = \frac{x_2}{y_2}$ means precisely that the cross-multiplied relation holds in $\mathbb{Z}$:
\begin{equation}\label{eq:cross-mult-z}
x_1 y_2 = x_2 y_1 \iff x_1 y_2 - x_2 y_1 = 0 \in \mathbb{Z}.
\end{equation}

\textbf{2. Algebraic Chain in $F$:}
We want to establish whether:
\[
g\left(\frac{x_1}{y_1}\right) \overset{?}{=} g\left(\frac{x_2}{y_2}\right) \iff \frac{f(x_1)}{f(y_1)} \overset{?}{=} \frac{f(x_2)}{f(y_2)}.
\]
Since $y_1, y_2 \ne 0$ and $f$ is injective, $f(y_1) \ne 0$ and $f(y_2) \ne 0$. Multiplying both sides by $f(y_1)f(y_2) \in F^\times$ (clearing denominators) yields an equivalent equality in $F$:
\begin{align}
\frac{f(x_1)}{f(y_1)} = \frac{f(x_2)}{f(y_2)} &\iff f(x_1) \cdot f(y_2) = f(x_2) \cdot f(y_1) \nonumber \\
&\iff f(x_1 y_2) = f(x_2 y_1) && \text{(since $f$ is a ring homomorphism)} \nonumber \\
&\iff f(x_1 y_2) - f(x_2 y_1) = 0_F \nonumber \\
&\iff f(x_1 y_2 - x_2 y_1) = 0_F. \label{eq:hom-zero-check}
\end{align}

\textbf{3. Conclusion:}
By \eqref{eq:cross-mult-z}, $x_1 y_2 - x_2 y_1 = 0 \in \mathbb{Z}$. Since any ring homomorphism maps $0$ to $0$ (\cref{prop:homomorphism-preserves-zero}):
\[
f(x_1 y_2 - x_2 y_1) = f(0) = 0_F.
\]
Thus, condition \eqref{eq:hom-zero-check} holds, which proves that $g\left(\frac{x_1}{y_1}\right) = g\left(\frac{x_2}{y_2}\right)$.
\end{proof}

\begin{steps}[Core Pedagogical Principle: Non-Unique Names of Fractions]
A rational number has infinitely many equivalent names (e.g., $\frac{1}{2} = \frac{2}{4} = \frac{-3}{-6}$). Any function defined on rational numbers using numerator and denominator symbols must be proven invariant under replacement of representatives. As shown above, this reduces entirely to the cross-multiplication identity in the underlying ring $\mathbb{Z}$.
\end{steps}
\end{content}

\subsection{Verification of Field Homomorphism Properties}

\begin{speech}[00:29:02 - 00:30:32]
I have now defined the map $g$, finally. Then you have to check that this map $g$ is actually... I mean, all of the properties of $g$, that it's a homomorphism of fields. You have to check that it respects addition, and you have to check that it respects multiplication.

All of it is slightly painful, because at first blush it's obvious: you just pick representatives for $\mathbb{Q}$ as $x_1/y_1$, and then you just check it. But then you have to check that all of that stuff respects the different choices, the different names of elements of $\mathbb{Q}$, just like we did here.

And the checks are all the same, so you can do it yourself. That's the end of what I'm going to give as the proof of this; you can check the rest. We're going to have to do many more things like this, so...
\end{speech}

\begin{speech}[00:30:32 - 00:32:18]
Okay, so I leave this to you. We're going to have to do things like this anyway, and it's all just unraveling the definitions. Meaning you have to check... \inlinemetanote{writes You have to check $g(\frac{x_1}{y_1}) + g(\frac{x_2}{y_2}) = g(\frac{x_1}{y_1} + \frac{x_2}{y_2})$} You have to check something like $g(x_1/y_1) + g(x_2/y_2) = g(x_1/y_1 + x_2/y_2)$.

You have to go check that; that's what you have to check. And you can see, when you write this, that this side... Okay, I'm going to do part of it for you. I don't know why I'm feeling compassionate today here.

This one is $f(x_1)/f(y_1) + f(x_2)/f(y_2)$; that's the definition here. And this has a different definition. First of all, how do we do it? We first have to add this fraction, right? Because our definition only applies to names of rational numbers, and this is a sum, so we first have to do this. Well, you know how to add fractions: this is equal to $g$ of $(x_1 y_2 + x_2 y_1)/(y_1 y_2)$. At least I hope you know how to add fractions. At least I hope I know how to add fractions, so.

Okay, so that's the addition of the fraction. And then this, by definition, is $f(x_1 y_2 + x_2 y_1)/f(y_1 y_2)$, okay?
\end{speech}

\begin{speech}[00:32:18 - 00:32:48]
So, that's the one side to check, and that was supposed to equal this side, all right? And how are you going to check this? Well, you use the homomorphism property, you move everything upstairs, and then finally you're going to use the thing I erased, okay?

So, I leave it to you. There's nothing really I can add to this; you have to do it. But I do say it's painful --- painful for me to write it on the board even.
\end{speech}

\begin{content}[Verification of the Field Homomorphism Axioms]
\begin{claim}[Preservation of Field Operations]\label[claim]{clm:g-field-hom}
The map $g \colon \mathbb{Q} \to F$ is a field homomorphism; that is, for all $q_1, q_2 \in \mathbb{Q}$:
\begin{align}
g(q_1 + q_2) &= g(q_1) + g(q_2), \label{eq:g-add} \\
g(q_1 \cdot q_2) &= g(q_1) \cdot g(q_2), \label{eq:g-mult} \\
g(1_\mathbb{Q}) &= 1_F. \label{eq:g-unit}
\end{align}
\end{claim}

\begin{proof}
Let $q_1 = \frac{x_1}{y_1}$ and $q_2 = \frac{x_2}{y_2}$ with $x_i, y_i \in \mathbb{Z}$ and $y_i \ne 0$.

\paragraph{1. Verification of Additivity \eqref{eq:g-add}:}
Adding fractions according to standard rational arithmetic:
\[
q_1 + q_2 = \frac{x_1}{y_1} + \frac{x_2}{y_2} = \frac{x_1 y_2 + x_2 y_1}{y_1 y_2}.
\]
Applying the definition of $g$ and the ring homomorphism properties of $f$:
\begin{align*}
g(q_1 + q_2) &= g\left(\frac{x_1 y_2 + x_2 y_1}{y_1 y_2}\right) \\
&= \frac{f(x_1 y_2 + x_2 y_1)}{f(y_1 y_2)} \\
&= \frac{f(x_1)f(y_2) + f(x_2)f(y_1)}{f(y_1)f(y_2)}.
\end{align*}
On the other hand, evaluating the sum of images in the field $F$:
\begin{align*}
g(q_1) + g(q_2) &= \frac{f(x_1)}{f(y_1)} + \frac{f(x_2)}{f(y_2)} \\
&= f(x_1)f(y_1)^{-1} + f(x_2)f(y_2)^{-1}.
\end{align*}
Bringing both terms over the common denominator $f(y_1)f(y_2)$ in the field $F$:
\[
\frac{f(x_1)}{f(y_1)} + \frac{f(x_2)}{f(y_2)} = \frac{f(x_1)f(y_2) + f(x_2)f(y_1)}{f(y_1)f(y_2)}.
\]
Both expressions coincide identically, establishing \eqref{eq:g-add}. By \cref{clm:g-well-defined}, this value is independent of the choice of fraction representatives.

\paragraph{2. Verification of Multiplicativity \eqref{eq:g-mult}:}
Multiplying fractions inside $\mathbb{Q}$:
\[
q_1 \cdot q_2 = \frac{x_1}{y_1} \cdot \frac{x_2}{y_2} = \frac{x_1 x_2}{y_1 y_2}.
\]
Applying the definition of $g$ and the multiplicativity of $f$:
\begin{align*}
g(q_1 \cdot q_2) &= g\left(\frac{x_1 x_2}{y_1 y_2}\right) \\
&= \frac{f(x_1 x_2)}{f(y_1 y_2)} \\
&= \frac{f(x_1) f(x_2)}{f(y_1) f(y_2)} \\
&= \frac{f(x_1)}{f(y_1)} \cdot \frac{f(x_2)}{f(y_2)} \\
&= g(q_1) \cdot g(q_2),
\end{align*}
establishing \eqref{eq:g-mult}.

\paragraph{3. Preservation of Multiplicative Identity \eqref{eq:g-unit}:}
Representing $1_\mathbb{Q} \in \mathbb{Q}$ as $\frac{1_\mathbb{Z}}{1_\mathbb{Z}}$:
\[
g(1_\mathbb{Q}) = g\left(\frac{1_\mathbb{Z}}{1_\mathbb{Z}}\right) = \frac{f(1_\mathbb{Z})}{f(1_\mathbb{Z})}.
\]
Since $f \colon \mathbb{Z} \to F$ is a unital ring homomorphism, $f(1_\mathbb{Z}) = 1_F$. Thus:
\[
g(1_\mathbb{Q}) = \frac{1_F}{1_F} = 1_F,
\]
establishing \eqref{eq:g-unit}. Thus $g$ is a field homomorphism.
\end{proof}
\end{content}

\section{The Field of Fractions of an Integral Domain}

\subsection{Motivation: Embedding an Integral Domain into its Smallest Field}

\begin{speech}[00:32:48 - 00:34:30]
Okay, so that's the story about $\mathbb{Z}$. Their fates are really closely tied, $\mathbb{Z}$ and $\mathbb{Q}$. $\mathbb{Q}$ is not just any field that $\mathbb{Z}$ is in; it's the smallest field it's in.

So, we could ask: can we construct this relationship for other rings? Because we really like this relationship. If you study the integers, it's very useful sometimes to be able to divide. So, that's the movement to the field. Let's try this with a ring in general.

Suppose $R$ is an integral domain. It turns out that this is kind of extremely important, and we will see in the definition --- it may not be clear from the discussion so far --- but let's say $R$ is an integral domain. Can we find the \qt{smallest} field containing $R$?

Already in this language, if we find a field containing $R$, that means $R$ has to be a domain (Recall: \cref{def:integral-domain}). Because domain means that when two elements are non-zero, their product stays non-zero; so if you're inside a field, that's always the case, okay? This is one of the reasons why we're talking only about integral domains.

So, that's the question: if you find on the street an integral domain, can you put it in its smallest field? That's the question. And the answer is yes, and that construction is called the field of fractions. And I will show you how to do it.
\end{speech}

\begin{speech}[00:34:30 - 00:36:15]
In some ways, there's lots of conceptual nice things about it, but this pain, you can imagine, we're not going to get away from. It's just going to be more abstract.

One of the things you will see here is that when we define the field of fractions of a ring, it will have this property that the elements are not single objects, but equivalence classes. And you should expect that answer, because that already happens for the rational numbers.

By the way, for a rational number, the objection --- and I'm surprised no one made it --- is that you could say that you could fix the name by taking the numerator and denominator and canceling all common factors. There's still a sign issue, but we leave that aside. If you try to do arithmetic with that, it's super painful, because you have to do a lot of cancellation. So, that is the argument against it; that's not really how you do arithmetic in the rational numbers. But that is the argument against what I was saying about the rational numbers.

In a sense, if you do that, it's kind of even more painful, because then when you write these expressions, you can't write these expressions --- you have to cancel everywhere, and you know...
\end{speech}

\begin{hidden}
% timestamp: 00:36:10
% topic: Transition from Z subset Q to general field of fractions of an integral domain R
% board_state: Z \subset Q, R an integral domain, question: smallest field containing R?
% next_goal: Define set of pairs S = {(x, y) \in R x R | y != 0} and equivalence relation
% open_loops: none
\end{hidden}

\begin{content}[Generalization from Integers to Integral Domains]
\begin{remark}[Why Integral Domains?]\label[remark]{rem:why-integral-domains}
If a commutative ring $R$ embeds into a field $F$ via an injective homomorphism $\iota \colon R \hookrightarrow F$, then $R$ must necessarily have no zero divisors:
\[
a, b \in R \setminus \{0\} \implies \iota(a), \iota(b) \ne 0 \implies \iota(ab) = \iota(a)\iota(b) \ne 0 \implies ab \ne 0.
\]
Moreover, $R$ is not the zero ring: $\iota(1) = 1 \neq 0 = \iota(0)$ in $F$, so $1 \neq 0$ in $R$. Thus, the property of being an \newterm{integral domain} (\cref{def:integral-domain}) is a strictly necessary prerequisite for any ring to be embedded into a field of fractions.
\end{remark}

\begin{highlight}{MidnightBlue}[The Goal: Field of Fractions]
Given an arbitrary integral domain $R$, our goal is to construct the minimal field containing $R$, called the \newterm{field of fractions} (or quotient field) $\operatorname{Frac}(R)$, generalizing the classical relationship:
\[
\mathbb{Z} \hookrightarrow \mathbb{Q} = \operatorname{Frac}(\mathbb{Z}).
\]
\end{highlight}
\end{content}

\subsection{The Set of Pairs and Equivalence Relations}

\begin{speech}[00:36:15 - 00:38:08]
Let's think about this in the abstract. The main thing is just the definition. You can imagine that here's our ring $R$ --- we don't know what it is, it's just any ring, but it's an integral domain. We want to construct a field containing $R$.

The field should have elements that are like $x/y$, where $x, y \in R$ and $y \ne 0$. This should be the element of that field, because if it contains $R$, I should be able to take $x$ and multiply it by the inverse of $y$ in the field, so it should have these elements in it. Is that clear? That's what happened with the rational numbers. So, we can just define it that way.

So, let $S$ be the set of pairs. This whole thing is going to be a big disappointment if you think something fun is going to happen. I just tell you ahead of time: this is going to be a big disappointment. But some things in life are like that.

Okay, that's the set of pairs. For the integers, if we started with integers, we'd just get the set of pairs of integers where the second one is not zero, okay? But we can think about the set of pairs for any $R$, is that clear? This is inside $R \times R$; it's just some subset of pairs.
\end{speech}

\begin{speech}[00:38:08 - 00:39:18]
All right, and now the thing that would come is that there is an equivalence relation on $S$. Maybe we call this equivalence relation three lines --- you can use various symbols for equivalence, that's kind of standard. How about three lines?

Do you know what an equivalence relation is? This was in your foundations class. An equivalence relation on a set $T$ tells you when two elements are related, and the rules for the equivalence relation are: $x$ is always related to itself (this is the reflexive property); if $x$ is related to $y$, then $y$ is related to $x$ (this is the symmetric property).

And if you know the definition, what's the third one? Yeah?
\end{speech}

\begin{student}[Student Answer]
Transitive.
\end{student}

\begin{speech}[00:39:20 - 00:40:21]
Transitive property: it says that if we have $x$ related to $y$ and $y$ related to $z$, then $x$ is related to $z$. That's the whole thing about equivalence relations.

It comes up many times in mathematics; it will come up again in this course, so if you have not seen it, think about it. And the main thing we use about this equivalence relation is that if I have $T$ with this equivalence relation, then I have another set, which is the set of equivalence classes: $T/\!\equiv$ is the set of equivalence classes.

And that's going to be very important for us. Whenever we have a set with an equivalence relation, that equivalence relation divides the set up into disjoint equivalence classes, okay? And this thing here is the set of equivalence classes.

There are three objects in this kind of abstract situation: there's the set that we start with, there's the equivalence relation, and then there's the set of equivalence classes.
\end{speech}

\begin{content}[The Set of Pairs and General Equivalence Relations]
\begin{definition}[The Set of Formal Fractions $S$]\label[definition]{def:set-of-pairs-s}
Let $R$ be an integral domain. We define the set of formal numerator-denominator pairs $S$ by:
\begin{equation}\label{eq:set-of-pairs}
S := \{ (x, y) \in R \times R \mid y \ne 0 \} \subset R \times R.
\end{equation}
\end{definition}

\begin{remark}[Review: Equivalence Relations and Partitions]\label[remark]{rem:equivalence-relation-review}
An \newterm{equivalence relation} $\equiv$ on a set $T$ is a binary relation satisfying three foundational axioms:
\begin{enumerate}[label=\textbf{(\roman*)}]
    \item \textbf{Reflexivity:} $\forall x \in T \colon x \equiv x$.
    \item \textbf{Symmetry:} $\forall x, y \in T \colon x \equiv y \implies y \equiv x$.
    \item \textbf{Transitivity:} $\forall x, y, z \in T \colon (x \equiv y \land y \equiv z) \implies x \equiv z$.
\end{enumerate}
The quotient set $T/\!\equiv$ denotes the set of \newterm{equivalence classes}, which partitions $T$ into mutually disjoint subsets.
\end{remark}

\begin{center}
\begin{tikzpicture}[scale=1.2]
% \begin{hidden}
% - Canvas: scale=1.2, set T as an outer rectangle [0,5] x [0,2.8], divided into 3 disjoint blocks
% - Subsets: T_1 at [0,1.7], T_2 at [1.7,3.4], T_3 at [3.4,5]
% - Elements: dots inside each class representing elements related by \equiv
% - Labels: T_1, T_2, T_3 above each partition, T overall
% \end{hidden}
    % Outer bounding box T
    \draw[thick, gray!80, fill=gray!5] (0,0) rectangle (5, 2.5);
    \node[above right, font=\bfseries] at (0, 2.5) {Set $T$};

    % Vertical partition lines
    \draw[thick, dashed, MidnightBlue] (1.65, 0) -- (1.65, 2.5);
    \draw[thick, dashed, MidnightBlue] (3.35, 0) -- (3.35, 2.5);

    % Equivalence Classes
    \node[MidnightBlue, font=\bfseries] at (0.8, 2.1) {$T_1 = [x_1]$};
    \node[MidnightBlue, font=\bfseries] at (2.5, 2.1) {$T_2 = [x_2]$};
    \node[MidnightBlue, font=\bfseries] at (4.2, 2.1) {$T_3 = [x_3]$};

    % Representative dots: labels left of the first and right of the second point, away from the arrows
    \fill (0.5, 1.0) circle (1.5pt) node[left] {\footnotesize $a$};
    \fill (1.1, 0.7) circle (1.5pt) node[right] {\footnotesize $b$};
    \draw[<->, shorten <=2pt, shorten >=2pt, thick, ForestGreen] (0.5, 1.0) -- (1.1, 0.7);

    \fill (2.2, 1.2) circle (1.5pt) node[left] {\footnotesize $c$};
    \fill (2.8, 0.9) circle (1.5pt) node[right] {\footnotesize $d$};
    \draw[<->, shorten <=2pt, shorten >=2pt, thick, ForestGreen] (2.2, 1.2) -- (2.8, 0.9);

    \fill (3.8, 1.1) circle (1.5pt) node[left] {\footnotesize $e$};
    \fill (4.5, 0.6) circle (1.5pt) node[right] {\footnotesize $f$};
    \draw[<->, shorten <=2pt, shorten >=2pt, thick, ForestGreen] (3.8, 1.1) -- (4.5, 0.6);

    \node[below, font=\small] at (2.5, -0.2) {Partition of $T$ into disjoint equivalence classes $T/\!\equiv$};
\end{tikzpicture}
\end{center}
\end{content}

\subsection{The Equivalence Relation on Pairs}

\begin{speech}[00:40:21 - 00:41:03]
We have to define the equivalence relation. So, we now define this equivalence relation on our set $S$, which is the set of pairs. You should think about these pairs as numerator and denominator, but those are just words here.

We need to define: when is $(x_1, y_1)$ equivalent to $(x_2, y_2)$? We need a definition for this.

Okay, so what's the answer going to be? Yeah?
\end{speech}

\begin{student}[Student Answer]
$x_1 y_2 = x_2 y_1$.
\end{student}

\begin{speech}[00:41:06 - 00:42:06]
Very good. So, the definition is just exactly what you do with integer arithmetic: if you have a fraction, you say they're equal if $x_1 y_2 = x_2 y_1$, okay? That's the definition. There's nothing more or less than that; that's the whole definition. 

And it's motivated by the integers. It's very important for this definition that this is only about $R$, right? This whole discussion starts with $R$, so we can only say things that are about $R$. We can make these pairs --- that's only about $R$. We want to define this equivalence relation; at the end, these things are only in $R$, and the definition is only about some equation in $R$, okay? So, we haven't gotten out of $R$, in some sense, in this whole discussion. Is that clear?
\end{speech}

\begin{hidden}
% timestamp: 00:42:00
% topic: Definition of the equivalence relation on S and verification of reflexivity, symmetry, transitivity
% board_state: def:equiv-relation-S, 3 properties on board
% next_goal: Verify transitivity using the integral domain property
% open_loops: none
\end{hidden}

\begin{content}[Definition of the Equivalence Relation on \texorpdfstring{$S$}{S}]
\begin{definition}[Equivalence Relation on $S$]\label[definition]{def:equiv-relation-S}
On the set $S = \{ (x, y) \in R \times R \mid y \ne 0 \}$, we define the relation $\equiv$ by:
\begin{equation}\label{eq:equiv-relation-formula}
(x_1, y_1) \equiv (x_2, y_2) :\iff x_1 y_2 = x_2 y_1 \quad (\text{in } R).
\end{equation}
\end{definition}

\begin{steps}[Conceptual Motivation: Cross-Multiplication]
This relation is the formal algebraic counterpart to equating cross-products of fractions:
\[
\frac{x_1}{y_1} = \frac{x_2}{y_2} \iff x_1 y_2 = x_2 y_1.
\]
Crucially, formula \eqref{eq:equiv-relation-formula} involves only multiplication within the ring $R$, without invoking any external concept of division or fraction.
\end{steps}
\end{content}

\subsection{Verification of Equivalence Relation Properties}

\begin{speech}[00:42:06 - 00:44:09]
Okay, so now we come to an issue: why is this an equivalence relation? This is the question that you should have asked your teacher in first grade or whatever it was. That's like the whole question about the rational numbers, in some sense --- I mean the foundations.

\inlinemetanote{erases the lower chalkboard}

This is the first part where we have to do a little bit of algebra. 

The first two properties: $x = x$, that's just identically true, okay? That's not a problem. We have to check three properties. 

The first is if I take $(x_1, y_1)$, this is equivalent to itself \inlinemetanote{equivalent to itself}. Why is this? Because by definition, this is equal to $x_1 y_1 = x_1 y_1$, so we're very happy about that; there was nothing to do here. 

You also have to check the symmetric property: that if we have $(x_1, y_1) \equiv (x_2, y_2)$, then the other way around is fine, and it'll just be the same equation. There's nothing to check about that.
\end{speech}

\begin{speech}[00:44:09 - 00:45:00]
So, the interesting one is the transitive property. This is the only fun one to check. We start with this --- these are the inputs: $(x_1, y_1) \equiv (x_2, y_2)$ and $(x_2, y_2) \equiv (x_3, y_3)$. And we must show --- this is for us to do --- that $(x_1, y_1) \equiv (x_3, y_3)$. 

The thing that's a little funny about this is that we have to write the first equation. This is what we have, that's what we've assumed, that's in the left column. What we have...
\inlinemetanote{bell chimes}

Okay, you'll have to wait until after the break for this excitement to continue.
\end{speech}

\begin{content}[Verification of Reflexivity and Symmetry on \texorpdfstring{$S$}{S}]
\begin{proposition}[Reflexivity and Symmetry of $\equiv$]\label[proposition]{prop:equiv-refl-sym}
The relation $\equiv$ defined on $S = \{ (x, y) \in R \times R \mid y \ne 0 \}$ by:
\begin{equation}\label{eq:equiv-def-recap}
(x_1, y_1) \equiv (x_2, y_2) \iff x_1 y_2 = x_2 y_1
\end{equation}
is reflexive and symmetric.
\end{proposition}

\begin{proof}
Let $(x_1, y_1), (x_2, y_2) \in S$:
\begin{itemize}
    \item \textbf{Reflexivity:} Trivially $x_1 y_1 = x_1 y_1$, which by \eqref{eq:equiv-def-recap} immediately implies:
    \[
    (x_1, y_1) \equiv (x_1, y_1).
    \]
    \item \textbf{Symmetry:} Assume $(x_1, y_1) \equiv (x_2, y_2)$. Then $x_1 y_2 = x_2 y_1$. By symmetry of equality in $R$:
    \[
    x_2 y_1 = x_1 y_2 \implies (x_2, y_2) \equiv (x_1, y_1).
    \]
\end{itemize}
\end{proof}
\end{content}

\begin{pause}[Mid-Lecture Break]
The lecturer announces the standard mid-lecture pause. The recording resumes after the break at timestamp 00:45:05.
\end{pause}

% --- TEIL 3 ---

\begin{speech}[00:45:05 - 00:46:37]
Okay, so let's start. I should say that very few people have found my office, and it turns out I'll be in this visitor office for most of the semester, because they're clearing some kind of nuclear waste from my current office. So, my office now is HG G 65.1, if you want to find me there. \inlinemetanote{writes \qt{HG G 65.1} on the chalkboard}

And yeah, I'm afraid I'll be there for most of the semester. 

We come to this point that's actually kind of fun and interesting. We started with an integral domain, and led by the analogy with the rational numbers, we come to defining this field as some equivalence class of pairs, where --- and I should have said this, and I did say it, but I didn't write it, so the way to think about this... \inlinemetanote{writes \qt{Think of as $x_1/y_1$} on the chalkboard}

This is the mathematical text: think of it as $x_1/y_1$, or $x/y$. You think of this pair in $S$ as $x/y$, okay? Just like with the integers. But as an object, it's this formal object. 

And then there was this proposal, which is correct, on how to define the equivalence relation. And now we need to actually check something: is it an equivalence relation? We have to prove the transitive property, and so we start with: what do we have?
\end{speech}

\begin{speech}[00:46:37 - 00:48:07]
What we have is that equivalence. If we translate that into equations, it says: $x_1 y_2$ --- and we have to be very careful of the subscripts, otherwise you'll confuse yourself --- $x_1 y_2 = x_2 y_1$.

And then what we also have here is $x_2 y_3 = x_3 y_2$.

Okay, and what we want is this thing, which is $x_1 y_3 = x_3 y_1$.

The puzzling thing about this is that we have these quadratic equations and we need to get another quadratic equation. This seems a little bit unfortunate, because you might think you have to do some linear algebra with quadratic equations, and that's not going to work --- you're not going to get a new quadratic equation from these.

So, the correct answer is that we're supposed to multiply here. The clever thing is to multiply the top equation by $y_3$, which is what we want. Multiply by $y_3$, and we get $x_1 y_2 y_3 = x_2 y_1 y_3$.

Okay, so this is kind of good: we see the $x_1 y_3$ here. Can you see what we have to do next? This is actually algebra. What?
\end{speech}

\begin{student}[Student Answer]
\inlinemetanote{suggests substituting from the second equation into the first}
\end{student}

\begin{speech}[00:48:08 - 00:49:55]
Yeah, you could do that, but why don't we just take this $y_2 y_3$ here and use this equation? It's the same thing. So, we take the $y_2 y_3$ and we use this equation to flip them.

This is all about subscripts, and I should just tell you that subscripts have always been difficult for me.

Do you accept that? Because I just took this and flipped them using this equation.

So, now what do we do? We have here $x_1 y_3$, that looks pretty good, and we have $y_3 y_1$, but they're both multiplied by $y_2$, right?

Now we use the fact that in these pairs, the second one is always born non-zero; that's part of being in this set: the second one has to be non-zero.

And this is an integral domain, so then we use the integral domain property.

As I said, this is a place where it's extremely important that it's an integral domain. Since $y_2 \ne 0$, we can cancel, okay? Because being an integral domain says that's the cancellation property. Is that okay?

If $y_2$ is not zero, we can cancel it. So, we cancel it. This is kind of dramatic, but we cancel it here. \inlinemetanote{cancels $y_2$ from both sides on the chalkboard}

And that's what we want! We have now proven transitivity. And you see there was a little bit of cleverness here; it's easy to confuse yourself.
\end{speech}

\begin{hidden}
% timestamp: 00:49:50
% topic: Transitivity of the equivalence relation on S
% board_state: prop:equiv-on-S, eq:trans-hyp1, eq:trans-hyp2, cancellation via integral domain
% next_goal: Formal definition of the field of fractions F = S/~ and canonical inclusion map i: R -> F
% open_loops: none
\end{hidden}

\begin{content}[Proof of Transitivity and the Role of the Integral Domain]
\begin{proposition}[Transitivity of $\equiv$ on $S$]\label[proposition]{prop:equiv-on-S}
Let $R$ be an integral domain. The relation $\equiv$ on $S$ defined by $(x_1, y_1) \equiv (x_2, y_2) \iff x_1 y_2 = x_2 y_1$ is transitive. Consequently, $\equiv$ is an equivalence relation on $S$.
\end{proposition}

\begin{proof}
Let $(x_1, y_1), (x_2, y_2), (x_3, y_3) \in S$, which implies $y_1, y_2, y_3 \in R \setminus \{0\}$.

\textbf{1. Assumptions:}
We assume the relations:
\begin{align}
(x_1, y_1) \equiv (x_2, y_2) &\implies x_1 y_2 = x_2 y_1, \label{eq:trans-hyp1} \\
(x_2, y_2) \equiv (x_3, y_3) &\implies x_2 y_3 = x_3 y_2. \label{eq:trans-hyp2}
\end{align}

\textbf{2. Target:}
We must prove that $(x_1, y_1) \equiv (x_3, y_3)$, which means showing that:
\begin{equation}\label{eq:trans-target}
x_1 y_3 = x_3 y_1.
\end{equation}

\textbf{3. Derivation via Multiplication and Substitution:}
Multiplying equation \eqref{eq:trans-hyp1} on both sides by $y_3 \in R$:
\begin{equation}\label{eq:trans-mult}
(x_1 y_2) y_3 = (x_2 y_1) y_3 \iff (x_1 y_3) y_2 = (x_2 y_3) y_1.
\end{equation}
Substituting $x_2 y_3 = x_3 y_2$ from \eqref{eq:trans-hyp2} into the right-hand side of \eqref{eq:trans-mult}:
\begin{equation}\label{eq:trans-subst}
(x_1 y_3) y_2 = (x_3 y_2) y_1 = (x_3 y_1) y_2.
\end{equation}
Subtracting the right-hand side gives:
\begin{equation}\label{eq:trans-factored}
(x_1 y_3 - x_3 y_1) y_2 = 0.
\end{equation}

\textbf{4. Cancellation via the Integral Domain Property:}
By definition of $S$, $y_2 \ne 0$. Because $R$ is an \textbf{integral domain}, it contains no zero divisors:
\[
ab = 0 \text{ with } b \ne 0 \implies a = 0.
\]
Applying this cancellation property to \eqref{eq:trans-factored} with $b = y_2 \ne 0$:
\[
x_1 y_3 - x_3 y_1 = 0 \implies x_1 y_3 = x_3 y_1.
\]
Thus, $(x_1, y_1) \equiv (x_3, y_3)$, which completes the proof.
\end{proof}

\begin{steps}[Why Integral Domains Are Essential for Transitivity]
Notice where the ring axioms alone would fail: equation \eqref{eq:trans-factored} holds in any commutative ring. However, to deduce $x_1 y_3 = x_3 y_1$ from $(x_1 y_3 - x_3 y_1)y_2 = 0$, we must divide out and cancel the non-zero element $y_2$. If $R$ has zero divisors, the relation $\equiv$ can \emph{fail} to be transitive. For instance, in $\mathbb{Z}/6\mathbb{Z}$ (with $y_2 = \bar{2}$ and $x_1 y_3 - x_3 y_1 = \bar{3}$): $(\bar{3}, \bar{1}) \equiv (\bar{0}, \bar{2})$ since $\bar{3} \cdot \bar{2} = \bar{0} = \bar{0} \cdot \bar{1}$, and $(\bar{0}, \bar{2}) \equiv (\bar{0}, \bar{1})$ since $\bar{0} \cdot \bar{1} = \bar{0} = \bar{0} \cdot \bar{2}$, but $(\bar{3}, \bar{1}) \not\equiv (\bar{0}, \bar{1})$ since $\bar{3} \cdot \bar{1} = \bar{3} \neq \bar{0} = \bar{0} \cdot \bar{1}$.
\end{steps}
\end{content}

\subsection{Definition of the Field of Fractions and Canonical Embedding}

\begin{speech}[00:49:55 - 00:51:06]
We have come so far as to define $S$ as a set with an equivalence relation. So, now we can define the points of the field that we're constructing, which will be the equivalence classes...

\inlinemetanote{The lecturer begins erasing the lower chalkboards}

This was almost like the first piece of algebra in this class, I think. I don't remember what we did; did we have any algebra before?

By that I mean playing with equations. Which, in the end... you know, I'm saying a lot of words, but in the end the subject is about equations, that's...
\end{speech}

\begin{speech}[00:51:06 - 00:52:37]
Okay, so where are we? We started with $R$ as an integral domain. We constructed $S$ with an equivalence relation, and from this we define $F$.

The name of this is the field of fractions of $R$ --- \qt{field of fractions of $R$}, that's the name of this object. $F$ is equal to the set of equivalence classes of $S$. Is it completely clear what this thing is?

We have defined $S$, we have defined an equivalence relation, we took great pains to check all of its properties, and so we get the set of equivalence classes, and that is $F$.

Then we need to do a lot of things. But let's start by defining the map which is the inclusion from $R$ to $F$, because we said this $F$ should be a field that contains $R$. That means we have to have some way of getting points from $R$ to $F$. And what's that way going to be?

This inclusion $i$: we have to put $R$ into $F$. In other words, if I have an element $r \in R$ --- normally people don't even write it, but I just... we have to put it into $F$. That means we have to assign it an equivalence class in this set. So, what are we going to do? 

Yeah?
\end{speech}

\begin{student}[Student Answer]
$(r, 1)$?
\end{student}

\begin{speech}[00:52:39 - 00:53:41]
Yeah, that's a good idea: $(r, 1)$.

It has some advantages: $1 \ne 0$, because an integral domain is never the zero ring (Recall: \cref{def:integral-domain}). So, this is legal, this is an element of $S$, and then we can take its equivalence class.

I'm not going to write equivalence classes like $[(r, 1)]$ everywhere because it becomes too complicated. You could say formally take its equivalence class like this, but if you start doing this, you get committed to writing things like $[(3, 4)]$ for rational numbers, and no one ever does that. So, we don't want to do it, but this would be kind of more logically correct.

When I write this, it's just the equivalence class. I looked at what the book does, and the book says that $(r, 1)$ is an element of $S$, and that is its equivalence class. That's a very clever solution. I wouldn't say it's totally standard, but anyway; standard is just to ignore the issue.

So, this is the equivalence class here. Is this reasonable? Well, it better at least be... yeah...
\end{speech}

\begin{student}[Student Question]
So, the set $S$, is it the whole $R \times R$, or is it just any subset of $R \times R$?
\end{student}

\begin{speech}[00:53:50 - 00:55:47]
No, the set $S$ was very precisely defined on the previous board: $S$ is equal to those elements in $R \times R$ where $y \ne 0$. That's $S$. \inlinemetanote{writes $S = \{ (x, y) \mid x, y \in R, \ y \neq 0 \}$ on the board}

And this is an element of $S$ because $1 \ne 0$; and $1 \ne 0$ because $R$ is an integral domain, and an integral domain is by definition not the zero ring. That's the full logical certificate, if you want it.

Still, we'd like more: we'd like this to be an injection. \inlinemetanote{writes \qt{Is $i$ injective?} on the board}

This means: when are $(r, 1)$ and $(\hat{r}, 1)$ in the same equivalence class of $S$, right? \inlinemetanote{writes When are $(r, 1)$ and $(\hat{r}, 1)$ in the same equivalence class of $S$} That's what injectivity is about.

How do we check this? The answer is that these are in the same equivalence class if and only if that equivalence class identity holds: if and only if $r \cdot 1 = \hat{r} \cdot 1$, which is if and only if $r = \hat{r}$. This is great for us, because it shows that this is an injection. \inlinemetanote{writes $r \cdot 1 = \hat{r} \cdot 1 \iff r = \hat{r}$ and puts a checkmark}

We've constructed... there's a lot of things to do in this, like a big house where you have to build all the pieces of it. This shows that now we have not only defined $F$, we've defined $i$, and we've put $R$ inside $F$. And this placement is injective, okay?
\end{speech}

\begin{hidden}
% timestamp: 00:55:40
% topic: Definition of the field of fractions F = S/~ and canonical inclusion map i: R -> F
% board_state: def:field-of-fractions, def:canonical-inclusion, prop:inclusion-injective
% next_goal: Define the field operations (zero, unit, addition, multiplication) on F
% open_loops: none
\end{hidden}

\begin{content}[Definition of the Field of Fractions and Canonical Embedding]
\begin{definition}[Field of Fractions]\label[definition]{def:field-of-fractions}
Let $R$ be an integral domain, and let $S := \{ (x, y) \in R \times R \mid y \ne 0 \}$. The \newterm{field of fractions} (or quotient field) of $R$, denoted by $\operatorname{Frac}(R)$ or $F$, is defined as the quotient set:
\begin{equation}\label{eq:field-of-fractions-def}
F := S/\!\equiv \, = \{ [(x, y)] \mid (x, y) \in S \}.
\end{equation}
The equivalence class $[(x, y)]$ is commonly written as $\frac{x}{y}$.
\end{definition}

\begin{definition}[Canonical Inclusion Map]\label[definition]{def:canonical-inclusion}
We define the canonical embedding $i \colon R \to F$ by assigning to each $r \in R$ the equivalence class of the pair $(r, 1)$:
\begin{equation}\label{eq:canonical-inclusion-def}
i(r) := [(r, 1)] \in F.
\end{equation}
\end{definition}

\begin{proposition}[Injectivity of the Canonical Inclusion]\label[proposition]{prop:inclusion-injective}
The canonical map $i \colon R \to F$ is injective.
\end{proposition}

\begin{proof}
Let $r, \hat{r} \in R$. By definition of equality in the quotient set $F = S/\!\equiv$:
\begin{align*}
i(r) = i(\hat{r}) &\iff [(r, 1)] = [(\hat{r}, 1)] \\
&\iff (r, 1) \equiv (\hat{r}, 1) \\
&\iff r \cdot 1 = \hat{r} \cdot 1 \\
&\iff r = \hat{r}.
\end{align*}
Therefore, $i$ is injective.
\end{proof}

\begin{steps}[Why $1_R \ne 0_R$ Is Guaranteed]
For the pair $(r, 1)$ to belong to $S$, the denominator must satisfy $1 \ne 0$. By condition (i) of \cref{def:integral-domain}, an integral domain must satisfy $1_R \ne 0_R$ (i.e., it cannot be the trivial zero ring). Hence, $(r, 1) \in S$ is unconditionally valid for every $r \in R$.
\end{steps}
\end{content}

\subsection{Field Operations on the Field of Fractions: Zero, Unit, and Addition}

\begin{speech}[00:55:50 - 00:57:04]
As I said, this is just the beginning of the story. $R$ already is born with all of the data: $(R, +, \cdot, 0, 1)$. We need to define all the structures on $F$ and show this is a homomorphism also.
\inlinemetanote{The lecturer pushes the upper chalkboards up and erases the middle boards}
\end{speech}

\begin{speech}[00:57:04 - 00:58:39]
We have to start talking about all the structures for $F$ as a field. And for many of them, we don't have any choice: the zero in $F$ has to correspond to the zero in $R$ for it to be a homomorphism. So, the zero in $F$ has to be $(0, 1)$. \inlinemetanote{writes Zero in $F$ has to be $(0, 1)$}

And also the unit: the unit in $F$ has to be $(1, 1)$, because of course it has to correspond to the unit in $R$ by the rule of getting $R$ to $F$. \inlinemetanote{writes unit in $F$ has to be $(1, 1)$}

We don't really have any choice in this matter. But we have to define plus and multiplication on $F$. \inlinemetanote{writes We have to define $+, \cdot$ on $F$} That's something that we have to construct: we have to define it, check it's compatible with $R$, check it satisfies the field axioms --- and it's a lot of long and boring work.

The interesting part to start with is that we have to define it. If we have an element here, $(x_1, y_1) \in S$, and we have another one, $(x_2, y_2) \in S$, then we have to define this operation. This is up to us; we must define it. \inlinemetanote{writes $(x_1, y_1) + (x_2, y_2)$ with \qt{We must define} under the $+$}

That's the problem. So, what's the solution? How are we going to define this? What do you think? 

Tell me.
\end{speech}

\begin{student}[Student Answer]
Take $x_1 y_2 + x_2 y_1$, and $y_1 y_2$.
\end{student}

\begin{speech}[00:58:43 - 01:00:00]
Yeah, so we just use the standard definition for how you add fractions. \inlinemetanote{writes $(x_1 y_2 + x_2 y_1, y_1 y_2)$ on the board} Meaning you pretend that it's this problem, which it kind of is: \inlinemetanote{writes $\frac{x_1}{y_1} + \frac{x_2}{y_2} = \frac{x_1 y_2 + x_2 y_1}{y_1 y_2}$ on the right} You know how to do this.

I hope I got the subscripts right; you should tell me if I get the subscripts wrong.

This is the definition, nothing more, nothing less.

Now, is it a legal definition? There's an easy part: for it to be legal, the denominator better not be zero, because in our set $S$, the one in the basement always has to be non-zero.

We have here that this guy's non-zero, that guy's non-zero, and that implies their product is non-zero, because it is an integral domain. \inlinemetanote{draws an arrow to $y_1 y_2$ and writes \qt{non zero, $R$ is an int dom}} Is that clear? Otherwise, if this were zero, this would not be legal.
\end{speech}

\begin{speech}[01:00:00 - 01:01:45]
That's the easy part. The hard part is you have to check that... this is a definition on $S$, but we have to check that it's well-defined on equivalence classes.

That means if I trade this person for another one in the same equivalence class, and I trade this for another person in the same equivalence class, the result has to be in the same equivalence class. Is that clear?

That's the check of rational arithmetic. That's the check that I claim, to understand the rational numbers, you should have done in school. Otherwise, what are you doing? You're doing some random manipulation.

So, you have to check this. \inlinemetanote{writes \qt{You have to check the indep. of choices of equiv. class} on the board}

We've already done this kind of thing twice in the lecture. If you just do it and follow your nose, you'll see it all unwind: you just use the definition, and you have to say formally that this result respects our definition of the equivalence class.

In various books they do it; mostly it's left as an exercise. It's really important, so you should do it. If I do it, it's just a bunch of indices --- it doesn't really help. You have to go through it yourself.

Is it clear what the problem is? I recommend trying this yourself, and if you get stuck, look at the book. The book does it. It'll give you a line of things, and you can imagine it's probably true.
\end{speech}

\begin{hidden}
% timestamp: 01:01:40
% topic: Definition of zero, unit, and addition on F and well-definedness
% board_state: zero in F is (0,1), unit in F is (1,1), addition formula on F, checking non-zero denominator via integral domain
% next_goal: Define multiplication on F and prove compatibility with equivalence classes
% open_loops: none
\end{hidden}

\begin{content}[Neutral Elements and Addition on the Field of Fractions]
\begin{definition}[Neutral Elements in $F$]\label[definition]{def:neutral-elements-F}
The zero element and multiplicative unit in $F$ are defined by:
\begin{align}
0_F &:= i(0_R) = [(0, 1)], \label{eq:zero-F} \\
1_F &:= i(1_R) = [(1, 1)]. \label{eq:unit-F}
\end{align}
\end{definition}

\begin{definition}[Addition on $F$]\label[definition]{def:addition-F}
For any two equivalence classes $[(x_1, y_1)]$, $[(x_2, y_2)] \in F$, addition is defined by:
\begin{equation}\label{eq:addition-F-formula}
[(x_1, y_1)] + [(x_2, y_2)] := [(x_1 y_2 + x_2 y_1, y_1 y_2)].
\end{equation}
\end{definition}

\begin{proposition}[Legitimacy and Well-Definedness of Addition]\label[proposition]{prop:addition-well-defined}
Addition on $F$ as given in \eqref{eq:addition-F-formula} is well-defined.
\end{proposition}

\begin{proof}
We verify two essential requirements:
\begin{enumerate}[label=\textbf{(\arabic*)}]
    \item \textbf{Legitimacy of the Denominator:} For $(x_1 y_2 + x_2 y_1, y_1 y_2)$ to lie in $S$, we must have $y_1 y_2 \ne 0$. Since $(x_1, y_1), (x_2, y_2) \in S$, we know $y_1 \ne 0$ and $y_2 \ne 0$. Because $R$ is an \textbf{integral domain}, the product of two non-zero elements is non-zero:
    \[
    y_1 \ne 0 \land y_2 \ne 0 \implies y_1 y_2 \ne 0.
    \]
    Thus, $(x_1 y_2 + x_2 y_1, y_1 y_2) \in S$.
    
    \item \textbf{Independence of Representatives:} Suppose $(x_1, y_1) \equiv (\hat{x}_1, \hat{y}_1)$ and $(x_2, y_2) \equiv (\hat{x}_2, \hat{y}_2)$, meaning:
    \begin{equation}\label{eq:add-well-def-hyp}
    x_1 \hat{y}_1 = \hat{x}_1 y_1 \quad \text{and} \quad x_2 \hat{y}_2 = \hat{x}_2 y_2.
    \end{equation}
    We must show that $(x_1 y_2 + x_2 y_1, y_1 y_2) \equiv (\hat{x}_1 \hat{y}_2 + \hat{x}_2 \hat{y}_1, \hat{y}_1 \hat{y}_2)$, which requires:
    \[
    (x_1 y_2 + x_2 y_1)(\hat{y}_1 \hat{y}_2) \overset{?}{=} (\hat{x}_1 \hat{y}_2 + \hat{x}_2 \hat{y}_1)(y_1 y_2).
    \]
    Expanding the left-hand side and applying \eqref{eq:add-well-def-hyp}:
    \begin{align*}
    (x_1 y_2 + x_2 y_1)\hat{y}_1 \hat{y}_2 &= (x_1 \hat{y}_1)(y_2 \hat{y}_2) + (x_2 \hat{y}_2)(y_1 \hat{y}_1) \\
    &= (\hat{x}_1 y_1)(y_2 \hat{y}_2) + (\hat{x}_2 y_2)(y_1 \hat{y}_1) \\
    &= (\hat{x}_1 \hat{y}_2 + \hat{x}_2 \hat{y}_1)y_1 y_2.
    \end{align*}
    This establishes equivalence, confirming that addition is independent of the choice of representative.
\end{enumerate}
\end{proof}
\end{content}

\subsection{Multiplication and Ring Embedding Properties}

\begin{speech}[01:01:45 - 01:03:17]
\inlinemetanote{The lecturer erases the left chalkboard}

You also have to do it with multiplication. Multiplication is a bit easier; addition is the most painful one. Also define multiplication:
\[
(x_1, y_1) \cdot (x_2, y_2) = (x_1 x_2, y_1 y_2),
\]
which will surprise no one, okay? \inlinemetanote{writes Also define $\cdot$: $(x_1, y_1) \cdot (x_2, y_2) \overset{\text{def}}{=} (x_1 x_2, y_1 y_2)$} That's the definition, right?

And again, next is the check of compatibility with the equivalence classes. This check is easy, so I'll do it since it's easy.
\end{speech}

\begin{speech}[01:03:17 - 01:05:40]
Suppose we have that $(x_1, y_1)$ is in the same equivalence class as $(\hat{x}_1, \hat{y}_1)$, and $(x_2, y_2)$ is in the same equivalence class as $(\hat{x}_2, \hat{y}_2)$. \inlinemetanote{writes $(x_1, y_1) \equiv (\hat{x}_1, \hat{y}_1)$, $(x_2, y_2) \equiv (\hat{x}_2, \hat{y}_2)$}

Then we have to show --- that's the input, right? --- we must show that $(x_1 x_2, y_1 y_2)$ is in the same equivalence class as if we chose the other representatives: $(\hat{x}_1 \hat{x}_2, \hat{y}_1 \hat{y}_2)$. \inlinemetanote{writes We must show $(x_1 x_2, y_1 y_2) \equiv (\hat{x}_1 \hat{x}_2, \hat{y}_1 \hat{y}_2)$} That's what we have to do.

What we know is this: the equation that these are in the same equivalence class is $x_1 \hat{y}_1 = \hat{x}_1 y_1$, and this equation is $x_2 \hat{y}_2 = \hat{x}_2 y_2$, all very confusing. \inlinemetanote{writes $x_1 \hat{y}_1 = \hat{x}_1 y_1$, $x_2 \hat{y}_2 = \hat{x}_2 y_2$}

And now we just multiply them together! These are what we know, and then we have these two equations: we can multiply them all together, and we get $x_1 x_2 \hat{y}_1 \hat{y}_2 = \hat{x}_1 \hat{x}_2 y_1 y_2$. \inlinemetanote{writes $\implies x_1 x_2 \hat{y}_1 \hat{y}_2 = \hat{x}_1 \hat{x}_2 y_1 y_2$}

That equation is exactly the certificate that these things are in the same equivalence class, okay? This implies that.

This shows that if I pick different equivalence classes here, the result doesn't matter. This one I did completely, because it's easy: all you do is multiply. For the other one, you have to do the kind of tricks we were doing before, where you have to clear the denominators, etc., so you can do that, so to speak, on your own time.
\end{speech}

\begin{speech}[01:05:40 - 01:06:50]
Then you have to check that these definitions are compatible with $R$. I mean, they're just all defined that way, meaning we have to check that this inclusion $i$ is a homomorphism of rings. \inlinemetanote{writes Check $i \colon R \to F$ is a hom of rings}

And again, there's really nothing to this: it means that if you take $r + \hat{r}$, and then you go by $i$, you get $(r + \hat{r}, 1)$. \inlinemetanote{writes $r + \hat{r} \xrightarrow{i} (r + \hat{r}, 1)$ and below it $= (r, 1) + (\hat{r}, 1)$}

Then you have to check that this thing is equal to $(r, 1) + (\hat{r}, 1)$, and that's the same by the definition that we used there. But you have to check it, and you have to also check it for multiplication. So, it's just a bunch of checking. Here, all the cards are played, and you just have to check everything, okay?
\end{speech}

\begin{hidden}
% timestamp: 01:06:40
% topic: Multiplication on F, verification of well-definedness, and homomorphism properties of i: R -> F
% board_state: def:mult-F, prop:mult-well-defined, thm:field-of-fractions-properties
% next_goal: Summary of the complete construction of the field of fractions
% open_loops: none
\end{hidden}

\begin{content}[Multiplication on $F$ and Ring Homomorphism Properties]
\begin{definition}[Multiplication on $F$]\label[definition]{def:multiplication-F}
For any $[(x_1, y_1)], [(x_2, y_2)] \in F$, multiplication is defined by:
\begin{equation}\label{eq:multiplication-F-formula}
[(x_1, y_1)] \cdot [(x_2, y_2)] := [(x_1 x_2, y_1 y_2)].
\end{equation}
\end{definition}

\begin{proposition}[Well-Definedness of Multiplication]\label[proposition]{prop:mult-well-defined}
Multiplication on $F$ is well-defined and independent of the choice of representative.
\end{proposition}

\begin{proof}
The pair $(x_1 x_2, y_1 y_2)$ lies in $S$, since $y_1 y_2 \neq 0$ in the integral domain $R$ (as for addition, \cref{prop:addition-well-defined}).

Let $(x_1, y_1) \equiv (\hat{x}_1, \hat{y}_1)$ and $(x_2, y_2) \equiv (\hat{x}_2, \hat{y}_2)$ in $S$. By definition of $\equiv$:
\begin{equation}\label{eq:mult-hyp}
x_1 \hat{y}_1 = \hat{x}_1 y_1 \quad \text{and} \quad x_2 \hat{y}_2 = \hat{x}_2 y_2.
\end{equation}
Multiplying these two equations in the commutative ring $R$:
\[
(x_1 \hat{y}_1)(x_2 \hat{y}_2) = (\hat{x}_1 y_1)(\hat{x}_2 y_2) \implies (x_1 x_2)(\hat{y}_1 \hat{y}_2) = (\hat{x}_1 \hat{x}_2)(y_1 y_2).
\]
By \eqref{eq:equiv-def-recap}, this algebraic identity is precisely the condition for equivalence:
\[
(x_1 x_2, y_1 y_2) \equiv (\hat{x}_1 \hat{x}_2, \hat{y}_1 \hat{y}_2).
\]
Thus, $[(x_1 x_2, y_1 y_2)] = [(\hat{x}_1 \hat{x}_2, \hat{y}_1 \hat{y}_2)]$, proving well-definedness.
\end{proof}

\begin{proposition}[$i$ is an Injective Ring Homomorphism]\label[proposition]{prop:i-is-homomorphism}
The canonical inclusion $i \colon R \hookrightarrow F$ is an injective ring homomorphism:
\begin{align}
i(r + \hat{r}) &= i(r) + i(\hat{r}), \label{eq:i-add} \\
i(r \cdot \hat{r}) &= i(r) \cdot i(\hat{r}), \label{eq:i-mult} \\
i(1_R) &= 1_F. \label{eq:i-unit}
\end{align}
\end{proposition}

\begin{proof}
Evaluating the definitions directly:
\begin{itemize}
    \item \textbf{Additivity:}
    \[
    i(r) + i(\hat{r}) = [(r, 1)] + [(\hat{r}, 1)] = [(r \cdot 1 + \hat{r} \cdot 1, 1 \cdot 1)] = [(r + \hat{r}, 1)] = i(r + \hat{r}).
    \]
    \item \textbf{Multiplicativity:}
    \[
    i(r) \cdot i(\hat{r}) = [(r, 1)] \cdot [(\hat{r}, 1)] = [(r \hat{r}, 1 \cdot 1)] = [(r \hat{r}, 1)] = i(r \hat{r}).
    \]
    \item \textbf{Identity:}
    \[
    i(1_R) = [(1, 1)] = 1_F.
    \]
\end{itemize}
Injectivity was already established in \cref{prop:inclusion-injective}.
\end{proof}

\begin{steps}[Multiplicative Inverses in $F$]
To confirm that $F$ is indeed a field, let $[(x, y)] \in F$ with $[(x, y)] \ne 0_F = [(0, 1)]$. Then $(x, y) \not\equiv (0, 1) \implies x \cdot 1 \ne 0 \cdot y \implies x \ne 0$. Thus, $(y, x) \in S$, and:
\[
[(x, y)] \cdot [(y, x)] = [(xy, yx)] = [(1, 1)] = 1_F.
\]
Therefore, every non-zero element in $F$ possesses a multiplicative inverse $[(x, y)]^{-1} = [(y, x)]$.
\end{steps}
\end{content}

\subsection{Summary of the Explicit Construction}

\begin{speech}[01:06:50 - 01:08:46]
That is the construction. 

As I said, when you understand it, you'll realize that actually you knew the whole thing anyway in your understanding of the rational numbers. You were subconsciously doing all of this --- not explicitly, but underneath the surface, you were assuming all of these properties were true; it made rational number arithmetic work. 

There are different ways to try to learn this material, but the way that's foundational is that you actually go yourself and check all the things that you're supposed to check. That takes a little bit of time, and it's somewhat painful, but you only have to do it once in your life. 

You can read the book, too. I think that doesn't help so much, because then you just read what the author has written, and it's just a bunch of things with equalities. You can convince yourself of anything if you read what someone else did. So, the best thing is to go through all of it yourself; then you'll understand it once in your life. 

And that's the construction, so we're done with this construction. I'll just summarize here: \inlinemetanote{writes on the lower left chalkboard: \qt{Construction}, \qt{Start with $R$ int domain}, \qt{Construct $F$ = field of fractions}, $R \hookrightarrow F$, $r \to (r, 1)$}

We start with $R$ as an integral domain, and then we construct explicitly --- that's as explicit a mathematical construction as you get. You just take $R$, and you make some sets and equivalence classes and relations; it's totally explicit. We construct this $F$, which is the field of fractions. 

And $R$ sits inside $F$ by this map: $r \mapsto (r, 1)$.
\end{speech}

\begin{content}[Summary of the Construction of the Field of Fractions]
\begin{theorem}[Existence and Structure of the Field of Fractions]\label[theorem]{thm:field-of-fractions-summary}
Let $R$ be an integral domain. There exists a field $F$, called the \newterm{field of fractions} of $R$ (denoted $\operatorname{Frac}(R)$), and an injective ring homomorphism:
\[
i \colon R \hookrightarrow F, \quad r \mapsto [(r, 1)],
\]
such that:
\begin{enumerate}[label=\textbf{(\arabic*)}]
    \item \textbf{Field Structure:} $F$ is a field with zero $0_F = [(0, 1)]$ and multiplicative identity $1_F = [(1, 1)]$.
    \item \textbf{Quotient Representation:} Every element of $F$ can be written as a quotient of elements in the image of $R$:
    \begin{equation}\label{eq:frac-quotient-representation}
    \forall q \in F \ \exists r, s \in R \ (s \ne 0) \quad \text{such that} \quad q = i(r) \cdot i(s)^{-1} = \frac{i(r)}{i(s)}.
    \end{equation}
\end{enumerate}
\end{theorem}
\end{content}

% --- TEIL 4 ---

\subsection{Algebraic Perspective: Fractions in an Integral Domain}

\begin{speech}[01:08:46 - 01:09:44]
When you get used to it, the way an algebraist always thinks about this is: you have elements in $R$, and the elements of $F$ are just $r/s$; you just view them as fractions in $R$. 

This is a perfectly reasonable way to think about it. In our language, when someone writes this, what they really mean is the equivalence class of $(r, s)$ in this set $S$. That's what it actually means. And it's only legal if the bottom one is not zero, which means you're not allowed to write this anyway if the bottom one's not zero. Then you can add fractions and multiply fractions just as you would normally do, and they correspond exactly to the rules that were prescribed here. That's the natural way to think about it. 

Once all of this is done, then you just start using it like you normally would anyway. But this is the underlying foundation. 

Are there any conceptual questions about that?
\end{speech}

\begin{content}[Algebraic View: Working with Fractions]
In mathematical practice, once the formal construction of $F = S/\!\equiv$ is established, we identify $R$ with its image $i(R) \subseteq F$ and adopt standard fraction notation:
\begin{equation}\label{eq:fraction-notation-standard}
\frac{r}{s} := [(r, s)] \in \operatorname{Frac}(R) \quad (r, s \in R, \ s \ne 0).
\end{equation}

\textbf{Arithmetic in the Field of Fractions:}
Under this notation, the field operations on $\operatorname{Frac}(R)$ take the standard form familiar from elementary arithmetic:
\begin{align}
\text{Equality:} && \frac{r_1}{s_1} &= \frac{r_2}{s_2} \iff r_1 s_2 = r_2 s_1, \label{eq:frac-eq} \\
\text{Addition:} && \frac{r_1}{s_1} + \frac{r_2}{s_2} &= \frac{r_1 s_2 + r_2 s_1}{s_1 s_2}, \label{eq:frac-add} \\
\text{Multiplication:} && \frac{r_1}{s_1} \cdot \frac{r_2}{s_2} &= \frac{r_1 r_2}{s_1 s_2}, \label{eq:frac-mult} \\
\text{Inverse:} && \left(\frac{r}{s}\right)^{-1} &= \frac{s}{r} \quad (\text{provided } r \ne 0). \label{eq:frac-inv}
\end{align}
\end{content}

\section{Examples of Fields of Fractions}

\subsection{The Rationals, Gaussian Rationals, and Subrings of \texorpdfstring{$\mathbb{Q}$}{Q}}

\begin{speech}[01:09:44 - 01:10:44]
If I take $\mathbb{Z}$, you can ask: what is the field of fractions? 

The answer is going to be $\mathbb{Q}$, which means that if you start with $\mathbb{Z}$ and you do this crazy construction, you get the rational numbers. I hope that's clear by now. 

All right, so what happens if I do something else? What happens if I start with the ring of Gaussian integers? What is the field of fractions? 

That's a strange question, because I've just explained what the field of fractions is: you take the Gaussian integers, you do this construction, and that's the field of fractions. So, when I say \qt{what is the field of fractions}, what the question is actually asking is: what's another way of thinking about the answer? 

Do you have a proposal for what that is? 

Yeah?
\end{speech}

\begin{student}[Student Answer]
Just $\mathbb{Q}$ adjoined with $i$.
\end{student}

\begin{hidden}
% timestamp: 01:10:45
% topic: Field of fractions examples: Z, Z[i], Z[1/3]
% board_state: Frac(Z) = Q, Frac(Z[i]) = Q(i), Frac(Z[1/3]) = Q
% next_goal: Discuss subrings of Q and field of fractions of a field
% open_loops: none
\end{hidden}

\begin{speech}[01:10:46 - 01:12:32]
Yeah, so the answer is this: $\mathbb{Q}(i)$, which is equal to the set of $a + bi$, where $a, b \in \mathbb{Q}$. 

It's correct, but now to check this is correct, you have to do a bunch of things: one is you have to show this is a field (you can invert), and then you have to show that this sits inside, and everything here is a ratio of that. It's all true, but that's the answer here, and it requires some checking. All right, this is correct. 

What about this? We saw last week that there are some funny rings that sit inside $\mathbb{Z}$ and $\mathbb{Q}$. This ring is the set of $a/3^b$, where $a \in \mathbb{Z}$ and $b$ is some integer greater than or equal to $0$, right? We discussed this last week (Recall: \cref{ex:dyadic-rationals-subring}): the subrings that sit in between. I hope you remember. 

Do you remember? 

This is an integral domain also, because it sits inside $\mathbb{Q}$. So, what is the field of fractions of... I just gave it a name, $\mathbb{Z}[1/3]$, because you're allowed to have threes in the bottom. There are two standard names for this: one is $\mathbb{Z}[1/3]$, and the other is... yeah, that's a better name, okay. 

What is the field of fractions of this integral domain?
\end{speech}

\begin{student}[Student Answer]
Also $\mathbb{Q}$.
\end{student}

\begin{speech}[01:12:34 - 01:13:33]
Also $\mathbb{Q}$! Nothing says it can't be $\mathbb{Q}$. It certainly sits inside $\mathbb{Q}$, and everything in $\mathbb{Q}$ is already a ratio of numbers here. So, it's just $\mathbb{Q}$. If you go through this definition, you get $\mathbb{Q}$. 

As I said, you think about this field of fractions as the smallest field that contains the integral domain. You can't really contain something smaller than $\mathbb{Q}$, because it already has $\mathbb{Z}$ in it, and $\mathbb{Z}$ already gives you everything in $\mathbb{Q}$. 

So, what is the field of fractions of this? The answer is also $\mathbb{Q}$. 

In fact, if you take any ring, any integral domain that sits in between $\mathbb{Z}$ and $\mathbb{Q}$, the field of fractions is just $\mathbb{Q}$.
\end{speech}

\begin{content}[Examples: Fractions of the Integers, Gaussian Integers, and Subrings of \texorpdfstring{$\mathbb{Q}$}{Q}]
\begin{example}[The Integers]\label[example]{ex:frac-integers}
For the ring of integers $\mathbb{Z}$:
\[
\operatorname{Frac}(\mathbb{Z}) = \mathbb{Q}.
\]
\end{example}

\begin{example}[The Gaussian Integers]\label[example]{ex:frac-gaussian-integers}
For the ring of Gaussian integers $\mathbb{Z}[i] = \{ a + bi \mid a, b \in \mathbb{Z} \}$:
\[
\operatorname{Frac}(\mathbb{Z}[i]) = \mathbb{Q}(i) := \{ a + bi \mid a, b \in \mathbb{Q} \}.
\]
The field $\mathbb{Q}(i)$ is called the field of \newterm{Gaussian rationals}.
\end{example}

\begin{example}[Triadic Rationals]\label[example]{ex:frac-triadic}
Consider the ring of triadic fractions $\mathbb{Z}[\frac{1}{3}]$ (the case $m = 3$ of \cref{ex:dyadic-rationals-subring}), defined by:
\[
\mathbb{Z}\left[\frac{1}{3}\right] := \left\{ \frac{a}{3^b} \;\middle|\; a \in \mathbb{Z}, \ b \in \mathbb{Z}_{\ge 0} \right\} \subset \mathbb{Q}.
\]
Since $\mathbb{Z} \subset \mathbb{Z}[\frac{1}{3}] \subset \mathbb{Q}$, any fraction of elements from $\mathbb{Z}[\frac{1}{3}]$ is already a rational number, and every rational number $\frac{p}{q}$ can be written as a quotient of elements in $\mathbb{Z} \subset \mathbb{Z}[\frac{1}{3}]$. Thus:
\[
\operatorname{Frac}\left(\mathbb{Z}\left[\frac{1}{3}\right]\right) = \mathbb{Q}.
\]
\end{example}

\begin{proposition}[Intermediate Subrings of $\mathbb{Q}$]\label[proposition]{prop:subrings-q-frac}
Let $R$ be any subring of $\mathbb{Q}$ with $\mathbb{Z} \subseteq R \subseteq \mathbb{Q}$. Then:
\begin{equation}\label{eq:intermediate-frac-q}
\operatorname{Frac}(R) = \mathbb{Q}.
\end{equation}
\end{proposition}

\begin{center}
\begin{tikzpicture}[scale=1.3, >=Stealth]
% \begin{hidden}
% - Canvas: scale=1.3, vertical tower of inclusions: Z at (0,0), Z[1/3] at (0,1.3), Q at (0,2.6)
% - Right side: Frac map arrows from Z, Z[1/3], and Q all pointing to Q at (3.5, 1.3)
% - Colors: MidnightBlue for rings, BrickRed for fields, ForestGreen for arrows
% \end{hidden}
    % Ring Tower on Left
    \node[MidnightBlue, font=\bfseries] (Z) at (0, 0) {$\mathbb{Z}$};
    \node[MidnightBlue, font=\bfseries] (Z3) at (0, 1.3) {$\mathbb{Z}[\frac{1}{3}]$};
    \node[BrickRed, font=\bfseries] (Qleft) at (0, 2.6) {$\mathbb{Q}$};

    \draw[right hook->, thick, MidnightBlue] (Z) -- (Z3);
    \draw[right hook->, thick, MidnightBlue] (Z3) -- (Qleft);

    % Common Field of Fractions on Right
    \node[draw=BrickRed, fill=BrickRed!10, rounded corners, inner sep=6pt, font=\bfseries] (Qtarget) at (3.5, 1.3) {$\operatorname{Frac} = \mathbb{Q}$};

    \draw[->, thick, ForestGreen] (Z.east) to[bend right=20] node[midway, below, sloped, font=\footnotesize] {$\operatorname{Frac}$} (Qtarget.south west);
    \draw[->, thick, ForestGreen] (Z3.east) -- node[midway, above, font=\footnotesize] {$\operatorname{Frac}$} (Qtarget.west);
    \draw[->, thick, ForestGreen] (Qleft.east) to[bend left=20] node[midway, above, sloped, font=\footnotesize] {$\operatorname{Frac}$} (Qtarget.north west);
\end{tikzpicture}
\end{center}

\begin{steps}[Why All Intermediate Subrings Yield $\mathbb{Q}$]
By categorical minimality, $\operatorname{Frac}(R)$ is the smallest field containing $R$. Since $\mathbb{Z} \subseteq R$, any field containing $R$ must contain $\mathbb{Z}$, and hence must contain a copy of $\operatorname{Frac}(\mathbb{Z}) = \mathbb{Q}$ (\cref{thm:universal-property-q}). Since $R \subseteq \mathbb{Q}$, $\mathbb{Q}$ is already a field containing $R$. Therefore, $\operatorname{Frac}(R)$ cannot be strictly larger or smaller than $\mathbb{Q}$, forcing $\operatorname{Frac}(R) = \mathbb{Q}$.
\end{steps}
\end{content}

\subsection{Field of Fractions of an Existing Field}

\begin{speech}[01:13:33 - 01:14:49]
Okay, so now a stranger question is: what happens if... well, yeah. 

You know, on the business of guessing: once I was at a conference, and the conference had some math researchers and some high school teachers. At dinner I was sitting next to some teachers, and I was trying to explain to them my theory of how you should teach children math. And the outcome was that they wanted me disbarred from academics. 

The theory, by the way, was that on exams you should make the kids just guess the answers: no calculation, just guess everything. That's, you know, maybe a little excessive, but you know... 

Maybe a compromise: on half the exam they should just guess everything. But I think that would be really, somehow, more in the spirit of mathematics than...
\end{speech}

\begin{speech}[01:14:49 - 01:15:25]
A kind of stranger question is: suppose I start with a field. Suppose $R$ is a field. \inlinemetanote{writes Suppose $R$ is field}

Then, of course, it's still an integral domain, because every field is an integral domain (Recall: \cref{ex:fields-are-domains}). Then what is the field of fractions? \inlinemetanote{writes $R$ is int domain, $\implies$ What is the field of fraction?}

What do you think the field of fractions is going to be here? 

Yeah?
\end{speech}

\begin{student}[Student Answer]
It's $R$ itself.
\end{student}

\begin{speech}[01:15:28 - 01:16:00]
Yes, $R$ itself! \inlinemetanote{writes \qt{$R$ itself} on the board}

And the answer is $R$ itself. 

To prove that: again, this fits in with this view that it's the smallest field. If it's already a field, how could there be a smaller field that contains it? It's already a field. 

To prove it, you have to go through this construction and show that in the case where $R$ is a field, that construction just reconstructs $R$. To do that, you have to go through that construction and see it, so that's also an exercise for you. 

Those are really the basic facts.
\end{speech}

\begin{hidden}
% timestamp: 01:15:55
% topic: Field of fractions of a field is the field itself
% board_state: R is a field -> Frac(R) = R itself
% next_goal: Begin discussion of polynomial rings R[x]
% open_loops: none
\end{hidden}

\begin{content}
\begin{proposition}[Field of Fractions of a Field]\label[proposition]{prop:frac-of-a-field}
Let $K$ be a field. Then the canonical inclusion $i \colon K \hookrightarrow \operatorname{Frac}(K)$ is an isomorphism of fields:
\begin{equation}\label{eq:frac-field-itself}
\operatorname{Frac}(K) \cong K.
\end{equation}
\end{proposition}

\begin{proof}
By \cref{prop:i-is-homomorphism}, $i \colon K \hookrightarrow \operatorname{Frac}(K)$ is an injective field homomorphism. To establish surjectivity, let $[(x, y)] \in \operatorname{Frac}(K)$ with $x, y \in K$ and $y \ne 0$.
Since $K$ is a field, $y$ has a multiplicative inverse $y^{-1} \in K$. Then:
\[
(x, y) \equiv (x y^{-1}, 1) \iff x \cdot 1 = (x y^{-1}) \cdot y = x.
\]
Consequently:
\[
[(x, y)] = [(x y^{-1}, 1)] = i(x y^{-1}).
\]
Every element of $\operatorname{Frac}(K)$ is the image of the element $x y^{-1} \in K$. Hence, $i$ is surjective and constitutes a field isomorphism.
\end{proof}
\end{content}

\section{Polynomial Rings and Power Series Rings}

\subsection{Definition of the Polynomial Ring in One Indeterminate}

\begin{speech}[01:16:00 - 01:18:16]
The next topic this week is about polynomial rings. These are really fun. We'll talk about polynomial rings \inlinemetanote{writes \qt{Polynomials rings, Power series rings}}, and in some sense even more fun than polynomial rings are power series rings. 

Again, this material, in some sense, you've been using all of this already in the first year, but here we take a formal approach. 

The idea is: if I take $R$ as a ring --- and now, no assumptions, meaning not necessarily an integral domain, so this just means commutative for us: a commutative ring with unit. That's our basic ring: when we say ring with no assumptions, we mean a commutative ring with unit (Recall: \cref{rem:ring-convention}). In particular, for everything today so far, $R$ was an integral domain. Now, I don't care if it's an integral domain; it's just a commutative ring with unit. 

Then I can make a new ring here: $R[x]$. 

Here, $x$ is the variable, or sometimes it's called the indeterminate. That's kind of a fancy name, but it's just $x$, right? The indeterminate or the variable. 

The name of this ring is the polynomial ring with coefficients in $R$. Or if you want to be even more precise: the polynomial ring in one variable with coefficients in $R$. That's kind of a long sentence, but if you meet a mathematician and say, \qt{I'm going to consider the polynomial ring in one variable with coefficients in $R$}, as long as you agree on what a ring is, that's absolutely unambiguous --- the person will know exactly what you're talking about. 

What is this? Well, I have to tell you what it is as a set. So, the first thing is: what's the set?
\end{speech}

\begin{speech}[01:18:16 - 01:20:17]
Here the definition is that it's the set of polynomials in $x$ with coefficients in $R$. Polynomials look like this: $a_0 + a_1 x + a_2 x^2 + \dots$ \inlinemetanote{writes $R[x] = \{ a_0 + a_1 x + a_2 x^2 + \dots + a_n x^n \mid a_i \in R, n \ge 0 \}$ on the board}

It's these symbols, up to eventually... One of the properties of a polynomial is that it's not infinite; a polynomial is finite. That's always part of the definition in mathematics of a polynomial. That means we stop somewhere, at some $n$, with $a_n x^n$. The elements of this ring are such symbols: formal combinations of $a_i x^i$, where these $a_i$ are in $R$, and possibly different $n$, where $n \in \mathbb{Z}_{\ge 0}$. 

If I have $\mathbb{Z}[x]$, the elements in here are sometimes written $f(x)$, because we like to write polynomials with the variable. It could be, for example, $f(x) = 3 + 2x + 15x^2 + 345x^{200}$. This is a totally legitimate element of it. Other times, you'll see people just write this as $f$ and suppress the variable. Those are two different ways you can write polynomials. 

For $R[x]$, it's just the same thing: you write an element here, $f(x)$, by picking some finite number of elements of $R$ --- the coefficients, all the way to $a_n$ --- and then you just put the powers of $x$ here. 

Here, $x$ is a formal variable: there's no relation; $x$ is in some sense just a placeholder. Those are the elements of $R[x]$.
\end{speech}

\begin{speech}[01:20:17 - 01:21:26]
$R[x]$ has a confusing element, just like the integers: there's a confusing polynomial, like the constant polynomial. If I put $17$ here, that is also in $R[x]$. 

That's the polynomial that only has a constant term; it just happens that there's no $x$ in that. It's weird to write this polynomial as a function of $x$, but it is a function of $x$ --- it's a constant function. In this ring, there are also the constant polynomials, the ones with no $x$. 

Is it clear what this set is? You take $R$, and then you take an $x$ that has nothing to do with $R$ --- it's just some kind of Platonic ideal indeterminate --- and then you consider these formal linear combinations. That's what a polynomial is. 

In particular, I don't really mean anything by this plus: it's just a formal linear combination. Some authors don't like to write this because they think you might think that this plus here means something, which is a legitimate thing to think, actually.
\end{speech}

\begin{content}[{The Polynomial Ring \texorpdfstring{$R[x]$}{R[x]}: Informal Formulation}]
\recap{Polynomial rings were already used informally: $\mathbb{Z}[x] \subseteq \mathbb{Q}[x] \subseteq \mathbb{R}[x] \subseteq \mathbb{C}[x]$ in Week~1 (\cref{ex:inclusions-standard-rings}) and $\mathbb{Z}[x^2] \subset \mathbb{Z}[x]$ in Week~2 (\cref{ex:even-polynomials-subring}). Here they are defined formally.}

\begin{definition}[Polynomial Ring over a Commutative Ring]\label[definition]{def:poly-ring-informal}
Let $R$ be a commutative ring with unit $1_R$. The \newterm{polynomial ring in one variable} (or indeterminate) $x$ with coefficients in $R$, denoted $R[x]$, is defined as the set of formal expressions:
\begin{equation}\label{eq:poly-set-def}
R[x] := \left\{ \sum_{i=0}^n a_i x^i = a_0 + a_1 x + a_2 x^2 + \dots + a_n x^n \;\middle|\; n \in \mathbb{N}_0, \ a_i \in R \right\}.
\end{equation}
\end{definition}

\begin{example}[{Polynomials in $\mathbb{Z}[x]$}]\label[example]{ex:poly-zx}
In the ring $\mathbb{Z}[x]$, an element may be denoted $f(x)$ or simply $f$:
\[
f(x) = 3 + 2x + 15x^2 + 345x^{200} \in \mathbb{Z}[x].
\]
Constant elements such as $17 \in \mathbb{Z}$ correspond to polynomials of degree $0$ with $n = 0$:
\[
17 = 17 \cdot x^0 \in \mathbb{Z}[x].
\]
\end{example}

\begin{steps}[The Indeterminate as a Formal Placeholder]
Crucially, $x$ is not an element of $R$, nor does it represent an unknown numerical value in this context. It is an \newterm{indeterminate} (formal variable) acting as a positional placeholder. The symbols $+$ and $x^i$ in \eqref{eq:poly-set-def} are purely typographical markers separating the coefficients $a_i \in R$, rather than evaluations of operations.
\end{steps}
\end{content}

\subsection{Rigorous Foundation: Polynomials as Almost-Zero Sequences}

\begin{speech}[01:21:26 - 01:22:28]
\inlinemetanote{The lecturer pushes the upper chalkboard up and begins erasing the lower board}
If you look at books where they're really formally correct, they will do some weird things just so you don't get confused. I will explain it for a moment, but I don't usually do that because I think it's a little dishonest. 

When I give these lectures, I try to actually give the lectures in the way I think about the subject. There's some kind of underlying honesty in the lecture. 

If I were to be dishonest and formally correct, the way you define these polynomials is the following way.
\end{speech}

\begin{speech}[01:22:28 - 01:24:00]
You say a polynomial, an element of $R[x]$, is a tuple of elements of $R$, $(a_0, a_1, a_2, \dots)$ --- an infinite tuple --- such that there exists some $N$ such that for all $m > N$, the element in this tuple is zero, okay? \inlinemetanote{writes \qt{an element of $R[x]$ is a tuple of elements of $R$} $(a_0, a_1, a_2, \dots)$ such that $\exists N \ \forall m > N, a_m = 0$}

That's a formal definition of a polynomial in one variable, where I just give you the coefficients. The coefficients are $a_0, a_1, a_2, \dots$; they're all elements of $R$. 

I just give you the coefficients, and there could be lots of them --- any finite number --- but eventually this sequence of coefficients is all zero. That's the formal definition of a polynomial. 

It's actually the same thing that I'm saying, but it's much more reasonable to think about it with the $x$'s. This way is just so that you don't get confused formally. 

That's the set. I hope there's no confusion about the set.
\end{speech}

\begin{hidden}
% timestamp: 01:23:55
% topic: Formal definition of R[x] as sequences with finite support (almost-zero tuples)
% board_state: R[x] as infinite tuple (a_0, a_1, ...) with a_m = 0 for m > N
% next_goal: Identification of the indeterminate x as (0, 1, 0, 0, ...)
% open_loops: none
\end{hidden}

\begin{speech}[01:24:00 - 01:24:24]
In this abstract way, what is the variable $x$? \inlinemetanote{pause}

You have five minutes to figure it out. 

What's the variable $x$? In my notation there's an $x$; here there's no $x$. So, what's $x$ here? 

Yeah?
\end{speech}

\begin{student}[Student Answer]
$(0, 1, 0, \dots)$
\end{student}

\begin{speech}[01:24:27 - 01:24:40]
Yeah. So, $x$ is $(0, 1, 0, 0, \dots)$; that's $x$ if you think about it this way, okay? \inlinemetanote{writes $x = (0, 1, 0, \dots)$ on the board}
\end{speech}

\begin{content}[Formal Definition of Polynomials as Sequences with Finite Support]
\begin{definition}[{Rigorous Definition of $R[x]$ via Coefficient Tuples}]\label[definition]{def:poly-rigorous-sequences}
Let $R$ be a commutative ring with unit. The polynomial ring $R[x]$ is formally defined as the set of sequences $(a_k)_{k \in \mathbb{N}_0}$ of elements in $R$ with \newterm{finite support} (almost-zero sequences):
\begin{equation}\label{eq:poly-formal-sequence}
R[x] := \left\{ (a_0, a_1, a_2, \dots) \in R^{\mathbb{N}_0} \;\middle|\; \exists N \in \mathbb{N}_0 \ \forall m > N \colon a_m = 0 \right\}.
\end{equation}
\end{definition}

\begin{proposition}[Identification of the Indeterminate $x$ and Constants]\label[proposition]{prop:x-identification}
Under the sequence model \eqref{eq:poly-formal-sequence}:
\begin{enumerate}[label=\textbf{(\arabic*)}]
    \item \textbf{Indeterminate $x$:} The indeterminate $x$ corresponds to the elementary sequence with $1$ in position $1$:
    \begin{equation}\label{eq:x-as-sequence}
    x := (0, 1, 0, 0, \dots).
    \end{equation}
    \item \textbf{Constant Elements:} For any $a \in R$, the constant polynomial corresponds to the sequence:
    \begin{equation}\label{eq:constant-as-sequence}
    a := (a, 0, 0, 0, \dots).
    \end{equation}
\end{enumerate}
\end{proposition}

\begin{center}
\begin{tikzpicture}[scale=1.2, >=Stealth]
% \begin{hidden}
% - Canvas: scale=1.2, horizontal sequence of boxes representing sequence indices k=0, 1, 2, ..., N, N+1, ...
% - Box 0 has a_0, Box 1 has a_1, Box N has a_N (highlighted MidnightBlue), Box N+1 to infty have 0 (gray!40)
% - Bracket below indicates "Finite Support / Vanishing Tail"
% \end{hidden}
    % Sequence slots
    \foreach \x/\label in {0/{$a_0$}, 1/{$a_1$}, 2/{$a_2$}, 3/{$\dots$}, 4/{$a_N$}} {
        \draw[thick, MidnightBlue, fill=MidnightBlue!10] (\x*1.2, 0) rectangle ++(1.0, 0.8);
        \node at (\x*1.2 + 0.5, 0.4) {\label};
    }
    % Zero tail slots
    \foreach \x/\label in {5/{$0$}, 6/{$0$}, 7/{$\dots$}} {
        \draw[thick, gray!60, fill=gray!10] (\x*1.2, 0) rectangle ++(1.0, 0.8);
        \node[gray!70] at (\x*1.2 + 0.5, 0.4) {\label};
    }

    % Indices: boxes are 1.0 wide with 0.2 gaps (1.2 tikz units = 1.44 cm per slot), so each label must stay
    % narrower than that; "k =" is written once at the left and the slots carry only the index
    \node[left, font=\footnotesize] at (0, -0.2) {$k =$};
    \node[below, font=\footnotesize] at (0.5, 0) {$0$};
    \node[below, font=\footnotesize] at (1.7, 0) {$1$};
    \node[below, font=\footnotesize] at (2.9, 0) {$2$};
    \node[below, font=\footnotesize] at (5.3, 0) {$N$};
    \node[below, font=\footnotesize] at (6.5, 0) {$N+1$};
    \node[below, font=\footnotesize] at (7.7, 0) {$N+2$};

    % Annotations: braces at y = -0.55, below the index labels (which end near y = -0.35)
    \draw[thick, BrickRed, decorate, decoration={brace, amplitude=6pt, mirror}] (0, -0.55) -- (5.8, -0.55)
        node[midway, below=8pt, text=BrickRed, font=\small\bfseries] {Arbitrary coefficients in $R$};

    \draw[thick, ForestGreen, decorate, decoration={brace, amplitude=6pt, mirror}] (6.0, -0.55) -- (9.4, -0.55)
        node[midway, below=8pt, text=ForestGreen, font=\small\bfseries] {All zero for $m > N$};
\end{tikzpicture}
\end{center}
\end{content}

\subsection{Ring Operations on \texorpdfstring{$R[x]$}{R[x]}: Addition and the Cauchy Product}

\begin{speech}[01:24:40 - 01:25:20]
That's the set. Once you agree on what the set is --- and I gave you two ways to think about it: one way which I think is very simple, and this way which is more formally correct. 

I need to give you the addition rule, the multiplication rule, $0$, and $1$. This is going to be easy: $0$ is just the zero polynomial. In my case, it's just the polynomial $0$; in this sequence case, you have to write an infinite number of zeros, which takes a long time. 

And $1$ is just the constant one polynomial; here it's $(1, 0, 0, 0, \dots)$.
\end{speech}

\begin{student}[Student Question]
Is the plus in that definition not the plus from the ring?
\end{student}

\begin{speech}[01:25:25 - 01:26:28]
It is eventually, but I haven't defined the plus for the ring yet. In this definition, I'm defining $R[x]$ as a set first, and then I'll define addition. When I finish, it will be. But when it's born, it's just a set. Is that clear? 

It's a good question, and that's the reason why thinking about it this way is slightly confusing. Although I encourage you to think about it this way, because you have a lot of experience with polynomials, and in the end, nothing will go wrong. 

You have to agree on what the set is, and I gave you two different ways to think about the set. Then you have to define addition. 

You can define addition in that informal way as the addition of polynomials, which is very good because I don't have to write it out, since you already know how to add polynomials. If I write it here, then I have to decide how to add polynomials, and that's not hard. How do you add polynomials here? 

Yeah?
\end{speech}

\begin{student}[Student Answer]
Componentwise addition.
\end{student}

\begin{speech}[01:26:30 - 01:26:48]
Componentwise addition, which is also how you add polynomials there: componentwise addition. And it will be the case that if I take the constant polynomial and I take this polynomial and add them, I get the thing I call $a_0 + x$, which is also kind of nice. \inlinemetanote{writes $a_0 + x$ on the chalkboard}

All right, addition is easy.
\end{speech}

\begin{speech}[01:26:48 - 01:27:43]
Then we have to define multiplication. Multiplication is more of a pain. If you think about it informally, you already know how to multiply polynomials: you do this thing where you distribute and whatever --- you know how to do it. 

How do you define polynomial multiplication here? Let's try to do that. 

If I use this list of coefficients: here's my first polynomial $f$, and my second polynomial $g$... \inlinemetanote{writes $f = (a_0, a_1, a_2, \dots)$ and $g = (b_0, b_1, b_2, \dots)$ on the board}

Then I have to define this polynomial which is $f \times g$. How do I do it? What's this first coefficient going to be? 

What do you think? 

Yeah?
\end{speech}

\begin{student}[Student Answer]
$a_0 b_0$.
\end{student}

\begin{speech}[01:27:45 - 01:28:54]
Yeah, because it's just the constant term. The secret to doing this is that you write this out as a polynomial. That's the easy way to think about it: write this as a polynomial, do polynomial multiplication, and the constant term is $a_0 b_0$. 

What is this next term? Well, it's going to be $a_1 b_0 + a_0 b_1$. The second term is going to be $a_2 b_0 + a_1 b_1 + a_0 b_2$. \inlinemetanote{writes $fg = (a_0 b_0, \ a_1 b_0 + a_0 b_1, \ a_2 b_0 + a_1 b_1 + a_0 b_2, \dots)$ on the chalkboard}

Formally, this is the definition of polynomial multiplication. 

Is that clear? 

If you want to be on the formal side: polynomials are these lists of coefficients that eventually become all zero; they contain the $0$ and $1$ by the elements here and all those zeros; you can add them by componentwise addition; and the interesting thing is multiplying them, where you use this rule.
\end{speech}

\begin{student}[Student Question]
This is essentially like if you do, like, in... what do you call that, the Cauchy product or convolution...
\end{student}

\begin{hidden}
% timestamp: 01:29:00
% topic: Ring structure on R[x]: addition, Cauchy product multiplication, and convolution interpretation
% board_state: def of fg as convolution, 0 = (0, 0, ...), 1 = (1, 0, ...)
% next_goal: Conclude the lecture on polynomial rings
% open_loops: none
\end{hidden}

\begin{speech}[01:29:06 - 01:30:00]
Yeah, you do this thing. That's what you're saying? Yeah, that's what you do: that's how you multiply polynomials. This is the same rule as you'd get here. Look: if you take the linear term, how do you get an $x$ here? You have to take $x$ times a constant; those are the two ways to get the $x$. 

So, you get this rule. 

Okay, that's the ring. We're going to talk about it more next time. But yeah, it's a little bit of a subtle matter, and it's good to make sure you completely understand it: given a ring, there's a polynomial ring. 

A polynomial is actually just a list of coefficients. It's very helpful to think about the polynomial that way, because that's how you normally think about it, and you won't go wrong thinking about it that way. 

Okay, stop. \inlinemetanote{applause; video fades out}
\end{speech}

\begin{content}[{The Ring Structure of \texorpdfstring{$R[x]$}{R[x]}}]
\begin{definition}[{Ring Operations on $R[x]$}]\label[definition]{def:poly-ring-operations}
Let $R$ be a commutative ring with unit $1_R$. The polynomial ring $(R[x], +, \cdot, 0_{R[x]}, 1_{R[x]})$ is equipped with the following structures:
\begin{enumerate}[label=\textbf{(\arabic*)}]
    \item \textbf{Neutral Elements:}
    \begin{align}
    0_{R[x]} &:= (0_R, 0_R, 0_R, \dots), \label{eq:poly-zero} \\
    1_{R[x]} &:= (1_R, 0_R, 0_R, \dots). \label{eq:poly-one}
    \end{align}
    \item \textbf{Addition (Componentwise):} For $f = (a_k)_{k \ge 0}$ and $g = (b_k)_{k \ge 0}$:
    \begin{equation}\label{eq:poly-addition}
    (a_0, a_1, a_2, \dots) + (b_0, b_1, b_2, \dots) := (a_0 + b_0, \ a_1 + b_1, \ a_2 + b_2, \ \dots).
    \end{equation}
    \item \textbf{Multiplication (Cauchy Product / Convolution):} The product $f \cdot g = (c_k)_{k \ge 0}$ is defined by:
    \begin{equation}\label{eq:poly-multiplication}
    c_k := \sum_{i=0}^k a_i b_{k-i} = a_0 b_k + a_1 b_{k-1} + \dots + a_k b_0.
    \end{equation}
    Explicitly:
    \begin{equation}\label{eq:poly-mult-explicit}
    f \cdot g = \bigl( a_0 b_0, \ a_0 b_1 + a_1 b_0, \ a_0 b_2 + a_1 b_1 + a_2 b_0, \ \dots \bigr).
    \end{equation}
\end{enumerate}
Both results are again almost-zero sequences: if $a_k = 0$ for $k > N$ and $b_k = 0$ for $k > M$, then $a_k + b_k = 0$ for $k > \max(N, M)$, and $c_k = 0$ for $k > N + M$ (each term $a_i b_{k-i}$ has $i > N$ or $k - i > M$).
\end{definition}

\begin{center}
\begin{tikzpicture}[scale=1.3, >=Stealth]
% \begin{hidden}
% - Canvas: scale=1.3, grid of terms a_i b_j with anti-diagonals i+j = k
% - Elements: coordinate lattice (i, j) for i, j in {0, 1, 2, 3}
% - Diagonal bands in BurntOrange and BrickRed showing terms summing to c_0, c_1, c_2, c_3
% \end{hidden}
    % Axes
    \draw[->, thick, gray!70] (-0.3, 0) -- (3.5, 0) node[right, text=black] {$i$ (Index of $a_i$)};
    \draw[->, thick, gray!70] (0, -0.3) -- (0, 3.5) node[above, text=black] {$j$ (Index of $b_j$)};

    % Points
    \foreach \i in {0, 1, 2, 3} {
        \foreach \j in {0, 1, 2, 3} {
            \fill (\i, \j) circle (1.8pt);
            \node[above right, font=\tiny, text=gray!80] at (\i, \j) {$a_{\i} b_{\j}$};
        }
    }

    % Anti-diagonal lines for Cauchy product: i + j = k. Each diagonal ends at (0, k) on the j-axis, and its
    % label c_k sits just left of that point (anchored east at (-0.35, k)), outside the grid, so no label
    % overlaps a lattice point, an a_i b_j label or another diagonal.
    \draw[ultra thick, BrickRed] (0, 0) circle (3pt);
    \node[text=BrickRed, font=\footnotesize, left] at (-0.35, 0) {$c_0 = a_0 b_0$};

    \draw[very thick, BurntOrange, dashed] (1, 0) -- (0, 1);
    \node[text=BurntOrange, font=\footnotesize, left] at (-0.35, 1) {$c_1 = a_1 b_0 + a_0 b_1$};

    \draw[very thick, ForestGreen, dashed] (2, 0) -- (0, 2);
    \node[text=ForestGreen, font=\footnotesize, left] at (-0.35, 2) {$c_2 = a_2 b_0 + a_1 b_1 + a_0 b_2$};

    \draw[very thick, MidnightBlue, dashed] (3, 0) -- (0, 3);
    \node[text=MidnightBlue, font=\footnotesize, left] at (-0.35, 3) {$c_3 = \sum_{i+j=3} a_i b_j$};
\end{tikzpicture}
\end{center}

\begin{steps}[Connection to the Cauchy Product and Discrete Convolution]
The formula $c_k = \sum_{i=0}^k a_i b_{k-i}$ exactly mirrors:
\begin{enumerate}[label=\textbf{(\arabic*)}]
    \item \textbf{Distributive Expansion:} In formal series notation:
    \[
    \left(\sum_{i=0}^n a_i x^i\right)\left(\sum_{j=0}^m b_j x^j\right) = \sum_{i=0}^n \sum_{j=0}^m a_i b_j x^{i+j} = \sum_{k=0}^{n+m} \left(\sum_{i+j=k} a_i b_j\right) x^k.
    \]
    The powers of the indeterminate $x$ naturally collect all terms where the sum of indices equals $k$.
    \item \textbf{Cauchy Product of Power Series:} In analysis, the product of two absolutely convergent power series $\sum a_n x^n$ and $\sum b_n x^n$ is given by the exact same convolution formula.
\end{enumerate}
\end{steps}
\end{content}

\begin{aiNote}[Outlook: Formal Power Series in the Next Lecture]
% [PERSISTENT]
In the lecture introduction, the professor remarked that \qt{in some sense even more fun than polynomial rings are power series rings}. This lecture formalized the polynomial ring $R[x]$ as sequences of coefficients with finite support (almost-zero sequences). Dropping the finite support requirement leads directly to the ring of \newterm{formal power series} $R[[x]]$, consisting of arbitrary infinite sequences $(a_0, a_1, a_2, \dots)$ equipped with the same Cauchy product multiplication, which is systematically investigated at the start of the next lecture (\cref{def:power-series-ring-formal}).
\end{aiNote}

\begin{overview}[Summary of Key Concepts]
\begin{itemize}
    \item \textbf{The Group of Units $\operatorname{Unit}(R)$ (also denoted $U(R)$, $R^\times$, or $R^*$):} The set of multiplicatively invertible elements $R^\times = \{x \in R \mid \exists y \in R \colon xy = 1\}$ forms an abelian group under ring multiplication. Key examples include $\mathbb{Z}^* = \{\pm 1\}$, $\mathbb{Z}[i]^* = \{\pm 1, \pm i\}$ (characterized via the norm $N(a+bi) = a^2 + b^2 = 1$), and $(\mathbb{Z}/m\mathbb{Z})^* = \{\bar{x} \mid \gcd(x, m) = 1\}$.
    \item \textbf{Euler's Totient Function $\phi(m)$:} Defined as the order $|(\mathbb{Z}/m\mathbb{Z})^*|$, with explicit values $\phi(p) = p - 1$ and $\phi(p^k) = p^k - p^{k-1} = p^k(1 - 1/p)$ for prime powers.
    \item \textbf{Universal Property of $\mathbb{Q}$:} As the minimal field containing the integral domain $\mathbb{Z}$, any injective ring homomorphism $f \colon \mathbb{Z} \to F$ into a field $F$ factors uniquely through an injective field homomorphism $g \colon \mathbb{Q} \to F$ defined by $g(x/y) = f(x)f(y)^{-1}$.
    \item \textbf{Field of Fractions $\operatorname{Frac}(R)$:} For any integral domain $R$, the field of fractions is constructed as the quotient of formal pairs $S = R \times (R \setminus \{0\})$ by the cross-multiplication equivalence relation $(x_1, y_1) \equiv (x_2, y_2) \iff x_1 y_2 = x_2 y_1$. Transitivity and the non-vanishing of the denominators $y_1 y_2 \neq 0$ crucially rely on the zero-divisor-free property of $R$.
    \item \textbf{Canonical Inclusion and Examples:} The map $i \colon R \hookrightarrow \operatorname{Frac}(R)$, $r \mapsto [(r, 1)]$, is an injective ring embedding. Examples include $\operatorname{Frac}(\mathbb{Z}) = \mathbb{Q}$, $\operatorname{Frac}(\mathbb{Z}[i]) = \mathbb{Q}(i)$, $\operatorname{Frac}(\mathbb{Z}[\frac{1}{3}]) = \mathbb{Q}$, and $\operatorname{Frac}(K) \cong K$ for any existing field $K$.
    \item \textbf{Polynomial Rings $R[x]$:} For a commutative ring with unity $R$, the polynomial ring $R[x]$ is formally defined as the ring of almost-zero sequences $(a_k)_{k \in \mathbb{N}_0}$ in $R^{\mathbb{N}_0}$. It is equipped with componentwise addition and Cauchy product (convolution) multiplication $c_k = \sum_{i=0}^k a_i b_{k-i}$, with indeterminate $x = (0, 1, 0, 0, \dots)$.
\end{itemize}
\end{overview}
% ==========================================
% LatexRefinement Step Output: step4-2026-09-30-wednesday-week3-algebra-1-offset-final.tex
% Input: step3-2026-09-30-wednesday-week3-algebra-1-offset-speech_refined.tex (gemini-3.8-flash)
% Step 4 pass: done manually by Claude (claude-opus-5-5) following last-refinement.md,
%              because the Gemini step 4 run did not complete
% Processed on: 2026-10-04
% ==========================================

% Primary Language: English
% End of the video: 00:42:50

\lecturechapter[3]{Wednesday}{30 Sep}{30 September 2026}{Field of Fractions, Polynomial Rings, and Formal Power Series}

\begin{overview}[Context]
Continuing from the previous lecture, we review the construction and intuitive arithmetic of the field of fractions $\operatorname{Frac}(R)$ for an integral domain $R$. We then return to the polynomial ring $R[x]$ over an arbitrary commutative ring with identity, establish its universal mapping property, and determine its units when $R$ is an integral domain. Finally, we introduce the formal power series ring $R[[x]]$, define its ring operations, and characterize its units via formal geometric series.
\end{overview}

\section{Field of Fractions and Ring Constructions}

\subsection{Field of Fractions of an Integral Domain}

\begin{speech}[00:00:00 - 00:01:02]
Did it start?

Okay, we start. The topic this week was... the first part was: $R$ is an integral domain, and then the construction of the field of fractions of $R$ (Recall: \cref{def:field-of-fractions}). Sometimes this is denoted by $\operatorname{Frac}(R)$, or...

And then the other two topics are: if $R$ is any ring --- which for us is commutative with unit ---, then there's the construction of the polynomials --- we started that, I'll say more --- and also of the power series.

This is the universal symbol for polynomial, and this is the universal mathematical way of describing power series, the power series ring.

I'll talk about these topics. I'll just say one more thing about the field of fractions.
\end{speech}

\begin{content}[Topic Overview: Integral Domains and Ring Constructions]
\begin{itemize}
    \item \textbf{Field of Fractions:} If $R$ is an integral domain, there exists an embedding into a field:
    \[
    R \subseteq F := \operatorname{Frac}(R),
    \]
    where $F$ is called the \newterm{field of fractions} of $R$.
    \item \textbf{Polynomial and Power Series Rings:} If $R$ is any commutative ring with identity $1$, we construct:
    \begin{itemize}
        \item $R[x]$: the \newterm{polynomial ring} in the indeterminate $x$ with coefficients in $R$,
        \item $R[[x]]$: the \newterm{formal power series ring} in the indeterminate $x$ with coefficients in $R$.
    \end{itemize}
\end{itemize}
\end{content}

\begin{speech}[00:01:02 - 00:01:46]
I gave you this very formal way to think about the field of fractions, which is: you take some set $S$ of pairs... I won't review the whole thing now. Set $S$ of pairs, then you have an equivalence relation, and you look at the objects in this field of fractions as being equivalence classes in this equivalence relation.

And then you can define everything at this formal level and check all the properties, and that is the underlying mathematical formalism of the field of fractions. Meaning that if you're ever thinking about the field of fractions and you're wondering about something, you can always resolve it foundationally by saying, \qt{I know what the definition is: it's this, this, and this, precisely.}
\end{speech}

\begin{content}[Formal Construction of the Field of Fractions]
\recap{The field of fractions was constructed in the previous lecture: the set of pairs (\cref{def:set-of-pairs-s}), the equivalence relation (\cref{def:equiv-relation-S}, \cref{prop:equiv-on-S}) and the quotient set (\cref{def:field-of-fractions}). The lecturer summarizes it again.}

Let $R$ be an integral domain. To construct the field of fractions formally:
\begin{enumerate}[label=\textbf{(\arabic*)}]
    \setcounter{enumi}{0} \item \textbf{Set of Pairs:} Consider the Cartesian product $S := R \times (R \setminus \{0\}) \subset R^2$.
    \setcounter{enumi}{1} \item \textbf{Equivalence Relation:} Define the relation $\equiv$ on $S$ by:
    \[
    (r, s) \equiv (t, w) \iff rw - st = 0.
    \]
    \setcounter{enumi}{2} \item \textbf{Equivalence Classes:} The field of fractions is defined as the quotient set:
    \[
    F = \operatorname{Frac}(R) := S / {\equiv}.
    \]
\end{enumerate}
\end{content}

\begin{speech}[00:01:46 - 00:03:24]
No one ever thinks about it that way. The way you think about it is: if $R$ is an integral domain, if you have some elements in $R$, $r$ and $s$, then for an element of the field of fractions, you just say this is $r$ over $s$ ($r/s$) --- that's it, you know? And when you do this, you make sure $s$ is not zero, because that was the second condition there, just like you would with rational numbers.

And if you want to add, multiply, all of this, you just use the rules you did with rational numbers. Like if you want to add $r$ over $s$ and $t$ over $w$, you know how to add fractions; you just do it here, it's just the right answer: it's this.

And so on. If you multiply, then you multiply fractions. Whatever rules you use to manipulate fractions, they're just correct here.

In particular, why is this a field? If I take $r$ over $s$, if this is not zero, that implies that $r$ is not zero. If $r$ were zero, then that's the zero element of this. And you can go check by the definitions that if $r$ is zero, this is the zero element. So, if this is not zero, $r$ must not be equal to zero, and then it has an inverse, and that's the multiplicative inverse. And it's legal because $r$ is not zero, so I can put him downstairs.

Anyway, the way you think about it --- I mean, you get to think about it however you want ---, the way typically people think about the field of fractions is you just write the fractions; that's why it's called the field of fractions! And you just go forward.

And the rules of this formalism are designed so that this actually works.

Okay.
\end{speech}

\begin{content}[Intuitive Arithmetic in the Field of Fractions]
\recap{These fraction rules were already listed in the previous lecture (\cref{eq:frac-eq,eq:frac-add,eq:frac-mult,eq:frac-inv}).}

In mathematical practice, elements of $\operatorname{Frac}(R)$ are represented as formal fractions $\frac{r}{s}$ with $r, s \in R$ and $s \neq 0$:
\begin{align}
    \operatorname{Frac}(R) &= \left\{ \frac{r}{s} \;\middle|\; r, s \in R, \ s \neq 0 \right\}. \label{eq:frac-field-elements}
\end{align}
The field operations are defined analogously to elementary fraction arithmetic:
\begin{align}
    \text{\textbf{Addition:}} && \frac{r}{s} + \frac{t}{w} &= \frac{rw + st}{sw}, \label{eq:frac-field-add} \\
    \text{\textbf{Multiplication:}} && \frac{r}{s} \cdot \frac{t}{w} &= \frac{rt}{sw}, \label{eq:frac-field-mult} \\
    \text{\textbf{Multiplicative Inverse:}} && \left( \frac{r}{s} \right)^{-1} &= \frac{s}{r} \qquad \left(\text{valid whenever } \frac{r}{s} \neq 0 \text{, i.e., } r \neq 0\right). \label{eq:frac-field-inv}
\end{align}

\begin{steps}[Why the Integral Domain Assumption Is Essential]
For the denominator $sw$ in \eqref{eq:frac-field-add} and \eqref{eq:frac-field-mult} to be non-zero, we must have:
\[
s \neq 0 \text{ and } w \neq 0 \implies sw \neq 0.
\]
This property holds if and only if $R$ has no zero divisors, which is precisely the defining condition of an \textbf{integral domain} (\cref{def:integral-domain}). Furthermore, the zero element of $\operatorname{Frac}(R)$ is $\frac{0}{1} = \frac{0}{s}$ (\cref{def:neutral-elements-F}). Thus, $\frac{r}{s} \neq 0$ implies $r \neq 0$, guaranteeing that $\left(\frac{r}{s}\right)^{-1} = \frac{s}{r}$ is a well-defined fraction with non-zero denominator.
\end{steps}
\end{content}

\begin{speech}[00:03:24 - 00:04:44]
This is a basic issue that comes up very often in math education about how to think about it. The first time this comes up --- which you've already passed --- is the real numbers.

To define the real numbers is complicated. You have to either define them as Cauchy sequences or Dedekind cuts, and that's pretty complicated to define. And then you have to check that it satisfies the properties. That's all fine; that's kind of like some kind of logical underlayer. But no one thinks about real numbers that way.

And it's just that that's part of learning mathematics: there's some kind of logical underlayer, and then there's the way you think about it, and usually they have nothing to do with each other.

Okay, that's maybe an exaggeration, but... Actually, about the real numbers, I can't help now but comment that these Dedekind cuts were discovered by Richard Dedekind, who had moved to Zurich in the second part of the 19th century, and he's an algebraist. He's a very famous mathematician. He's, of course, passed away by now.

And he was assigned to teach analysis. And he writes it very beautifully in this book that he wrote about the real numbers. He was assigned to teach analysis, and he's an algebraist, and he said, \qt{Well, I have this difficult time teaching analysis because I don't know what a real number is!} And then he spends the class constructing the real numbers. It's kind of amazing to read the introduction to that. You can find it.
\end{speech}

\begin{insight}[The Logical Underlayer vs. Working Mathematical Intuition]
The lecturer highlights a fundamental pedagogical dichotomy in abstract mathematics: the distinction between the \emph{foundational underlayer} and the \emph{working mathematical intuition}. 

Whether constructing $\mathbb{R}$ via Dedekind cuts or Cauchy sequences, or constructing $\operatorname{Frac}(R)$ via quotient sets of ordered pairs $S / {\equiv}$, the foundational machinery serves solely to establish rigorous existence and soundness. Once established, mathematicians operate with the intuitive rules (e.g., standard fraction manipulations $\frac{r}{s}$) without thinking of the underlying equivalence classes.
\end{insight}

\begin{aiNote}[Beyond the Lecture: Dedekind's Book on the Real Numbers]
% [PERSISTENT, ADDITION]
The book the lecturer refers to is Richard Dedekind's \emph{Stetigkeit und irrationale Zahlen} (\qt{Continuity and Irrational Numbers}, 1872). In its preface Dedekind recalls teaching the differential calculus at the Polytechnikum in Zurich (today's ETH) in 1858 and feeling the lack of a truly rigorous foundation of arithmetic; the book introduces what are now called Dedekind cuts.
\end{aiNote}

\section{Polynomial Rings}

\subsection{Definition and Structure of \texorpdfstring{$R[x]$}{R[x]}}

\begin{speech}[00:04:45 - 00:06:23]
Okay, so we go to polynomials.

I explained yesterday (Recall: \cref{def:poly-ring-informal}) that the way to think about polynomials is just the way you always think about polynomials. Elements of this ring... I advise you to think about them this way. But, in all things, you get your own choice how to think about them.

They're polynomials that have this formal variable $x$ that's denoted here, and you consider polynomials in $x$ with coefficients in $R$.

And this polynomial ring, that's the set, and it has addition. The addition is: you add it like polynomials. So, whatever power, you keep the power the same, and you add the coefficients of that power. You know how to add polynomials; you just do it here.

Multiplying polynomials is more complicated. But you already know how to multiply polynomials, and you just do it here. You can define it formally, which I did last time (Recall: \cref{def:poly-ring-operations}), but you just do it. And it has the zero polynomial and the one polynomial, and this is the ring of polynomials --- the ring of polynomials in the variable $x$ with coefficients in $R$.

So, we can do this for any ring. I said that. This is important to keep track of: sometimes we make assumptions on rings, sometimes not. Here we needed the assumption (i.e., $R$ an integral domain) \inlinemetanote{points to the left board for the field of fractions} --- that's kind of crucial. Here we don't need the assumption.
\end{speech}

\begin{content}[The Polynomial Ring over a Commutative Ring]
\recap{The polynomial ring was introduced in the previous lecture, informally (\cref{def:poly-ring-informal}) and as sequences of coefficients (\cref{def:poly-rigorous-sequences}, \cref{def:poly-ring-operations}). The lecturer recalls it.}

\begin{definition}[Polynomial Ring]\label[definition]{def:polynomial-ring-w3wed}
Let $R$ be an arbitrary commutative ring with identity $1$. The \newterm{polynomial ring} in the variable $x$ with coefficients in $R$, denoted $R[x]$, is the set of formal sums:
\begin{equation}\label{eq:polynomial-set-def}
R[x] := \left\{ a_0 + a_1 x + a_2 x^2 + \dots + a_n x^n \;\middle|\; n \in \mathbb{N}, \ a_i \in R \right\}.
\end{equation}
Equipped with the standard polynomial addition and Cauchy multiplication:
\begin{align}
    \left( \sum_{i=0}^n a_i x^i \right) + \left( \sum_{i=0}^n b_i x^i \right) &:= \sum_{i=0}^n (a_i + b_i) x^i, \label{eq:poly-addition-w3wed} \\
    \left( \sum_{i=0}^n a_i x^i \right) \cdot \left( \sum_{j=0}^m b_j x^j \right) &:= \sum_{k=0}^{n+m} \left( \sum_{i+j=k} a_i b_j \right) x^k, \label{eq:poly-multiplication-w3wed}
\end{align}
the structure $(R[x], +, \cdot, 0, 1)$ forms a commutative ring with identity $1 \in R \subset R[x]$ and zero element $0 \in R \subset R[x]$.
\end{definition}

\begin{steps}[Generality of the Polynomial Construction]
Unlike the construction of the field of fractions $\operatorname{Frac}(R)$, which strictly requires $R$ to be an integral domain, the polynomial ring $R[x]$ is well-defined for \textbf{any} commutative ring with identity $R$. Zero divisors in $R$ do not obstruct the definition of polynomial addition or multiplication.
\end{steps}
\end{content}

\begin{hidden}
% timestamp: 00:06:23
% topic: Review of field of fractions and definition of polynomial rings
% board_state: def:field-of-fractions, def:polynomial-ring-w3wed
% next_goal: Universal property of polynomial rings (extension theorem)
% open_loops: none
\end{hidden}

\subsection{The Universal Property of Polynomial Rings}

\begin{speech}[00:06:29 - 00:08:59]
Okay, so, properties of this polynomial ring.

The first is --- this is somehow the most basic property of polynomial rings: Suppose $R$ and $S$ are rings --- which, as I say, are commutative rings with unit, just for us --- and I have an element $y$ in $S$.

Then, given any homomorphism --- so you have some homomorphism $\phi$ from $R$ to $S$, a homomorphism of rings ---, there exists a unique extension --- I will say what this means in a moment: an extension from the polynomial ring to $S$ --- which satisfies these properties. Two properties.

Property \textbf{(1)} is that if I take a constant polynomial --- so this is just the constant polynomial, which is to say that $R$ sits inside $R[x]$ as the constants, okay? ---, if I take a constant polynomial, this should just be the value that I've already specified here: $\phi(r)$.

And the second property is: I have to say what happens to $x$, because that's not specified here. And I just say $x$ goes to $y$.

So, that's the claim.

If I start with these two rings --- there's the ring where my maps are starting, and the ring where the target map, where the homomorphism is ending --- and I pick an element $y \in S$, I can take any homomorphism from $R$ to $S$ and I can extend it to the polynomial ring by just setting $x$ to $y$.
\end{speech}

\begin{content}[The Universal Property of Polynomial Rings]
\begin{theorem}[{Universal Mapping Property of $R[x]$}]\label[theorem]{thm:universal-property-polynomial-ring}
Let $R$ and $S$ be commutative rings with identity, let $\phi \colon R \to S$ be a ring homomorphism, and let $y \in S$.
Then there exists a unique ring homomorphism $\hat{\phi} \colon R[x] \to S$ extending $\phi$ such that:
\begin{enumerate}[label=\textbf{(\arabic*)}]
    \setcounter{enumi}{0} \item \textbf{Extension of Constants:} For every constant polynomial $r \in R \subset R[x]$:
    \begin{equation}\label{eq:univ-prop-constants}
    \hat{\phi}(r) = \phi(r).
    \end{equation}
    \setcounter{enumi}{1} \item \textbf{Evaluation at $y$:} The indeterminate $x$ is mapped to $y$:
    \begin{equation}\label{eq:univ-prop-indeterminate}
    \hat{\phi}(x) = y.
    \end{equation}
\end{enumerate}
\end{theorem}

\begin{steps}[Evaluation Homomorphism and Factorization]
Going beyond the lecture: the map $\hat{\phi}$ is often called the \newterm{evaluation homomorphism} associated with $\phi$ and $y$. In the special case where $S = R$ and $\phi = \operatorname{id}_R$, $\hat{\phi}$ is simply the evaluation map $\operatorname{ev}_y \colon R[x] \to R$ given by $f(x) \mapsto f(y)$. 

The universal property states that $R[x]$ is the \qt{freest} ring containing $R$ and an additional generator $x$: mapping $R[x]$ to any ring $S$ requires only choosing a base ring homomorphism $\phi \colon R \to S$ and choosing where to send $x$. It is the same kind of statement as the universal property of $\mathbb{Q}$ from the previous lecture (\cref{thm:universal-property-q}): a homomorphism defined on a smaller ring extends uniquely to the bigger one (for $R[x]$, once the image of $x$ has been chosen).
\end{steps}
\end{content}

\begin{proofWrapper}[Proof of \cref{thm:universal-property-polynomial-ring}]

\begin{speech}[00:09:00 - 00:09:41]
Okay, and the proof of this is just by the formula.

I have to define $\hat{\phi}$ on any polynomial. A polynomial looks like this...

And I have to define it. Formally, I would have to first define it, then I have to check that it's a homomorphism --- because I say it extends to a homomorphism ---: define it, check it's a homomorphism, and then show it satisfies these two properties. That's the task at hand, so to speak.

So, how am I going to define this?

Yeah?
\end{speech}

\begin{student}[Student Answer]
Apply $\phi$ to the coefficients, and then send each $x$ to $y$?
\end{student}

\begin{speech}[00:09:53 - 00:12:19]
This is the definition: I have to somehow take this data and get to $S$.

Any element here --- these $a_i$'s, that's part of the definition of a polynomial --- is in $R$. And I only have on the board one way to ship elements from $R$ to $S$: it's here \inlinemetanote{points to $\phi \colon R \to S$}. I can't make up another one.

So, all these coefficients I just move by the only way I know how to get to $S$. That's the constant coefficient, then I move this coefficient to $S$, because this result has to be in $S$. And then the question is: how do I move $x$ to $S$? Well, I know one element of $S$, that's $y$.

And then I do this: $\phi(a_2) y^2$, and so on, $\phi(a_n) y^n$. So, I just define that. The first thing you have to accept is that this definition makes sense.

When you define a map, usually, if I define it so that whenever you give me anything, I tell you what the answer is, that's the definition of the map --- as long as it takes any possibility here, which it does.

In the case where you have equivalence classes, the matter is a little subtle, because in equivalence classes things have different names, and you have to be careful. But that's not what's going on here! This is not equivalence classes. The field of fractions was equivalence classes, but the polynomial ring is not an equivalence class. I tell you specifically who is a polynomial, and everyone has just one name.

So, that's the definition.

But is that the end? No, I have to check that it satisfies this property on constant polynomials. But for constant polynomials, well, it does satisfy that property --- that would just mean there's none of this higher part, so just by definition it satisfies that property.

And I have to check $x$. What is $\hat{\phi}(x)$ by this definition? Well, that would be just $1 \cdot x$ here, and everything else is $0$ or $1$. Formally, I'd have to put $\phi(0)$ and $\phi(1)$. I'm not going to write that, but you could write that. But $\phi(0)$ is $0$ and $\phi(1)$ is $1$, so you don't have to write it. This is just $y$.

So, I have now checked these things, but the one thing I haven't checked is that it's a homomorphism. Let's think about how you'd check that.
\end{speech}

\begin{content}[Definition of \texorpdfstring{$\hat{\phi}$}{phi-hat} and Verification of Conditions]
\textbf{1. Definition of the Candidate Map:}
Let $f(x) = a_0 + a_1 x + a_2 x^2 + \dots + a_n x^n \in R[x]$ with $a_i \in R$. We define the map $\hat{\phi} \colon R[x] \to S$ by:
\begin{equation}\label{eq:phi-hat-def}
\hat{\phi}(a_0 + a_1 x + \dots + a_n x^n) \stackrel{\text{def}}{:=} \phi(a_0) + \phi(a_1) y + \phi(a_2) y^2 + \dots + \phi(a_n) y^n.
\end{equation}
Since every polynomial in $R[x]$ has a unique representation as a formal sum of powers of $x$ (unlike equivalence classes in $\operatorname{Frac}(R)$, there are no ambiguous representatives), $\hat{\phi}$ is unconditionally well-defined.

\textbf{2. Verification of Properties \textbf{(1)} and \textbf{(2)}:}
\begin{itemize}
    \item \textbf{Constant Polynomials:} If $f(x) = r \in R$, then $a_0 = r$ and $a_i = 0$ for all $i \ge 1$. Hence:
    \[
    \hat{\phi}(r) = \phi(r).
    \]
    \item \textbf{Image of $x$:} For $f(x) = x$, we have $a_0 = 0$, $a_1 = 1$, and $a_i = 0$ for all $i \ge 2$. Since $\phi(0) = 0$ and $\phi(1) = 1$:
    \[
    \hat{\phi}(x) = \phi(0) + \phi(1) y = 0 + 1 \cdot y = y.
    \]
\end{itemize}
\end{content}

\begin{hidden}
% timestamp: 00:12:19
% topic: Proving the universal property of polynomial rings
% board_state: thm:universal-property-polynomial-ring, eq:phi-hat-def
% next_goal: Verify additivity and multiplicativity of the extension homomorphism
% open_loops: none
\end{hidden}

\begin{speech}[00:12:24 - 00:13:30]
\inlinemetanote{The lecturer pushes the middle blackboard up and continues on the lower board.}

To check it's a homomorphism, you have to check some things about zero going to zero, and one goes to one, which it certainly does. So, the main thing to check is that $\hat{\phi}$ of two polynomials --- so $f$ and $g$ now are my polynomials --- equals $\hat{\phi}(f) + \hat{\phi}(g)$.

For this, you need to know how to... sorry, these are $f$, right? This is $g$. So, you need to know how to add the polynomials here, and then you need to know how to add elements of $S$. This addition is addition of polynomials, and this addition is addition in $S$, right?
\end{speech}

\begin{speech}[00:13:30 - 00:15:10]
And the way this thing works is: what are we going to do here? Let's say that we take this $\hat{\phi}$... If $f$ is given by $a$'s, it's $a_0 + a_1 x$, let's say --- let's say that's $f$, just a linear one --- plus $b_0 + b_1 x$.

If we add them first in polynomials, then this is $(a_0 + b_0) + (a_1 + b_1) x$, okay? And then, if we go over by the definition, this is $\phi(a_0 + b_0) + \phi(a_1 + b_1) y$, right?

And here, I can now use the homomorphism property of $R$, and write this as $\phi(a_0) + \phi(a_1) y + \phi(b_0) + \phi(b_1) y$ --- I don't know if this is really worth doing all out, but: $\phi(b_0) + \phi(b_1) y$. And that, at the end, is $\hat{\phi}(a_0 + a_1 x) + \hat{\phi}(b_0 + b_1 x)$.
\end{speech}

\begin{content}[Verification of Additivity]
\textbf{3. Verification of Additivity ($\hat{\phi}(f + g) = \hat{\phi}(f) + \hat{\phi}(g)$):}
Let $f, g \in R[x]$. We verify additivity explicitly for linear polynomials $f(x) = a_0 + a_1 x$ and $g(x) = b_0 + b_1 x$:
\begin{align}
    \hat{\phi}(f + g) &= \hat{\phi}\bigl( (a_0 + a_1 x) + (b_0 + b_1 x) \bigr) \nonumber \\
    &= \hat{\phi}\bigl( (a_0 + b_0) + (a_1 + b_1) x \bigr) && \text{(addition in $R[x]$)} \nonumber \\
    &= \phi(a_0 + b_0) + \phi(a_1 + b_1) y && \text{(by definition of $\hat{\phi}$)} \nonumber \\
    &= \bigl(\phi(a_0) + \phi(b_0)\bigr) + \bigl(\phi(a_1) + \phi(b_1)\bigr) y && \text{($\phi \colon R \to S$ is a homomorphism)} \nonumber \\
    &= \bigl(\phi(a_0) + \phi(a_1) y\bigr) + \bigl(\phi(b_0) + \phi(b_1) y\bigr) && \text{(reordering in ring $S$)} \nonumber \\
    &= \hat{\phi}(a_0 + a_1 x) + \hat{\phi}(b_0 + b_1 x) \nonumber \\
    &= \hat{\phi}(f) + \hat{\phi}(g). \label{eq:additivity-check}
\end{align}
The same argument works for polynomials of arbitrary degree: writing $f = \sum_{i=0}^n a_i x^i$ and $g = \sum_{i=0}^n b_i x^i$ (padding the shorter one with zero coefficients),
\[
\hat{\phi}(f + g) = \sum_{i=0}^n \phi(a_i + b_i) \, y^i = \sum_{i=0}^n \phi(a_i) \, y^i + \sum_{i=0}^n \phi(b_i) \, y^i = \hat{\phi}(f) + \hat{\phi}(g).
\]
\end{content}

\begin{speech}[00:15:11 - 00:16:43]
What did I do here? I just expanded out by the definition, and at a certain point I have to use the homomorphism property of $R$.

Nothing happened here, actually; it was just a lot of waste of chalk. But at a certain point, I have to use the homomorphism property of $R$. There.

And I did this with linear polynomials, but it's just the same with more coefficients.

And then you have to check something that's more complicated with the indices: that if I take $f \cdot g$, this is equal to $\hat{\phi}(f) \cdot \hat{\phi}(g)$. Again, this follows from the same expansion.

I tried to do some checking for you, but in the end, I don't know how successful that was. The fact of the matter is: I give you the definition, and then you have to check the compatibility. And when you check the compatibility, all you use is the homomorphism property of this.

That's one abstract property of polynomial rings: if you can map $R$ to something, then you can pick any landing point for $x$, and it gives you a homomorphism.

Sometimes this is called the universal property of polynomial rings.
\end{speech}

\begin{content}[Verification of Multiplicativity and Uniqueness]
\textbf{4. Verification of Multiplicativity ($\hat{\phi}(f \cdot g) = \hat{\phi}(f) \cdot \hat{\phi}(g)$):}
Let $f(x) = \sum_{i=0}^n a_i x^i$ and $g(x) = \sum_{j=0}^m b_j x^j$. Using the Cauchy product in $R[x]$:
\begin{align}
    \hat{\phi}(f \cdot g) &= \hat{\phi}\left( \sum_{k=0}^{n+m} \left( \sum_{i+j=k} a_i b_j \right) x^k \right) \nonumber \\
    &= \sum_{k=0}^{n+m} \phi\left( \sum_{i+j=k} a_i b_j \right) y^k && \text{(by definition of $\hat{\phi}$)} \nonumber \\
    &= \sum_{k=0}^{n+m} \left( \sum_{i+j=k} \phi(a_i) \phi(b_j) \right) y^k && \text{($\phi$ is a ring homomorphism)} \nonumber \\
    &= \left( \sum_{i=0}^n \phi(a_i) y^i \right) \cdot \left( \sum_{j=0}^m \phi(b_j) y^j \right) && \text{(distributivity and Cauchy product in $S$)} \nonumber \\
    &= \hat{\phi}(f) \cdot \hat{\phi}(g). \label{eq:multiplicativity-check}
\end{align}
Furthermore, $\hat{\phi}(1) = \phi(1) = 1_S$. Thus, $\hat{\phi}$ is a ring homomorphism.

\textbf{5. Uniqueness:}
Any ring homomorphism $\psi \colon R[x] \to S$ satisfying conditions \textbf{(1)} and \textbf{(2)} of \cref{thm:universal-property-polynomial-ring} must preserve sums and products:
\[
\psi\left( \sum_{i=0}^n a_i x^i \right) = \sum_{i=0}^n \psi(a_i) \psi(x)^i = \sum_{i=0}^n \phi(a_i) y^i = \hat{\phi}\left( \sum_{i=0}^n a_i x^i \right).
\]
Hence $\psi = \hat{\phi}$, proving uniqueness.
\end{content}

\end{proofWrapper}

\subsection{Units in the Polynomial Ring over an Integral Domain}

\begin{speech}[00:16:45 - 00:18:16]
We know that if you have rings, you can think about their units. So, we can ask: what are the units in a polynomial ring? That's a very reasonable question.

And that question actually turns out to have some wrinkles. Let's start with: Suppose $R$ is an integral domain.

Then the question is: what are the units in the polynomial ring? Which is to say: which elements of this polynomial ring have a multiplicative inverse?

What do you think the answer to this is? If I have some $f$ in the polynomial ring, $f = a_0 + a_1 x + \dots + a_n x^n$, and if it's a unit, that means there exists a multiplicative inverse. There exists $g = b_0 + b_1 x + \dots + b_m x^m$, and I multiply them so that I get $1$, the constant polynomial. That would mean that if $f$ is a unit, then there exists $g$ such that when I multiply them, I get the constant polynomial $1$.

How could this happen?

Yeah?
\end{speech}

\begin{student}[Student Question]
If there is an element in $R$, you can plug it in and then it goes to zero?
\end{student}

\begin{speech}[00:18:21 - 00:18:25]
There exists an element of... Sorry, you have to speak louder.
\end{speech}

\begin{student}[Student Question, continued]
Like, when there exists an element in $R$...
\end{student}

\begin{hidden}
% timestamp: 00:18:27
% topic: Determining the units in the polynomial ring R[x]
% board_state: def:units-polynomial-ring, eq:poly-invertible-condition
% next_goal: Prove that Unit(R[x]) = Unit(R) when R is an integral domain
% open_loops: student suggested evaluating at elements of R; lecturer clarifies x cannot be replaced
\end{hidden}

\begin{speech}[00:18:27 - 00:18:56]
This $x$ stays! You can't replace $x$. We're talking about this polynomial ring.

Actually, the discussion is kind of simpler: if I multiply this $f$ with $g$, I get another polynomial. And for it to be $1$, all of its higher terms have to somehow cancel, right? So, the question is: can you arrange it so all the higher terms cancel? That's the real question.

Yeah?
\end{speech}

\begin{student}[Student Answer]
I suppose if it's a constant term?
\end{student}

\begin{speech}[00:18:59 - 00:20:09]
Yeah, so one way to do it is: the answer is that a unit... if I take the polynomial to be a constant, $f$ equals a constant polynomial and this happens to be a unit for $R$, right?

He's saying: we could have just the constants --- take that to be a constant polynomial, and if I further choose that to be a unit, then I know it has a friend, because it has an inverse in $R$.

So, this part I hope is good: the units of the original ring sit inside the units of this polynomial ring, and they sit inside as the constant polynomials.

I hope that's clear.

Could there be any more units here?

And you should, in the way I presented this, be sensitive to the fact that I didn't take any ring: I took an integral domain. So, you might want to think: why did I do that?

Yeah?
\end{speech}

\begin{content}[Constant Units in the Polynomial Ring]
Let $R$ be a commutative ring with identity, and consider $R[x]$:
\begin{itemize}
    \item A polynomial $f(x) \in R[x]$ is a \newterm{unit} (\cref{def:unit-ring}) if there exists $g(x) \in R[x]$ such that:
    \begin{equation}\label{eq:poly-unit-condition}
    f(x) \cdot g(x) = 1.
    \end{equation}
    \item Any unit $u \in \operatorname{Unit}(R)$ has an inverse $u^{-1} \in R$. Considered as constant polynomials in $R[x]$:
    \[
    u \cdot u^{-1} = 1 \implies u \in \operatorname{Unit}(R[x]).
    \]
    Thus, the units of $R$ always embed into the units of $R[x]$:
    \begin{equation}\label{eq:units-embedding}
    \operatorname{Unit}(R) \subseteq \operatorname{Unit}(R[x]) \quad \text{(as constant polynomials)}.
    \end{equation}
\end{itemize}
\end{content}

\begin{student}[Student Answer]
Because there are no more units?
\end{student}

\begin{speech}[00:20:11 - 00:20:37]
Yes! The answer here is that there are no more. It's actually equal.

The answer is that there are no more units: the units in $R$ equal the units in $R[x]$.

And I just remind you, because otherwise it's going to be wrong: $R$ is an integral domain.

How are you going to use this condition?
\end{speech}

\begin{student}[Student Answer]
So, you have the general form of $f$ and $g$, and if we multiply them, then the term with the highest power...
\end{student}

\begin{speech}[00:20:45 - 00:22:13]
Yeah, that is the right answer. When you multiply polynomials --- and as I said, you must have had some years of experience multiplying polynomials, and you should not forget those years of experience.

If you multiply polynomials, in the middle is a complicated mess of things. You have to add up all the terms; it's some complicated mess. But there are two terms which are very simple.

If you want the constant polynomial of the product, that's the product of the constants, right? And if you want the highest term in the product, that's the product of the highest terms.

With polynomials, if you're multiplying them, before you do all the multiplication, you should think: is what I'm thinking about resolved by understanding what happens to the constants or the highest term? Because those are the easiest ones, since there's not a big sum there.

The constant term of $fg$ is just equal to $a_0 b_0$. And the highest term --- meaning the highest degree term in $x$ --- comes from the highest term of $f$ times the highest term of $g$.
\end{speech}

\begin{content}[Leading and Constant Terms in Polynomial Multiplication]
Let $R$ be a commutative ring with identity, and let $f(x), g(x) \in R[x]$ be non-zero polynomials represented in canonical form:
\begin{align}
    f(x) &= a_0 + a_1 x + \dots + a_n x^n, \quad a_n \neq 0 \ (\deg(f) = n), \label{eq:poly-f-rep} \\
    g(x) &= b_0 + b_1 x + \dots + b_m x^m, \quad b_m \neq 0 \ (\deg(g) = m). \label{eq:poly-g-rep}
\end{align}
When multiplying $f(x)$ and $g(x)$, the extremal coefficients take an uncomplicated form:
\begin{align}
    \text{\textbf{Constant Term:}} && \operatorname{const}(fg) &= a_0 b_0, \label{eq:const-term-fg} \\
    \text{\textbf{Highest Degree Term:}} && \operatorname{highest}(fg) &= (a_n x^n)(b_m x^m) = a_n b_m x^{n+m}. \label{eq:highest-term-fg}
\end{align}
\end{content}

\begin{speech}[00:22:13 - 00:23:23]
Usually, when people write polynomials --- but you have to be careful about the conventions of whatever author you're reading ---, the typical thing is: when you write a polynomial like this, you could say $a_n$ is not zero, because if it were zero, you wouldn't have written it. So, it's kind of standard that when you write a polynomial like this, the highest term you say is not zero. And this one is not zero, because otherwise I wouldn't have to write it.

If I take these polynomials, then this is the highest term, and that's the highest term. This would be the highest term here: $a_n x^n$; that would be the highest term there: $b_m x^m$. And the highest term would be $a_n \cdot b_m \cdot x^{n+m}$. That would be the highest term.

And if we're in an integral domain, this product --- if they start out to be non-zero --- ends up being non-zero. So, this would still be non-zero if $R$ is an integral domain.
\end{speech}

\begin{content}[Units of the Polynomial Ring over an Integral Domain]
\begin{theorem}[{Units in $R[x]$ for an Integral Domain $R$}]\label[theorem]{thm:units-polynomial-ring-integral-domain}
Let $R$ be an integral domain. Then the units of the polynomial ring $R[x]$ are precisely the units of the base ring $R$, regarded as constant polynomials:
\begin{equation}\label{eq:units-rx-equals-units-r}
\operatorname{Unit}(R[x]) = \operatorname{Unit}(R).
\end{equation}
\end{theorem}

\begin{proof}
The inclusion $\operatorname{Unit}(R) \subseteq \operatorname{Unit}(R[x])$ holds trivially because any unit $u \in \operatorname{Unit}(R)$ has an inverse $u^{-1} \in R$, which satisfies $u \cdot u^{-1} = 1$ in $R[x]$ as constant polynomials.

Conversely, suppose $f(x) \in \operatorname{Unit}(R[x])$. Then there exists an inverse polynomial $g(x) \in R[x]$ such that:
\begin{equation}\label{eq:unit-inverse-rel}
f(x) \cdot g(x) = 1.
\end{equation}
Both $f$ and $g$ are non-zero, since $f g = 1 \neq 0$ (an integral domain has $1 \neq 0$). Write $f(x) = \sum_{i=0}^n a_i x^i$ with $a_n \neq 0$ and $g(x) = \sum_{j=0}^m b_j x^j$ with $b_m \neq 0$.
The leading term of the product is given by:
\[
\operatorname{highest}(f \cdot g) = a_n b_m x^{n+m}.
\]
Because $R$ is an \textbf{integral domain}, the absence of non-trivial zero divisors ensures:
\[
a_n \neq 0 \text{ and } b_m \neq 0 \implies a_n b_m \neq 0 \quad \text{in } R.
\]
On the other hand, the constant polynomial $1$ has degree $0$. Thus, comparing degrees on both sides of \eqref{eq:unit-inverse-rel} yields:
\[
n + m = \deg(fg) = \deg(1) = 0.
\]
Since $n, m \in \mathbb{N} = \{0, 1, 2, \dots\}$, this forces:
\[
n = 0 \quad \text{and} \quad m = 0.
\]
Therefore, $f(x) = a_0$ is a constant polynomial. The constant term relation gives:
\[
\operatorname{const}(fg) = a_0 b_0 = 1 \implies a_0 \in \operatorname{Unit}(R).
\]
This proves that every unit in $R[x]$ is a constant polynomial coming from $\operatorname{Unit}(R)$.
\end{proof}

\begin{steps}[Degree Formula in Integral Domains]
The proof relies on the fundamental degree formula:
\[
\deg(f \cdot g) = \deg(f) + \deg(g),
\]
which holds for all non-zero polynomials $f, g \in R[x]$ if and only if $R$ has no zero divisors. If $R$ possessed non-trivial zero divisors, the leading coefficient $a_n b_m$ could vanish, causing $\deg(f \cdot g) < \deg(f) + \deg(g)$.
\end{steps}
\end{content}

\begin{speech}[00:23:23 - 00:24:13]
This means that if you ever take a polynomial which has a positive-degree part --- meaning it's not the constant polynomial, meaning one of these coefficients with $x$'s is not zero ---, you will never be able to find an inverse. Because when you multiply it by somebody else, you'll always have a term in $n+m$, and if one of them is at least one, this will never be zero. This will never be a constant term.

Does that make sense? In other words, if your ring is an integral domain, when you multiply two polynomials with $x$'s, the $x$'s will never go away, because you can always look at the highest one, and that will persist.

So, the answer is: if it's an integral domain, you'll have no more units.
\end{speech}

\begin{speech}[00:24:13 - 00:25:00]
Of course, this begs the question: what happens when it's not?

\inlinemetanote{The lecturer writes on the lower right board: \qt{What happens when $R$ is not an integral domain?}}

I leave this to you to think about. That's kind of a fun thing. Well, maybe we'll discuss it later.

You should either prove that this argument somehow can be fixed, or you should find a ring that's not an integral domain where there are more units. Those are the two possible solutions.
\end{speech}

\begin{content}[The Non-Integral Domain Question]
\begin{center}
\textbf{Question posed by the lecturer:}
\emph{What happens when $R$ is not an integral domain?}
\end{center}
\end{content}

\begin{aiNote}[{Beyond the Lecture: Non-Constant Units in $R[x]$}]
% [PERSISTENT, ADDITION]
When $R$ is not an integral domain, $R[x]$ can have non-constant units. In general, a polynomial $f(x) = a_0 + a_1 x + \dots + a_n x^n \in R[x]$ is a unit if and only if $a_0 \in \operatorname{Unit}(R)$ and all higher coefficients $a_1, \dots, a_n \in R$ are nilpotent. 

For a concrete counterexample, consider the ring $R = \mathbb{Z}/4\mathbb{Z}$ (which is not an integral domain since $2 \cdot 2 = 0$). In $(\mathbb{Z}/4\mathbb{Z})[x]$, consider the linear polynomial $f(x) = 1 + 2x$. We compute:
\[
(1 + 2x)^2 = 1 + 4x + 4x^2 \equiv 1 \pmod 4.
\]
Thus $f(x) = 1 + 2x$ is a unit of degree $1$ with $f(x)^{-1} = 1 + 2x$.
\end{aiNote}

\begin{speech}[00:25:01 - 00:25:19]
My memory from school --- which was really a long time ago --- was that there were a lot of polynomials in school, weren't there?

You should have some experience adding, subtracting, and multiplying polynomials, and as I said, all that experience is directly applicable here.
\end{speech}

\begin{hidden}
% timestamp: 00:25:20
% topic: Units of polynomial rings R[x] over integral domains
% board_state: thm:units-polynomial-ring-integral-domain, eq:units-rx-equals-units-r
% next_goal: Formal definition and operations of formal power series rings R[[x]]
% open_loops: student reflection on units in R[x] when R is not an integral domain
\end{hidden}

\section{Formal Power Series Rings}

\subsection{Definition and Ring Operations in \texorpdfstring{$R[[x]]$}{R[[x]]}}

\begin{speech}[00:25:20 - 00:27:01]
\inlinemetanote{The lecturer pushes the lower board up and starts writing on the newly revealed lower board.}

Okay, so the next thing is power series. And you may have less experience manipulating power series, I don't know. But it's a beautiful thing.

Here, I have to define for you the power series ring. And the reason there are two brackets here somehow is because there are a lot more $x$'s, you know?

An element of the power series ring... As with polynomial rings, there are two ways to think about it: the formal way, or the way people actually think about it.

In the formal way, before, a polynomial was given by a sequence of elements of $R$, and that sequence had a highest term and an infinite number of zeros afterwards. That was the definition of a polynomial, if you remember from yesterday (Recall: \cref{def:poly-rigorous-sequences}). You could add these, subtract these, multiply them, and it would just be the same thing as the $x$'s; that's the abstract discussion.

From this point of view, power series are easy: you just don't put this condition of zeros. It's just infinite elements.

From this abstract point of view, elements of the power series ring are infinite lists --- infinite vectors of elements of $R$. Before, we had lists that eventually become all zero, but now we just lift that condition. You can do anything you want.
\end{speech}

\begin{content}[Formal Definition of Power Series as Sequences]
\begin{definition}[Formal Power Series Ring]\label[definition]{def:power-series-ring-formal}
Let $R$ be a commutative ring with identity $1$. 
\begin{enumerate}[label=\textbf{(\arabic*)}]
    \setcounter{enumi}{0} \item \textbf{Sequence Perspective:} Formally, a \newterm{formal power series} with coefficients in $R$ is an infinite sequence of elements of $R$:
    \begin{equation}\label{eq:power-series-sequence}
    (a_0, a_1, a_2, \dots, a_n, \dots) \in R^{\mathbb{N}}.
    \end{equation}
    In contrast to polynomials in $R[x]$ (which require $a_k = 0$ for all sufficiently large $k$), no vanishing condition is imposed on the terms of a formal power series.
    \setcounter{enumi}{1} \item \textbf{Notation:} The set of all formal power series in the indeterminate $x$ with coefficients in $R$ is denoted by:
    \begin{equation}\label{eq:power-series-ring-symbol}
    R[[x]].
    \end{equation}
\end{enumerate}
\end{definition}
\end{content}

\begin{speech}[00:27:02 - 00:27:57]
In the way to think about it as power series, all this means is that before, for polynomials, we had the condition that eventually you stop --- that's a polynomial ---; and for power series, you just don't stop. You can go forever.

It's hard to go forever, so typically you would write the following: the elements of the power series ring are the sum from $i = 0$ to infinity of $a_i x^i$, where the $a_i$'s are in $R$.

You can see this is exactly the same data as that. The $x$, as I said, is just the variable. It's not influencing the $a$'s. Like here, it just makes clear that $x$ is just a placeholder, right? And that's all it's doing here also.

So, that's a power series. And you must have had a whole class on power series.
\end{speech}

\begin{content}[Representation as Formal Sums]
Every formal power series $(a_n)_{n \in \mathbb{N}} \in R[[x]]$ is represented as a formal infinite sum:
\begin{equation}\label{eq:power-series-sum-notation}
a_0 + a_1 x + a_2 x^2 + \dots = \sum_{i=0}^\infty a_i x^i, \quad a_i \in R.
\end{equation}

\begin{steps}[Algebraic Meaning of the Infinite Sum]
In algebra, the expression $\sum_{i=0}^\infty a_i x^i$ is purely a formal symbol representing the sequence of coefficients $(a_0, a_1, a_2, \dots)$. There are no questions of convergence, metrics, or limits involved, and the symbol $x$ serves solely as an algebraic placeholder tracking positions.
\end{steps}
\end{content}

\begin{speech}[00:27:58 - 00:29:02]
Which is good, because I'm not going to teach you about them in that way.

Do you know how to add power series, multiply power series?

The zero is just the zero power series. That's where every coefficient forever is zero. The one power series is just a one and then everybody else is zero.

\inlinemetanote{The lecturer writes $(R[[x]], +, \cdot, 0, 1)$ on the board, followed by $(0, 0, 0, \dots) = 0$ and $(1, 0, 0, \dots) = 1$.}

You know how to add power series, and I say that you do know how to add power series. In every analysis class, I think you add power series.

Here's the rule, which will surprise no one:

\inlinemetanote{The lecturer writes the addition formula on the board.}

It's a lot of work to write all of these things. That's the exciting rule: it's componentwise addition.
\end{speech}

\begin{content}[{Ring Structure and Addition in \texorpdfstring{$R[[x]]$}{R[[x]]}}]
The formal power series ring is equipped with a ring structure $(R[[x]], +, \cdot, 0, 1)$:
\begin{itemize}
    \item \textbf{Additive Identity ($0$):}
    \begin{equation}\label{eq:power-series-zero}
    0_{R[[x]]} := (0, 0, 0, \dots) = \sum_{i=0}^\infty 0 \cdot x^i.
    \end{equation}
    \item \textbf{Multiplicative Identity ($1$):}
    \begin{equation}\label{eq:power-series-one}
    1_{R[[x]]} := (1, 0, 0, \dots) = 1 + \sum_{i=1}^\infty 0 \cdot x^i.
    \end{equation}
    \item \textbf{Addition:} Defined componentwise:
    \begin{equation}\label{eq:power-series-addition}
    \sum_{i=0}^\infty a_i x^i + \sum_{i=0}^\infty b_i x^i \stackrel{\text{def}}{=} \sum_{i=0}^\infty (a_i + b_i) x^i.
    \end{equation}
\end{itemize}
\end{content}

\begin{speech}[00:29:03 - 00:30:12]
More complicated is: how do you multiply power series?

That's a definition, by the way. I must say that I really dislike indices in mathematics somehow.

We have to define this. How are we going to define this? You define it to be the sum from $i = 0$ to infinity of $c_i x^i$. And then I have to define what $c_i$ is.

Let's say $c_k$ is... you can write this as a sum from $i = 0$ to $k$, and I write $a_{k-i} b_i$.

\inlinemetanote{The lecturer writes $c_k = \sum_{i=0}^k a_{k-i} b_i$ on the board.}

All right, I did it. But you know this. This is how you multiply power series.
\end{speech}

\begin{content}[{Cauchy Multiplication in \texorpdfstring{$R[[x]]$}{R[[x]]}}]
\begin{definition}[Cauchy Product of Power Series]\label[definition]{def:power-series-multiplication}
Let $\sum_{i=0}^\infty a_i x^i$ and $\sum_{i=0}^\infty b_i x^i$ be formal power series in $R[[x]]$. Their product is defined by:
\begin{equation}\label{eq:power-series-mult}
\left( \sum_{i=0}^\infty a_i x^i \right) \cdot \left( \sum_{i=0}^\infty b_i x^i \right) \stackrel{\text{def}}{=} \sum_{k=0}^\infty c_k x^k,
\end{equation}
where each coefficient $c_k \in R$ is given by the Cauchy convolution formula:
\begin{equation}\label{eq:cauchy-coeff-formula}
c_k := \sum_{i=0}^k a_{k-i} b_i = \sum_{i+j=k} a_i b_j.
\end{equation}
\end{definition}

\begin{steps}[Finiteness of Each Coefficient Calculation]
Notice that to determine any individual coefficient $c_k$ of the product, only the finite collection of coefficients $\{a_0, \dots, a_k\}$ and $\{b_0, \dots, b_k\}$ is involved. The sum in \eqref{eq:cauchy-coeff-formula} has exactly $k+1$ terms in $R$, meaning that no infinite sums in the base ring $R$ are required to define the product.
\end{steps}
\end{content}

\begin{speech}[00:30:13 - 00:31:25]
That is a ring. Is it a ring? Formally, you have to check all the axioms.

The specific axiom that probably you don't see in your head immediately --- unless you've thought about it --- is you have to check the associativity. It's kind of obvious this is associative, commutative, everything for addition, because all we're doing is componentwise; really nothing's happening there.

This one may be more interesting if you haven't thought about it. This is somehow a slightly strange rule that tells you how to get the new power series from the old power series. You have to check that this is associative, etcetera.

Commutative is kind of obvious, even though I've broken the symmetry here; but if you think about it, it's kind of clear that it's commutative. But then you have to check it's associative.

Formally, you have to check it, unless somehow you already know that it's true. The formal way to check it is: you have to multiply three things in two different ways, look at the coefficients, and it's an identity. And it will work out.

You have to decide how you interact with math, but formally, if you want to say you understand what a power series ring is, you have to check that the multiplication is associative.
\end{speech}

\begin{note}[Board Transition and Cleaning]
The lecturer takes a wet sponge and squeegee and thoroughly cleans the lower blackboard to prepare for the discussion of units in $R[[x]]$, while the upper board displays the definitions and operations of $R[[x]]$.
\end{note}

\begin{speech}[00:31:29 - 00:32:34]
One of the issues about power series rings is that if you want to communicate to someone else an element of the power series ring, it's a lot harder than in the polynomial ring. In the polynomial ring, you just write it; it's a finite amount of data. You just write seven coefficients, and that's it.

In a power series ring, a priori, there's an infinite amount of data there. So, if you want to communicate it, you have to give some rule or something like that. We'll do some examples eventually, but it's somehow a more complicated ring.

When you encounter power series in analysis, they could be like the Taylor expansion of trigonometric functions, and they have rules; that's a way to explain it.

Anyway, if you have not encountered these before, you should make yourself comfortable with the idea of polynomials in $R$, power series in $R$, how to add, multiply, and how to manipulate them.
\end{speech}

\begin{hidden}
% timestamp: 00:32:34
% topic: Ring structure and operations of formal power series R[[x]]
% board_state: def:power-series-ring-formal, def:power-series-multiplication
% next_goal: Determine the units of the formal power series ring R[[x]]
% open_loops: none
\end{hidden}

\subsection{Units in the Formal Power Series Ring}

\begin{speech}[00:32:38 - 00:33:50]
Let's ask this question: what are the units of this ring?

\inlinemetanote{The lecturer writes on the clean lower board: \qt{What are unit $R[[x]]$?}}

Again, that means we start with a particular power series $f$:
\[
f = a_0 + a_1 x + a_2 x^2 + \dots
\]
and it goes forever. In principle... of course, it could also all be zero. A polynomial is a power series; it's just one that doesn't last too long. Because zero is fine.
\[
g = b_0 + b_1 x + b_2 x^2 + \dots
\]
For this to be a unit, if $f$ is a unit, then there exists $g$ such that $f \cdot g = 1$.

And to be $1$ is dramatic. It means when you multiply this, the rule in principle gives you infinitely many terms; if it's $1$, it means almost all of them are zero.
\end{speech}

\begin{content}[{Setting up the Invertibility Condition in \texorpdfstring{$R[[x]]$}{R[[x]]}}]
Let $f(x) \in R[[x]]$ be a formal power series:
\begin{equation}\label{eq:f-series-def}
f(x) = a_0 + a_1 x + a_2 x^2 + \dots = \sum_{i=0}^\infty a_i x^i.
\end{equation}
By definition, $f(x)$ is a \newterm{unit} in $R[[x]]$ if and only if there exists an inverse series $g(x) = \sum_{j=0}^\infty b_j x^j \in R[[x]]$ such that:
\begin{equation}\label{eq:fg-equals-one}
f(x) \cdot g(x) = 1_{R[[x]]} = 1 + 0 \cdot x + 0 \cdot x^2 + \dots
\end{equation}
\end{content}

\begin{speech}[00:33:51 - 00:35:06]
If we multiply them, $f$ times $g$, we get the first term: $a_0 b_0$.

The second term is $a_1 b_0 + a_0 b_1$, times $x$.

\inlinemetanote{The lecturer writes $fg = a_0 b_0 + (a_1 b_0 + a_0 b_1) x + (a_2 b_0 + a_1 b_1 + a_0 b_2) x^2 + \dots$ on the board.}

You can tell I've been spending too much time in Europe, because I say \qt{nought} all the time, right?

Here: $a_2$... That was a complicated joke.

That's how the power series multiplication works, right?

How is it possible that this could be $1$? Could this be $1$?

$1$, I remind you, is $1 + 0 \cdot x + 0 \cdot x^2$, and so on.

The only way it can be $1$ is if this first term is $1$. This would imply that $a_0 b_0 = 1$.

So, this means that the constant term of $f$ is a unit, right?
\end{speech}

\begin{content}[The Constant Term Invertibility Condition]
Multiplying $f(x)$ and $g(x)$ via Cauchy multiplication:
\begin{align}
    f(x) \cdot g(x) &= a_0 b_0 + (a_1 b_0 + a_0 b_1) x + (a_2 b_0 + a_1 b_1 + a_0 b_2) x^2 + \dots \nonumber \\
    &\stackrel{!}{=} 1 + 0 \cdot x + 0 \cdot x^2 + \dots \label{eq:fg-expansion-matched}
\end{align}
Equating coefficients power by power, the constant term yields:
\begin{equation}\label{eq:const-coeff-relation}
a_0 b_0 = 1 \implies a_0 = \operatorname{const}(f) \in \operatorname{Unit}(R).
\end{equation}
Thus, a necessary condition for $f(x)$ to be invertible in $R[[x]]$ is that its constant term $a_0$ is a unit in the base ring $R$.
\end{content}

\begin{speech}[00:35:07 - 00:36:18]
What about all this other stuff? What's your feeling about that?

Of course, it's also true that the constants... if we take the units of $R$, they sit inside the units of the power series by just sending the unit $a$ to the constant power series $a$.

\inlinemetanote{The lecturer writes $\operatorname{Unit}(R) \subset \operatorname{Unit}(R[[x]])$, $a \mapsto a$ on the board.}

That's also true, right?

So, the same question as above: is this an equality?

And now we come to something interesting. What was the argument before (i.e., the proof of \cref{thm:units-polynomial-ring-integral-domain})? The argument before was that --- at least when $R$ was an integral domain --- you go to the two polynomials, take their highest piece, and multiply them together.

That was a great argument for polynomials; it absolutely fails for power series, because power series have no highest piece. So, this is problematic from the point of view of using that argument.
\end{speech}

\begin{speech}[00:36:18 - 00:36:51]
Since we don't have so much time, maybe I'll tell you the answer, which I'm always reluctant to do. It is part of my job actually, but I'm still reluctant.

Something that's kind of surprising here is:

\inlinemetanote{The lecturer writes on the board: \qt{$f$ is a unit in $R[[x]]$ iff $\operatorname{const}(f)$ is a unit in $R$.}}

$f$ is a unit in $R[[x]]$ if and only if the constant term is a unit in $R$.

All I have to do to check whether it's a unit is check whether the constant term is a unit.
\end{speech}

\begin{content}[{Characterization of Units in \texorpdfstring{$R[[x]]$}{R[[x]]}}]
\begin{theorem}[Units in the Formal Power Series Ring]\label[theorem]{thm:units-power-series-ring}
Let $R$ be an arbitrary commutative ring with identity $1$. A formal power series $f(x) = \sum_{i=0}^\infty a_i x^i \in R[[x]]$ is a unit if and only if its constant term $a_0 = \operatorname{const}(f)$ is a unit in $R$:
\begin{equation}\label{eq:units-power-series-characterization}
f(x) \in \operatorname{Unit}(R[[x]]) \iff a_0 \in \operatorname{Unit}(R).
\end{equation}
\end{theorem}

\begin{steps}[Striking Contrast with Polynomial Rings]
Recall from \cref{thm:units-polynomial-ring-integral-domain} that in $R[x]$ (over an integral domain), a polynomial is a unit if and only if it is a constant unit: no positive-degree terms may appear. 

In radical contrast, in $R[[x]]$, a power series can have completely arbitrary non-zero coefficients $a_1, a_2, a_3, \dots$ for all higher powers of $x$ without restriction; as long as the single scalar $a_0$ is invertible in $R$, the entire infinite series $f(x)$ possesses an exact multiplicative inverse in $R[[x]]$!
\end{steps}
\end{content}

\begin{speech}[00:36:51 - 00:37:44]
That's a weird thing, isn't it? But it shouldn't surprise you if you played with polynomials.

I wanted to give two proofs. We will see. You should always be suspicious when a mathematician says that they will give you two proofs: it usually means both will be sloppy.

By the way, since people didn't seem to recognize this Euler note (Recall: the old 10-franc note from the previous lecture), I brought some copies, so you can see if you're interested: the Swiss franc Euler note from the previous century.
\inlinemetanote{The lecturer holds up and shows a historical Swiss 10-franc banknote featuring Leonhard Euler.}

When I came to Switzerland, I went to a used coin shop and bought a booklet of them.
\end{speech}

\begin{hidden}
% timestamp: 00:37:45
% topic: Proving the unit criterion for formal power series R[[x]]
% board_state: thm:units-power-series-ring, eq:units-power-series-characterization
% next_goal: Construct the inverse series via reduction to 1 and geometric series expansion
% open_loops: none
\end{hidden}

\begin{proofWrapper}[Proof of \cref{thm:units-power-series-ring}]

\begin{speech}[00:37:45 - 00:39:06]
Okay, so I will try to prove this. I want to prove this statement at the end.

Suppose $f$ has the form $a_0 + a_1 x + a_2 x^2 + \dots$, where $a_0 \in R$ is a unit.

\inlinemetanote{The lecturer writes: \qt{Suppose $f = a_0 + a_1 x + a_2 x^2 + \dots$}, \qt{$a_0 \in R$ is a unit}.}

That means there exists some $y \in R$ such that $a_0 y = 1$ in $R$, right?

\inlinemetanote{The lecturer writes $\exists y \in R, \ a_0 y = 1 \text{ in } R$.}

The first thing I can do is just multiply $f$ by $y$.

Whenever I think about this, $y$ is now the constant power series. Are you happy with this?

And who knows what I get here, but I know one thing: it starts with $1$. Because when I multiply by a constant power series, it just multiplies all these coefficients by $y$, right?

So, whatever it is, it starts with $1$, and then it's like $(a_1 y) x$...

\inlinemetanote{The lecturer writes $fy = 1 + (a_1 y) x + (a_2 y) x^2 + \dots$ on the board.}

$y$ is an element of the ring; there's nothing to do with $x$. Then $(a_2 y) x^2$, and so on.

The first step is: I use the fact it's a unit, multiply by something in $R$, and make this constant term $1$.
\end{speech}

\begin{content}[Step 1: Normalization to Constant Term 1]
\textbf{1. Forward Implication ($\Rightarrow$):}
As established in \eqref{eq:const-coeff-relation}, if $f(x) \cdot g(x) = 1$, then $\operatorname{const}(f) \cdot \operatorname{const}(g) = a_0 b_0 = 1$, forcing $a_0 \in \operatorname{Unit}(R)$.

\textbf{2. Reverse Implication ($\Leftarrow$) --- Normalization:}
Suppose $a_0 \in \operatorname{Unit}(R)$. Then there exists an inverse element $y \in R$ such that:
\begin{equation}\label{eq:y-inverse-a0}
a_0 y = 1 \quad \text{in } R.
\end{equation}
Multiplying $f(x)$ by the constant power series $y \in R \subset R[[x]]$:
\begin{align}
    f(x) \cdot y &= (a_0 + a_1 x + a_2 x^2 + \dots) \cdot y \nonumber \\
    &= a_0 y + (a_1 y) x + (a_2 y) x^2 + \dots \nonumber \\
    &= 1 + (a_1 y) x + (a_2 y) x^2 + \dots \label{eq:normalized-series-fy}
\end{align}
Because $y \in \operatorname{Unit}(R) \subset \operatorname{Unit}(R[[x]])$, $f(x)$ is invertible if and only if the normalized series $f(x) \cdot y$ is invertible. Indeed, if $(f(x) \cdot y)^{-1}$ exists, then $f(x)^{-1} = y \cdot (f(x) \cdot y)^{-1}$.
\end{content}

\begin{speech}[00:39:07 - 00:40:42]
Then it's enough to invert this.

Inversion means that there's somebody else I multiply to get $1$. It's sufficient to find $g \in R[[x]]$ such that $g$ times this result --- $1$ plus... and now I just make these coefficients $b$'s, because I don't want to write $a_1 y$, and I don't know anything about it anyway, so I just make it $b$:
\[
1 + b_1 x + b_2 x^2 + \dots
\]
equals $1$.

\inlinemetanote{The lecturer writes: \qt{Sufficient to find $g \in R[[x]]$, $g(1 + b_1 x + b_2 x^2 + \dots) = 1$}.}

It's sufficient to find $g$ such that whenever I have something of this form --- which is what I have there ---, this is equal to $1$.

Meaning, I only have to invert such things. All I've done is use the fact it's a unit to change the first coefficient to $1$. And then it's sufficient to invert that, because once I invert that, I can just multiply it by $y$, and then we're done.

So, I have to invert this:

\inlinemetanote{The lecturer writes $(1 + b_1 x + b_2 x^2 + \dots)^{-1}$.}

You never finish writing, so maybe it's like this: I have to invert this.

Now, you know how to do this, right? You definitely know how to do this. It's somehow in your brain; you have to bring it out.
\end{speech}

\begin{content}[Step 2: Reduction to Inverting $1 + \gamma$]
Setting $b_k := a_k y \in R$ for all $k \ge 1$, it suffices to demonstrate that any power series of the form:
\begin{equation}\label{eq:one-plus-b-series}
1 + b_1 x + b_2 x^2 + b_3 x^3 + \dots
\end{equation}
possesses a multiplicative inverse in $R[[x]]$.
\end{content}

\begin{speech}[00:40:43 - 00:41:21]
You write this whole thing as $\gamma$.

$\gamma$ is equal to this power series:

\inlinemetanote{The lecturer brackets $b_1 x + b_2 x^2 + \dots$ and labels it $\gamma = b_1 x + b_2 x^2 + \dots$.}

And this power series has $x$'s in it.

So, you can say: I want to invert this: $1 + \gamma$.

\inlinemetanote{The lecturer writes $(1 + \gamma)^{-1}$.}

And if I write this as $1 / (1 + \gamma)$, surely you know how to invert this:
\[
1 - \gamma + \gamma^2 - \gamma^3 + \dots
\]

\inlinemetanote{The lecturer writes $= \frac{1}{1+\gamma} = 1 - \gamma + \gamma^2 - \gamma^3 + \dots$.}

That's the answer! This actually proves it's invertible.
\end{speech}

\begin{content}[Step 3: The Formal Geometric Series]
Define the auxiliary power series $\gamma(x) \in R[[x]]$ consisting of all higher-degree terms:
\begin{equation}\label{eq:gamma-def}
\gamma := b_1 x + b_2 x^2 + b_3 x^3 + \dots = \sum_{k=1}^\infty b_k x^k.
\end{equation}
The expression to invert is simply $1 + \gamma$. Inspired by the classical geometric series $\frac{1}{1 - z} = \sum_{k=0}^\infty z^k$, we substitute $z = -\gamma$ to propose the candidate inverse formula:
\begin{equation}\label{eq:geometric-series-gamma}
(1 + \gamma)^{-1} \stackrel{\text{formal}}{=} \frac{1}{1 + \gamma} = 1 - \gamma + \gamma^2 - \gamma^3 + \gamma^4 - \dots = \sum_{k=0}^\infty (-1)^k \gamma^k.
\end{equation}
\end{content}

\begin{speech}[00:41:22 - 00:42:16]
Why does it prove it? You know this; this is some identity from your previous life.

I just write it, and I claim this is the formula for the inverse.

The reason is that this thing makes sense. It looks like an infinite sum, but this $\gamma$ has $x$'s in it. $\gamma$ is divisible by $x$: $x$ divides $\gamma$. That means $x^2$ divides $\gamma^2$, and $x^k$ divides $\gamma^k$.

\inlinemetanote{The lecturer writes $x \mid \gamma \implies x^2 \mid \gamma^2 \dots x^k \mid \gamma^k$.}

When I expand this, if I want the coefficient in the power series of any particular power of $x$, only finitely many of these terms enter.

The whole point of this argument is that this thing is a well-defined element of the power series ring.

\inlinemetanote{The lecturer writes \qt{well defined elem of $R[[x]]$} on the board.}
\end{speech}

\begin{content}[Step 4: Well-Definedness via Divisibility by Powers of $x$]
\textbf{Why the Infinite Sum of Power Series Is Algebraically Sound:}
In an arbitrary ring, infinite sums are not defined. However, because $\gamma$ has no constant term (its lowest degree term is $b_1 x$), $x$ divides $\gamma$:
\begin{equation}\label{eq:x-divisibility-chain}
x \mid \gamma \implies x^2 \mid \gamma^2 \implies \dots \implies x^k \mid \gamma^k.
\end{equation}
Consequently, the lowest power of $x$ appearing with non-zero coefficient in $\gamma^k$ is at least $x^k$:
\[
\gamma^k = (b_1 x + b_2 x^2 + \dots)^k = b_1^k x^k + (\text{terms of degree } \ge k+1).
\]
Now consider the coefficient of $x^N$ in the formal sum $\sum_{k=0}^\infty (-1)^k \gamma^k$:
\begin{itemize}
    \item For every $k > N$, every term in $\gamma^k$ has degree at least $k > N$.
    \item Therefore, terms with $k > N$ contribute strictly $0$ to the coefficient of $x^N$!
\end{itemize}
Hence, for any fixed power $x^N$, only the finitely many terms $k \in \{0, 1, \dots, N\}$ contribute to its coefficient:
\begin{equation}\label{eq:finite-coeff-contribution}
[x^N]\left( \sum_{k=0}^\infty (-1)^k \gamma^k \right) = [x^N]\left( \sum_{k=0}^N (-1)^k \gamma^k \right) \in R.
\end{equation}
Because every individual coefficient of the resulting power series is computed as a finite sum of elements in $R$, the series:
\[
g(x) := \sum_{k=0}^\infty (-1)^k \gamma(x)^k
\]
is a rigorous, well-defined element of $R[[x]]$. Multiplying by $(1 + \gamma)$ explicitly cancels all terms by telescoping:
\[
(1 + \gamma) \cdot \sum_{k=0}^\infty (-1)^k \gamma^k = \sum_{k=0}^\infty (-1)^k \gamma^k + \sum_{k=0}^\infty (-1)^k \gamma^{k+1} = 1.
\]
This manipulation of infinite sums is legitimate because it can be checked coefficient by coefficient: for the coefficient of $x^N$, only the finitely many terms with $k \le N$ contribute, and for these the identity is the finite telescoping sum $(1 + \gamma) \sum_{k=0}^{N} (-1)^k \gamma^k = 1 + (-1)^N \gamma^{N+1}$, whose last term has no $x^N$ component.
Thus, $(1 + \gamma)$ is invertible, completing the proof that $f(x)$ is a unit in $R[[x]]$.
\end{content}

\end{proofWrapper}

\begin{speech}[00:42:17 - 00:42:50]
That is what you have to convince yourself of: it's just that it makes sense to write this. I gave you the argument, but that's different than you accepting the argument.

It just says that if you do this, every one of these terms is well-defined. You're not allowed to add infinitely many things in algebra. But since the powers of $x$ get bigger and bigger, you don't have to do it.

Maybe we'll talk a little bit more about that next time.
\inlinemetanote{lecturer takes off microphone, class applauds, video ends}
\end{speech}

\begin{overview}[Summary of Key Concepts]
\begin{itemize}
    \item \textbf{Field of Fractions:} For an integral domain $R$, $\operatorname{Frac}(R)$ is formally the set of equivalence classes of pairs $(r, s) \in R \times (R \setminus \{0\})$ with $(r, s) \equiv (t, w) \iff rw = st$. In practice one computes with fractions $\frac{r}{s}$ by the usual rules; the integral domain property guarantees that the denominators $sw$ stay non-zero, and every $\frac{r}{s} \neq 0$ has the inverse $\frac{s}{r}$.
    \item \textbf{Polynomial Rings:} $R[x]$ exists for every commutative ring $R$ with unit. Its universal property: for every ring homomorphism $\phi \colon R \to S$ and every $y \in S$, there is a unique ring homomorphism $\hat{\phi} \colon R[x] \to S$ with $\hat{\phi}|_R = \phi$ and $\hat{\phi}(x) = y$, namely $\hat{\phi}\bigl(\sum a_i x^i\bigr) = \sum \phi(a_i) y^i$.
    \item \textbf{Units of $R[x]$:} If $R$ is an integral domain, then $\operatorname{Unit}(R[x]) = \operatorname{Unit}(R)$, because the leading coefficient of a product is the product of the leading coefficients. Without this assumption there can be more units, e.g., $1 + 2x$ in $(\mathbb{Z}/4\mathbb{Z})[x]$.
    \item \textbf{Formal Power Series:} $R[[x]]$ consists of all sequences $(a_0, a_1, a_2, \dots)$ of elements of $R$, written $\sum_{i \ge 0} a_i x^i$, with componentwise addition and the Cauchy product $c_k = \sum_{i=0}^k a_{k-i} b_i$; every coefficient of a product is a finite sum.
    \item \textbf{Units of $R[[x]]$:} For every commutative ring $R$ with unit, $f \in R[[x]]$ is a unit if and only if its constant term is a unit in $R$. After normalizing the constant term to $1$, the inverse of $1 + \gamma$ is the formal geometric series $\sum_{k \ge 0} (-1)^k \gamma^k$, which is well-defined because $x^k \mid \gamma^k$.
\end{itemize}
\end{overview}

% week 4:
% ==========================================
% LatexRefinement Step Output: step4-2026-10-06-tuesday-week4-algebra-1-offset-final.tex
% Model: gemini-3.8-flash
% Temperature: 0.35
% TopP: 0.9
% TopK: 10
% MaxOutputTokens: 65535
% ThinkingBudget: 4096
% ThinkingLevel: HIGH
% Prompt Tokens: 79’172
% Candidates Tokens: 38’383 (inkl. Thinking Tokens)
% Cached Tokens: 28’918
% Processed on: 2026-10-09 00:30:29
% ==========================================

% Primary Language: English
% End of the video: 01:30:14

\lecturechapter[4]{Tuesday}{06 Oct}{6 October 2026}{Polynomial Rings, Laurent Series, Quotient Rings, and the First Isomorphism Theorem}
% Old chapter title: Polynomial Rings, Laurent Series, and Ideals

\begin{overview}[Course Content Overview]
This lecture explores fundamental algebraic ring constructions, bridging polynomial extensions, localized series, and residue structures:
\begin{itemize}
    \item \textbf{Polynomial and Power Series Rings:} Multivariable polynomial rings, the binomial theorem in commutative rings, and the Frobenius identity in characteristic $p$.
    \item \textbf{Fields of Fractions and Laurent Series:} The field of rational functions $R(x)$ and formal Laurent series $R((x)) \cong \operatorname{Frac}(R[[x]])$.
    \item \textbf{Ideals and Quotient Rings:} Formal definition of ideals, the linear algebra analogy, and the construction of the quotient ring $R/I$ from scratch (highlighting how ideal absorption ensures well-defined coset multiplication).
    \item \textbf{The First Isomorphism Theorem for Rings:} Formulation and complete three-step proof of $R/\ker(f) \cong \operatorname{Im}(f)$.
    \item \textbf{Ideals Generated by Elements and Polynomial Quotients:} Contrasting the quotients $\mathbb{C}[x]/(x^2+1)$ (which possesses zero divisors) and $\mathbb{Q}[x]/(x^2+1) \cong \mathbb{Q}(i)$ (which forms the field of Gaussian rationals).
\end{itemize}
\end{overview}

\section{Review: Polynomial Rings and Power Series}

\begin{speech}[00:00:00 - 00:01:54]
Okay, so we start, and I wanted to start with some things from last week. If $R$ is a ring, we had defined this polynomial ring in one variable and also the power series ring. \inlinemetanote{writes $R \to R[x], R[[x]]$} And just some comments: the polynomial ring is a subring of the power series ring. \inlinemetanote{writes $R[x] \subset R[[x]]$} That's important to note, because every polynomial you can view as a finite power series.

So, this is one comment. The second comment --- and people were asking me after class --- is: what happens if I have two variables or three variables? There's two ways to approach this. One is: you could just define this... \inlinemetanote{writes $\to R[x,y] \to \text{defined}$} defined in the same way as we defined this, which is that it's the set of finite sums \inlinemetanote{writes $\{ \sum_{\text{finite}} a_{i,j} x^i y^j \mid a_{i,j} \in R \}$} $\sum a_{i,j} x^i y^j$, where these indices go over some finite set, and the coefficients are in $R$; and it's just the same definition as before. So, that's one way to do it. I hope that's clear --- it's just multiplying with more polynomials.

But there's a more formal way, which is that $R$ is a ring, so we've already defined this --- that's where we add one variable ---, and so, this is a ring, so we can also define this: this is a ring and we have one more variable. \inlinemetanote{writes $R \to R[x] \to (R[x])[y]$} And you can take this to be the definition of $R[x,y]$. \inlinemetanote{writes $= R[x,y]$}

So, we just add the variables one at a time. This is a polynomial ring over the polynomial ring over $R$, so it's another way. And, of course, if you have two different paths, then you have the burden of proving they're both the same, etc., and I leave this burden for you to carry, all right? \inlinemetanote{draws brace and writes \qt{same}} And these are the same; maybe I just say that: same.
\end{speech}

\begin{content}[Polynomial Rings in One and Several Variables]
Let $R$ be a commutative ring with identity. Recall the construction of the \newterm{polynomial ring} $R[x]$ and the \newterm{formal power series ring} $R[[x]]$. Every polynomial can be naturally identified with a formal power series having only finitely many non-zero coefficients, yielding the subring inclusion:
\[
R[x] \subseteq R[[x]].
\]

\begin{definition}[Polynomial Ring in Two Variables]\label[definition]{def:two-variable-polynomial-ring}
For two indeterminates $x$ and $y$, the polynomial ring $R[x, y]$ can be defined in two equivalent ways:
\begin{enumerate}[label=\textbf{(\arabic*)}]
    \item \textbf{Explicit Definition (Finite Linear Combinations):}
    \[
    R[x, y] := \left\{ \sum_{i, j \ge 0}^{\text{finite}} a_{i,j} x^i y^j \;\middle|\; a_{i,j} \in R \right\}.
    \]
    \item \textbf{Iterative Definition (Polynomials over Polynomials):}
    Viewing $R[x]$ as the coefficient ring, we form the univariate polynomial ring in $y$:
    \[
    R \longrightarrow R[x] \longrightarrow (R[x])[y] =: R[x, y].
    \]
\end{enumerate}
\end{definition}

\begin{steps}[Equivalence of Definitions]
Both constructions yield canonically isomorphic rings:
\[
(R[x])[y] \cong R[x, y] \cong (R[y])[x].
\]
In the iterative perspective, an element is a polynomial $P(y) = \sum_{j=0}^m p_j(x) y^j$, where each coefficient $p_j(x) = \sum_{i=0}^{k_j} a_{i,j} x^i \in R[x]$. Expanding via the distributive law expresses $P(y)$ uniquely as the finite double sum $\sum_{i, j} a_{i,j} x^i y^j$. By induction, this extends to $n$ variables:
\[
R[x_1, \dots, x_n] \cong (R[x_1, \dots, x_{n-1}])[x_n].
\]
\end{steps}
\end{content}

\subsection{The Binomial Formula in Commutative Rings}

\begin{speech}[00:01:54 - 00:03:47]
Okay, so then there's some comments. Essentially, well, something you should know is: if I have this polynomial ring in two variables and I take $x + y$, that's in this polynomial ring in two variables. And I can raise it to the $n$, and this is something you know --- this is this binomial expansion, or I don't know how it's called, binomial formula. \inlinemetanote{writes $R[x,y]$ and $(x+y)^n = \sum_{a=0}^n x^{n-a} y^a \binom{n}{a}$} This is equal to the sum --- so, from $a = 0$ to $n$ --- of $x^{n-a} y^a \binom{n}{a}$.

Okay? That's the binomial formula. Are you familiar with this binomial formula?

So, for example, $(x+y)^5$ is equal to $1 x^5 + 5 x^4 y + 10 x^3 y^2 + 10 x^2 y^3 + 5 x y^4 + y^5$. \inlinemetanote{writes $(x+y)^5 = x^5 + 5x^4y + 10x^3y^2 + 10x^2y^3 + 5xy^4 + y^5$} I guess I should have done a smaller one.

That's an example. And this holds in every ring, because these elements $x$ and $y$ are here, and we've already discussed that these numbers make sense in any ring --- they're just how you add up $1$. Is that clear? So, this formula holds in any ring. And to prove it, it's just the same as by induction: if you wanted to prove this formula in any ring, you just write it out as $(x+y)(x+y)\dots$ and use the distributive law, and you get $x^2$, and you regroup everything, and you get $xy + xy + y^2$, \inlinemetanote{writes $(x+y)^2 = (x+y)(x+y) = x^2 + xy + xy + y^2$} and that gives you the formula; and then you prove it by induction. Okay.
\end{speech}

\begin{content}[The Binomial Expansion]
Let $R$ be any commutative ring with identity $1_R$. For integers $k \in \mathbb{Z}$, the integer scalar $k \cdot 1_R$ is canonically defined by repeated addition or negation of $1_R$. Consequently, the binomial coefficients $\binom{n}{a} \in \mathbb{N}$ act canonically as coefficients in $R$.

\begin{theorem}[Binomial Formula in Commutative Rings]\label[theorem]{thm:binomial-formula}
In the polynomial ring $R[x, y]$, or for any two commuting elements $x, y \in R$ ($x y = y x$), the following identity holds for all $n \in \mathbb{N}$:
\begin{equation}\label{eq:binomial-formula}
(x + y)^n = \sum_{a=0}^n \binom{n}{a} x^{n-a} y^a.
\end{equation}
\end{theorem}

\begin{example}[Expansion for $n=5$ and $n=2$]\label[example]{ex:binomial-5}
For $n = 5$:
\[
(x + y)^5 = x^5 + 5 x^4 y + 10 x^3 y^2 + 10 x^2 y^3 + 5 x y^4 + y^5.
\]
For $n = 2$, distributivity and commutativity give:
\[
(x + y)^2 = (x + y)(x + y) = x^2 + x y + y x + y^2 = x^2 + 2 x y + y^2.
\]
\end{example}

\begin{steps}[Proof by Induction]
The general identity follows by induction on $n$:
\begin{align*}
(x + y)^{n+1} &= (x + y) \sum_{a=0}^n \binom{n}{a} x^{n-a} y^a \\
&= \sum_{a=0}^n \binom{n}{a} x^{n+1-a} y^a + \sum_{a=0}^n \binom{n}{a} x^{n-a} y^{a+1} \\
&= x^{n+1} + \sum_{k=1}^n \left[ \binom{n}{k} + \binom{n}{k-1} \right] x^{n+1-k} y^k + y^{n+1} \\
&= \sum_{k=0}^{n+1} \binom{n+1}{k} x^{n+1-k} y^k,
\end{align*}
using Pascal's identity $\binom{n}{k} + \binom{n}{k-1} = \binom{n+1}{k}$.
\end{steps}
\end{content}

\subsection{Characteristic \texorpdfstring{$p$}{p} and the Frobenius Identity}

\begin{speech}[00:03:47 - 00:05:43]
There's one comment here, which is kind of a fun comment. That holds in any ring, that binomial formula. But what happens if we take this ring $\mathbb{Z}/p\mathbb{Z}$, where $p$ is a prime? \inlinemetanote{writes $\mathbb{Z}/p\mathbb{Z}$, $p$ prime} And then in this ring, you could take this ring in $x, y$, \inlinemetanote{writes $(\mathbb{Z}/p\mathbb{Z})[x,y]$} and the fun thing to consider there is $(x+y)^p$. \inlinemetanote{writes $(x+y)^p$}

So, that's an example. And if I take this formula in $\mathbb{Z}/p\mathbb{Z}$ --- $\mathbb{Z}/5\mathbb{Z}$ in this case, where it's the same exponent as here --- something kind of nice happens: all these internal coefficients here, they all have a $5$ in them, so they all go away. And this is true just in general: in $\mathbb{Z}/p\mathbb{Z}$, if you take $(x+y)^p$, you get this formula that you secretly wanted anyway, which is this: \inlinemetanote{writes $(x+y)^p = x^p + y^p$} It's a simple formula.

And the reason is because for all of the binomial coefficients $\binom{p}{a}$, $p$ divides $\binom{p}{a}$ when $a$ goes from $1$ to $p-1$. \inlinemetanote{writes $p \mid \binom{p}{a}$ when $a \in \{1, \dots, p-1\}$} For all the internal coefficients, $p$ divides that. And why is that? It's because this binomial starts with a $p!\dots$ \inlinemetanote{writes $\binom{p}{a} = \frac{p!}{a!(p-a)!}$} and there are no $p$'s here to cancel the $p$ in the $p!$, because these are all lesser numbers. Is that clear?

So, you'll always have, if you look at the exponent being prime in this binomial theorem, all of the middle coefficients here will all have that prime in it. So, if I look at it in $\mathbb{Z}/p\mathbb{Z}$, I get this great formula... This is a well-known formula; it will come up. Okay, so that's an aside about the binomial formula.
\end{speech}

\begin{content}[The Freshman's Dream in Characteristic \texorpdfstring{$p$}{p}]
\begin{lemma}[Divisibility of Binomial Coefficients by a Prime]\label[lemma]{lem:prime-divides-binomial}
Let $p$ be a prime number. For every integer $a \in \{1, 2, \dots, p-1\}$, the prime $p$ divides the binomial coefficient $\binom{p}{a}$:
\begin{equation}\label{eq:prime-divides-binom}
p \mathrel{\Big|} \binom{p}{a}.
\end{equation}
\end{lemma}

\begin{proof}
Recall the algebraic formula:
\[
\binom{p}{a} = \frac{p!}{a! (p-a)!} \implies a! (p-a)! \binom{p}{a} = p! = p \cdot (p-1)!.
\]
Thus, $p$ divides the product $a! (p-a)! \binom{p}{a}$. Since $1 \le a \le p-1$, every factor occurring in the factorials $a! = 1 \cdot 2 \dots a$ and $(p-a)! = 1 \cdot 2 \dots (p-a)$ is strictly less than $p$. Because $p$ is prime, $\operatorname{gcd}(a! (p-a)!, p) = 1$. By Euclid's lemma, $p$ must divide the remaining factor $\binom{p}{a}$.
\end{proof}

\begin{proposition}[Frobenius Endomorphism Identity]\label[proposition]{prop:frobenius-identity}
Let $R$ be a commutative ring of characteristic $p$ (e.g., $R = \mathbb{Z}/p\mathbb{Z}$). In the polynomial ring $R[x, y]$, or for any commuting elements $x, y \in R$, we have:
\begin{equation}\label{eq:frobenius-freshmans-dream}
(x + y)^p = x^p + y^p.
\end{equation}
\end{proposition}

\begin{proof}
Expanding $(x + y)^p$ via \cref{thm:binomial-formula}:
\[
(x + y)^p = x^p + \sum_{a=1}^{p-1} \binom{p}{a} x^{p-a} y^a + y^p.
\]
By \cref{lem:prime-divides-binomial}, $\binom{p}{a} \equiv 0 \pmod p$ for all $1 \le a \le p-1$. In characteristic $p$, each intermediate term vanishes, leaving $(x + y)^p = x^p + y^p$.
\end{proof}
\end{content}

\begin{hidden}
% timestamp: 00:05:43
% topic: Review of multivariable polynomials, binomial theorem, and Frobenius identity in characteristic p.
% board_state: def:two-variable-polynomial-ring, thm:binomial-formula, lem:prime-divides-binomial, prop:frobenius-identity
% next_goal: Discuss integral domains for R[x] and R[[x]], and fields of fractions.
% open_loops: none
\end{hidden}

\section{Integral Domains, Fields of Fractions, and Laurent Series}

\subsection{Polynomial Rings and Power Series Rings as Integral Domains}

\begin{speech}[00:05:43 - 00:06:38]
And then there were some definitions: that if I take $R$ now to be an integral domain, \inlinemetanote{writes $R$ is domain} we had defined here $R[x]$ --- that's polynomials in one variable ---; this is also an integral domain. \inlinemetanote{writes $\implies R[x]$ also int.\ domain} We discussed that: you use the highest term, and if $R$ is an integral domain, when you multiply polynomials, the highest term lives, so it stays an integral domain.

What I didn't say is: this is also an integral domain. \inlinemetanote{writes $R[[x]]$ also int.\ domain} And one needs a different argument here, because power series don't have highest terms. And so, the answer is... well, what do you think the answer is? You can't use the highest term here, so what should you use?
\end{speech}

\begin{student}[Student Answer]
Lowest term.
\end{student}

\begin{speech}[00:06:40 - 00:07:07]
Lowest term, exactly! You could have used the lowest term here, but psychologically, with polynomials, you go to the highest term. But it turns out lowest term is a better idea here, because a power series does have a lowest term. And when you multiply the two lowest terms of power series, it's the lowest term of the product. So, it's the same argument here: it's just using the lowest term. Which is, you know, psychologically a little strange --- you could have used the lowest term here, but you didn't, right? Okay, so this is also an integral domain.
\end{speech}

\begin{content}[Integral Domain Properties: Degrees vs. Orders]
\begin{theorem}[Preservation of the Integral Domain Property]\label[theorem]{thm:domain-poly-power-series}
If $R$ is an integral domain, then:
\begin{enumerate}[label=\textbf{(\arabic*)}]
    \item The polynomial ring $R[x]$ is an integral domain.
    \item The formal power series ring $R[[x]]$ is an integral domain.
\end{enumerate}
\end{theorem}

\begin{proof}
Let $f, g$ be non-zero elements in the respective ring.
\begin{enumerate}[label=\textbf{(\arabic*)}]
    \item \textbf{For $R[x]$ (Highest Degree Term):}
    Let $f = a_n x^n + \dots + a_0$ and $g = b_m x^m + \dots + b_0$ with $a_n \neq 0$ and $b_m \neq 0$. The leading term of the product $f g$ is:
    \[
    a_n b_m x^{n+m}.
    \]
    Since $R$ is an integral domain, $a_n b_m \neq 0$. Thus, $f g \neq 0$ and $\deg(f g) = \deg(f) + \deg(g)$.
    \item \textbf{For $R[[x]]$ (Lowest Degree / Order Term):}
    A formal power series does not have a maximal degree, but every non-zero series $f = \sum_{k=0}^\infty a_k x^k$ has a well-defined \newterm{order} (minimal non-zero index):
    \[
    \operatorname{ord}(f) := \min \{ k \in \mathbb{N} \mid a_k \neq 0 \}.
    \]
    We write $f = x^n (a_n + a_{n+1} x + \dots)$ with $a_n \neq 0$, and $g = x^m (b_m + b_{m+1} x + \dots)$ with $b_m \neq 0$. The lowest non-vanishing term of the Cauchy product is:
    \[
    a_n b_m x^{n+m}.
    \]
    Since $R$ is an integral domain, $a_n b_m \neq 0$. Consequently, $f g \neq 0$ with $\operatorname{ord}(f g) = \operatorname{ord}(f) + \operatorname{ord}(g)$.
\end{enumerate}
\end{proof}
\end{content}

\subsection{Fields of Fractions: Rational Functions and Laurent Series}

\begin{speech}[00:07:07 - 00:08:35]
And the point I wanted to make here is that the main construction last week was the field of fractions. \inlinemetanote{draws arrow and writes \qt{field of fractions}} We'll have a different construction this week, but that was the field of fractions; and the important thing about a field of fractions is: you can only build that on an integral domain.

So, if I have an integral domain, I could take the field of fractions of that integral domain --- that was the construction. But I could also take the field of fractions of polynomials, because that's an integral domain; and that is called... well, the symbol for that is $R$ with round parentheses --- that's round brackets ---; that's the name of the field of fractions of the polynomial ring when $R$ is an integral domain. And it's called the field of rational functions. \inlinemetanote{writes \qt{field of rational functions}}

So, rational functions with coefficients in $R$. This is polynomials with coefficients in $R$; that's the field of rational functions. And you know how to think about this because we've thought about fields of fractions: an element here in rational functions is a polynomial divided by another polynomial. \inlinemetanote{writes $R(x) \ni \frac{f(x)}{g(x)}$} So, both of these, $f$ and $g$, are polynomials, \inlinemetanote{writes $f, g \in R[x]$} and, of course, $g$ is not allowed to be zero. \inlinemetanote{writes $g \neq 0$} So, that's how to think about an element here, and that's...
\end{speech}

\begin{content}[The Field of Rational Functions]
Recall that for any integral domain $D$, its \newterm{field of fractions} $\operatorname{Frac}(D)$ consists of equivalence classes of formal quotients $\frac{a}{b}$ with $a, b \in D$ and $b \neq 0$.

\begin{definition}[Field of Rational Functions $R(x)$]\label[definition]{def:rational-functions}
Let $R$ be an integral domain. The field of fractions of the polynomial ring $R[x]$ is denoted by round parentheses $R(x)$ and is called the \newterm{field of rational functions}:
\begin{equation}\label{eq:field-rational-functions}
R(x) := \operatorname{Frac}(R[x]) = \left\{ \frac{f(x)}{g(x)} \;\middle|\; f(x), g(x) \in R[x], \ g(x) \neq 0 \right\}.
\end{equation}
\end{definition}
\end{content}

\begin{speech}[00:08:35 - 00:10:20]
Now, power series are also integral domains, so we can again consider the field of fractions, and the symbol for that is two parentheses. \inlinemetanote{writes $R((x))$} All right, so this is kind of universal math notation: you have an integral domain, you have the polynomial ring, power series, field of fractions, and this one has a special name. Does anyone know what the name of this thing is? It's called the field of Laurent series. \inlinemetanote{writes \qt{Field of Laurent Series}} That's the name of this.

Actually, I'm having some doubts about how to spell Laurent. Is it with an \qt{a}? No, I think it's with an \qt{e}. I'm not $100\%$ sure about that.

So, this is the field of Laurent series, and this will come up in your complex analysis classes for sure, but with a $\mathbb{C}$ here. And usually the argument is called $z$ in complex analysis.

So, the field of Laurent series... you could say: we've already constructed the field of fractions, so why do I have to talk about it here? You could say this field of Laurent series we think about in the same way, as $f(x)$ over $g(x)$, where $f$ and $g$ are power series. \inlinemetanote{writes $R((x)) \ni \frac{f(x)}{g(x)}$} And this would be absolutely correct. \inlinemetanote{writes $f, g \in R[[x]], g \neq 0$} That would be correct, with $g$ not equal to zero.

It's just there's a nicer way to think about it, and that has to do with the last thing I said --- not yesterday, last week. I think it's with \qt{e}.
\end{speech}

\begin{content}[The Field of Fractions of Formal Power Series]
\begin{definition}[Field of Fractions $R((x))$]\label[definition]{def:frac-power-series}
Let $R$ be an integral domain. Since $R[[x]]$ is an integral domain (\cref{thm:domain-poly-power-series}), we define its field of fractions, denoted with double round parentheses $R((x))$:
\begin{equation}\label{eq:frac-power-series}
R((x)) := \operatorname{Frac}(R[[x]]) = \left\{ \frac{f(x)}{g(x)} \;\middle|\; f(x), g(x) \in R[[x]], \ g(x) \neq 0 \right\}.
\end{equation}
This field is known algebraically as the \newterm{field of formal Laurent series}.
\end{definition}
\end{content}

\begin{note}[Board Cleaning Transition]
\inlinemetanote{The lecturer pauses to erase the middle and right boards, asking the class if they have already encountered Laurent series in complex analysis.}
\end{note}

\subsection{Formal Laurent Series with Finitely Many Negative Powers}

\begin{speech}[00:11:02 - 00:13:20]
You won't go wrong if you think about it this way: that's the definition, that's how we constructed it. But there's another way, which is why it's called the field of Laurent series. So, what is a Laurent series? \inlinemetanote{writes \qt{Def: Laurent series in $x$}}

A Laurent series is a power series where you are allowed to have finitely many negative coefficients --- that's the definition. So, I can say this: a Laurent series in a variable $x$ is a power series with --- and this is the crucial thing --- finitely many negative powers. \inlinemetanote{writes \qt{is a power series with finitely many negative powers}} You can have as many positive-degree coefficients as you want, it's a power series, but only finitely many negative powers.

So, that's the crucial thing for a Laurent series. What is a power series? A power series might be like the power series for $e^x$; we've seen these anyway: \inlinemetanote{writes $1 + x + \frac{x^2}{2!} + \frac{x^3}{3!} + \dots \in R[[x]]$} That's a power series. These are elements of the power series ring.

In a Laurent series, you're allowed to have some negative powers, so you could have $x^{-12} + 3x^{-2} + x^{-1}$, then the constant term... \inlinemetanote{writes $x^{-12} + 3x^{-2} + x^{-1} + 17 - x + x^{15} + \dots$} I got too excited about the constant term: $17 - x + x^{15} + \dots$ That is a Laurent series.

Have you seen these objects in mathematics? Okay, well, maybe it's the first time, but you can imagine: you just take power series, and you're allowed to have negative terms, but only finitely many. And it's extremely important that you have only finitely many. You'll find that if you have only finitely many, you can add them, and when you add them you'll still only have finitely many, because each of the two only had finitely many. When you multiply them, the lowest term now... if I multiply this with something with lowest term $-7$, its lowest term will be $x^{-19}$, so it will still have only finitely many. Is that clear?
\end{speech}

\begin{content}[Formal Definition and Arithmetic of Laurent Series]
\begin{definition}[Formal Laurent Series]\label[definition]{def:laurent-series}
A \newterm{formal Laurent series} in the indeterminate $x$ over a ring $R$ is a formal sum of the form:
\begin{equation}\label{eq:laurent-series-def}
f(x) = \sum_{n = -N}^\infty a_n x^n = a_{-N} x^{-N} + a_{-N+1} x^{-N+1} + \dots + a_0 + a_1 x + a_2 x^2 + \dots
\end{equation}
where $N \in \mathbb{Z}$ and $a_n \in R$ for all $n \ge -N$. That is, $f(x)$ contains at most \textbf{finitely many negative powers} of $x$.
\end{definition}

\begin{example}[Power Series vs. Laurent Series]\label[example]{ex:laurent-series}
\begin{itemize}
    \item \textbf{Power Series:} Elements of $R[[x]]$ contain only non-negative powers:
    \[
    1 + x + \frac{x^2}{2!} + \frac{x^3}{3!} + \dots \in R[[x]].
    \]
    \item \textbf{Laurent Series:} Elements can have negative exponents up to some lowest power:
    \[
    x^{-12} + 3 x^{-2} + x^{-1} + 17 - x + x^{15} + \dots \in R((x)).
    \]
\end{itemize}
\end{example}
\end{content}

\begin{speech}[00:13:20 - 00:15:20]
And here's the great thing: if you're non-zero here, you can divide. That's why the Laurent series turns out to be a field. And this is something you should think about --- we did discuss this last time.

So, if you have a power series... maybe I just say this, because this is maybe now departing a little bit from the course, but it is the case that if I take this field of Laurent series... \inlinemetanote{writes \qt{field of Laurent series $\cong R((x))$}} this field of Laurent series is isomorphic to $R((x))$, the field of fractions of the power series. And what you have to show there is that if you have Laurent series, you can divide; that's the interesting thing. And that has to do with the last thing I said last week, which is that it's pretty easy to divide power series: all you have to do is find the first non-zero coefficient, then use this rule I explained last week, and then you can play with this power to take it out. So, I leave this for you.

I'll say one thing: why don't we allow power series with infinitely many negative terms? Because if you try to multiply them, then you have to do infinite sums. And we don't do infinite sums in algebra; we only do finite sums. Because if you do infinite sums, you need convergence, and you can try to do that --- that's an interesting subject ---, but in this class we're not doing infinite sums. And if you have finitely many negative terms, when you multiply, you see you'll never have to do infinite sums; it's just the same.

So, anyway, that's just this comment here: this field of Laurent series, it turns out, is just a nice way of thinking about this field of fractions for power series. And you can imagine it's nicer to think about this kind of object than the ratio. And this second way of looking at things is not available for polynomials: there's only one way to think about it there, as ratios of polynomials. Okay.
\end{speech}

\begin{content}[Isomorphism between Laurent Series and the Field of Fractions]
\begin{theorem}[Algebraic Characterization of $R((x))$]\label[theorem]{thm:laurent-field-isomorphism}
When $K$ is a field, the set of formal Laurent series with finitely many negative powers forms a field under formal addition and multiplication, and is naturally isomorphic to the field of fractions of $K[[x]]$:
\begin{equation}\label{eq:laurent-isomorphism}
\operatorname{Frac}(K[[x]]) \cong K((x)) = K[[x]][x^{-1}].
\end{equation}
\end{theorem}

\begin{steps}[Why Invertibility Holds and Why Infinitely Many Negative Terms Fail]
\begin{enumerate}[label=\textbf{(\arabic*)}]
    \item \textbf{Invertibility in $K((x))$:}
    Any non-zero Laurent series can be factored as $f(x) = x^k u(x)$, where $k \in \mathbb{Z}$ is the minimal index with $a_k \neq 0$, and $u(x) = a_k + a_{k+1} x + \dots \in K[[x]]$ with $a_k \neq 0$. Since $K$ is a field, the leading coefficient $a_k$ is a unit in $K$, so $u(x)$ is an invertible unit in the power series ring $K[[x]]$. Its multiplicative inverse is:
    \[
    f(x)^{-1} = x^{-k} u(x)^{-1} \in K((x)).
    \]
    \item \textbf{Failure of Infinitely Many Negative Terms:}
    If series with infinitely many negative powers were permitted ($\sum_{n=-\infty}^\infty a_n x^n$), the coefficient of $x^k$ in the product of two such series would require evaluating the infinite sum:
    \[
    c_k = \sum_{j=-\infty}^\infty a_j b_{k-j}.
    \]
    Without topological convergence, such an infinite sum is undefined in pure algebra. Restricting to finitely many negative powers ensures that for each $k$, only finitely many non-zero products $a_j b_{k-j}$ occur, keeping all algebraic operations purely finite.
\end{enumerate}
\end{steps}
\end{content}

\begin{hidden}
% timestamp: 00:15:20
% topic: Laurent series as field of fractions R((x)); introduction of ideals and quotient rings.
% board_state: def:rational-functions, def:frac-power-series, def:laurent-series, thm:laurent-field-isomorphism
% next_goal: Formal definition of ideals and ring homomorphisms.
% open_loops: none
\end{hidden}

\section{Ideals and Ring Homomorphisms}

\subsection{Definition and Axioms of an Ideal}

\begin{speech}[00:15:33 - 00:16:29]
Okay, so the main topic this week is about ideals and quotient rings, and the main construction is that of a quotient ring.

If you like polynomial rings, the class will talk about results of polynomial rings, but mostly polynomial rings in one variable. The subject of polynomial rings in several variables is very important in mathematics, and there's a whole class about it, more or less: the class is called commutative algebra.
\end{speech}

\begin{note}[Board Cleaning and Transition]
\inlinemetanote{The lecturer finishes erasing the boards to begin the new chapter on ideals and quotient rings.}
\end{note}

\begin{speech}[00:16:30 - 00:18:33]
Okay, so now we finish these constructions and we start something new. So, $R$ is our ring, \inlinemetanote{writes $R$ ring} and the first thing, which we already defined: $I$ is an ideal. \inlinemetanote{writes $I \subset R$ is an ideal} We already defined this, but I'll just review it briefly, because now it's crucial for this week.

$I$ is an ideal if it satisfies that if I take $(I, +, 0)$, this is a subgroup of the additive part of the ring. \inlinemetanote{writes $\bullet \ (I, +, 0) \subset (R, +, 0)$ subgroup} And the second is: if I take anybody in the ring and anybody in the ideal, then their product lies in the ideal. \inlinemetanote{writes $\bullet \ \forall r \in R, \forall x \in I \text{ then } r \cdot x \in I$} This says that the ideal is an additive subgroup under addition, but there's another condition. The ideal has two constraints: one is the addition constraint, and the other is the multiplication constraint. They're just different, because they're talking about two different operations.

And our example we had seen --- we did some examples anyway ---, but a place where we know we always get ideals: if I take a homomorphism... \inlinemetanote{writes $f \colon R \to S$} Maybe I should put it here: so, $f$ is a homomorphism of rings, \inlinemetanote{writes $f$ is a hom of rings} then whenever you have a homomorphism of rings, you get an ideal for free. Sorry. \inlinemetanote{drops chalk}

And we discussed this also before: then the kernel of $f$ --- which is defined to be all the elements in the ring such that $f(a)$ is equal to $0$ \inlinemetanote{writes $\ker(f) := \{ a \in R \mid f(a) = 0 \}$} --- is an ideal. \inlinemetanote{writes \qt{is an ideal}} We proved this in one of the early weeks.
\end{speech}

\begin{content}[Definition of Ideals and the Kernel of a Homomorphism]
Let $R$ be a commutative ring with identity.

\begin{definition}[Ideal]\label[definition]{def:ideal}
A subset $I \subseteq R$ is called an \newterm{ideal} of $R$ (denoted $I \trianglelefteq R$) if it satisfies the following two conditions:
\begin{enumerate}[label=\textbf{(\arabic*)}]
    \item \textbf{Additive Subgroup Condition:}
    $(I, +, 0)$ is a subgroup of the additive group $(R, +, 0)$. That is, $0 \in I$, and for all $x, y \in I$, $x - y \in I$.
    \item \textbf{Absorption / Multiplication Constraint:}
    $I$ absorbs multiplication by arbitrary elements of the ambient ring $R$:
    \[
    \forall r \in R, \quad \forall x \in I \implies r \cdot x \in I.
    \]
\end{enumerate}
\end{definition}

\begin{proposition}[Kernel of a Ring Homomorphism as an Ideal]\label[proposition]{prop:kernel-is-ideal}
Let $f \colon R \to S$ be a homomorphism of commutative rings. The \newterm{kernel} of $f$, defined by:
\begin{equation}\label{eq:kernel-definition}
\ker(f) := \{ a \in R \mid f(a) = 0_S \},
\end{equation}
is an ideal of $R$.
\end{proposition}

\begin{proof}
We verify the two defining conditions of \cref{def:ideal}:
\begin{enumerate}[label=\textbf{(\arabic*)}]
    \item \textbf{Additive Subgroup:} Since $f(0_R) = 0_S$, $0_R \in \ker(f)$. For $x, y \in \ker(f)$, additivity of $f$ yields $f(x - y) = f(x) - f(y) = 0_S - 0_S = 0_S$, so $x - y \in \ker(f)$.
    \item \textbf{Multiplicative Absorption:} For any $r \in R$ and $x \in \ker(f)$, multiplicativity gives:
    \[
    f(r \cdot x) = f(r) \cdot f(x) = f(r) \cdot 0_S = 0_S.
    \]
    Thus, $r \cdot x \in \ker(f)$.
\end{enumerate}
\end{proof}
\end{content}

\subsection{The Linear Algebra Analogy: Kernels, Images, and Isomorphisms}

\begin{speech}[00:18:33 - 00:20:20]
So, in general, when you study algebra, even in linear algebra --- I remind you in the study of linear algebra, this is a reminder: \inlinemetanote{writes \qt{Reminder:}}

If you have a vector space $V$ and you have a linear transformation, let's call it $\phi$. \inlinemetanote{writes $\phi \colon V \to W$} So, $V$ and $W$ are vector spaces \inlinemetanote{writes $V, W$ vector spaces} over some field --- it's not relevant for this example, say the complex numbers ---, and $\phi$ is a linear transformation. \inlinemetanote{writes $\phi$ linear transformation} I don't know what you called homomorphisms of vector spaces: linear transformations.

In the abstract study of linear algebra, when you have such a linear transformation, you automatically get two other vector spaces, so to speak for free. The two vector spaces are the kernel \inlinemetanote{writes $\ker(\phi) \subseteq V$} --- this is a subspace of $V$ ---, and the image, \inlinemetanote{writes $\operatorname{Im}(\phi) \subseteq W$} which is a subspace of $W$. And there is a theorem that links this data. If you remember, it says that if I take $V$ and I quotient by the kernel --- I hope you remember this, it's kind of crucial here ---, if you take $V$ and you quotient by the kernel, this is isomorphic as vector spaces to the image. \inlinemetanote{writes $V / \ker(\phi) \cong \operatorname{Im}(\phi)$}

All right, so this is kind of a model. Did you learn it in these terms? I've never taught linear algebra here. Okay. So, we would like to think a little bit about this kind of construction for rings, and we'll see a lot of it to start.
\end{speech}

\begin{overview}[Linear Algebra Analogy: The First Isomorphism Theorem]
For a linear transformation $\phi \colon V \to W$ between vector spaces over a field $\mathbb{F}$:
\begin{itemize}
    \item \textbf{Kernel:} $\ker(\phi) = \{ v \in V \mid \phi(v) = 0 \} \le V$ is a vector subspace of $V$.
    \item \textbf{Image:} $\operatorname{Im}(\phi) = \{ \phi(v) \mid v \in V \} \le W$ is a vector subspace of $W$.
    \item \textbf{First Isomorphism Theorem for Vector Spaces:}
    \[
    \frac{V}{\ker(\phi)} \cong \operatorname{Im}(\phi).
    \]
\end{itemize}
In ring theory, we seek an identical structural correspondence for a ring homomorphism $f \colon R \to S$.
\end{overview}

\begin{speech}[00:20:20 - 00:21:11]
Yeah, I haven't taught linear algebra for... I don't know, it's possible not in this century. But it's a great class, linear algebra. I used to tell students it's like the most important math class, linear algebra, because humans are actually pretty good at solving linear equations and pretty terrible at solving any other kind of equation.

That's not true, I did teach it. \inlinemetanote{cleans and wipes the lower board}
\end{speech}

\subsection{Ring Homomorphisms vs. Linear Transformations}

\begin{speech}[00:21:11 - 00:23:08]
Okay, so if I have this ring homomorphism... \inlinemetanote{writes $f \colon R \to S$} So, these are both rings now --- I try to make the parallel here: $R$ and $S$ are rings, \inlinemetanote{writes $R, S$ are rings} and I remind you they're commutative rings with unit, and $f$ always takes unit: $f$ takes $1$ to $1$, that's the definition of a ring homomorphism. And $f$ is a homomorphism, \inlinemetanote{writes $f$ hom} so we just want to complete that board.

Now, in this situation, we get two objects for free --- or, I don't know, \qt{free} might be the wrong word, but anyway, this gives us two objects. One is the kernel inside $R$, but this is not... there, everything was vector spaces, but this is not a ring now: this is an ideal. \inlinemetanote{writes $\ker(f) \subseteq R$ ideal} But we also get the image of $f$ inside $S$, and that is a subring. \inlinemetanote{writes $\operatorname{Im}(f) \subseteq S$ subring} We saw that the image was a subring.

So, that's kind of the parallel there. And then we're going to work on this quotient business. But you see, the matter is slightly... in this vector space world, everything is a vector space. But in the ring world, sometimes you get rings and sometimes you get ideals, and they're slightly different objects. So, we have to come to terms with that.

Let me just say some basic statements: \inlinemetanote{writes \qt{statements:}} One is that $f$ is surjective if and only if this subring is the whole ring. \inlinemetanote{writes \textbf{(1)} $f$ is surjective $\iff \operatorname{Im}(f) = S$} That's almost like a tautology, right? That's what surjective means: you get everybody here, and if the image is everybody, it's surjective. That's essentially a tautology.
\end{speech}

\begin{content}[The Morphism Data for Rings: Ideals vs. Subrings]
When comparing linear transformations $\phi \colon V \to W$ and ring homomorphisms $f \colon R \to S$, a crucial categorical divergence arises:
\begin{center}
\begin{tabular}{l l l}
    \toprule
    \textbf{Structure} & \textbf{Linear Algebra ($\phi \colon V \to W$)} & \textbf{Ring Theory ($f \colon R \to S$)} \\
    \midrule
    \textbf{Domain Sub-object} & $\ker(\phi) \le V$ (Vector Subspace) & $\ker(f) \trianglelefteq R$ (\textbf{Ideal}, not a subring with $1_R$) \\
    \textbf{Codomain Sub-object} & $\operatorname{Im}(\phi) \le W$ (Vector Subspace) & $\operatorname{Im}(f) \le S$ (\textbf{Subring}, contains $1_S = f(1_R)$) \\
    \bottomrule
\end{tabular}
\end{center}

\begin{proposition}[Surjectivity Criterion]\label[proposition]{prop:surjective-image}
Let $f \colon R \to S$ be a ring homomorphism. Then $f$ is surjective if and only if:
\begin{equation}\label{eq:surjective-criterion}
\operatorname{Im}(f) = S.
\end{equation}
\end{proposition}
\end{content}

\begin{hidden}
% timestamp: 00:23:08
% topic: Parallel between linear transformations and ring homomorphisms; properties of kernel and image.
% board_state: def:ideal, prop:kernel-is-ideal, prop:surjective-image
% next_goal: Characterization of injectivity via the zero ideal (statement 2).
% open_loops: none
\end{hidden}

\subsection{Injectivity and the Zero Ideal}

\begin{speech}[00:23:08 - 00:24:24]
The second statement --- and that's also true there --- is: $f$ is injective if and only if... \inlinemetanote{writes \textbf{(2)} $f$ is injective $\iff \ker(f) = (0)$}

So, what's the statement going to be? It's going to be about the kernel: the kernel of $f$ is the zero ideal. This is the standard terminology for the zero ideal, the standard notation for the zero ideal. The zero ideal is the ideal that has $0$ and nothing else. If you look at the definition of an ideal, the ideal that has the single element $0$... This is going to be used a lot, so it should be completely clear: \inlinemetanote{writes \qt{Zero ideal: $(0)$, $0$, $\{0\}$}}

The zero ideal: sometimes this is denoted like this, $(0)$; sometimes it's just denoted $0$; but the zero ideal is the ideal that contains one element: $\{0\}$. \inlinemetanote{writes $\{0\}$} It's definitely a subgroup, because it's just zero: zero, when you add it, stays zero, it has inverses, there's nothing wrong with zero. And the second condition is: it's closed under multiplication, because if you take zero in the ideal and multiply by anybody, you stay zero. So, that's the zero ideal.

And this last statement says: $f$ is injective if and only if the kernel is the zero ideal.
\end{speech}

\begin{content}[Injectivity and the Zero Ideal]
\begin{definition}[The Zero Ideal]\label[definition]{def:zero-ideal}
In any ring $R$, the singleton set consisting solely of the additive neutral element:
\begin{equation}\label{eq:zero-ideal-def}
(0) = \{0\}
\end{equation}
is an ideal of $R$, called the \newterm{zero ideal}.
\end{definition}

\begin{proposition}[Injectivity and the Kernel]\label[proposition]{prop:injective-kernel}
Let $f \colon R \to S$ be a ring homomorphism. Then $f$ is injective if and only if its kernel is trivial:
\begin{equation}\label{eq:injective-kernel-criterion}
f \text{ is injective} \iff \ker(f) = (0).
\end{equation}
\end{proposition}
\end{content}

\begin{proofWrapper}[Proof of \cref{prop:injective-kernel}]

\begin{speech}[00:24:24 - 00:25:39]
So, let's try to see what the argument for that is --- that's almost a tautology. Proof of \textbf{(2)}: \inlinemetanote{writes \qt{Proof of (2):}}

Suppose... we have to do two things. Suppose $f$ is injective. \inlinemetanote{writes $(\implies)$ Suppose $f$ is injective} Then let $a \in R$ be non-zero. \inlinemetanote{writes let $a \in R, a \neq 0$} Then $f(a)$ is not equal to $f(0)$, \inlinemetanote{writes then $f(a) \neq f(0)$} because injective means if you take different elements, they go to different places. And $f(0) = 0$, \inlinemetanote{writes $= 0$} so that means $f(a)$ is not zero, which implies $a$ is not in the kernel of $f$. \inlinemetanote{writes $\implies a \notin \ker(f)$}

So, if $f$ is injective, nobody who is not zero is in the kernel of $f$, so this implies the kernel of $f$ equals $(0)$. \inlinemetanote{writes $\implies \ker(f) = (0)$} That's one half of it: I assume it's injective, and I deduce the kernel of $f$ is zero --- it can only have one element in it.
\end{speech}

\begin{content}[Forward Direction: Injectivity Implies Trivial Kernel]
\begin{proof}[Forward Implication ($\implies$)]
Suppose $f$ is injective. Let $a \in R$ with $a \neq 0$.
Since $f$ is injective and $a \neq 0$:
\[
f(a) \neq f(0).
\]
Since $f$ is a ring homomorphism, $f(0) = 0_S$. Thus:
\[
f(a) \neq 0_S \implies a \notin \ker(f).
\]
Consequently, no non-zero element of $R$ belongs to $\ker(f)$. Since $0 \in \ker(f)$ always holds, we conclude:
\[
\ker(f) = (0).
\]
\end{proof}
\end{content}

\begin{speech}[00:25:39 - 00:26:39]
The other way around is: suppose the kernel of $f$ is $(0)$ --- \inlinemetanote{writes $(\impliedby)$ Suppose $\ker(f) = (0)$} and as I said, it's standard to use both notations, depending on what you want to emphasize.

Suppose the kernel of $f$ is zero, and let $f(a) = f(b)$. \inlinemetanote{writes then let $f(a) = f(b)$} We're assuming it's zero, and then we pick two elements $a, b$ in the ring that go to the same place. This implies, by the homomorphism property, $f(a) - f(b) = f(a - b) = 0$, because I'm assuming they're the same. \inlinemetanote{writes $\implies f(a) - f(b) = f(a-b) = 0$}

So, this implies that $a - b$ is in the kernel; \inlinemetanote{writes $\implies a - b \in \ker(f)$} but we're assuming the kernel is zero, so this implies $a - b = 0$, which implies $a = b$. \inlinemetanote{writes $\implies a - b = 0 \implies a = b$}

That's the second implication. And actually, if you look at this proof, it's the same proof as you had in linear algebra; there's actually no difference.
\end{speech}

\begin{content}[Reverse Direction: Trivial Kernel Implies Injectivity]
\begin{proof}[Reverse Implication ($\impliedby$)]
Suppose $\ker(f) = (0)$. Let $a, b \in R$ such that $f(a) = f(b)$.
Using the additive homomorphism property of $f$:
\[
f(a) - f(b) = 0_S \implies f(a - b) = 0_S.
\]
By definition of the kernel, this implies:
\[
a - b \in \ker(f).
\]
Since $\ker(f) = (0) = \{0\}$ by hypothesis, it follows that:
\[
a - b = 0_R \implies a = b.
\]
Thus, $f$ is injective, completing the proof of \cref{prop:injective-kernel}.
\end{proof}
\end{content}

\end{proofWrapper}

\begin{hidden}
% timestamp: 00:26:39
% topic: Equivalence of injectivity and trivial kernel; parallel to linear algebra.
% board_state: prop:surjective-image, def:zero-ideal, prop:injective-kernel
% next_goal: Transition to the quotient ring construction R/I.
% open_loops: none
\end{hidden}

\begin{speech}[00:26:39 - 00:28:05]
For this linear algebra picture, we've recovered all the pieces except this one, \inlinemetanote{points to $V/\ker(\phi) \cong \operatorname{Im}(\phi)$ on the upper right board} which says that $V$ mod the kernel is the image, and we now want to recover that. And for that, we need a construction. There was a construction already in linear algebra you needed, which was the notion of a quotient vector space; if you remember that, it's going to be very helpful.

The idea of a math education is that, actually, you're not allowed to forget anything you've learned. \inlinemetanote{erases the board} That's one theory. And the other theory is that education is what's left after you've forgotten everything you've learned. You can decide which theory of education you'd like to subscribe to. \inlinemetanote{cleans the lower board completely with a sponge}
\end{speech}

\begin{note}[Board Cleaning and Pedagogical Aside]
The lecturer erases the lower blackboard to prepare for the construction of quotient rings, sharing an aside on two contrasting theories of education.
\end{note}

\section{The Construction of Quotient Rings}

\begin{speech}[00:28:05 - 00:29:35]
Okay, so the task at hand is to understand this quotient. And now we do this not just for the kernel --- we'll get back to the kernel ---, but for any ideal. So, we let $I \subseteq R$ be an ideal. \inlinemetanote{writes Let $I \subseteq R$ be an ideal} And now this is the main construction for this week: we will construct the quotient ring. \inlinemetanote{writes We will construct the quotient ring} That's the name: quotient ring. And the symbol is $R/I$. \inlinemetanote{writes $R/I$} That's the symbol.

We have to construct it, and the parallel to linear algebra, of course, is $R$ mod the kernel --- that's the quotient ring. We'll get back to that in a moment. So, we forget about this linear algebra analogy for now, and we go for this construction: $R$ is a ring, $I$ is an ideal, and we need to construct the quotient.

We need to construct this object as a ring. That means, as usual, it's a multi-step construction: we have to first construct it as a set; then we have to find zero in it; we have to find one in it, \inlinemetanote{writes $(R/I, +, \cdot, 0, 1)$} zero and one; and we have to define addition and multiplication. You have to construct the whole thing --- that's just what you have to do; there's no other way around it.
\end{speech}

\begin{content}[The Quotient Ring Construction: Roadmap]
Let $R$ be a commutative ring with identity $1$, and let $I \trianglelefteq R$ be an ideal. Our objective is to construct a new ring:
\begin{equation}\label{eq:quotient-ring-symbol}
R/I = (R/I, +, \cdot, 0_{R/I}, 1_{R/I}),
\end{equation}
called the \newterm{quotient ring} (or residue class ring) of $R$ modulo $I$.

\begin{steps}[The Five Stages of Constructing a Ring from Scratch]
To endow the quotient $R/I$ with a complete ring structure, we must rigorously execute the following stages:
\begin{enumerate}[label=\textbf{Step \arabic*:}]
    \setcounter{enumi}{0}
    \item \textbf{Underlying Set:} Define an equivalence relation on $R$ and take $R/I$ to be the set of its equivalence classes.
    \item \textbf{Distinguished Elements:} Identify the additive zero $0_{R/I}$ and multiplicative identity $1_{R/I}$.
    \item \textbf{Addition Operation:} Define binary addition $+$ on equivalence classes and prove it is \emph{well-defined}.
    \item \textbf{Multiplication Operation:} Define binary multiplication $\cdot$ on equivalence classes and prove it is \emph{well-defined} using the ideal absorption property.
    \item \textbf{Ring Axioms Verification:} Verify associativity, commutativity, distributivity, and existence of inverses.
\end{enumerate}
\end{steps}
\end{content}

\subsection{Step 1: The Underlying Set and Equivalence Relation Modulo an Ideal}

\begin{speech}[00:29:35 - 00:31:26]
Step one --- there's a lot of steps and, as usual, I don't do all of them: \inlinemetanote{writes Step 1: What is the set $R/I$?} What is the set $R/I$? And here the notation is totally standard: everybody uses just this notation. What is the set?

If you remember linear algebra, that should help you now. We define an equivalence relation \inlinemetanote{writes We define an equivalence relation $\equiv$ on $R$} on $R$. Sometimes this equivalence relation could be denoted by this, \inlinemetanote{writes $\equiv_I$} but you have to carry around a lot of stuff, so I'm just going to leave that $I$ out. We're going to define an equivalence relation on $R$.

Here is the definition: if I have $x$ and $y$ in $R$, then $x$ is equivalent to $y$ \inlinemetanote{writes Def: $x, y \in R$, $x \equiv y$} --- and as I said, if you like, you can put $I$ here, because the equivalence relation depends on $I$ --- if and only if $x - y$ is in the ideal $I$. \inlinemetanote{writes $\iff x - y \in I$} That's the definition.

In words --- as an aside: \inlinemetanote{writes \qt{Side:}} if you want to be really precise, you could write $\equiv_I$, \inlinemetanote{writes $\equiv_I$} because this depends on $I$ --- this is usually put: \qt{$x$ is equivalent to $y \bmod I$}. \inlinemetanote{writes $x$ is equivalent to $y \bmod I$} This is very standard language: $x$ is equivalent to $y \bmod I$. That's what this equivalence relation is; it's just like for the integers.
\end{speech}

\begin{content}[The Equivalence Relation Modulo an Ideal]
\begin{definition}[Congruence Modulo an Ideal]\label[definition]{def:congruence-modulo-ideal}
Let $R$ be a ring and $I \trianglelefteq R$ an ideal. For elements $x, y \in R$, we define the relation of \newterm{congruence modulo $I$} (denoted $x \equiv_I y$ or simply $x \equiv y$) by:
\begin{equation}\label{eq:congruence-def}
x \equiv y \iff x - y \in I.
\end{equation}
In words, we say that \qt{$x$ is equivalent to $y$ modulo $I$}.
\end{definition}

\begin{steps}[Connection to Modular Arithmetic in $\mathbb{Z}$]
This directly generalizes congruence modulo $n$ in the integers:
\[
x \equiv y \pmod n \iff x - y \in n\mathbb{Z},
\]
where $I = n\mathbb{Z} = (n)$ is the principal ideal generated by $n$.
\end{steps}
\end{content}

\begin{speech}[00:31:26 - 00:33:28]
Formally, you now have to check that it's an equivalence relation. To do that, you have to remember the three properties of an equivalence relation: $x \equiv x$, which is zero being in the ideal; the second property is if $x \equiv y$, then $y \equiv x$, which would just put a minus sign here, but we know that $I$ is a subgroup under addition, so it's closed under taking negatives; and the last one is transitivity. Maybe I'll do the last one: the transitive property. \inlinemetanote{writes Transitivity:}

The transitive property says: if $x \equiv y$ and $y \equiv z$, then $x \equiv z$. \inlinemetanote{writes if $x \equiv y$ and $y \equiv z$ then $x \equiv z$} Now, $x \equiv y$ means $x - y \in I$, \inlinemetanote{writes $x - y \in I$} and $y \equiv z$ means $y - z \in I$. \inlinemetanote{writes $y - z \in I$} What I need to show is $x - z \in I$. But these two elements are in $I$, so I can add them: \inlinemetanote{writes $(x-y) + (y-z) \in I$} since $I$ is closed under addition, I can add them. Is that clear? And then, happily, the $y$ cancels out, so you get $x - z \in I$, \inlinemetanote{writes $x - z \in I$} and we're happy about that. That's the proof of transitivity.

If you think about what we've used: we've used zero is in $I$, we've used the minus sign is in $I$, and we've used that $I$ is closed under addition; so, we've used all the group conditions for $I$. We couldn't get away with less. What we have not used is the multiplicative condition --- that's going to come later. An ideal has two assumptions: one is the additive structure, it's a subgroup, and we've used it all here; but we have not touched the multiplicative condition. That'll come in a moment.
\end{speech}

\begin{content}[Verification: Equivalence Relation Properties]
\begin{proposition}[$\equiv_I$ is an Equivalence Relation]\label[proposition]{prop:equiv-rel-modulo-ideal}
Let $I \trianglelefteq R$ be an ideal. The relation $\equiv_I$ is an equivalence relation on $R$.
\end{proposition}

\begin{proof}
We verify the three defining axioms:
\begin{enumerate}[label=\textbf{(\arabic*)}]
    \item \textbf{Reflexivity ($x \equiv x$):}
    For any $x \in R$, $x - x = 0 \in I$ because $(I, +, 0)$ is a subgroup. Thus, $x \equiv x$.
    \item \textbf{Symmetry ($x \equiv y \implies y \equiv x$):}
    If $x \equiv y$, then $x - y \in I$. Since $(I, +, 0)$ contains additive inverses:
    \[
    y - x = -(x - y) \in I \implies y \equiv x.
    \]
    \item \textbf{Transitivity ($x \equiv y \land y \equiv z \implies x \equiv z$):}
    Suppose $x \equiv y$ and $y \equiv z$. By definition:
    \[
    x - y \in I \quad \text{and} \quad y - z \in I.
    \]
    Since $I$ is closed under addition, their sum lies in $I$:
    \[
    (x - y) + (y - z) \in I \implies x - z \in I \implies x \equiv z.
    \]
\end{enumerate}
\end{proof}

\begin{steps}[Subgroup Property Sufficiency]
Notice that verifying that $\equiv_I$ is an equivalence relation relies \textbf{exclusively} on the fact that $(I, +, 0)$ is an additive subgroup of $(R, +, 0)$. The multiplicative absorption condition $R \cdot I \subseteq I$ has not yet been utilized!
\end{steps}
\end{content}

\begin{hidden}
% timestamp: 00:33:28
% topic: Definition of congruence modulo an ideal and proof that it is an equivalence relation.
% board_state: def:congruence-modulo-ideal, prop:equiv-rel-modulo-ideal
% next_goal: Define the quotient set R/I as equivalence classes and examine naming conventions.
% open_loops: none
\end{hidden}

\begin{speech}[00:33:28 - 00:34:48]
If it's an equivalence relation, then you can imagine that we're going to take its equivalence classes. \inlinemetanote{erases the right portion of the lower board}

Now we can define the set: $R/I$ as a set is equal to the set of equivalence classes \inlinemetanote{writes $R/I = \text{set of equivalence classes}$} of this relation on $R$. \inlinemetanote{writes of $\equiv_I$ on $R$} So, I've defined this as a set.
\end{speech}

\begin{content}[The Quotient Set $R/I$]
\begin{definition}[The Quotient Set $R/I$]\label[definition]{def:quotient-set}
The underlying set of the quotient ring $R/I$ is defined as the set of all equivalence classes of $R$ with respect to the equivalence relation $\equiv_I$:
\begin{equation}\label{eq:quotient-set-def}
R/I := R / {\equiv_I} = \{ [x] \mid x \in R \}.
\end{equation}
\end{definition}
\end{content}

\subsection{Naming Conventions for Equivalence Classes (Cosets)}

\begin{speech}[00:34:48 - 00:36:05]
And now, the question of names: how do I name things? That's the same discussion we had about the rational numbers, the same for the field of fractions. One of the conceptual issues about the quotient is that elements do not have unique names. And this is just something you have to live with.

The question is: how do I name an element in $R/I$? \inlinemetanote{writes How do I name an element in $R/I$?}

One way to do it is: you take $x \in R$. \inlinemetanote{writes $x \in R$} Since $x$ is in $R$, it determines an equivalence class: the equivalence class which $x$ lies in. And that is usually called $\bar{x}$. \inlinemetanote{writes $\bar{x} \in R/I$ the equivalence class that $x$ lies in} This is the equivalence class that $x$ lies in. You take $x$, you get an equivalence class, and that's usually denoted $\bar{x}$. That's what we do with integers also. Is that okay?
\end{speech}

\begin{speech}[00:36:05 - 00:37:33]
There are three standard ways to denote this. If $x \in R$, we can write $\bar{x} \in R/I$. \inlinemetanote{writes $x \in R \implies \bar{x} \in R/I$}

If you really want to write things out explicitly, you can write it like that: $x + I \in R/I$, \inlinemetanote{writes $x + I \in R/I$} because the equivalence class of $x$ is equal to all those elements $x + i$ where $i \in I$. \inlinemetanote{writes $[x] = \{x + i \mid i \in I\}$} Because if you're equivalent to $x$, the difference is in $I$. That is explicitly what the equivalence class is. And since it's $x$ plus something in $I$, some people will write it like this: $x + I$. \inlinemetanote{points to $x + I$} It's kind of messy to keep that around, but in some arguments it's nice. That's the second way to write it.

The third way to write it --- which is by far the most common --- is: you just call it $x$ or $[x]$. \inlinemetanote{writes $[x]$ or just $x$} That's by far the most common. But if you use this, you have to be very careful: when $x$ was born, it was an element, and then when you put it here, it becomes an equivalence class; you have to keep track of the \qt{data type} of $x$ in your head. Those are the three ways to do it. They all have their own advantages --- well, this one has limited advantage, but okay.
\end{speech}

\begin{content}[Three Standard Notations for Residue Classes (Cosets)]
An element of $R/I$ is an equivalence class containing elements of $R$. Because representations are not unique (different representatives can specify the exact same class), three standard notations are commonly employed:

\begin{enumerate}[label=\textbf{(\arabic*)}]
    \item \textbf{Overline Notation ($\bar{x}$):}
    Denoted $\bar{x} \in R/I$. This is compact and mirrors the notation used for residue classes $\mathbb{Z}/n\mathbb{Z}$.
    \item \textbf{Coset Notation ($x + I$):}
    Explicitly writing the equivalence class as an affine shift of the ideal:
    \begin{equation}\label{eq:coset-explicit-def}
    [x] = x + I := \{ x + i \mid i \in I \} \subseteq R.
    \end{equation}
    Since $y \equiv x \iff y - x = i \in I \iff y = x + i$, the equivalence class of $x$ is literally the subset $x + I$.
    \item \textbf{Bare Representative Notation ($x$ or $[x]$):}
    Writing simply $[x]$ or, by abuse of notation, $x \in R/I$. While by far the most common in advanced algebra, this requires carefully tracking whether $x$ represents an individual ring element in $R$ or an entire equivalence class in $R/I$ (the \qt{data type} in one's head).
\end{enumerate}
\end{content}

\subsection{Step 2 \& 3: Distinguished Elements and the Well-Definedness of Addition}

\begin{speech}[00:37:33 - 00:39:08]
All right, so since we've defined the set, we need to now define addition. First, let's do zero and one: we define zero in $R/I$. \inlinemetanote{writes $0 \in R/I$, Def: $0 := \bar{0}$} By definition, $0 := \bar{0}$, because this is the equivalence class of zero. Zeros are always called zeros. So, that's the zero here: it's the equivalence class of zero, which is not a surprise. Similarly, the definition of one is that it's the equivalence class of one: $1 := \bar{1}$. \inlinemetanote{writes Def: $1 := \bar{1}$} Now we have at least this set $R/I$, and we have some distinguished elements: zero and one.

Now, let us define addition. \inlinemetanote{writes Let us define $+$} And it's still called plus. Very formally, you shouldn't call the binary operations in two different rings the same thing, because they're kind of different. But no one does that; everyone calls it plus, and it makes you have to concentrate to follow the definitions.

What does that mean? To define plus, I take an element of $R/I$ and I take another element of $R/I$, and I have to define this plus: \inlinemetanote{writes $\bar{x} + \bar{y} \overset{\text{def}}{=}$}
\end{speech}

\begin{content}[Distinguished Elements and the Definition of Addition]
\begin{definition}[Zero and Identity in $R/I$]\label[definition]{def:distinguished-elements}
The distinguished neutral elements in the quotient ring $R/I$ are defined by the classes of the respective identities in $R$:
\begin{align}
0_{R/I} &:= \bar{0} = 0 + I = I, \label{eq:zero-quotient-def} \\
1_{R/I} &:= \bar{1} = 1 + I. \label{eq:one-quotient-def}
\end{align}
\end{definition}

\begin{definition}[Addition on $R/I$]\label[definition]{def:addition-quotient}
For any two equivalence classes $\bar{x}, \bar{y} \in R/I$, their sum is defined by adding representatives in $R$:
\begin{equation}\label{eq:addition-quotient-formula}
\bar{x} + \bar{y} := \overline{x + y}, \quad \text{or equivalently} \quad (x + I) + (y + I) := (x + y) + I.
\end{equation}
\end{definition}
\end{content}

\begin{speech}[00:39:08 - 00:40:48]
This equation here is the definition. And how am I going to define it? Well, I only know how to add $x$ and $y$ one way, because they were born in $R$, and it's the only thing I know how to do. So, this is the only thing possibly I could write: $\overline{x + y}$. \inlinemetanote{writes $\overline{x + y}$} To read it, you have to concentrate so that you know what's being defined: the left-hand side is being defined, while the right-hand side was already previously known in $R$. This whole equation is the definition of plus.

And it's not a full definition yet, because you have to check that it's well-defined: these elements could have different names. Let's check that it's well-defined in $R/I$: \inlinemetanote{writes Let us check it is well-defined in $R/I$}

What could we do? We could have $\bar{x} = \bar{x}'$, where $x' = x + i$ with $i \in I$. \inlinemetanote{writes $\bar{x} = \bar{x}', x' = x + i, i \in I$} We could have just had a different name for the same equivalence class. And if we had a different name, we know it will just be some shift of $x$ by an element of $I$, because their difference is in $I$ by definition. Similarly, for $y$, we could have a different name: $y' = y + j$, where $j \in I$. \inlinemetanote{writes $\bar{y} = \bar{y}', y' = y + j, j \in I$}
\end{speech}

\begin{hidden}
% timestamp: 00:40:48
% topic: Defining addition in R/I and setting up the check for well-definedness.
% board_state: def:distinguished-elements, def:addition-quotient
% next_goal: Complete the proof of well-definedness of addition in R/I.
% open_loops: none
\end{hidden}

\begin{speech}[00:40:48 - 00:42:38]
Okay, so now we have to test that this definition is the same. We already have the formula for the names $x$ and $y$. Let's try to see what happens if we use these names: we'd better get the same answer! If we use these names, \inlinemetanote{writes $\bar{x}' + \bar{y}'$} this definition would give us $\overline{x' + y'}$. \inlinemetanote{writes $= \overline{x' + y'}$} Then we use the formulas for $x'$ and $y'$, and then I regroup them --- well, okay, maybe I don't: \inlinemetanote{writes $= \overline{(x+i) + (y+j)}$} $(x + i) + (y + j)$; and then I regroup it: $(x + y) + (i + j)$. \inlinemetanote{writes $= \overline{(x+y) + (i+j)}$} But now $i + j$ has to be in the ideal $I$, because $I$ is closed under addition. And so, this equivalence class is actually the same as that equivalence class, because it's the equivalence class of the same thing shifted by an element of $I$.

So, the answer is: we have checked that if I use different names and combine them with addition, we get the same equivalence class on the right-hand side. This proves it's well-defined. And if you remember, as I said, from linear algebra, you've done all of this before.

Then, formally, now that I've defined addition and zero, you have to formally check that this addition is associative, that zero plays the role of a zero, and that there are additive inverses; and I leave all of this to you. \inlinemetanote{erases the middle board}

But there is multiplication, and that's worth doing carefully, because so far we have not used the multiplicative part of the hypothesis on the ideal. And it comes out now, and it's kind of crucial. So, I will try to do this carefully, maybe after the break, but let me just say for...
\end{speech}

\begin{content}[Proof that Addition is Well-Defined on $R/I$]
\begin{proposition}[Well-Definedness of Addition]\label[proposition]{prop:addition-well-defined}
The addition operation defined in \cref{def:addition-quotient} is well-defined: it is independent of the choice of representatives.
\end{proposition}

\begin{proof}
Suppose $\bar{x} = \bar{x}'$ and $\bar{y} = \bar{y}'$ in $R/I$. By \cref{def:congruence-modulo-ideal}, there exist elements $i, j \in I$ such that:
\begin{align}
x' &= x + i, \label{eq:x-prime-shift} \\
y' &= y + j. \label{eq:y-prime-shift}
\end{align}
Evaluating the sum using the primed representatives yields:
\begin{align*}
\bar{x}' + \bar{y}' &= \overline{x' + y'} \\
&= \overline{(x + i) + (y + j)} \\
&= \overline{(x + y) + (i + j)}.
\end{align*}
Since $(I, +, 0)$ is an additive subgroup, $i, j \in I$ implies $i + j \in I$. Therefore, $(x + y) + (i + j) \equiv (x + y) \pmod I$, which implies:
\[
\overline{(x + y) + (i + j)} = \overline{x + y} = \bar{x} + \bar{y}.
\]
Thus, the sum is independent of representatives.
\end{proof}

\begin{steps}[The Quotient Additive Group]
Since addition is well-defined, associativity, commutativity, the identity $\bar{0}$, and inverses $-\bar{x} = \overline{-x}$ all inherit directly from $(R, +, 0)$. Hence, $(R/I, +, \bar{0})$ forms an abelian group.
\end{steps}
\end{content}

\subsection{Step 4: Multiplication and the Crucial Role of the Ideal Absorption Property}

\begin{speech}[00:42:38 - 00:43:48]
The next step is to define multiplication on $R/I$. \inlinemetanote{writes Define $\cdot$ on $R/I$} That means, if we have $x$ and $y$, it's exactly as before: we need to define this, \inlinemetanote{writes $\bar{x} \cdot \bar{y} \overset{\text{def}}{=}$} and this is going to be the definition. And how are you going to define it? You only know how to do one thing to define things, so you really have no choice except this: \inlinemetanote{writes $\overline{xy}$} $\overline{x y}$. I could put the dot there too if I want.

This just says that if I want to multiply two equivalence classes, I take one element of one equivalence class, another element of the other equivalence class, I multiply them in the ring, and then take the resulting equivalence class. Is that clear?
\end{speech}

\begin{content}[Definition of Multiplication on $R/I$]
\begin{definition}[Multiplication on $R/I$]\label[definition]{def:multiplication-quotient}
For equivalence classes $\bar{x}, \bar{y} \in R/I$, their product is defined by multiplying representatives in $R$:
\begin{equation}\label{eq:multiplication-quotient-def}
\bar{x} \cdot \bar{y} := \overline{x \cdot y}, \quad \text{or equivalently} \quad (x + I) \cdot (y + I) := (x \cdot y) + I.
\end{equation}
\end{definition}
\end{content}

\begin{speech}[00:43:48 - 00:45:11]
And then I have to prove this is well-defined. \inlinemetanote{writes Prove well-defined} So, I use the same choices: suppose I pick these different names and I use the definition, \inlinemetanote{writes $\bar{x}' \cdot \bar{y}' = \overline{x' y'}$} then I get $\overline{x' y'}$, and then I use their formulas: \inlinemetanote{writes $= \overline{(x+i)(y+j)}$} $(x + i)(y + j)$ \inlinemetanote{corrects to $(x+i)(y+j)$}.

And when I multiply them out --- this is the crucial part ---, I get... sorry, there's no primes here. I get $x y$, and then I get $i \cdot y + j \cdot x$ --- it's commutative, so I can write it either way --- plus $i j$. \inlinemetanote{writes $= \overline{xy + iy + jx + ij}$ feedback on terms}

With one name, I get $\overline{x y}$. With the other name, I get this. And the difference is this. So, the only way this is well-defined is if this whole thing is in $I$. \inlinemetanote{brackets $iy + jx + ij$ and writes $\in I$} That's the only way it could be well-defined. And it is in $I$! Here, finally, we use the fact that if you have an element of $I$ and multiply it by anything, you stay in $I$. So, that's why it's here. Okay, so...
\end{speech}

\begin{content}[Proof that Multiplication is Well-Defined on $R/I$]
\begin{proposition}[Well-Definedness of Multiplication]\label[proposition]{prop:multiplication-well-defined}
The multiplication operation defined in \cref{def:multiplication-quotient} is well-defined.
\end{proposition}

\begin{proof}
Let $\bar{x}' = \bar{x}$ and $\bar{y}' = \bar{y}$. Then $x' = x + i$ and $y' = y + j$ for some $i, j \in I$. Expanding the product of the alternative representatives:
\begin{align*}
\bar{x}' \cdot \bar{y}' &= \overline{x' \cdot y'} \\
&= \overline{(x + i)(y + j)} \\
&= \overline{x y + i y + j x + i j}.
\end{align*}
To establish that $\overline{(x + i)(y + j)} = \overline{x y}$, we must verify that the difference belongs to $I$:
\[
(x y + i y + j x + i j) - x y = i y + j x + i j \stackrel{?}{\in} I.
\]
We examine each of the three cross-terms:
\begin{itemize}
    \item $i \in I$ and $y \in R \implies i \cdot y \in I$ by the \textbf{multiplicative absorption property} of ideals ($I \cdot R \subseteq I$).
    \item $j \in I$ and $x \in R \implies j \cdot x \in I$ by the \textbf{multiplicative absorption property} ($R \cdot I \subseteq I$).
    \item $i, j \in I \implies i \cdot j \in I$ (in fact, $I \cdot I \subseteq I$).
\end{itemize}
Since $(I, +, 0)$ is closed under addition:
\[
i y + j x + i j \in I \implies \overline{x y + i y + j x + i j} = \overline{x y} = \bar{x} \cdot \bar{y}.
\]
Therefore, multiplication on $R/I$ is strictly well-defined.
\end{proof}

\begin{steps}[Why Subrings Fail and Ideals are Necessary]
If $I$ were merely an additive subgroup or a subring (closed under multiplication within itself, but not absorbing arbitrary elements of $R$), the terms $i y$ and $j x$ would generally fail to lie in $I$, causing the multiplication on equivalence classes to collapse! \textbf{This step is the sole reason why ideals, rather than general subrings, are required to define quotient rings.}
\end{steps}
\end{content}

\begin{hidden}
% timestamp: 00:45:11
% topic: Well-definedness of multiplication in R/I and mid-lecture break.
% board_state: def:multiplication-quotient, prop:multiplication-well-defined
% next_goal: Recap of ideal axioms and why absorption is crucial after the break.
% open_loops: none
\end{hidden}

\begin{pause}[Mid-Lecture Break]
The lecturer takes a short pause for the mid-lecture break. The lecture resumes at timestamp 00:45:11 with a pedagogical recap of the essential role played by the ideal axioms.
\end{pause}

\subsection{Recap: The Crucial Role of the Ideal Axioms}

\begin{speech}[00:45:11 - 00:46:13]
Okay, so we start, and the last moment in the first half of the class was about this observation, and it's a really important thing. So, I repeat it now: this is the only time...

If you remember, the definition of an ideal had two clauses. The first is that $(I, +, 0)$ is an additive subgroup, and we've used that repeatedly: in the definition of the equivalence relation, in the definition of addition; and if you wanted to prove associativity and all of those properties that I didn't prove, you'd be using it here.

But it has a second property, which was the multiplicative property: that for every element of the ring and for every $a$ in the ideal, the product is in the ideal. \inlinemetanote{writes on the lower right board: Def of Ideal: (1) $(I, +, 0) < (R, +, 0)$ subgroup, (2) $\forall r \in R, \forall a \in I \implies ra \in I$} That's just a totally separate property, because the first one doesn't have anything to do with multiplication. This one has to do with multiplication, and this is the first time we use it.
\end{speech}

\begin{content}[Review: The Two Defining Axioms of an Ideal]
\begin{definition}[Ideal Axioms Reviewed]\label[definition]{def:ideal-axioms-reviewed}
A subset $I \subseteq R$ is an ideal if and only if:
\begin{enumerate}[label=\textbf{(\arabic*)}]
    \item \textbf{Additive Subgroup Axiom:}
    \[
    (I, +, 0) \le (R, +, 0) \quad \text{is an additive subgroup.}
    \]
    \item \textbf{Multiplicative Absorption Axiom:}
    \[
    \forall r \in R, \quad \forall a \in I \implies r \cdot a \in I.
    \]
\end{enumerate}
\end{definition}

\begin{steps}[Division of Labor Between the Two Axioms]
\begin{itemize}
    \item Axiom \textbf{(1)} governs all additive aspects: transitivity of equivalence classes, well-definedness of addition, additive identity, and additive inverses in $R/I$.
    \item Axiom \textbf{(2)} is reserved solely for multiplication: ensuring that products of cosets do not depend on the representative chosen.
\end{itemize}
\end{steps}
\end{content}

\begin{speech}[00:46:13 - 00:47:25]
And it's used to show that this definition of multiplication --- which, as I tried to say, philosophically you didn't have very much choice about: if you want to multiply two equivalence classes, you pick an element in one equivalence class, you pick an element in the other, you multiply those elements, and you take the equivalence class; that's pretty much the only choice you have to define it, because you start with just an abstract ring.

And to show this is well-defined, which I did here quickly at the end with different elements in those equivalence classes, you come to this issue: you're taking the equivalence class of this element, and it has to be the same as the equivalence class of that element. And the only way it is the same is if the extra terms, when you combine them, are all in $I$.

And to see that: this is a product of two elements of $I$. That's okay, because for any element of $R$ and any element of $I$, you only need one element of $I$. So, this one is in $I$; and these two have one element of $I$ and one element of $R$ each, because $x$ and $y$ are arbitrary elements of $R$. So, these are three times we apply this rule.
\end{speech}

\subsection{Completion of the Quotient Ring Structure}

\begin{speech}[00:47:25 - 00:48:47]
Okay, so that proves...

And this is one aspect of the construction which is not in your linear algebra class. Although it's not far from what happened in your linear algebra class when you considered the problem of scalar multiplication on the quotient, which I won't repeat, but anyway. So, even that shouldn't be totally foreign to you.

So, now we've constructed this: \inlinemetanote{writes $((R/I), +, \cdot, \bar{0}, \bar{1})$ on the lower left board} we have plus, we have multiplication, we have $0$ and $1$. \inlinemetanote{adds checkmarks above the symbols}

The tricky part --- it's not even that tricky, but the part that requires the foundational work --- is the definition as equivalence classes, and making sure these are well-defined. For the rest, you just follow your nose. The identity: you can imagine, you have to check $\bar{1} \cdot \bar{x}$, \inlinemetanote{writes $\bar{1} \cdot \bar{x} = \overline{1 \cdot x} = \bar{x}$ with a checkmark} so by definition this is $\overline{1 \cdot x}$, and so that's $\bar{x}$. That checks one of the identity properties, and there are about $10$ other properties you have to check, and they're all checked like this.

So, that's the construction, and that's the main construction this week. Last week we had the construction of the field of fractions, which you should understand completely; and now there is this construction, which is the construction of the quotient ring, which is your task to understand really completely this week: the construction.
\end{speech}

\begin{content}[The Complete Quotient Ring Structure]
\begin{theorem}[The Quotient Ring $(R/I, +, \cdot, \bar{0}, \bar{1})$]\label[theorem]{thm:quotient-ring-structure}
Let $R$ be a commutative ring with identity, and let $I \trianglelefteq R$ be an ideal. The quintuple:
\begin{equation}\label{eq:quotient-ring-quintuple}
(R/I, +, \cdot, \bar{0}, \bar{1})
\end{equation}
equipped with the coset operations:
\[
\bar{x} + \bar{y} := \overline{x + y}, \qquad \bar{x} \cdot \bar{y} := \overline{x \cdot y}, \qquad 0_{R/I} := \bar{0}, \qquad 1_{R/I} := \bar{1},
\]
forms a commutative ring with identity.
\end{theorem}

\begin{proof}
Having proven that addition and multiplication are well-defined, the remaining ring axioms follow directly from the ring axioms of $R$ applied to representatives:
\begin{itemize}
    \item \textbf{Multiplicative Identity:} $\bar{1} \cdot \bar{x} = \overline{1 \cdot x} = \bar{x}$ for all $\bar{x} \in R/I$.
    \item \textbf{Associativity:} $(\bar{x} \cdot \bar{y}) \cdot \bar{z} = \overline{(x y) z} = \overline{x (y z)} = \bar{x} \cdot (\bar{y} \cdot \bar{z})$.
    \item \textbf{Distributivity:} $\bar{x} \cdot (\bar{y} + \bar{z}) = \overline{x (y + z)} = \overline{x y + x z} = (\bar{x} \cdot \bar{y}) + (\bar{x} \cdot \bar{z})$.
    \item \textbf{Commutativity:} $\bar{x} \cdot \bar{y} = \overline{x y} = \overline{y x} = \bar{y} \cdot \bar{x}$.
\end{itemize}
\end{proof}
\end{content}

\begin{hidden}
% timestamp: 00:48:47
% topic: Completion of the quotient ring construction R/I.
% board_state: def:multiplication-quotient, prop:multiplication-well-defined, thm:quotient-ring-structure
% next_goal: Address student question regarding relationship of lectures to textbook and exam scope.
% open_loops: none
\end{hidden}

\begin{speech}[00:48:47 - 00:51:00]
Are there any questions about this?

People were asking me: what's the relationship of this class to the book? That's a really good question.

Every week I pick out some pages in that book --- it's a huge book ---, and somehow the class covers a very small part of the book, so it's not your goal to read that book from cover to cover, although it's not a bad thing to do!

Every week I pick out some topics there, and my lectures are roughly on those topics. Now, my notation is sometimes different, and the presentation is sometimes different in the book. The book is useful if you want... I skip some details in lectures; if you want all the details, sometimes they're in the book. That selection for the whole semester also covers the actual curriculum I'm required to cover, more or less, for this course, so it's helpful for me. But I would say: use the book if it's helpful for you; if not, you can use the lectures.

What happened when I taught this course at ETH in 2018 was that some students wrote some notes, and they're really great notes. Of course, they sent them to me afterwards, but it's $8$ years ago and I couldn't find them. But probably some of your colleagues write some good notes; you can also use that. The lectures are recorded, so you can also look at that.

Another option is that I gave some really good references --- some really classic algebra references. They have the unfortunate quality that their PDF file is not available through the library, but you can look at those books also. But if you look at the other books, the material won't be exactly in the order I'm giving you. If you want a book where it's in the order I'm giving you, then you should follow those page numbers in the actual reference; and that book is freely available on ITS, on the ETH website. So, that's the discussion about the book: use it if it's helpful.

For the exam material and things, I will only use material that I lectured on in class, not some exotic topic in the book that I didn't lecture on.

Okay.
\end{speech}

\begin{note}[Course References and Exam Scope]
The lecturer answers student inquiries regarding the course literature:
\begin{itemize}
    \item \textbf{Textbook Alignment:} Weekly reading assignments correspond to specific sections of the comprehensive course textbook (freely accessible to students via the ETH website). The textbook provides supplementary proofs and alternative presentations.
    \item \textbf{Exam Scope:} The examinations will strictly test the concepts, definitions, and theorems presented during the lectures, rather than extraneous textbook topics.
\end{itemize}
\end{note}

\section{The First Isomorphism Theorem for Rings}

\begin{speech}[00:51:00 - 00:52:37]
All right, so we go back to this issue now. We've made some progress. \inlinemetanote{writes $f \colon R \to S$, $R, S$ rings, $f$ hom on the upper middle board} If I have $f \colon R \to S$ --- where $R$ and $S$ are rings and $f$ is a homomorphism ---, then this is sometimes called the First Isomorphism Theorem. \inlinemetanote{writes First Isom. Theorem}

And it plays the role of the missing piece that we had in linear algebra that I wanted to develop here: that says that this ring, $R/\ker(f)$ \inlinemetanote{writes $R/\ker(f)$} --- because the kernel is an ideal, and if we have an ideal we can take the quotient, which is what the last $45$ minutes were about ---, this ring is canonically isomorphic to the image: $\cong \operatorname{Im}(f)$, \inlinemetanote{writes $\cong \operatorname{Im}(f)$} okay?

That plays the role of the statement in vector spaces. Is that clear? That's the replacement. So, we have to prove this.

Let us see what we have to do: we start with $f \colon R \to S$, \inlinemetanote{writes We start $f \colon R \to S$} and we will define a homomorphism $\bar{f} \colon R/\ker(f) \to S$. \inlinemetanote{writes We will define a hom $\bar{f} \colon R/\ker(f) \to S$} Just to start. And how do you think we're going to define this?
\end{speech}

\begin{content}[Statement of the First Isomorphism Theorem for Rings]
\begin{theorem}[First Isomorphism Theorem for Rings]\label[theorem]{thm:first-isomorphism-theorem-statement}
Let $f \colon R \to S$ be a homomorphism of commutative rings. Then the kernel $\ker(f)$ is an ideal of $R$, the image $\operatorname{Im}(f)$ is a subring of $S$, and there exists a canonical ring isomorphism:
\begin{equation}\label{eq:first-isomorphism-ring}
\bar{f} \colon R/\ker(f) \overset{\cong}{\longrightarrow} \operatorname{Im}(f), \qquad \bar{x} \mapsto f(x).
\end{equation}
In particular, if $f$ is surjective, then $R/\ker(f) \cong S$.
\end{theorem}
\end{content}

\begin{note}[Board Cleaning Transition]
\inlinemetanote{The lecturer washes the lower blackboard completely with a damp sponge from 00:52:37 to 00:53:30 to prepare the space for the three-step proof of the First Isomorphism Theorem.}
\end{note}

\begin{proofWrapper}[Proof of the First Isomorphism Theorem]

\subsection{Step 1: Construction and Well-Definedness of the Induced Map \texorpdfstring{$\bar{f}$}{fbar}}

\begin{speech}[00:53:30 - 00:54:07]
Okay, so we have to define $\bar{f}$. What does that mean? We have to define this: \inlinemetanote{writes $\bar{f}(\bar{x}) \overset{\text{def}}{=} f(x) \in S$ on the lower board} $\bar{f}$ is a homomorphism that starts from $R/\ker(f)$. What this thing eats is elements of $R/\ker(f)$ \inlinemetanote{writes $\bar{x} \in R/\ker(f)$} --- it eats equivalence classes.

What do you think we're going to do here? Yeah?
\end{speech}

\begin{student}[Student Answer]
Send it to $f(x)$?
\end{student}

\begin{speech}[00:54:08 - 00:54:34]
Yeah. In this kind of generality, there's very little you can do, because you can imagine $R$ is like one country, $S$ is another, and there's only one flight between them, and you've got to take it! There's just no other way. To get in here, you've got to get on $f$ somehow, and that's the only way. So, that's the definition.
\end{speech}

\begin{speech}[00:54:34 - 00:55:47]
When you define it like this, as I said, it's the standard pain about this naming business. And, you know, I say it slightly facetiously, but naming is an incredibly crucial thing. As you study mathematics more and more, you'll see this idea of what you can name, what you can name precisely, and how you can name it turns out to be an incredibly important, subtle point. So, it's good to get used to thinking about that.

Anyway, the point is: we have to show this is well-defined. \inlinemetanote{writes Show $\bar{f}$ is well defined}

What does that mean? Suppose we have somebody else in the equivalence class: $x' = x + i$, where $i \in \ker(f)$. \inlinemetanote{writes $x' = x + i, i \in \ker(f)$} If we take another representative in the same equivalence class, $\bar{x}' = \bar{x}$, \inlinemetanote{writes $\bar{x}' = \bar{x}$} it satisfies this equation. Then we must check that $\bar{f}(\bar{x}') = \bar{f}(\bar{x})$. \inlinemetanote{writes We must check $\bar{f}(\bar{x}') = \bar{f}(\bar{x})$} We have to check it.
\end{speech}

\begin{speech}[00:55:47 - 00:56:42]
Let's see: what is this? Here we use the definition: that's $f(x')$, \inlinemetanote{writes $\bar{f}(\bar{x}') \overset{\text{def}}{=} f(x')$} and then we use the definition of $x'$, so that's $f(x + i)$. \inlinemetanote{writes $= f(x+i)$}

I should have said here very carefully that $i \in \ker(f)$. \inlinemetanote{writes $i \in \ker(f)$} What do we do here? Well, we can expand this as $f(x) + f(i)$, \inlinemetanote{writes $= f(x) + f(i)$} because $f$ itself is a homomorphism. And now $i$ is in the kernel, so this is equal to $f(x)$, \inlinemetanote{writes $= f(x)$} because $f(i) = 0$, \inlinemetanote{writes $f(i) = 0$} and that's what we wanted, right? And this is equal to $\bar{f}(\bar{x})$. \inlinemetanote{writes $= \bar{f}(\bar{x})$} Okay.
\end{speech}

\begin{student}[Student Question]
Don't you need to first show that $\bar{f}$ is a homomorphism, because we are using the argument that $\bar{f}$ of $i$ bar, which is $\bar{f}$ of $0$ bar is $0$, but...
\end{student}

\begin{speech}[00:56:55 - 00:57:33]
No, I'm just showing it's well-defined: I only want to show this is equal to that.

All this $f$ stuff is fine, because $f$ is a homomorphism. Remember: $f$ is a homomorphism.

We have to do all the things you're saying, but there's a certain order you do it in. First, I define $\bar{f}$ --- this is part of the definition of $\bar{f}$. I propose a definition here. The problem with this definition is that it may depend on which element of the equivalence class I pick. And now I check that it doesn't depend on that. But I haven't proven this is a homomorphism yet; I will have to do that.
\end{speech}

\begin{student}[Student Question]
Where do we know that $i$ is in the kernel?
\end{student}

\begin{speech}[00:57:37 - 00:57:48]
Well, that's right here: for another representative, that's what it means, that their difference is in the kernel.
\end{speech}

\begin{speech}[00:57:48 - 00:58:45]
If that was confusing, I'll say it again, because it's kind of important.

I proposed this as a definition, but it has the flaw that I have to check that this definition does not depend on which element of the equivalence class I chose. And this is another element of the equivalence class. Any two elements of the equivalence class differ by an element of the kernel. So, that's the equation there, and I use this equation to show this map is well-defined.

Actually, I wrote out all the steps. In my experience in teaching, it doesn't always help writing out all the steps: there's a lot of steps, and every time you have to check what I'm using at every equality. You have to do that: I already know, but you have to verify it, and if I just write it, it doesn't always help.

So, that was a little painful. But we're not done with the pain!
\end{speech}

\begin{content}[Definition of $\bar{f}$ and Proof of Well-Definedness]
We define the induced map:
\begin{equation}\label{eq:induced-map-def}
\bar{f} \colon R/\ker(f) \longrightarrow S, \qquad \bar{f}(\bar{x}) := f(x).
\end{equation}

\begin{claim}[$\bar{f}$ is Well-Defined]\label[claim]{clm:fbar-well-defined}
The map $\bar{f}$ is independent of the choice of representative in the coset $\bar{x}$.
\end{claim}

\begin{proof}
Suppose $\bar{x}' = \bar{x}$ in $R/\ker(f)$. By definition of congruence modulo $\ker(f)$, there exists an element $i \in \ker(f)$ such that:
\[
x' = x + i.
\]
Applying the definition of $\bar{f}$ and the additive homomorphism property of $f$:
\begin{align*}
\bar{f}(\bar{x}') &= f(x') && \text{(by definition of $\bar{f}$)} \\
&= f(x + i) && \text{(since $x' = x + i$)} \\
&= f(x) + f(i) && \text{(since $f$ is a ring homomorphism)} \\
&= f(x) + 0_S && \text{(since $i \in \ker(f) \implies f(i) = 0_S$)} \\
&= f(x) \\
&= \bar{f}(\bar{x}).
\end{align*}
Thus, $\bar{f}(\bar{x}') = \bar{f}(\bar{x})$, proving that $\bar{f}$ is well-defined as a set-theoretic function.
\end{proof}
\end{content}

\begin{speech}[00:58:45 - 01:00:22]
This just shows it's well-defined, so we have $\bar{f} \colon R/\ker(f) \to S$ as a function now. \inlinemetanote{writes $\bar{f} \colon R/\ker(f) \to S$ function}

And we have to show all the properties to show it's a ring homomorphism.

We would have to show, for example, $\bar{f}(\bar{0})$: \inlinemetanote{writes $\bar{f}(\bar{0})$} what is this? We use our definition, and that's equal to $f(0)$, \inlinemetanote{writes $= f(0)$} and $f(0) = 0$ \inlinemetanote{writes $= 0$} because $f$ itself is a homomorphism. So, that's good: it's sending $0$ to $0$. That's maybe what you were worried about; now we've shown it.

We also have to show $\bar{f}(\bar{x} + \bar{y}) = \bar{f}(\bar{x}) + \bar{f}(\bar{y})$. \inlinemetanote{writes $\bar{f}(\bar{x}+\bar{y}) = \bar{f}(\bar{x}) + \bar{f}(\bar{y})$} You just have to expand these sides completely, and if you do that, you'll get a chain more or less similar to that, and you'll find that it's compatible with the definition.

So, expand it yourself! \inlinemetanote{writes expand yourself} As I said, you'll get a chain...
\end{speech}

\begin{speech}[01:00:22 - 01:02:03]
That's the additive structure. Then you also have to show that $\bar{f}$ respects the multiplicative structure: you have to check that $\bar{f}(\bar{x} \cdot \bar{y}) = \bar{f}(\bar{x}) \cdot \bar{f}(\bar{y})$. \inlinemetanote{writes check $\bar{f}(\bar{x} \cdot \bar{y}) = \bar{f}(\bar{x}) \cdot \bar{f}(\bar{y})$ on the lower left board}

If we do that: first, we have the definition of the product in the quotient, that's $\bar{f}(\overline{x y})$; \inlinemetanote{writes $\bar{f}(\overline{xy})$} and then, by definition of $\bar{f}$, that's $f(x y)$; \inlinemetanote{writes $= f(xy)$} and expanding this, that's $f(x) f(y)$. \inlinemetanote{writes $= f(x) f(y)$} And that's true because $f$ is a homomorphism --- true because $f$ was born a homomorphism! \inlinemetanote{writes True because $f$ is hom $R \to S$} Okay.

You have to concentrate with the bars, but really, the proof of this is just unraveling the definitions.
\end{speech}

\begin{content}[Proof that $\bar{f}$ is a Ring Homomorphism]
\begin{claim}[$\bar{f}$ is a Ring Homomorphism]\label[claim]{clm:fbar-is-homomorphism}
The well-defined map $\bar{f} \colon R/\ker(f) \to S$ preserves zero, identity, addition, and multiplication.
\end{claim}

\begin{proof}
Using the coset definitions in $R/\ker(f)$ and the ring homomorphism properties of $f \colon R \to S$:
\begin{enumerate}[label=\textbf{(\arabic*)}]
    \item \textbf{Preservation of Zero:}
    \[
    \bar{f}(\bar{0}) = f(0_R) = 0_S.
    \]
    \item \textbf{Preservation of Identity:}
    \[
    \bar{f}(\bar{1}) = f(1_R) = 1_S.
    \]
    \item \textbf{Preservation of Addition:}
    \begin{align*}
    \bar{f}(\bar{x} + \bar{y}) &= \bar{f}(\overline{x + y}) && \text{(by definition of addition in $R/\ker(f)$)} \\
    &= f(x + y) && \text{(by definition of $\bar{f}$)} \\
    &= f(x) + f(y) && \text{(since $f$ is a ring homomorphism)} \\
    &= \bar{f}(\bar{x}) + \bar{f}(\bar{y}). && \text{(by definition of $\bar{f}$)}
    \end{align*}
    \item \textbf{Preservation of Multiplication:}
    \begin{align*}
    \bar{f}(\bar{x} \cdot \bar{y}) &= \bar{f}(\overline{x \cdot y}) && \text{(by definition of multiplication in $R/\ker(f)$)} \\
    &= f(x \cdot y) && \text{(by definition of $\bar{f}$)} \\
    &= f(x) \cdot f(y) && \text{(since $f$ is a ring homomorphism)} \\
    &= \bar{f}(\bar{x}) \cdot \bar{f}(\bar{y}). && \text{(by definition of $\bar{f}$)}
    \end{align*}
\end{enumerate}
Therefore, $\bar{f}$ is a homomorphism of commutative rings.
\end{proof}
\end{content}

\begin{hidden}
% timestamp: 01:02:03
% topic: Well-definedness of induced map \bar{f} and verification that it is a ring homomorphism (Step 1).
% board_state: thm:first-isomorphism-theorem-statement, clm:fbar-well-defined, clm:fbar-is-homomorphism
% next_goal: Step 2: Surjectivity onto the image of f; Step 3: Injectivity via kernel of \bar{f}.
% open_loops: none
\end{hidden}

\subsection{Step 2: Surjectivity of \texorpdfstring{$\bar{f}$}{fbar} onto \texorpdfstring{$\operatorname{Im}(f)$}{Im(f)}}

\begin{speech}[01:02:03 - 01:03:46]
That's the first step in the proof. So, I remind you: we have $f \colon R \to S$, the ideal is $\ker(f)$, and we've defined $\bar{f} \colon R/\ker(f) \to S$. \inlinemetanote{writes $f \colon R \to S$, $\ker(f) \trianglelefteq R$, $\bar{f} \colon R/\ker(f) \to S$ is hom} And this is a ring homomorphism. We know the domain is a ring --- that was the first part of this class ---, we've defined $\bar{f}$, and we proved it's a ring homomorphism.

All right, step two: what is the image of $\bar{f}$? \inlinemetanote{writes Step 2: What is $\operatorname{Im}(\bar{f})$?} Where do we go here? The answer is: it's the same as the image of $f$. So, actually, $\bar{f}$ goes from $R/\ker(f)$ to the image of $f$. \inlinemetanote{writes $\bar{f} \colon R/\ker(f) \to \operatorname{Im}(f)$} This is some subring of $S$.

Why? Because the definition just says that if I take an equivalence class and choose a representative, I go to $f(x)$. So, everything I get is in the image of $f$. On the other hand, if I'm in the image of $f$, that means that I come from some element of $R$, and I can take its equivalence class. So, I get exactly the image of $f$.

That's step two, and this is a surjection. So, $\bar{f}$ is surjective onto the image of $f$, \inlinemetanote{writes $\bar{f}$ is surjective onto the image of $f$} okay?
\end{speech}

\begin{speech}[01:03:46 - 01:04:46]
Sometimes in math you write two arrows for surjectivity; \inlinemetanote{draws a two-headed arrow on the board} do you follow that kind of stuff or no? Okay, I won't give you more notation now.

First, we defined $\bar{f}$, and $\bar{f}$ is defined in the only possible way you can. The second step is: we realize it's surjective onto the image of $f$.
\end{speech}

\begin{content}[Surjectivity of the Induced Homomorphism onto the Image]
\begin{claim}[Surjectivity of $\bar{f}$ onto $\operatorname{Im}(f)$]\label[claim]{clm:fbar-surjective}
The image of $\bar{f}$ is precisely the subring $\operatorname{Im}(f) \subseteq S$. Consequently, the restricted codomain map:
\begin{equation}\label{eq:fbar-surjective-map}
\bar{f} \colon R/\ker(f) \twoheadrightarrow \operatorname{Im}(f)
\end{equation}
is a surjective ring homomorphism.
\end{claim}

\begin{proof}
We verify mutual inclusion of sets:
\begin{itemize}
    \item \textbf{$\operatorname{Im}(\bar{f}) \subseteq \operatorname{Im}(f)$:} For any coset $\bar{x} \in R/\ker(f)$, $\bar{f}(\bar{x}) = f(x) \in \operatorname{Im}(f)$ by definition of the image of $f$.
    \item \textbf{$\operatorname{Im}(f) \subseteq \operatorname{Im}(\bar{f})$:} Let $y \in \operatorname{Im}(f)$. By definition of the image, there exists an element $x \in R$ such that $f(x) = y$. Considering the coset $\bar{x} \in R/\ker(f)$, we have:
    \[
    \bar{f}(\bar{x}) = f(x) = y.
    \]
    Thus, $y \in \operatorname{Im}(\bar{f})$.
\end{itemize}
Hence, $\operatorname{Im}(\bar{f}) = \operatorname{Im}(f)$, proving that $\bar{f}$ maps $R/\ker(f)$ surjectively onto the subring $\operatorname{Im}(f)$.
\end{proof}
\end{content}

\subsection{Step 3: Injectivity and Trivial Kernel of \texorpdfstring{$\bar{f}$}{fbar}}

\begin{speech}[01:04:46 - 01:06:32]
And the third step --- this one requires some real concentration --- is: what is the kernel of $\bar{f}$? \inlinemetanote{writes Step 3: What is the $\ker(\bar{f})$?}

Let's try to think about that. Suppose I have $\bar{x} \in R/\ker(f)$, \inlinemetanote{writes Suppose $\bar{x} \in R/\ker(f)$} and $\bar{f}(\bar{x}) = 0 \in S$. \inlinemetanote{writes and $\bar{f}(\bar{x}) = 0 \in S$} This is another way of saying $\bar{x} \in \ker(\bar{f})$. \inlinemetanote{writes $(\bar{x} \in \ker(\bar{f}))$}

What can I conclude? Well, by definition, $\bar{f}(\bar{x}) = f(x)$, so this implies that $f(x) = 0 \in S$, \inlinemetanote{writes Then $f(x) = 0 \in S$} which implies $x \in \ker(f)$. \inlinemetanote{writes $\implies x \in \ker(f)$}

What have we done here? We said: we start with any equivalence class; if it's in the kernel of $\bar{f}$, then that element has to be in the kernel of $f$. And that means that equivalence class is the zero equivalence class, because it's equivalent to zero modulo the kernel: $x \equiv_{\ker(f)} 0$. \inlinemetanote{writes $\implies x \equiv_{\ker(f)} 0$} And this means that $\bar{x} = \bar{0} \in R/\ker(f)$. \inlinemetanote{writes $\implies \bar{x} = \bar{0} \in R/\ker(f)$}

So, that proves that the kernel of $\bar{f}$ is zero! That's step three: the conclusion is that the kernel of $\bar{f}$ is zero.
\end{speech}

\begin{content}[Injectivity of $\bar{f}$ via Triviality of the Kernel]
\begin{claim}[$\bar{f}$ is Injective]\label[claim]{clm:fbar-injective}
The kernel of the homomorphism $\bar{f} \colon R/\ker(f) \to \operatorname{Im}(f)$ is the zero ideal of the quotient ring:
\begin{equation}\label{eq:kernel-fbar-zero}
\ker(\bar{f}) = (\bar{0}).
\end{equation}
Consequently, $\bar{f}$ is injective.
\end{claim}

\begin{proof}
Let $\bar{x} \in \ker(\bar{f}) \subseteq R/\ker(f)$. By definition of the kernel:
\[
\bar{f}(\bar{x}) = 0_S.
\]
Using the definition of the map $\bar{f}(\bar{x}) := f(x)$, this means:
\[
f(x) = 0_S \implies x \in \ker(f).
\]
By definition of cosets modulo $\ker(f)$, an element belongs to the ideal if and only if its coset coincides with the zero coset:
\[
x \in \ker(f) \iff x \equiv_{\ker(f)} 0 \iff \bar{x} = \bar{0} \in R/\ker(f).
\]
Hence, the only element in $\ker(\bar{f})$ is the zero coset $\bar{0}$, establishing:
\[
\ker(\bar{f}) = (\bar{0}).
\]
By the injectivity criterion established earlier in \cref{prop:injective-kernel}, a ring homomorphism is injective if and only if its kernel is trivial. Therefore, $\bar{f}$ is injective.
\end{proof}

\begin{steps}[Conclusion of the Isomorphism]
Combining the three steps:
\begin{itemize}
    \item $\bar{f} \colon R/\ker(f) \to \operatorname{Im}(f)$ is a well-defined ring homomorphism (Step 1).
    \item $\bar{f}$ is surjective onto $\operatorname{Im}(f)$ (Step 2).
    \item $\bar{f}$ is injective, having trivial kernel (Step 3).
\end{itemize}
A bijective ring homomorphism is an isomorphism. Thus, $R/\ker(f) \cong \operatorname{Im}(f)$, which completes the proof of the First Isomorphism Theorem (\cref{thm:first-isomorphism-theorem-statement}).
\end{steps}
\end{content}

\end{proofWrapper}

\subsection{Recapitulation and the Anatomy of the Theorem}

\begin{speech}[01:06:32 - 01:07:37]
Those are the three steps: you define $\bar{f}$, you show it's surjective onto the image, and you show its kernel is zero.

The reason why the last one is confusing is you have to understand who is zero where. The kernel of $f$ is not zero; the kernel of $f$ is just whatever the kernel of $f$ is. But we quotient by the kernel of $f$, so the kernel of $\bar{f}$ is zero.

That's an admittedly confusing state of affairs, but you have to somehow come to terms with that.
\end{speech}

\begin{note}[Board Cleaning and Re-statement]
\inlinemetanote{The lecturer wipes the lower blackboard clean from 01:07:00 to 01:07:37 to write out the clean summary of the three objects and their respective zero elements.}
\end{note}

\begin{speech}[01:07:37 - 01:09:56]
I'll say it in words, because it's confusing. We start with $f \colon R \to S$. \inlinemetanote{writes $f \colon R \to S$ on the lower board} This has a kernel, $\ker(f) \trianglelefteq R$, \inlinemetanote{writes has kernel $\ker(f) \trianglelefteq R$} which could be anything --- maybe zero, maybe not, could be anything in the world.

Then, from this, we make $\bar{f} \colon R/\ker(f) \to \operatorname{Im}(f) \subseteq S$. \inlinemetanote{writes $\bar{f} \colon R/\ker(f) \to \operatorname{Im}(f) \subseteq S$} So, we study this. And the last step is that the kernel of $\bar{f}$ --- which is an ideal in $R/\ker(f)$ \inlinemetanote{writes $\ker(\bar{f}) \subseteq R/\ker(f)$} --- is the zero ideal in this quotient ring: $\ker(\bar{f}) = (0) \subseteq R/\ker(f)$. \inlinemetanote{writes $\ker(\bar{f}) = (0) \subseteq R/\ker(f)$}

The kernel of $f$ may not be zero, but when I consider the kernel of $\bar{f}$ in the quotient ring, it's zero!

And the way to think about that --- we'll discuss this more tomorrow --- is that you can write the kernel of $\bar{f}$ as kind of whatever was in the kernel before, but now quotiented out by $\ker(f)$, so that's just the zero group: \inlinemetanote{writes $\ker(\bar{f}) = \frac{\ker(f)}{\ker(f)} = 0$} $\frac{\ker(f)}{\ker(f)} = 0$. That's a nice way to think about it; I'll explain tomorrow how this kind of logic makes sense. Roughly speaking, there was a kernel, but when we quotient out, we quotient out by that kernel, so after we quotient out, we have no kernel left.

These three steps, taken together: first we defined $\bar{f}$, we showed it's surjective onto the image of $f$, and now we showed the kernel is zero, which means it's also injective. Kernel being zero implies that $\bar{f}$ is injective --- we started the lecture with that, I think. \inlinemetanote{writes $\implies \bar{f}$ is injective}

So, steps 1 to 3 show that $\bar{f} \colon R/\ker(f) \to \operatorname{Im}(f)$ \inlinemetanote{writes Steps 1--3: $\bar{f} \colon R/\ker(f) \to \operatorname{Im}(f)$} is well-defined, injective, and surjective, so $\bar{f}$ is an isomorphism of rings. \inlinemetanote{writes is isom of rings}

That's the end of the proof of the First Isomorphism Theorem.
\end{speech}

\begin{content}[Summary: Anatomy of the First Isomorphism Theorem]
The architecture of the First Isomorphism Theorem relies on distinguishing the three rings and their respective kernels:
\begin{center}
\begin{tabular}{l l l}
    \toprule
    \textbf{Mapping} & \textbf{Domain / Codomain} & \textbf{Kernel} \\
    \midrule
    $f$ & $R \longrightarrow S$ & $\ker(f) \trianglelefteq R$ (Arbitrary ideal of $R$) \\
    $\bar{f}$ & $R/\ker(f) \overset{\cong}{\longrightarrow} \operatorname{Im}(f)$ & $\ker(\bar{f}) = (\bar{0}) \trianglelefteq R/\ker(f)$ (The zero ideal in $R/\ker(f)$) \\
    \bottomrule
\end{tabular}
\end{center}

\begin{align}
f \colon R &\longrightarrow S && \text{with kernel } \ker(f) \trianglelefteq R, \label{eq:f-kernel-summary} \\
\bar{f} \colon R/\ker(f) &\overset{\sim}{\longrightarrow} \operatorname{Im}(f) \subseteq S && \text{with kernel } \ker(\bar{f}) = (0) \subseteq R/\ker(f). \label{eq:fbar-kernel-summary}
\end{align}
\end{content}

\begin{speech}[01:09:56 - 01:10:46]
There was a lot of chalk on the board, but if you really understand this, you'll see there's actually nothing happening in this proof: you just have to follow the definitions at every stage. It's similar to when you really understand that linear algebra result: the fact that the image is the space mod the kernel has almost nothing happening in that proof --- the definitions almost make it so, and that's what happens here also.

So, let's do an example.
\end{speech}

\section{Ideals Generated by Elements}

\begin{speech}[01:10:46 - 01:12:23]
I have some room here, so I can say: one of the things to think about rings, their ideals, and quotient rings is that we need to have some way of writing down ideals. Otherwise, we can't have too much fun with examples!

A standard way of writing an ideal is: let $x_1, \dots, x_n$ be elements of the ring $R$. \inlinemetanote{writes Let $x_1, \dots, x_n \in R$} They could also be infinite, but in this class typically it's finite.

If you have some list of elements in $R$, the standard notation is you put parentheses around them, and you define an ideal: \inlinemetanote{writes then $(x_1, \dots, x_n) \subseteq R$ is the ideal defined by $(x_1, \dots, x_n) = \{ \sum_{i=1}^n a_i x_i \mid a_i \in R \}$} this is the ideal defined by this equation, so I just have to tell you who is in this ideal.

This is equal to those elements that are of the form $a_1 x_1 + \dots + a_n x_n$. In words, you could say these are $R$-linear combinations of these elements: \inlinemetanote{writes $R$-linear combs of the $x_i$} $a_1 x_1 + \dots + a_n x_n$, where the $a_i$ are in $R$.

So, this is a particular set.
\end{speech}

\begin{content}[Definition of Ideals Generated by Elements]
Let $R$ be a commutative ring with identity $1$.

\begin{definition}[Ideal Generated by Finitely Many Elements]\label[definition]{def:ideal-generated-by-elements}
Let $x_1, \dots, x_n \in R$. The \newterm{ideal generated by} $x_1, \dots, x_n$, denoted by $(x_1, \dots, x_n)$, is the set of all $R$-linear combinations of the generators:
\begin{equation}\label{eq:finitely-generated-ideal-def}
(x_1, \dots, x_n) := \left\{ \sum_{i=1}^n a_i x_i \;\middle|\; a_1, \dots, a_n \in R \right\}.
\end{equation}
When $n = 1$, the ideal $(x) = \{ a x \mid a \in R \} = x R$ is called a \newterm{principal ideal}.
\end{definition}
\end{content}

\begin{speech}[01:12:23 - 01:13:39]
As I said, in words, you could say they're $R$-linear combinations of the $x_i$.

Another phrasing is that this is the \qt{ideal generated by the $x_i$}. \inlinemetanote{writes \qt{ideal generated by the $x_i$}} That's also language used: ideal generated by the $x_i$. Those are two different ways to express the same thing.

And it's just this set. You have to accept that this is an ideal, which means you have to show that zero is here --- of course zero is here if I take all these coefficients to be zero. You have to show that if something is here, then its negative is here: I just flip all the signs of all the $a_i$. You have to show that if something's here, when I add it to something else, it's still here: then I just add the coefficients of the linear combination. And finally, I have to show the multiplicative property: if something's here and I multiply it by $r$, then it's still there; but you just view that $r$ as a scalar, and it multiplies all these coefficients.

So, you have to check it --- actually, I just checked it for you!

This is a very nice way: if you take a ring, you can just pick any elements and you get an ideal. It's the ideal generated by that.

Are there questions about this? This is a really fundamental way of generating ideals.
\end{speech}

\begin{content}[Verification of the Ideal Axioms for \texorpdfstring{$(x_1, \dots, x_n)$}{(x1, ..., xn)}]
\begin{proposition}\label[proposition]{prop:finitely-generated-is-ideal}
For any elements $x_1, \dots, x_n \in R$, the set $(x_1, \dots, x_n)$ defined in \cref{def:ideal-generated-by-elements} is indeed an ideal of $R$. Moreover, it is the smallest ideal of $R$ containing the set $\{x_1, \dots, x_n\}$.
\end{proposition}

\begin{proof}
We check the two defining conditions of an ideal:
\begin{enumerate}[label=\textbf{(\arabic*)}]
    \item \textbf{Additive Subgroup:}
    \begin{itemize}
        \item Setting $a_1 = \dots = a_n = 0$ yields $0 = \sum_{i=1}^n 0 \cdot x_i \in (x_1, \dots, x_n)$.
        \item For any element $u = \sum_{i=1}^n a_i x_i$, its additive inverse is $-u = \sum_{i=1}^n (-a_i) x_i \in (x_1, \dots, x_n)$.
        \item If $u = \sum_{i=1}^n a_i x_i$ and $v = \sum_{i=1}^n b_i x_i$, their sum is:
        \[
        u + v = \sum_{i=1}^n (a_i + b_i) x_i \in (x_1, \dots, x_n).
        \]
    \end{itemize}
    \item \textbf{Multiplicative Absorption:}
    For any ring element $r \in R$ and $u = \sum_{i=1}^n a_i x_i \in (x_1, \dots, x_n)$:
    \[
    r \cdot u = r \sum_{i=1}^n a_i x_i = \sum_{i=1}^n (r a_i) x_i.
    \]
    Since $R$ is a ring, each coefficient $r a_i \in R$, so $r \cdot u \in (x_1, \dots, x_n)$.
\end{enumerate}
\end{proof}
\end{content}

\subsection{Quotients of Univariate Polynomial Rings}

\begin{speech}[01:13:39 - 01:14:56]
I have a few minutes, so let's do some examples.

Now you have a very flexible way of generating a lot of rings. And the rings that you have the power to generate already have all sorts of interesting, different behavior. Part of the course is exploring some of that structure --- by far not all of it!
\end{speech}

\begin{note}[Board Cleaning and Transition to Examples]
\inlinemetanote{The lecturer washes the left and middle blackboards with a sponge to prepare for concrete polynomial examples.}
\end{note}

\begin{speech}[01:14:56 - 01:16:26]
For example, I can take my ring $R$ to be polynomials in one variable: $R = \mathbb{C}[x]$. \inlinemetanote{writes $R = \mathbb{C}[x]$} You should be totally happy with this ring: it's polynomials with complex coefficients.

And I can take an element here, $f$, to be $x^2 + 1$. That's a polynomial. \inlinemetanote{writes $f = x^2 + 1 \in R$}

And then I can take the ideal --- I'm using the notation from before: if you have one element, that's an element; if I put parentheses around it, this is the ideal: $(f) = (x^2 + 1) \subseteq R$. \inlinemetanote{writes $(f) = (x^2+1) \subseteq R$ ideal}

Very concretely, using that definition, this is just the ideal of $f$ times a polynomial, where the polynomial is in the ring: \inlinemetanote{writes $(f) = \{ f \cdot p \mid p \in R \}$} it's just one linear combination. All this is, is polynomials that have an $x^2 + 1$ factor (i.e., multiples of $x^2 + 1$).

What we've done today is: we have a new ring now, $R$ modulo this ideal. \inlinemetanote{writes $R/(f) = \mathbb{C}[x]/(x^2+1)$} This is the quotient ring, $\mathbb{C}[x]/(x^2 + 1)$.

Are you happy to think about that? We take this polynomial ring, we take this ideal, and we have this quotient ring. Do you accept that all the symbols on this board have been defined?

So, I can ask you now any question about this ring. For example: is this an integral domain? \inlinemetanote{writes Is this an integral domain?}

What's your opinion about that?
\end{speech}

\begin{student}[Student Answer]
No, because $x-i$ times $x+i$... $x-i$ is not zero and $x+i$ is not zero.
\end{student}

\begin{hidden}
% timestamp: 01:16:26
% topic: Example: R = C[x] and the quotient C[x]/(x^2+1); student identifies zero divisors.
% board_state: def:ideal-generated-by-elements, prop:finitely-generated-is-ideal
% next_goal: Formalize the zero divisor argument in C[x]/(x^2+1); contrast with Q[x]/(x^2+1).
% open_loops: none
\end{hidden}

\begin{speech}[01:16:26 - 01:18:39]
Yeah, the answer is no! \inlinemetanote{writes No!}

And the reason is: I can take $\overline{x + i} \cdot \overline{x - i}$, \inlinemetanote{writes $\overline{x+i} \cdot \overline{x-i}$} because $x + i$ and $x - i$ are in the polynomial ring, so their classes are in the quotient ring. And by definition, this is equal to $\overline{(x + i)(x - i)} = \overline{x^2 + 1}$, \inlinemetanote{writes $= \overline{(x+i)(x-i)} = \overline{x^2+1}$} and this is zero: $\bar{0}$, right? \inlinemetanote{writes $= \bar{0}$}

This is the beginning of your argument. That's a totally solid start!

To prove that it's not an integral domain, you have to find two non-zero things that multiply to zero. You found two things that multiply to zero. Why are these non-zero?
\end{speech}

\begin{student}[Student Explanation]
Because the difference to zero... if you multiply $x^2 + 1$ by any polynomial, the degree is at least $2$.
\end{student}

\begin{speech}[01:18:39 - 01:18:56]
That's the correct argument! We have to show that $\overline{x + i} \neq \bar{0}$ and $\overline{x - i} \neq \bar{0}$. \inlinemetanote{writes But $\overline{x+i} \neq \bar{0}$ and $\overline{x-i} \neq \bar{0}$} If we have these three statements --- the product equals zero, and they're both not equal to zero ---, then it's not an integral domain. You agree with that?

And he gave the correct argument for that: because if $x + i$ were equivalent to zero, this would mean that $x + i$ is in this ideal generated by $x^2 + 1$. \inlinemetanote{writes $\implies x+i \in (x^2+1)$}

But the ideal generated by $x^2 + 1$ has $(x^2 + 1)$ times polynomials in it. And you can never lower the degree to a linear polynomial by taking a quadratic polynomial and multiplying it by another polynomial. So, the full argument says: the answer is no.
\end{speech}

\begin{content}[Zero Divisors in the Quotient Ring \texorpdfstring{$\mathbb{C}[x]/(x^2+1)$}{C[x]/(x^2+1)}]
\begin{example}[Non-Domain Quotient Ring]\label[example]{ex:cx-quotient-not-domain}
Consider the polynomial ring $R = \mathbb{C}[x]$ and the principal ideal generated by $f(x) = x^2 + 1$:
\[
I = (x^2 + 1) = \{ (x^2 + 1) p(x) \mid p(x) \in \mathbb{C}[x] \}.
\]
The quotient ring $\mathbb{C}[x]/(x^2 + 1)$ is \textbf{not an integral domain}.
\end{example}

\begin{proof}
In $\mathbb{C}[x]$, the polynomial $x^2 + 1$ factors into linear factors:
\[
x^2 + 1 = (x + i)(x - i).
\]
In the quotient ring $\mathbb{C}[x]/(x^2 + 1)$, consider the cosets $\overline{x + i}$ and $\overline{x - i}$:
\begin{enumerate}[label=\textbf{(\arabic*)}]
    \item \textbf{Their Product is Zero:}
    \[
    \overline{x + i} \cdot \overline{x - i} = \overline{(x + i)(x - i)} = \overline{x^2 + 1} = \bar{0} \in \mathbb{C}[x]/(x^2 + 1).
    \]
    \item \textbf{Neither Factor is Zero:}
    Suppose $\overline{x + i} = \bar{0}$. Then $x + i \in (x^2 + 1)$, which means there exists $p(x) \in \mathbb{C}[x]$ such that:
    \[
    x + i = (x^2 + 1) p(x).
    \]
    If $p(x) = 0$, the right-hand side is $0 \neq x + i$. If $p(x) \neq 0$, comparing degrees gives:
    \[
    1 = \deg(x + i) = \deg(x^2 + 1) + \deg(p) = 2 + \deg(p) \ge 2,
    \]
    a contradiction. Thus, $\overline{x + i} \neq \bar{0}$. By identical degree considerations, $\overline{x - i} \neq \bar{0}$.
\end{enumerate}
Since $\overline{x + i}$ and $\overline{x - i}$ are non-zero elements whose product is zero, they are zero divisors. Hence, $\mathbb{C}[x]/(x^2 + 1)$ is not an integral domain.
\end{proof}
\end{content}

\subsection{The Quotient \texorpdfstring{$\mathbb{Q}[x]/(x^2+1)$}{Q[x]/(x^2+1)} and the Gaussian Rationals}

\begin{speech}[01:18:56 - 01:20:56]
What about this one? This is kind of cutting closer to what this class is about, actually: $\mathbb{Q}[x]/(x^2+1)$. \inlinemetanote{writes $\mathbb{Q}[x]/(x^2+1)$}

Is this an integral domain? \inlinemetanote{writes Is this an integral domain?}

His argument you cannot use here, because $i$ is not a rational number, right? You agree you can't use the same argument.

So, I can ask the same question, though: what do you think the answer to that is?

This is really a question that the last part of the class is about, but I'm going to give you a proof of this. The answer is: this is an integral domain! \inlinemetanote{writes is an int. domain}

Actually, more than that! \inlinemetanote{laughs} It's a field! \inlinemetanote{writes Field}

It doesn't look like it immediately, but I will try to give you a proof in the next few minutes. There are a lot of ways to prove it, and a certain part of the class is about such things called field extensions. If you take the successor course in the spring, that's entirely about this kind of stuff: Galois theory, etc. But it's a very good example to think about.

So far, the only thing you're supposed to understand is that the argument proposed for the complex numbers does not work here. But let's try to find an argument that does work, that's related to this lecture.

In some sense, that question is kind of like searching for the complex numbers: because you know the roots of $x^2 + 1$, built into that question is this quest for the complex numbers.
\end{speech}

\begin{content}[The Rational Quotient \texorpdfstring{$\mathbb{Q}[x]/(x^2+1)$}{Q[x]/(x^2+1)}]
\begin{question}
Is the quotient ring $\mathbb{Q}[x]/(x^2 + 1)$ an integral domain?
\end{question}

\begin{answer}
\textbf{Yes! In fact, it is a field.}
Notice that the factorization $x^2 + 1 = (x + i)(x - i)$ used over $\mathbb{C}$ is unavailable in $\mathbb{Q}[x]$, because $i \notin \mathbb{Q}$. In fact, $x^2 + 1$ is \newterm{irreducible} over $\mathbb{Q}$. We will prove that $\mathbb{Q}[x]/(x^2 + 1)$ is isomorphic to the field of Gaussian rationals $\mathbb{Q}(i)$.
\end{answer}
\end{content}

\subsection{Construction of the Evaluation Homomorphism at \texorpdfstring{$i$}{i}}

\begin{speech}[01:20:56 - 01:22:36]
But we know the complex numbers, so we write that down. \inlinemetanote{writes $\mathbb{C}$} And I take $\mathbb{Q}$: \inlinemetanote{writes $\mathbb{Q} \hookrightarrow \mathbb{C}$} $\mathbb{Q}$ embeds in the complex numbers; there's no doubt about that.

And now I look at the polynomial ring: $\mathbb{Q}[x]$. \inlinemetanote{writes $\mathbb{Q}[x]$} And we proved that if I choose any value for $x$, I get a homomorphism from the polynomial ring. Do you remember that?

So, I'm going to choose a homomorphism $f$ that on constants just acts as the inclusion, but $f(x) = i$. \inlinemetanote{writes $f \colon \mathbb{Q}[x] \to \mathbb{C}$, $f(x) = i$} I get to do it! Because I told you: for the polynomial ring, if I know how the base ring maps to any other ring, I can always extend the map sending $x$ to anything I want. So, I could just do this --- it's my right from previous lectures.

In other words, I define a homomorphism $f$ from polynomials to the complex numbers, \inlinemetanote{writes $f \colon \mathbb{Q}[x] \to \mathbb{C}$} where $f(p(x)) = p(i)$, \inlinemetanote{writes $f(p(x)) = p(i)$} okay? Do you accept this?

Even if you forgot that theorem, you can just check that this is a legal definition: given any polynomial, evaluating at $i$ gives a complex number, and it's a ring homomorphism. \inlinemetanote{writes Ring hom} So, you can just prove that from scratch in your head because I gave you the formula, or you can remember it has to be true because you can send $x$ anywhere.
\end{speech}

\begin{hidden}
% timestamp: 01:22:36
% topic: Defining the evaluation homomorphism f: Q[x] -> C, p(x) -> p(i).
% board_state: def:evaluation-homomorphism-i
% next_goal: Compute Im(f) = Q(i) and apply the First Isomorphism Theorem.
% open_loops: none
\end{hidden}

\begin{content}[The Evaluation Homomorphism at \texorpdfstring{$i$}{i}]
Let $\iota \colon \mathbb{Q} \hookrightarrow \mathbb{C}$ denote the canonical inclusion of fields. By the universal property of polynomial rings, specifying the image of the indeterminate $x$ uniquely determines a ring homomorphism extending $\iota$.

\begin{center}
\begin{tikzpicture}[scale=1.2, >=Stealth]
    % \begin{hidden}
    % - Canvas: scale 1.2, commutative triangle with Q at (0,1.5), C at (3.5,1.5), Q[x] at (0,0)
    % - Elements: inclusion Q into C, canonical inclusion Q into Q[x], evaluation homomorphism f from Q[x] to C
    % - Labels: iota on top, inc on left, f: p(x) -> p(i) on hypotenuse
    % \end{hidden}
    \node (Q) at (0, 1.5) {$\mathbb{Q}$};
    \node (C) at (3.5, 1.5) {$\mathbb{C}$};
    \node (Qx) at (0, 0) {$\mathbb{Q}[x]$};
    
    \draw[right hook->, thick] (Q) -- node[above] {$\iota$} (C);
    \draw[right hook->, thick] (Q) -- node[left] {$\mathrm{inc}$} (Qx);
    \draw[->, thick, MidnightBlue] (Qx) -- node[below right] {$f \colon p(x) \mapsto p(i)$} (C);
\end{tikzpicture}
\end{center}

\begin{definition}[Evaluation Homomorphism at $i$]\label[definition]{def:evaluation-homomorphism-i}
We define the \newterm{evaluation homomorphism} at the imaginary unit $i \in \mathbb{C}$ by:
\begin{equation}\label{eq:eval-homomorphism-def}
f \colon \mathbb{Q}[x] \longrightarrow \mathbb{C}, \qquad f(p(x)) := p(i).
\end{equation}
Explicitly, if $p(x) = \sum_{k=0}^m a_k x^k \in \mathbb{Q}[x]$, then $f(p(x)) = \sum_{k=0}^m a_k i^k \in \mathbb{C}$.
\end{definition}

\begin{steps}[Homomorphism Property Verification]
For any $p(x), q(x) \in \mathbb{Q}[x]$:
\begin{align*}
f(p(x) + q(x)) &= (p + q)(i) = p(i) + q(i) = f(p(x)) + f(q(x)), \\
f(p(x) \cdot q(x)) &= (p \cdot q)(i) = p(i) \cdot q(i) = f(p(x)) \cdot f(q(x)), \\
f(1) &= 1.
\end{align*}
Thus, $f$ is a well-defined homomorphism of commutative rings with identity.
\end{steps}
\end{content}

\subsection{Image of the Evaluation Homomorphism and the First Isomorphism Theorem}

\begin{speech}[01:22:36 - 01:23:54]
What is the image of this? \inlinemetanote{writes What is $\operatorname{Im}(f) = ?$ on the middle board} What is the image? Do you know what that is? Yeah?
\end{speech}

\begin{student}[Student Answer]
I mean, if we consider it like polynomials of degree $1$ with a constant, we get the Gaussian rationals?
\end{student}

\begin{speech}[01:23:54 - 01:24:51]
Yeah, and you don't get any more than that! Because if you take a polynomial, all the quadratic terms give you $-1$, and all you get is $i$. So, the answer is: the image of this is what we'd call the Gaussian integers (i.e., Gaussian rationals $\mathbb{Q}(i)$). \inlinemetanote{writes $\mathbb{Q}(i)$}

If you take a polynomial with rational coefficients and you evaluate it at $i$, all the even terms just give you rational numbers, and all the odd terms give you $i$'s. And the coefficients are all rational numbers, and you can get any one you want here.

So, now we use our First Isomorphism Theorem! \inlinemetanote{writes First Isom. Theorem} That says that $\mathbb{Q}[x]$ modulo the kernel of $f$ is isomorphic to the image of $f$, which is the Gaussian integers (i.e., Gaussian rationals $\mathbb{Q}(i)$) --- which, by the way, is a field! It's a domain and it's a field: \inlinemetanote{writes $\mathbb{Q}[x]/\ker(f) \cong \operatorname{Im}(f) = \mathbb{Q}(i)$} $\mathbb{Q}[x]/\ker(f) \cong \operatorname{Im}(f) = \mathbb{Q}(i)$.
\end{speech}

\begin{content}[Image of $f$ and Application of the First Isomorphism Theorem]
\begin{proposition}[Image of the Evaluation Map]\label[proposition]{prop:image-eval-map}
The image of the evaluation homomorphism $f \colon \mathbb{Q}[x] \to \mathbb{C}$ is the field of Gaussian rationals:
\begin{equation}\label{eq:image-gaussian-rationals}
\operatorname{Im}(f) = \mathbb{Q}(i) := \{ a + b i \mid a, b \in \mathbb{Q} \} \subset \mathbb{C}.
\end{equation}
\end{proposition}

\begin{proof}
Let $p(x) = a_n x^n + \dots + a_1 x + a_0 \in \mathbb{Q}[x]$. Since $i^2 = -1$, every even power $i^{2k} = (-1)^k \in \mathbb{Q}$ and every odd power $i^{2k+1} = (-1)^k i$. Collecting the even and odd terms yields:
\[
p(i) = \left( \sum_{k \ge 0} a_{2k} (-1)^k \right) + \left( \sum_{k \ge 0} a_{2k+1} (-1)^k \right) i = a + b i,
\]
where $a, b \in \mathbb{Q}$. Conversely, for any $a + b i \in \mathbb{Q}(i)$, the linear polynomial $p(x) = b x + a \in \mathbb{Q}[x]$ satisfies $f(p(x)) = a + b i$. Thus, $\operatorname{Im}(f) = \mathbb{Q}(i)$.
\end{proof}

\begin{theorem}[First Isomorphism Theorem Application]\label[theorem]{thm:first-isom-applied-qi}
Applying the First Isomorphism Theorem (\cref{thm:first-isomorphism-theorem-statement}) to $f \colon \mathbb{Q}[x] \to \mathbb{C}$:
\begin{equation}\label{eq:isom-qi-statement}
\frac{\mathbb{Q}[x]}{\ker(f)} \cong \operatorname{Im}(f) = \mathbb{Q}(i).
\end{equation}
Since $\mathbb{Q}(i)$ is a subfield of $\mathbb{C}$, the quotient $\mathbb{Q}[x]/\ker(f)$ is an integral domain and a field.
\end{theorem}
\end{content}

\begin{aiNote}[Terminology Note: Gaussian Rationals vs.\ Gaussian Integers]
% [PERSISTENT]
In the lecture audio, the professor refers to $\mathbb{Q}(i) = \{ a + b i \mid a, b \in \mathbb{Q} \}$ as the \qt{Gaussian integers}. In standard algebraic terminology, the Gaussian integers are $\mathbb{Z}[i] = \{ a + b i \mid a, b \in \mathbb{Z} \}$, while $\mathbb{Q}(i) = \operatorname{Frac}(\mathbb{Z}[i])$ is the field of \newterm{Gaussian rationals}. The blackboard notation $\mathbb{Q}(i)$ is preserved verbatim.
\end{aiNote}

\subsection{Computing the Kernel via Polynomial Long Division}

\begin{speech}[01:24:51 - 01:27:11]
Now it would be nice to know what this kernel is: what is the kernel of $f$? \inlinemetanote{writes What is the ker of $f$? on the lower right board}

Well, one thing that's definitely in the kernel is this polynomial: $x^2 + 1$ is definitely in the kernel, \inlinemetanote{writes $x^2 + 1 \in \ker(f)$} because $f(x^2 + 1)$, by definition, is equal to $i^2 + 1 = 0$, right? \inlinemetanote{writes $f(x^2+1) = i^2 + 1 = 0$} So, whatever the kernel is, it definitely contains this polynomial.

Once this polynomial is in the kernel, then the whole ideal generated by this polynomial is also in the kernel: \inlinemetanote{writes $(x^2+1) \subseteq \ker(f)$} $(x^2 + 1) \subseteq \ker(f)$, because if I take this polynomial and multiply it by anything else, it will still be in the kernel.

The question now is: is there anything else? Is the ideal generated by this polynomial equal to the kernel of $f$? \inlinemetanote{writes Is $(x^2+1) = \ker(f)$?}

If the answer is yes, then I've answered that question! Spoiler: it's going to be yes. If it is yes, then I will have answered this question, because I've computed the kernel here and shown it's isomorphic to the field of Gaussian integers (i.e., rationals $\mathbb{Q}(i)$). That is a way of writing the field of Gaussian rationals in a slightly hidden form.

So, all I have to do is prove that $(x^2 + 1)$ is the kernel.
\end{speech}

\begin{content}[Containment of the Principal Ideal $(x^2+1)$ in the Kernel]
\begin{lemma}[Ideal Containment in the Kernel]\label[lemma]{lem:ideal-contained-in-kernel}
The principal ideal $(x^2 + 1)$ is contained in $\ker(f)$:
\begin{equation}\label{eq:ideal-contained-ker}
(x^2 + 1) \subseteq \ker(f).
\end{equation}
\end{lemma}

\begin{proof}
Evaluating $f$ on the generator $x^2 + 1$:
\[
f(x^2 + 1) = i^2 + 1 = -1 + 1 = 0 \implies x^2 + 1 \in \ker(f).
\]
Since $\ker(f)$ is an ideal of $\mathbb{Q}[x]$ (\cref{prop:kernel-is-ideal}), it must absorb multiplication by any polynomial $q(x) \in \mathbb{Q}[x]$:
\[
\forall q(x) \in \mathbb{Q}[x] \implies q(x)(x^2 + 1) \in \ker(f).
\]
Hence, $(x^2 + 1) \subseteq \ker(f)$.
\end{proof}
\end{content}

\begin{speech}[01:27:11 - 01:28:36]
We have $f \colon \mathbb{Q}[x] \to \mathbb{C}$, where $f$ eats a polynomial and gives back a complex number, $p(i)$. A polynomial is in the kernel of $f$ if and only if $p(i) = 0$. \inlinemetanote{writes $p \in \ker(f) \iff p(i) = 0$} In real life, you say $i$ is a root of $p$, right? \inlinemetanote{writes \qt{$i$ is a root of $p$}} So, the polynomials in the kernel are those where $i$ is a root.

We have to see how that's related to $x^2 + 1$. We could do the following: we try to divide $p(x)$ by $x^2 + 1$, \inlinemetanote{writes $x^2+1 \mid p(x)$ on the lower board} and everything is happening in this ring of polynomials $\mathbb{Q}[x]$. \inlinemetanote{writes in $\mathbb{Q}[x]$}

This is what's called polynomial long division, which you should have learned in school. If you have a polynomial, you can start dividing it: it will be equal to some polynomial $q(x)$ times $x^2 + 1$ plus a remainder term. \inlinemetanote{writes $p(x) = q(x)(x^2+1) + r(x)$}

Do you remember polynomial long division? You just keep taking out the highest powers all the way down. If you do polynomial long division, you start dividing and you'll eventually get some polynomial times $(x^2 + 1)$ plus a remainder term. And since this is a quadratic polynomial, this remainder term is a linear polynomial: \inlinemetanote{writes linear poly $r_1 x + r_2$, $r_1, r_2 \in \mathbb{Q}$} some $r_1 x + r_2$, where $r_1, r_2 \in \mathbb{Q}$. Of course, these could be zero.
\end{speech}

\begin{content}[Polynomial Division by \texorpdfstring{$x^2+1$}{x^2+1}]
\begin{lemma}[{Division Algorithm in $\mathbb{Q}[x]$}]\label[lemma]{lem:poly-division-algorithm}
For any polynomial $p(x) \in \mathbb{Q}[x]$, there exist unique polynomials $q(x), r(x) \in \mathbb{Q}[x]$ such that:
\begin{equation}\label{eq:poly-division-x2-plus-1}
p(x) = q(x)(x^2 + 1) + r(x), \quad \text{with } \deg(r) < 2.
\end{equation}
In particular, the remainder $r(x)$ is a linear polynomial with rational coefficients:
\begin{equation}\label{eq:remainder-linear-form}
r(x) = r_1 x + r_2, \qquad r_1, r_2 \in \mathbb{Q}.
\end{equation}
\end{lemma}
\end{content}

\begin{speech}[01:28:36 - 01:29:21]
That's polynomial long division; that's something you learned in school, and you should remember it now.

Now, I just stick $i$ into this equation: \inlinemetanote{writes $p(i) = q(i)(i^2+1) + r(i)$} I put $i$'s everywhere: $p(i) = q(i)(i^2 + 1) + r(i)$.

And now, the assumption is $p(i) = 0$, so the left-hand side is zero. This first term is also zero because $i^2 + 1 = 0$. That means the conclusion is: $r(i) = 0$. \inlinemetanote{writes $0 = q(i) \cdot 0 + r(i) \implies r(i) = 0$}

How could that be? The remainder $r$ is a linear polynomial. How could $r(i)$ be zero? There's only one way it could be zero, which is what? Yeah?
\end{speech}

\begin{student}[Student Answer]
Both coefficients are zero.
\end{student}

\begin{speech}[01:29:21 - 01:30:14]
Yeah, it's the only way! If this were a non-zero polynomial, since these coefficients are rational numbers, the only way this is zero when you stick in $i$ is if this is the zero polynomial, because $i$ is not a rational number.

So, this implies that $r$ is the zero polynomial! \inlinemetanote{writes $\implies r = 0$} So, actually it's not here, and $p(x)$ is in the ideal! \inlinemetanote{points to $p(x) = q(x)(x^2+1)$} And that finishes the argument, actually.

The conclusion is that $p(i) = 0$ implies $p(x) \in (x^2 + 1)$. \inlinemetanote{writes $p(i) = 0 \implies p(x) \in (x^2+1)$}

And the final conclusion, using the First Isomorphism Theorem, is that $\mathbb{Q}[x]/(x^2 + 1)$ \inlinemetanote{writes $\mathbb{Q}[x]/(x^2+1)$} is isomorphic \inlinemetanote{writes $\cong \mathbb{Q}(i)$} to this ring of Gaussian integers (i.e., Gaussian rationals $\mathbb{Q}(i)$).

Okay, made it in time! That's it. \inlinemetanote{lecturer finishes and packs up}
\end{speech}

\begin{content}[Proof that \texorpdfstring{$\ker(f) = (x^2+1)$}{ker(f) = (x^2+1)} and Final Isomorphism]
\begin{theorem}[Equality of Kernel and Principal Ideal]\label[theorem]{thm:kernel-equals-ideal}
The kernel of the evaluation homomorphism $f \colon \mathbb{Q}[x] \to \mathbb{C}$ at $i$ is precisely the principal ideal generated by $x^2 + 1$:
\begin{equation}\label{eq:kernel-equals-ideal}
\ker(f) = (x^2 + 1).
\end{equation}
\end{theorem}

\begin{proof}
We established in \cref{lem:ideal-contained-in-kernel} that $(x^2 + 1) \subseteq \ker(f)$. We now prove the reverse inclusion $\ker(f) \subseteq (x^2 + 1)$.

Let $p(x) \in \ker(f)$. By definition, $p(i) = 0$. By \cref{lem:poly-division-algorithm}, divide $p(x)$ by $x^2 + 1$:
\[
p(x) = q(x)(x^2 + 1) + r(x), \quad \text{where } r(x) = r_1 x + r_2 \text{ with } r_1, r_2 \in \mathbb{Q}.
\]
Evaluating both sides at $x = i$:
\[
p(i) = q(i)(i^2 + 1) + r(i).
\]
Since $p(i) = 0$ and $i^2 + 1 = -1 + 1 = 0$, we obtain:
\[
0 = q(i) \cdot 0 + r(i) \implies r(i) = 0.
\]
Substituting the linear form of $r(x)$:
\[
r(i) = r_1 i + r_2 = 0 \implies r_2 + r_1 i = 0.
\]
Since $r_1, r_2 \in \mathbb{Q}$ and $i \notin \mathbb{Q}$ (so $1$ and $i$ are linearly independent over $\mathbb{Q}$), this forces:
\[
r_1 = 0 \quad \text{and} \quad r_2 = 0 \implies r(x) = 0.
\]
Therefore, the remainder vanishes identically, and:
\[
p(x) = q(x)(x^2 + 1) \in (x^2 + 1).
\]
Thus, $\ker(f) \subseteq (x^2 + 1)$, which completes the proof that $\ker(f) = (x^2 + 1)$.
\end{proof}

\begin{corollary}[Realization of the Gaussian Rationals as a Quotient Ring]\label[corollary]{cor:gaussian-rationals-quotient}
Combining \cref{thm:first-isom-applied-qi} and \cref{thm:kernel-equals-ideal}, there is a canonical isomorphism of fields:
\begin{equation}\label{eq:gaussian-rationals-quotient-isom}
\frac{\mathbb{Q}[x]}{(x^2 + 1)} \overset{\cong}{\longrightarrow} \mathbb{Q}(i), \qquad \overline{p(x)} \mapsto p(i).
\end{equation}
In particular, $\mathbb{Q}[x]/(x^2 + 1)$ is an integral domain and a field.
\end{corollary}
\end{content}

\begin{overview}[Summary of Key Concepts]
\begin{itemize}
    \item \textbf{Polynomial and Power Series Rings:} The rings $R[x]$ and $R[[x]]$ preserve the integral domain property of $R$, relying on the highest-degree terms and lowest-degree orders, respectively.
    \item \textbf{Laurent Series and Rational Functions:} The field of rational functions $R(x)$ is $\operatorname{Frac}(R[x])$, while the field of formal Laurent series $K((x))$ (finite negative powers) is naturally isomorphic to $\operatorname{Frac}(K[[x]])$.
    \item \textbf{Ideals and Well-Defined Quotients:} An ideal $I \trianglelefteq R$ requires additive closure and multiplicative absorption ($R \cdot I \subseteq I$). In constructing the quotient ring $R/I$, the additive subgroup condition guarantees well-defined addition, while the absorption property is strictly required to ensure well-defined coset multiplication.
    \item \textbf{The First Isomorphism Theorem for Rings:} Every ring homomorphism $f \colon R \to S$ induces a canonical ring isomorphism $\bar{f} \colon R/\ker(f) \overset{\cong}{\to} \operatorname{Im}(f)$ via $\bar{x} \mapsto f(x)$.
    \item \textbf{Polynomial Quotients:} Modding out by $x^2+1$ over $\mathbb{C}$ produces zero divisors via $(x+i)(x-i) \equiv 0$, whereas over $\mathbb{Q}$, polynomial irreducibility and evaluation at $i$ yield the field of Gaussian rationals $\mathbb{Q}[x]/(x^2+1) \cong \mathbb{Q}(i)$.
\end{itemize}
\end{overview}

\end{document}
